\documentclass[review,11pt,3p,authoryear,nopreprintline]{elsarticle}

\usepackage[T1]{fontenc}
\usepackage{amsmath,amssymb,amsthm}
\usepackage{mathrsfs}
\usepackage{bm}
\usepackage{needspace}
\usepackage{algorithm}
\usepackage{algpseudocode}
\usepackage{xcolor}
\usepackage[colorlinks=true,citecolor=blue,linkcolor=blue,urlcolor=blue]{hyperref}

\newtheorem{theorem}{Theorem}
\newtheorem{proposition}{Proposition}
\newtheorem{corollary}{Corollary}
\newtheorem{lemma}{Lemma}
\newdefinition{example}{Example}
\newdefinition{definition}{Definition}
\newdefinition{remark}{Remark}
\newdefinition{property}{Property}
\newdefinition{condition}{Condition}

\begin{document}

\begin{frontmatter}
  \title{Failure Probability Bounds with Uncertain Surrogate Residual Distributions}

  \author[inst1]{Renlong Wang}
  \ead{rw106@rice.edu}
  \affiliation[inst1]{organization={Department of Civil and Environmental Engineering, Rice University}, city={Houston}, postcode={TX 77005}, country={USA}}

  \begin{abstract}
  \end{abstract}

  \begin{keyword}
  \end{keyword}
\end{frontmatter}

\section{Introduction}

\section{Problem Description}
\label{sec-problem}

Let $\bm X\sim\mathbb P$ be a random input taking values in $\mathcal X\subseteq\mathbb R^n$, and let $g:\mathcal X\to\mathbb R$ be a measurable limit-state function. The failure probability is defined as
\[
  p_f:=\Pr\{g(\bm X)\le0\}=\mathbb E_{\bm X\sim\mathbb P}\left[\mathbf 1_{\{g(\bm X)\le0\}}\right]=\int_{\mathcal X}\mathbf 1_{\{g(\bm x)\le0\}}\,\mathbb P(\mathrm d\bm x).
\]

Evaluating $g$ can require an expensive high-fidelity simulation. We use the training dataset
\[
  \mathcal D^{\mathrm{tr}}:=\left\{\left(\bm x_i^{\mathrm{tr}},g(\bm x_i^{\mathrm{tr}})\right)\right\}_{i=1}^{N^{\mathrm{tr}}}
\]
to fit a computationally efficient, measurable surrogate $\hat g:=\hat g_{\mathcal D^{\mathrm{tr}}}$. Model fitting and selection use only the training data and are completed before the test data are used. We condition on $\mathcal D^{\mathrm{tr}}$, treat $\hat g$ as fixed, and suppress this dependence in the notation.

Define the signed residual function
\[
  \xi(\bm x):=g(\bm x)-\hat g(\bm x),\qquad \forall \bm x\in\mathcal X.
\]
For a new input $\bm X\sim\mathbb P$ independent of the training data, let $\mathbb Q^\star$ be the distribution of $\xi^\star:=\xi(\bm X)$. Let $\Xi\subseteq\mathbb R$ be a known Borel set with $\mathbb Q^\star(\Xi)=1$; the default choice is $\Xi=\mathbb R$. For any Borel set $\mathcal S$ in a Euclidean space, write $\mathcal B(\mathcal S)$ for its Borel sigma-algebra. Let $\mathcal M_+(\mathcal S)$ denote the finite nonnegative Borel measures concentrated on $\mathcal S$, and let $\mathcal M_1(\mathcal S)$ denote the probability measures.

For a candidate residual law $\mathbb Q\in\mathcal M_1(\Xi)$, let $\Pi$ be a joint distribution on $\mathcal X\times\Xi$ with marginals $\Pi_{\bm X}=\mathbb P$ and $\Pi_\xi=\mathbb Q$. Its surrogate-based failure probability is
\begin{equation}
  p_{\hat f}(\Pi):=\Pr\{\hat g(\bm X)+\xi\le0\}=\int_{\mathcal X\times\Xi}\mathbf 1_{\{\hat g(\bm x)+\xi\le0\}}\,\Pi(\mathrm d\bm x,\mathrm d\xi).
  \label{eq-surrogate-failure-probability}
\end{equation}
For fixed $\hat g$, this functional is linear in $\Pi$. In particular, the true joint distribution $\Pi^\star$ satisfies $p_{\hat f}(\Pi^\star)=p_f$.

The residual marginal and its dependence on the input jointly determine the failure probability. We represent residual-marginal uncertainty by an ambiguity set $\mathcal Q$ calibrated to contain $\mathbb Q^\star$ with high confidence. For a given $\mathcal Q$, define the joint ambiguity set
\[
  \mathcal P(\mathcal Q) :=\left\{\Pi\in\mathcal M_1(\mathcal X\times\Xi):\Pi_{\bm X}=\mathbb P,\ \Pi_\xi\in\mathcal Q\right\}.
\]
The true residual is determined by the input, but its unknown functional relationship is not encoded in $\mathcal Q$. The set $\mathcal P(\mathcal Q)$ therefore allows all couplings with the specified marginals, including randomized conditional residual laws. The upper and lower failure probability bounds are
\begin{equation}
  \overline p_{\hat f}:=\sup_{\Pi\in\mathcal P(\mathcal Q)}p_{\hat f}(\Pi),
  \label{eq-upper-bound}
\end{equation}
and
\begin{equation}
  \underline p_{\hat f}:=\inf_{\Pi\in\mathcal P(\mathcal Q)}p_{\hat f}(\Pi),
  \label{eq-lower-bound}
\end{equation}
respectively. If $\mathbb Q^\star\in\mathcal Q$, then $\Pi^\star\in\mathcal P(\mathcal Q)$ and hence $\underline p_{\hat f}\le p_f\le\overline p_{\hat f}$.


Section~\ref{sec-ambiguity-set} constructs the residual ambiguity set and establishes finite-sample coverage of the failure probability bounds, simultaneously across response thresholds. Section~\ref{sec-reformulation} then derives their Lagrangian duals, proves strong duality through a finite reformulation, and obtains closed-form expressions. Section~\ref{sec-implementation} integrates these results into an existing surrogate modeling workflow.


\section{Ambiguity Set Design}
\label{sec-ambiguity-set}

Failure occurs when $\xi\le-\hat g(\bm X)$, so the residual cumulative distribution function (CDF) gives the marginal mass at or below each failure threshold. We use the D\"umbgen--Wellner (DW) construction \citep{L.J23} to control empirical CDF deviations simultaneously over all thresholds, with tolerances that vary across probability levels.

Let $q$ and $p$ denote an empirical probability and a candidate probability. Following \citet{B.J79} and the CDF confidence bands of \citet{O95}, we measure their local discrepancy by the Bernoulli Kullback--Leibler divergence
\begin{equation}
  \phi(q,p):=q\log\frac{q}{p}+(1-q)\log\frac{1-q}{1-p},\qquad \forall q\in[0,1],\, p\in(0,1),
  \label{eq-bernoulli-kl}
\end{equation}
where logarithms are natural and a term with zero numerator is interpreted as zero. For $q$ close to an interior $p$, the approximation $\phi(q,p)\approx(q-p)^2/[2p(1-p)]$ scales the deviation by the Bernoulli variance. Thus, the same absolute deviation receives a larger weight near the probability boundaries. The full divergence also covers zero and full empirical counts through $\phi(0,p)=-\log(1-p)$ and $\phi(1,p)=-\log p$. At the boundaries of the candidate probability, we set $\phi(0,0)=\phi(1,1)=0$, $\phi(q,0)=+\infty$ for $q>0$, and $\phi(q,1)=+\infty$ for $q<1$.

To compare deviations simultaneously across probability levels, we use the additive correction of \citet{L.J23}. For $v\in(0,1)$, define
\begin{equation*}
  \rho(v):=\log\log\frac{e}{4v(1-v)},\qquad \omega(v):=\log\!\left(1+\rho(v)^2\right).
\end{equation*}
Before sampling, we fix the divergence parameter $s=1$, which gives \eqref{eq-bernoulli-kl}, and the correction parameter $\nu=1$. Set $\psi_0(v):=\rho(v)+\omega(v)$, with $\psi_0(0)=\psi_0(1)=+\infty$, and define the bivariate correction
\begin{equation*}
  \psi(q,p):=\min_{v\in[\min\{q,p\},\,\max\{q,p\}]}\psi_0(v).
\end{equation*}
The functions $\rho$ and $\omega$ vanish at $v=1/2$ and increase toward the probability boundaries. The minimum over the interval between $q$ and $p$ is therefore large only when both probabilities lie in the same tail. This correction balances fluctuations across probability scales in the simultaneous CDF constraint.

We construct the empirical CDF from a test dataset
\[
  \mathcal D^{\mathrm{te}}:=\left\{\left(\bm x_j^{\mathrm{te}},g(\bm x_j^{\mathrm{te}})\right)\right\}_{j=1}^{N^{\mathrm{te}}}.
\]
The test inputs are independent draws from $\mathbb P$ and are independent of $\mathcal D^{\mathrm{tr}}$. We fix $N^{\mathrm{te}}\ge1$ before sampling and use these data to calibrate the residual ambiguity set.

Define the signed test residuals
\[
  \xi_j^{\mathrm{te}}:=\xi(\bm x_j^{\mathrm{te}})=g(\bm x_j^{\mathrm{te}})-\hat g(\bm x_j^{\mathrm{te}}),\qquad j=1,\ldots,N^{\mathrm{te}}.
\]
Conditional on $\mathcal D^{\mathrm{tr}}$, these residuals are independent and identically distributed according to $\mathbb Q^\star$. Their empirical CDF is
\begin{equation}
  \widehat F^{\mathrm{te}}(t):=\frac{1}{N^{\mathrm{te}}}\sum_{j=1}^{N^{\mathrm{te}}}\mathbf 1_{\{\xi_j^{\mathrm{te}}\le t\}},\qquad \forall t\in\mathbb R.
  \label{eq-empirical-cdf}
\end{equation}
For $\mathbb Q\in\mathcal M_1(\Xi)$, write $F_{\mathbb Q}(t):=\mathbb Q\left(({-\infty},t]\right)$. At a threshold $t$, $q=\widehat F^{\mathrm{te}}(t)$ is the mean of the Bernoulli indicators in \eqref{eq-empirical-cdf}, and $p=F_{\mathbb Q}(t)$ is the candidate probability. Thus $N^{\mathrm{te}}\phi(q,p)$ is the Bernoulli log-likelihood difference between the empirical and candidate probabilities. Combining it with the multiscale correction gives $N^{\mathrm{te}}\phi(q,p)-\psi(q,p)$.

At the probability boundaries, the corrected discrepancy is $-\infty$ for $q=p\in\{0,1\}$ and $+\infty$ for $p=0<q$ or $p=1>q$. These conventions extend the constraint to every residual threshold.

The following definition constrains the residual law through the corrected CDF discrepancy.

\begin{definition}[DW-based CDF ambiguity set]
  \label{definition-dw-ambiguity-set}
  Given the empirical CDF $\widehat F^{\mathrm{te}}$ and a threshold $\theta\ge0$, define
  \begin{equation} 
    \label{eq-cdf-ambiguity-set}
      \mathcal Q^{\mathrm{DW}}(\theta):=\left\{\mathbb Q\in\mathcal M_1(\Xi):\ 
      \sup_{t\in\mathbb R}\left[ N^{\mathrm{te}}\phi\!\left(\widehat F^{\mathrm{te}}(t),F_{\mathbb Q}(t)\right) 
       {}-\psi\!\left(\widehat F^{\mathrm{te}}(t),F_{\mathbb Q}(t)\right)\right]\le\theta\right\}. 
  \end{equation} 
\end{definition}

We calibrate $\theta$ against a distribution-free reference law. Following \citet{L.J23}, let $V_1,\ldots,V_{N^{\mathrm{te}}}$ be independent standard uniform variables, independent of the training and test data. Their empirical CDF is
\begin{equation*}
  G_{N^{\mathrm{te}}}(u):=\frac{1}{N^{\mathrm{te}}}\sum_{j=1}^{N^{\mathrm{te}}}\mathbf 1_{\{V_j\le u\}},\qquad \forall u\in[0,1].
\end{equation*}
Define the uniform reference statistic
\begin{equation}
  T_{N^{\mathrm{te}}}^{\mathrm U}:=\sup_{u\in(0,1)}\left[N^{\mathrm{te}}\phi\!\left(G_{N^{\mathrm{te}}}(u),u\right)-\psi\!\left(G_{N^{\mathrm{te}}}(u),u\right)\right].
  \label{eq-dw-reference-statistic}
\end{equation}

The following proposition calibrates $\theta$ to obtain finite-sample coverage of the residual law.

\begin{proposition}[Finite-sample threshold calibration]
  \label{proposition-01}
  Under the sampling conditions stated above, let $\delta\in(0,1)$ and define
  \begin{equation}
    \theta_{N^{\mathrm{te}}}(\delta):=\inf\left\{\tau\ge0:\Pr\left\{T_{N^{\mathrm{te}}}^{\mathrm U}\le\tau\right\}\ge1-\delta\right\}.
    \label{eq-dw-calibrated-threshold}
  \end{equation}

  Then, for every $\theta\ge\theta_{N^{\mathrm{te}}}(\delta)$, the DW-based CDF ambiguity set satisfies
  \begin{equation}
    \Pr\left\{\mathbb Q^\star\in\mathcal Q^{\mathrm{DW}}(\theta)\,\middle|\,\mathcal D^{\mathrm{tr}}\right\}\ge1-\delta.
    \label{eq-cdf-coverage}
  \end{equation}
\end{proposition}
\begin{proof}
  The reference statistic is finite almost surely, and its distribution function is right-continuous. Hence \eqref{eq-dw-calibrated-threshold} gives
  \[
    \Pr\left\{T_{N^{\mathrm{te}}}^{\mathrm U}>\theta_{N^{\mathrm{te}}}(\delta)\right\}\le\delta.
  \]
  Conditioning on $\mathcal D^{\mathrm{tr}}$ fixes $\hat g$ and $\mathbb Q^\star$. Define the generalized quantile function $F_{\mathbb Q^\star}^{-1}(u):=\inf\{t\in\mathbb R:F_{\mathbb Q^\star}(t)\ge u\}$ for all $u\in(0,1)$. The sample $F_{\mathbb Q^\star}^{-1}(V_j)$ for $j=1,\ldots,N^{\mathrm{te}}$, has the same conditional joint law as the test residuals. Thus, given $\mathcal D^{\mathrm{tr}}$,
  \[
    \left(\widehat F^{\mathrm{te}}(t)\right)_{t\in\mathbb R}\overset{\mathrm d}{=}\left(G_{N^{\mathrm{te}}}\!\left(F_{\mathbb Q^\star}(t)\right)\right)_{t\in\mathbb R}.
  \]
  For $0<F_{\mathbb Q^\star}(t)<1$, the discrepancy at $\mathbb Q^\star$ is evaluated at a probability level included in \eqref{eq-dw-reference-statistic}. When $F_{\mathbb Q^\star}(t)$ is zero or one, the empirical CDF equals that boundary almost surely, giving discrepancy $-\infty$. The supremum in \eqref{eq-cdf-ambiguity-set}, evaluated at $\mathbb Q^\star$, is therefore bounded by $T_{N^{\mathrm{te}}}^{\mathrm U}$ under this coupling. Since $\mathbb Q^\star(\Xi)=1$, we obtain
  \[
    \Pr\left\{\mathbb Q^\star\notin\mathcal Q^{\mathrm{DW}}\!\left(\theta_{N^{\mathrm{te}}}(\delta)\right)\,\middle|\,\mathcal D^{\mathrm{tr}}\right\}
    \le\Pr\left\{T_{N^{\mathrm{te}}}^{\mathrm U}>\theta_{N^{\mathrm{te}}}(\delta)\right\}\le\delta.
  \]
  For every test sample, $\theta\ge\theta_{N^{\mathrm{te}}}(\delta)$ implies $\mathcal Q^{\mathrm{DW}}\!\left(\theta_{N^{\mathrm{te}}}(\delta)\right)\subseteq\mathcal Q^{\mathrm{DW}}(\theta)$. This inclusion proves \eqref{eq-cdf-coverage} for every such $\theta$.
\end{proof}

Proposition~\ref{proposition-01} requires no continuity or moment assumptions on the residual law. A conservative upper bound on the reference quantile preserves \eqref{eq-cdf-coverage}; its numerical evaluation is described in Section~\ref{sec-implementation}.
% If calibration is performed by simulation, quantile estimation error must be accounted for in the confidence guarantee. A threshold chosen from the test data also preserves \eqref{eq-cdf-coverage}, provided $\theta\ge\theta_{N^{\mathrm{te}}}(\delta)$ almost surely, by the same nesting argument.

The following theorem transfers this coverage guarantee to the failure probability interval.

\begin{theorem}[Finite-sample failure probability bounds]
  \label{theorem-01}
  Under the sampling conditions of Proposition~\ref{proposition-01}, let $\delta\in(0,1)$ and $\theta\ge\theta_{N^{\mathrm{te}}}(\delta)$, with the calibrated threshold defined in \eqref{eq-dw-calibrated-threshold}. Then the bounds defined in \eqref{eq-upper-bound} and \eqref{eq-lower-bound}, evaluated with $\mathcal Q=\mathcal Q^{\mathrm{DW}}(\theta)$, satisfy
  \begin{equation}
    \Pr\left\{\underline p_{\hat f}\le p_f\le\overline p_{\hat f}\,\middle|\,\mathcal D^{\mathrm{tr}}\right\}\ge1-\delta.
    \label{eq-failure-probability-coverage}
  \end{equation}
\end{theorem}

\begin{proof}
  On the event $\mathbb Q^\star\in\mathcal Q^{\mathrm{DW}}(\theta)$, the true joint distribution $\Pi^\star$ belongs to $\mathcal P\left(\mathcal Q^{\mathrm{DW}}(\theta)\right)$. Hence
  \[
    \underline p_{\hat f}\le p_{\hat f}(\Pi^\star)=p_f\le\overline p_{\hat f}.
  \]
  By Proposition~\ref{proposition-01}, this event has conditional probability at least $1-\delta$ given $\mathcal D^{\mathrm{tr}}$, proving \eqref{eq-failure-probability-coverage}.
\end{proof}

Theorem~\ref{theorem-01} uses one residual-coverage event for both endpoints, so the confidence level applies to the entire interval. For fixed test data, $0\le\theta_1\le\theta_2$ implies $\mathcal Q^{\mathrm{DW}}(\theta_1)\subseteq\mathcal Q^{\mathrm{DW}}(\theta_2)$. The lower bound is nonincreasing and the upper bound nondecreasing in $\theta$, so higher confidence gives a wider or unchanged interval.

The following corollary extends the guarantee to all response thresholds using the same residual ambiguity set.

\begin{corollary}[Simultaneous bounds across response thresholds]
  \label{cor-simultaneous-threshold-bounds}
  Under the conditions of Theorem~\ref{theorem-01}, let $\mathcal Q=\mathcal Q^{\mathrm{DW}}(\theta)$. For $c\in\mathbb R$, write $p_f(c):=\Pr \left\{g(\bm X)\le c\right\}$ and define
  \[
    p_{\hat f}(\Pi;c):=\int_{\mathcal X\times\Xi}
      \mathbf1_{\{\hat g(\bm x)+\xi\le c\}}\,\Pi(\mathrm d\bm x,\mathrm d\xi),
  \]
  with lower and upper bounds
  \[
    \underline p_{\hat f}(c):=\inf_{\Pi\in\mathcal P(\mathcal Q)}p_{\hat f}(\Pi;c),
    \qquad
    \overline p_{\hat f}(c):=\sup_{\Pi\in\mathcal P(\mathcal Q)}p_{\hat f}(\Pi;c).
  \]
  Then
  \begin{equation}
    \Pr\left\{\underline p_{\hat f}(c)\le p_f(c)\le\overline p_{\hat f}(c),
      \ \forall c\in\mathbb R\,\middle|\,\mathcal D^{\mathrm{tr}}\right\}\ge1-\delta.
    \label{eq-threshold-curve-coverage}
  \end{equation}
  At $c=0$, these definitions recover the original failure probability and its bounds.
\end{corollary}
\begin{proof}
  On $\mathbb Q^\star\in\mathcal Q^{\mathrm{DW}}(\theta)$, the true joint law $\Pi^\star$ belongs to the same set $\mathcal P(\mathcal Q)$ for every $c$. Consequently,
  \[ 
    \underline p_{\hat f}(c) =\inf_{\Pi\in\mathcal P(\mathcal Q)} p_{\hat f}(\Pi;c) \le p_{\hat f}(\Pi^\star;c) =p_f(c) \le\sup_{\Pi\in\mathcal P(\mathcal Q)} p_{\hat f}(\Pi;c) =\overline p_{\hat f}(c), \qquad\forall c\in\mathbb R.
  \]
  Both bound curves are nondecreasing. Monotonicity and the CDF limits of $p_f$ imply that any violation is detected at a rational threshold or an atom of $g(\bm X)$, a countable set. The simultaneous event is therefore measurable. Proposition~\ref{proposition-01} gives the common residual-coverage event conditional probability at least $1-\delta$, proving \eqref{eq-threshold-curve-coverage}.
\end{proof}

Corollary~\ref{cor-simultaneous-threshold-bounds} permits response thresholds to be selected after observing the test data. Complementing the endpoints gives simultaneous bounds on $\Pr \left\{g(\bm X)>c\right\}$. The functions $g$, $\hat g$, and the law $\mathbb P$ remain fixed; fragility curves conditioned on hazard intensity require coverage for the corresponding conditional residual laws. The coverage argument also applies to any residual ambiguity set containing $\mathbb Q^\star$ with conditional probability at least $1-\delta$. It remains to evaluate the DW extrema in \eqref{eq-upper-bound} and \eqref{eq-lower-bound} and their response-threshold extensions.


\section{Tractable Reformulation}
\label{sec-reformulation}

Throughout this section, the training data, test residuals, and $\theta\ge0$ are fixed. Input and residual cell probabilities give finite primal and dual problems and closed-form bounds. The main results use $\Xi=\mathbb R$; Remark~\ref{rem-interval-support} treats known interval support.

For each empirical probability level $q_j:=j/N^{\mathrm{te}}$, define the scalar acceptance interval
\begin{equation}
  [\alpha_j(\theta),\beta_j(\theta)]
  :=\left\{p\in[0,1]:
  N^{\mathrm{te}}\phi(q_j,p)-\psi(q_j,p)\le\theta\right\},
  \qquad j=0,\ldots,N^{\mathrm{te}}.
  \label{eq-dw-rank-acceptance}
\end{equation}
Let $\xi_{(1)}^{\mathrm{te}}<\cdots<\xi_{(K)}^{\mathrm{te}}$ be the distinct ordered test residuals, and set
\begin{equation}
  c_0:=0,\quad
  c_k:=\sum_{j=1}^{N^{\mathrm{te}}}\mathbf1_{\{\xi_j^{\mathrm{te}}\le\xi_{(k)}^{\mathrm{te}}\}},\quad
  l_k:=\alpha_{c_k}(\theta),\quad
  u_k:=\beta_{c_{k-1}}(\theta),\qquad k=1,\ldots,K.
  \label{eq-dw-counts-and-levels}
\end{equation}
The counts retain all multiplicities, so $c_K=N^{\mathrm{te}}$ even when test residuals coincide. Set $\xi_{(0)}^{\mathrm{te}}:=-\infty$ and $\xi_{(K+1)}^{\mathrm{te}}:=+\infty$. Separating the observed values from the intervening intervals gives $J:=2K+1$ residual cells:
\begin{equation}
  \Xi_{2k}:=\left\{\xi_{(k)}^{\mathrm{te}}\right\}
    \quad\left(k=1,\ldots,K\right),
  \qquad
  \Xi_{2k+1}:=\left(\xi_{(k)}^{\mathrm{te}},\xi_{(k+1)}^{\mathrm{te}}\right)
    \quad\left(k=0,\ldots,K\right).
  \label{eq-residual-cells}
\end{equation}
For any candidate residual law $\mathbb Q\in\mathcal M_1(\Xi)$, the open rank probability measures mass strictly below an observed residual, while the closed rank probability also includes mass at that value. For $k=1,\ldots,K$, the partition gives
\[
  \mathbb Q\left\{\xi<\xi_{(k)}^{\mathrm{te}}\right\}
    =\sum_{j<2k}\mathbb Q(\Xi_j),
  \qquad
  \mathbb Q\left\{\xi\le\xi_{(k)}^{\mathrm{te}}\right\}
    =\sum_{j\le2k}\mathbb Q(\Xi_j).
\]
The two probabilities differ by the atom mass $\mathbb Q(\Xi_{2k})$.

For $i=0,\ldots,K$, the observed residual levels classify inputs by the location of the failure threshold $-\hat g(\bm x)$:
\begin{equation}
  \mathcal X_i:=\left\{\bm x\in\mathcal X:
    \xi_{(i)}^{\mathrm{te}}\le-\hat g(\bm x)<\xi_{(i+1)}^{\mathrm{te}}\right\}.
  \label{eq-input-cells}
\end{equation}
The input cell probabilities give the surrogate-response threshold probabilities through
\[
  \mathbb P\left\{\hat g(\bm X)\le-\xi_{(k)}^{\mathrm{te}}\right\}
    =\sum_{i=k}^{K}\mathbb P(\mathcal X_i),
  \qquad k=1,\ldots,K.
\]
The half-open convention assigns $-\hat g(\bm x)=\xi_{(k)}^{\mathrm{te}}$ to $\mathcal X_k$. The tail sum therefore includes threshold equalities even when the surrogate response has atoms. 
The two partitions satisfy
\[
  \mathcal X=\bigsqcup_{i=0}^{K}\mathcal X_i,
  \qquad
  \Xi=\bigsqcup_{j=1}^{J}\Xi_j.
\]
Figure~\ref{fig-residual-input-cells} illustrates the two partitions for $K=3$.
\begin{figure}[!htbp]
  \centering
  \includegraphics[width=0.5\columnwidth]{fig_residual_input_cells.png}
  \caption{Illustration of the residual and input-threshold partitions for $K=3$. The residual partition consists of $K$ singleton cells and $K+1$ open intervals, whereas the input partition is induced by half-open intervals of the failure threshold $-\hat g(\bm x)$. Here, $w_i=\mathbb P(\mathcal X_i)$ denotes the probability mass of input cell $\mathcal X_i$.}
  \label{fig-residual-input-cells}
\end{figure}

The following proposition expresses the DW ambiguity set through linear constraints on cumulative residual cell probabilities.

\begin{proposition}[Finite representation of the DW ambiguity set]
  \label{prop-finite-rank-constraints}
  Each acceptance set in \eqref{eq-dw-rank-acceptance} is a nonempty closed interval containing $q_j$. The DW ambiguity set has the finite representation
  \begin{equation}
    \mathcal Q^{\mathrm{DW}}(\theta)
    =\left\{\mathbb Q\in\mathcal M_1(\Xi):
      \sum_{j<2k}\mathbb Q(\Xi_j)\le u_k,\quad
      \sum_{j\le2k}\mathbb Q(\Xi_j)\ge l_k,\quad k=1,\ldots,K\right\}.
    \label{eq-finite-rank-constraints}
  \end{equation}
  The sequences $l_k$ and $u_k$ are nondecreasing and satisfy
  \[
    0\le l_k,u_k\le1,\qquad l_{k-1}\le u_k,\qquad l_k\le u_{k+1},
    \qquad k=1,\ldots,K,
  \]
  where $l_0:=0$ and $u_{K+1}:=1$.
\end{proposition}
\begin{proof}
  For fixed $q$, the Bernoulli divergence increases as $p$ moves away from $q$, while $\psi(q,p)$ cannot increase as the interval between $q$ and $p$ expands. The corrected discrepancy therefore decreases toward $q$ from below and increases above $q$. It is continuous for $p\in(0,1)$ and equals $-\psi_0(q)\le0$ at $p=q$. The boundary conventions in Section~\ref{sec-ambiguity-set} give a nonempty closed acceptance interval also for $q\in\{0,1\}$. Scalar inversion of the DW constraint therefore gives the CDF band \citep[Supplement, Section~S.6]{L.J23}
  \[
    \alpha_{N^{\mathrm{te}}\widehat F^{\mathrm{te}}(t)}(\theta)
      \le F_{\mathbb Q}(t)
      \le\beta_{N^{\mathrm{te}}\widehat F^{\mathrm{te}}(t)}(\theta),
    \qquad\forall t\in\mathbb R.
  \]

  For $k=1,\ldots,K-1$, the empirical CDF on $\left[\xi_{(k)}^{\mathrm{te}},\xi_{(k+1)}^{\mathrm{te}}\right)$ equals $c_k/N^{\mathrm{te}}$. Monotonicity makes the lower band constraint on this interval equivalent to $\mathbb Q\left\{\xi\le\xi_{(k)}^{\mathrm{te}}\right\}\ge\alpha_{c_k}(\theta)$. By continuity of probability from below, the upper band constraint is equivalent to $\mathbb Q\left\{\xi<\xi_{(k+1)}^{\mathrm{te}}\right\}\le\beta_{c_k}(\theta)$. The exterior intervals give $\mathbb Q\left\{\xi<\xi_{(1)}^{\mathrm{te}}\right\}\le\beta_0(\theta)$ and $\mathbb Q\left\{\xi\le\xi_{(K)}^{\mathrm{te}}\right\}\ge\alpha_{N^{\mathrm{te}}}(\theta)$; their other band constraints are automatic because $\alpha_0(\theta)=0$ and $\beta_{N^{\mathrm{te}}}(\theta)=1$. The partition identities above then yield \eqref{eq-finite-rank-constraints}.

  With $p$ fixed, the corrected discrepancy decreases as $q$ approaches $p$ from either side. For empirical levels $q<q'$, an accepted value $p<q$ at level $q'$ remains accepted at level $q$; an accepted value $p>q'$ at level $q$ remains accepted at level $q'$. Since each acceptance interval contains its empirical level, both endpoints are nondecreasing. The inequalities $\alpha_{c_k}(\theta)\le c_k/N^{\mathrm{te}}\le\beta_{c_k}(\theta)$ give the stated compatibility conditions. Singleton intervals are allowed: if $q_j=1/2$ and $\theta=0$, the interval is $\{1/2\}$ because $\psi(1/2,p)=0$.
\end{proof}

Proposition~\ref{prop-finite-rank-constraints} restricts only the residual cell masses, leaving the distribution within each cell free.

For a feasible joint law, define its matrix of joint cell probabilities by $\pi_{ij}:=\Pi(\mathcal X_i\times\Xi_j)$. The fixed input marginal and the residual restrictions imply $\bm\pi=(\pi_{ij})\in\mathcal C_\theta$, where
\begin{equation}
  \mathcal C_\theta:=\left\{\bm\pi\in\mathbb R_+^{(K+1)\times J}:
  \begin{array}{ll}
    \displaystyle\sum_{j=1}^{J}\pi_{ij}=\mathbb P(\mathcal X_i),&i=0,\ldots,K,\\[1mm]
    \displaystyle\sum_{i=0}^{K}\sum_{j<2k}\pi_{ij}\le u_k,&k=1,\ldots,K,\\[1mm]
    \displaystyle\sum_{i=0}^{K}\sum_{j\le2k}\pi_{ij}\ge l_k,&k=1,\ldots,K
  \end{array}\right\}.
  \label{eq-cell-mass-polytope}
\end{equation}
To bound the failure objective, compare the indicator on each product cell:
\begin{equation}
  \underline h_{ij}:=\mathbf1_{\{j\le2i\}}
  \le\mathbf1_{\{\xi\le-\hat g(\bm x)\}}
  \le\overline h_{ij}:=\mathbf1_{\{j\le2i+1\}},
  \qquad \forall (\bm x,\xi)\in\mathcal X_i\times\Xi_j.
  \label{eq-cell-costs}
\end{equation}
For inputs in $\mathcal X_i$, the failure status can vary only within $\Xi_{2i+1}$. The other cells are uniformly failing or safe, as illustrated in Figure~\ref{fig-joint-partition} for $K=3$.
\begin{figure}[!htbp]
  \centering
  \includegraphics[width=0.5\columnwidth]{fig_joint_partition.png}
  \caption{Joint partition of the input and residual spaces for $K=3$. Each product cell has probability mass $\pi_{ij}=\Pi(\mathcal X_i\times\Xi_j)$. The blue, yellow, and red cells correspond to $(\underline h_{ij},\overline h_{ij})=(0,0)$, $(0,1)$, and $(1,1)$, respectively.}
  \label{fig-joint-partition}
\end{figure}

The upper problem remains a linear program over joint measures, with the full input marginal fixed:
\begin{equation}
  \begin{aligned}
    \overline p_{\hat f}
    =\sup_{\Pi\in\mathcal M_+(\mathcal X\times\Xi)}\quad&
      \int_{\mathcal X\times\Xi}\mathbf1_{\{\xi\le-\hat g(\bm x)\}}\,\Pi(\mathrm d\bm x,\mathrm d\xi)\\
    \text{s.t.}\quad&
      \int_{\mathcal X\times\Xi}\mathbf1_{\mathcal U}(\bm x)\,\Pi(\mathrm d\bm x,\mathrm d\xi)
        =\mathbb P(\bm X\in\mathcal U),&&\forall\mathcal U\in\mathcal B(\mathcal X),\\
    &\int_{\mathcal X\times\Xi}\mathbf1_{\{\xi<\xi_{(k)}^{\mathrm{te}}\}}\,\Pi(\mathrm d\bm x,\mathrm d\xi)\le u_k,
        &&k=1,\ldots,K,\\
    &\int_{\mathcal X\times\Xi}\mathbf1_{\{\xi\le\xi_{(k)}^{\mathrm{te}}\}}\,\Pi(\mathrm d\bm x,\mathrm d\xi)\ge l_k,
        &&k=1,\ldots,K.
  \end{aligned}
  \label{eq-measure-upper}
\end{equation}
The lower problem minimizes the same integral under the same constraints. The marginal equalities fix both the input law and total mass one. The product of $\mathbb P$ and the empirical residual law is feasible for every $\theta\ge0$.

The following lemma gives finite-dimensional Lagrangian duals and establishes weak duality.

\begin{lemma}[Finite-dimensional Lagrangian duals and weak duality]
  \label{lem-weak-duality}
  Write $\bm u=(u_1,\ldots,u_K)^\top$ and $\bm l=(l_1,\ldots,l_K)^\top$. The nonnegative multipliers $\bm\lambda,\bm\gamma\in\mathbb R_+^K$ correspond to the upper and lower rank constraints. The input marginal equality has an unrestricted multiplier $\eta\in\mathcal F_{\mathrm b}(\mathcal X)$, the space of bounded Borel measurable real-valued functions. Define the residual penalty
  \begin{equation}
    \zeta_{\bm\lambda,\bm\gamma}(\xi)
    :=\sum_{k=1}^{K}\lambda_k\mathbf1_{\{\xi<\xi_{(k)}^{\mathrm{te}}\}}
      -\sum_{k=1}^{K}\gamma_k\mathbf1_{\{\xi\le\xi_{(k)}^{\mathrm{te}}\}},
    \label{eq-dual-step-function}
  \end{equation}
  whose value on $\Xi_j$ is
  \begin{equation}
    \zeta_j:=\sum_{k=1}^{K}\lambda_k\mathbf1_{\{j<2k\}}
      -\sum_{k=1}^{K}\gamma_k\mathbf1_{\{j\le2k\}},
    \qquad j=1,\ldots,J.
    \label{eq-dual-cell-values}
  \end{equation}
  Eliminating the input marginal multiplier gives the upper and lower Lagrangian dual problems
  \begin{align}
    \overline d_\theta
      &:=\inf_{\bm\lambda,\bm\gamma\ge\bm0}\left\{
        \bm\lambda^\top\bm u-\bm\gamma^\top\bm l
        +\sum_{i=0}^{K}\mathbb P(\mathcal X_i)\max_{1\le j\le J}
          \left\{\overline h_{ij}-\zeta_j\right\}\right\},
      \label{eq-upper-measure-dual}\\
    \underline d_\theta
      &:=\sup_{\bm\lambda,\bm\gamma\ge\bm0}\left\{
        -\bm\lambda^\top\bm u+\bm\gamma^\top\bm l
        +\sum_{i=0}^{K}\mathbb P(\mathcal X_i)\min_{1\le j\le J}
          \left\{\underline h_{ij}+\zeta_j\right\}\right\}.
      \label{eq-lower-measure-dual}
  \end{align}
  Their optimal values satisfy
  \begin{equation}
    \underline d_\theta\le\underline p_{\hat f}
      \le\overline p_{\hat f}\le\overline d_\theta.
    \label{eq-weak-duality}
  \end{equation}
\end{lemma}
\begin{proof}
  The input marginal equality is equivalent to
  \[
    \int_{\mathcal X\times\Xi}\eta(\bm x)\,\Pi(\mathrm d\bm x,\mathrm d\xi)
      =\mathbb E_{\mathbb P}\left[\eta(\bm X)\right],
    \qquad\forall\eta\in\mathcal F_{\mathrm b}(\mathcal X).
  \]
  Relaxing this equality and the rank inequalities gives the upper Lagrangian \citep[Section~5.1]{B.V04}
  \[
    \begin{aligned}
      \mathcal L^+(\Pi,\eta,\bm\lambda,\bm\gamma)
      ={}&\int_{\mathcal X\times\Xi}\mathbf1_{\{\xi\le-\hat g(\bm x)\}}\,\Pi(\mathrm d\bm x,\mathrm d\xi)
        +\mathbb E_{\mathbb P}\left[\eta(\bm X)\right]
        -\int_{\mathcal X\times\Xi}\eta(\bm x)\,\Pi(\mathrm d\bm x,\mathrm d\xi)\\
      &+\sum_{k=1}^{K}\lambda_k\left(u_k-\int_{\mathcal X\times\Xi}
          \mathbf1_{\{\xi<\xi_{(k)}^{\mathrm{te}}\}}\,\Pi(\mathrm d\bm x,\mathrm d\xi)\right)
        +\sum_{k=1}^{K}\gamma_k\left(\int_{\mathcal X\times\Xi}
          \mathbf1_{\{\xi\le\xi_{(k)}^{\mathrm{te}}\}}\,\Pi(\mathrm d\bm x,\mathrm d\xi)-l_k\right)\\
      ={}&\mathbb E_{\mathbb P}\left[\eta(\bm X)\right]
        +\bm\lambda^\top\bm u-\bm\gamma^\top\bm l
        +\int_{\mathcal X\times\Xi}\left[
          \mathbf1_{\{\xi\le-\hat g(\bm x)\}}-\eta(\bm x)-\zeta_{\bm\lambda,\bm\gamma}(\xi)
        \right]\,\Pi(\mathrm d\bm x,\mathrm d\xi).
    \end{aligned}
  \]
  Its supremum over finite nonnegative measures is finite exactly when the integrand is nonpositive at every $(\bm x,\xi)$. A positive value at any point permits an arbitrarily large point mass there; otherwise the supremum integral is zero, attained by the zero measure. The functional upper dual therefore minimizes $\mathbb E_{\mathbb P}[\eta(\bm X)]+\bm\lambda^\top\bm u-\bm\gamma^\top\bm l$ subject to $\eta(\bm x)\ge\mathbf1_{\{\xi\le-\hat g(\bm x)\}}-\zeta_{\bm\lambda,\bm\gamma}(\xi)$ at every pair. For the lower problem, reversing the rank penalties gives
  \[
    \mathcal L^-
      =\mathbb E_{\mathbb P}\left[\eta(\bm X)\right]
        -\bm\lambda^\top\bm u+\bm\gamma^\top\bm l
        +\int_{\mathcal X\times\Xi}\left[
          \mathbf1_{\{\xi\le-\hat g(\bm x)\}}-\eta(\bm x)+\zeta_{\bm\lambda,\bm\gamma}(\xi)
        \right]\,\Pi(\mathrm d\bm x,\mathrm d\xi).
  \]
  Its infimum over finite nonnegative measures is finite exactly when the integrand is nonnegative everywhere, in which case the infimum integral is zero. The lower functional dual maximizes $\mathbb E_{\mathbb P}[\eta(\bm X)]-\bm\lambda^\top\bm u+\bm\gamma^\top\bm l$ subject to $\eta(\bm x)\le\mathbf1_{\{\xi\le-\hat g(\bm x)\}}+\zeta_{\bm\lambda,\bm\gamma}(\xi)$ at every pair.

  For fixed rank multipliers, the best marginal functions are the pointwise envelopes. For $\bm x\in\mathcal X_i$, they satisfy the identities
  \begin{equation}
    \begin{aligned}
      \eta^+_{\bm\lambda,\bm\gamma}(\bm x)
        &:=\sup_{\xi\in\Xi}\left\{\mathbf1_{\{\xi\le-\hat g(\bm x)\}}
          -\zeta_{\bm\lambda,\bm\gamma}(\xi)\right\}
          =\max_{1\le j\le J}\left\{\overline h_{ij}-\zeta_j\right\},\\
      \eta^-_{\bm\lambda,\bm\gamma}(\bm x)
        &:=\inf_{\xi\in\Xi}\left\{\mathbf1_{\{\xi\le-\hat g(\bm x)\}}
          +\zeta_{\bm\lambda,\bm\gamma}(\xi)\right\}
          =\min_{1\le j\le J}\left\{\underline h_{ij}+\zeta_j\right\}.
    \end{aligned}
    \label{eq-dual-envelope-identities}
  \end{equation}
  To verify them, note that outside $\Xi_{2i+1}$ the two cell coefficients equal the actual failure indicator. In the overlapping interval, the lower coefficient zero is attained by $\xi=\left(\xi_{(i+1)}^{\mathrm{te}}-\hat g(\bm x)\right)/2$ if $i<K$, and by $\xi=-\hat g(\bm x)+1$ if $i=K$. The upper coefficient one is attained by $\xi=-\hat g(\bm x)$ whenever $-\hat g(\bm x)>\xi_{(i)}^{\mathrm{te}}$. The only remaining case is $-\hat g(\bm x)=\xi_{(i)}^{\mathrm{te}}$ with $i\ge1$. Then the overlapping interval is safe, but
  \[
    \zeta_{2i}=\zeta_{2i+1}-\gamma_i\le\zeta_{2i+1},
    \qquad
    1-\zeta_{2i}\ge1-\zeta_{2i+1}.
  \]
  Choosing the left endpoint therefore attains an upper-envelope value at least as large as the candidate from the overlapping interval. Together with \eqref{eq-cell-costs}, this proves both identities, including threshold equalities.

  The envelopes are bounded and constant on the Borel cells $\mathcal X_i$, so they belong to $\mathcal F_{\mathrm b}(\mathcal X)$. They respectively minimize and maximize the expectation over feasible marginal functions. Substituting them into the functional dual objectives gives \eqref{eq-upper-measure-dual} and \eqref{eq-lower-measure-dual}. Introducing epigraph variables $z_i$ for the upper maxima and hypograph variables $z_i$ for the lower minima gives the equivalent LPs
  \begin{equation}
    \begin{aligned}
      \overline d_\theta
      =\inf_{\bm\lambda,\bm\gamma\ge\bm0,\ \bm z\in\mathbb R^{K+1}}\quad&
        \bm\lambda^\top\bm u-\bm\gamma^\top\bm l+\sum_{i=0}^{K}\mathbb P(\mathcal X_i) z_i\\
      \text{s.t.}\quad&z_i+\zeta_j\ge\overline h_{ij},
        \qquad i=0,\ldots,K,\quad j=1,\ldots,J,
    \end{aligned}
    \label{eq-finite-upper-dual}
  \end{equation}
  and
  \begin{equation}
    \begin{aligned}
      \underline d_\theta
      =\sup_{\bm\lambda,\bm\gamma\ge\bm0,\ \bm z\in\mathbb R^{K+1}}\quad&
        -\bm\lambda^\top\bm u+\bm\gamma^\top\bm l+\sum_{i=0}^{K}\mathbb P(\mathcal X_i) z_i\\
      \text{s.t.}\quad&z_i-\zeta_j\le\underline h_{ij},
        \qquad i=0,\ldots,K,\quad j=1,\ldots,J.
    \end{aligned}
    \label{eq-finite-lower-dual}
  \end{equation}
  Each $z_i$ can be chosen as its corresponding finite envelope, including on input cells with $\mathbb P(\mathcal X_i)=0$. Thus introducing these variables preserves both optimal values.

  For a feasible joint law, the signs of the rank penalties give $\mathcal L^-\le p_{\hat f}(\Pi)\le\mathcal L^+$. Feasible lower dual objectives therefore bound the primal objectives from below, and feasible upper dual objectives bound them from above. Taking the respective optima over the common nonempty primal feasible set proves \eqref{eq-weak-duality}.
\end{proof}

Lemma~\ref{lem-weak-duality} eliminates the function-valued marginal multiplier, leaving $2K$ nonnegative multipliers. Equations~\eqref{eq-finite-upper-dual} and \eqref{eq-finite-lower-dual} give equivalent LP forms.

The following theorem proves that the finite cell allocations and their duals recover the original failure probability bounds.

\begin{theorem}[Exact finite reformulations and strong duality]
  \label{thm-finite-dual-reformulations}
  For $\mathcal Q=\mathcal Q^{\mathrm{DW}}(\theta)$ and every $\theta\ge0$, the failure probability bounds admit the exact representations
  \begin{equation}
    \overline p_{\hat f}
      =\max_{\bm\pi\in\mathcal C_\theta}\sum_{i=0}^{K}\sum_{j=1}^{J}\overline h_{ij}\pi_{ij},
    \qquad
    \underline p_{\hat f}
      =\min_{\bm\pi\in\mathcal C_\theta}\sum_{i=0}^{K}\sum_{j=1}^{J}\underline h_{ij}\pi_{ij}.
    \label{eq-finite-primal-values}
  \end{equation}
  Their Lagrangian dual values in Lemma~\ref{lem-weak-duality} satisfy
  \begin{equation}
    \overline p_{\hat f}=\overline d_\theta,
    \qquad
    \underline p_{\hat f}=\underline d_\theta.
    \label{eq-strong-duality}
  \end{equation}
  The finite primal and dual optima, the measure primal optima, and the functional dual optima are attained.
\end{theorem}
\begin{proof}
  Any feasible joint law induces a matrix in $\mathcal C_\theta$. Integrating \eqref{eq-cell-costs} therefore bounds the upper measure problem by the finite maximization from above and the lower measure problem by the finite minimization from below.

  For the reverse inequalities, fix $\bm\pi\in\mathcal C_\theta$. Draw $\bm x$ from $\mathbb P$. On $\mathcal X_i$ with $\mathbb P(\mathcal X_i)>0$, select a provisional residual cell $j$ with probability $\pi_{ij}/\mathbb P(\mathcal X_i)$; the row equality makes these probabilities sum to one. For $j\ne2i+1$, choose a fixed representative of $\Xi_j$, where the actual failure indicator equals both coefficients. In the overlapping cell, use the following assignments to define joint laws $\Pi^+$ and $\Pi^-$.

  For the upper problem, assign $\xi=-\hat g(\bm x)$. This causes failure and remains in $\Xi_{2i+1}$ unless $-\hat g(\bm x)=\xi_{(i)}^{\mathrm{te}}$. At equality, which requires $i\ge1$, the assignment transfers the provisional mass to the left endpoint. For any $\widetilde\xi\in\left(\xi_{(i)}^{\mathrm{te}},\xi_{(i+1)}^{\mathrm{te}}\right)$ and every observed threshold,
  \[
    \mathbf1_{\{\xi_{(i)}^{\mathrm{te}}<\xi_{(k)}^{\mathrm{te}}\}}
      =\mathbf1_{\{\widetilde\xi<\xi_{(k)}^{\mathrm{te}}\}},
    \qquad
    \mathbf1_{\{\xi_{(i)}^{\mathrm{te}}\le\xi_{(k)}^{\mathrm{te}}\}}
      \ge\mathbf1_{\{\widetilde\xi\le\xi_{(k)}^{\mathrm{te}}\}}.
  \]
  Thus every open upper probability is unchanged and every closed lower probability can only increase. The upper assignment preserves all rank inequalities and attains the coefficient $\overline h_{i,2i+1}=1$.

  For the lower problem, assign $\xi=\left(\xi_{(i+1)}^{\mathrm{te}}-\hat g(\bm x)\right)/2$ if $i<K$, and $\xi=-\hat g(\bm x)+1$ if $i=K$. These values remain in $\Xi_{2i+1}$ and strictly exceed the failure threshold, including at threshold equality. Hence the lower assignment retains the cell matrix and attains $\underline h_{i,2i+1}=0$.

  If $\mathbb P(\mathcal X_i)=0$, the corresponding row is zero; set $\xi=\xi_{(1)}^{\mathrm{te}}$ on that input cell. All assignments are Borel measurable, and their finite mixtures define probability kernels. Combining them with $\mathbb P$ gives joint laws with the full prescribed input marginal, including randomization at input atoms. Both laws are feasible and satisfy
  \[
    p_{\hat f}(\Pi^+)=\sum_{i=0}^{K}\sum_{j=1}^{J}\overline h_{ij}\pi_{ij},
    \qquad
    p_{\hat f}(\Pi^-)=\sum_{i=0}^{K}\sum_{j=1}^{J}\underline h_{ij}\pi_{ij}.
  \]
  The lower law retains every provisional cell mass. The upper law may change the cell matrix at threshold equalities, but it preserves feasibility and the assigned upper objective. These identities establish the reverse inequalities required for \eqref{eq-finite-primal-values}.

  For every $\theta\ge0$, the empirical product law induces the feasible matrix
  \[
    \pi_{i,2k}^{\mathrm{emp}}:=\mathbb P(\mathcal X_i)\frac{c_k-c_{k-1}}{N^{\mathrm{te}}},
    \qquad i=0,\ldots,K,\quad k=1,\ldots,K,
  \]
  with all odd columns zero. Its open and closed rank probabilities are $c_{k-1}/N^{\mathrm{te}}$ and $c_k/N^{\mathrm{te}}$, respectively, so feasibility includes $\theta=0$. The nonnegative entries of a matrix in $\mathcal C_\theta$ sum to one, and the constraints are closed. Thus $\mathcal C_\theta$ is nonempty and compact. The finite objectives lie in $[0,1]$ and attain their optima, which the joint-law constructions realize in the measure problems.

  The finite primal LPs in \eqref{eq-finite-primal-values} have row multipliers $z_i$ and rank multipliers $\lambda_k,\gamma_k$. Their dual constraints are precisely $z_i+\zeta_j\ge\overline h_{ij}$ and $z_i-\zeta_j\le\underline h_{ij}$. Finite LP strong duality therefore gives \eqref{eq-strong-duality} and dual attainment \citep[Section~5.2.3]{B.V04}. By \eqref{eq-dual-envelope-identities}, optimal rank multipliers also yield optimal bounded marginal functions for the original functional duals. Feasibility and finite optimal values suffice; strict rank inequalities are unnecessary, including at $\theta=0$.
\end{proof}

Theorem~\ref{thm-finite-dual-reformulations} reduces the optimization to joint cell masses while retaining the full input law through the within-cell construction.

The following corollary evaluates both bounds in closed form and identifies optimal dual certificates.

\begin{corollary}[Closed-form bounds and optimal certificates]
  \label{thm-closed-form-bounds}
  Under the conditions of Theorem~\ref{thm-finite-dual-reformulations},
  \begin{align}
    \overline p_{\hat f}
      &=\min\left\{1,\min_{k=1,\ldots,K}\left[u_k+\sum_{i=k}^{K}\mathbb P(\mathcal X_i)\right]\right\},
      \label{eq-dw-closed-upper}\\
    \underline p_{\hat f}
      &=\max\left\{0,\max_{k=1,\ldots,K}\left[l_k+\sum_{i=k}^{K}\mathbb P(\mathcal X_i)-1\right]\right\}.
      \label{eq-dw-closed-lower}
  \end{align}
  Each dual optimum admits a certificate with at most one nonzero rank multiplier, of value one.
\end{corollary}
\begin{proof}
  For a candidate residual law $\mathbb Q\in\mathcal M_1(\Xi)$, define the finite transportation set
  \[
    \mathcal T(\mathbb Q):=\left\{\bm\pi\in\mathbb R_+^{(K+1)\times J}:
    \begin{array}{ll}
      \displaystyle\sum_{j=1}^{J}\pi_{ij}=\mathbb P(\mathcal X_i),&i=0,\ldots,K,\\[1mm]
      \displaystyle\sum_{i=0}^{K}\pi_{ij}=\mathbb Q(\Xi_j),&j=1,\ldots,J
    \end{array}\right\}.
  \]
  It contains the product matrix $\pi_{ij}=\mathbb P(\mathcal X_i)\mathbb Q(\Xi_j)$. Every residual cell is nonempty, so any nonnegative column allocation with total mass one is realized by a residual law. Proposition~\ref{prop-finite-rank-constraints} consequently gives
  \[
    \mathcal C_\theta
      =\bigcup_{\mathbb Q\in\mathcal Q^{\mathrm{DW}}(\theta)}\mathcal T(\mathbb Q).
  \]
  The parameter $\mathbb Q$ fixes only residual partition masses in this finite matrix problem. The joint-law construction in Theorem~\ref{thm-finite-dual-reformulations} preserves feasibility and the assigned objective while allowing the residual marginal to change, including endpoint transfers in the upper problem.

  For fixed cell masses $\mathbb Q(\Xi_j)$, the upper problem maximizes mass on pairs $j\le2i+1$; the lower problem maximizes safe mass on pairs $j>2i$ and subtracts it from one. Both are maximum-flow problems with source-to-row capacities $\mathbb P(\mathcal X_i)$, column-to-sink capacities $\mathbb Q(\Xi_j)$, and capacity two on each permitted row-to-column edge. A transportation matrix restricts to a feasible flow. Conversely, any flow can be completed to a transportation matrix because the remaining row and column masses are nonnegative with equal totals. Thus the maximum flow equals the maximum permitted mass.

  A cut of capacity one is available, so a minimum cut crosses no edge of capacity two. Its source-side columns can be taken to be exactly the neighbors of its source-side rows. Since the neighborhoods are nested, adding rows whose neighbors are already included cannot increase the capacity. Apart from empty and full row sets, upper-problem cuts reduce to rows $i<k$ and columns $j<2k$, while safe-mass cuts reduce to rows $i\ge k$ and columns $j>2k$. For $k=1,\ldots,K$, their capacities are respectively
  \[
    \sum_{i=k}^{K}\mathbb P(\mathcal X_i)+\sum_{j<2k}\mathbb Q(\Xi_j),
    \qquad
    \sum_{i=0}^{k-1}\mathbb P(\mathcal X_i)+\sum_{j>2k}\mathbb Q(\Xi_j).
  \]
  Empty and full row sets give capacity one. The max-flow/min-cut theorem \citep{F.F56} therefore yields
  \begin{align*}
    \max_{\bm\pi\in\mathcal T(\mathbb Q)}\sum_{i,j}\overline h_{ij}\pi_{ij}
      &=\min\left\{1,\min_k\left[\sum_{i=k}^{K}\mathbb P(\mathcal X_i)
        +\sum_{j<2k}\mathbb Q(\Xi_j)\right]\right\},\\
    \min_{\bm\pi\in\mathcal T(\mathbb Q)}\sum_{i,j}\underline h_{ij}\pi_{ij}
      &=\max\left\{0,\max_k\left[\sum_{i=k}^{K}\mathbb P(\mathcal X_i)
        +\sum_{j\le2k}\mathbb Q(\Xi_j)-1\right]\right\}.
  \end{align*}

  These values are nondecreasing in the respective open and closed cumulative residual probabilities. The rank constraints bound the former from above by $u_k$ and the latter from below by $l_k$. To attain all relevant limits simultaneously, choose laws $\mathbb Q^+$ and $\mathbb Q^-$ with cell probabilities
  \begin{align*}
    &\mathbb Q^+(\Xi_1):=u_1,\quad
      \mathbb Q^+(\Xi_{2k}):=u_{k+1}-u_k,\quad
      \mathbb Q^+(\Xi_{2k+1}):=0,\\
    &\mathbb Q^-(\Xi_{2k-1}):=0,\quad
      \mathbb Q^-(\Xi_{2k}):=l_k-l_{k-1},\quad
      \mathbb Q^-(\Xi_J):=1-l_K,
      \qquad k=1,\ldots,K.
  \end{align*}
  By Proposition~\ref{prop-finite-rank-constraints}, the prescribed masses are nonnegative and sum to one. Choosing a representative in each residual cell realizes both laws. Their cumulative probabilities satisfy
  \begin{align*}
    &\sum_{j<2k}\mathbb Q^+(\Xi_j)=u_k,
      \qquad \sum_{j\le2k}\mathbb Q^+(\Xi_j)=u_{k+1}\ge l_k,\\
    &\sum_{j\le2k}\mathbb Q^-(\Xi_j)=l_k,
      \qquad \sum_{j<2k}\mathbb Q^-(\Xi_j)=l_{k-1}\le u_k.
  \end{align*}
  Both laws therefore belong to $\mathcal Q^{\mathrm{DW}}(\theta)$. The union representation of $\mathcal C_\theta$ and \eqref{eq-finite-primal-values} show that these allocations attain the upper and lower bounds in \eqref{eq-dw-closed-upper} and \eqref{eq-dw-closed-lower}. Zero masses and coincident cumulative limits are allowed, including at $\theta=0$.

  Finally, let $\bm e_k$ be the $k$th unit vector. Certificates for the upper and lower LPs \eqref{eq-finite-upper-dual} and \eqref{eq-finite-lower-dual} are respectively
  \begin{align*}
    &\bm\lambda=\bm e_k,\qquad\bm\gamma=\bm0,\qquad
      z_i=\mathbf1_{\{i\ge k\}},\\
    &\bm\lambda=\bm0,\qquad\bm\gamma=\bm e_k,\qquad
      z_i=\mathbf1_{\{i\ge k\}}-1.
  \end{align*}
  Their feasibility follows from
  \[
    \mathbf1_{\{i\ge k\}}+\mathbf1_{\{j<2k\}}\ge\overline h_{ij},
    \qquad
    \mathbf1_{\{i\ge k\}}-1+\mathbf1_{\{j\le2k\}}\le\underline h_{ij}.
  \]
  Their objectives are $u_k+\sum_{i=k}^{K}\mathbb P(\mathcal X_i)$ and $l_k+\sum_{i=k}^{K}\mathbb P(\mathcal X_i)-1$, the inner candidates in the closed-form expressions. Zero rank multipliers with $z_i=1$ or $z_i=0$ supply the outer candidates. Selecting the best candidate gives an optimal full dual solution, whose functional marginal multiplier follows from \eqref{eq-dual-envelope-identities}. The inclusion of $i=k$ in the tail sum retains inputs whose failure threshold equals the observed residual.
\end{proof}

Corollary~\ref{thm-closed-form-bounds} extends the classical fixed-marginal sum problem studied by \citet{M82} and \citet{F.N.S87} to the setting where the residual law varies within a calibrated DW band, with open and closed rank constraints accounting for possible atoms. Even when both marginals are known, dependence uncertainty may lead to substantially different failure probabilities. For equally likely surrogate responses $0$ and $2$ and equally likely residuals $-1$ and $1$, the failure probability ranges from $0$ to $1/2$ over all couplings, whereas independence gives $1/4$. Thus, further calibration of the residual marginal alone cannot eliminate dependence uncertainty.

For a response threshold $c$, applying the construction to $g-c$ and $\hat g-c$ preserves the residuals and rank bounds. Corollary~\ref{thm-closed-form-bounds} therefore gives
\begin{align}
  \overline p_{\hat f}(c)
    &=\min\left\{1,\min_{k=1,\ldots,K}\left[u_k
      +\mathbb P\left\{\hat g(\bm X)\le
      c-\xi_{(k)}^{\mathrm{te}}\right\}\right]\right\},
    \label{eq-threshold-closed-upper}\\
  \underline p_{\hat f}(c)
    &=\max\left\{0,\max_{k=1,\ldots,K}\left[l_k
      +\mathbb P\left\{\hat g(\bm X)\le
      c-\xi_{(k)}^{\mathrm{te}}\right\}-1\right]\right\}.
    \label{eq-threshold-closed-lower}
\end{align}
These identities hold for every $c\in\mathbb R$ and $\theta\ge0$. When $\theta\ge\theta_{N^{\mathrm{te}}}(\delta)$, the coverage is simultaneous by \eqref{eq-threshold-curve-coverage}, even though the extremizing couplings may vary with $c$. In particular, for a target probability $p_{\mathrm{tar}}\in(0,1)$, the condition $\overline p_{\hat f}(c)\le p_{\mathrm{tar}}$ holds if and only if
\[
  u_k+\mathbb P\left\{\hat g(\bm X)\le
  c-\xi_{(k)}^{\mathrm{te}}\right\}
  \le p_{\mathrm{tar}}
\]
for at least one $k\in\{1,\ldots,K\}$. 

The following remark incorporates known residual support into the bounds.

\begin{remark}[Known interval support]
  \label{rem-interval-support}
  Suppose $\Xi=[a,b]$ with $-\infty<a\le b<\infty$ and the observed residuals lie in this interval. Then
  \begin{align*}
    \overline p_{\hat f}
      &=\min\left\{\mathbb P\left\{\hat g(\bm X)\le-a\right\},
        \min_{k=1,\ldots,K}\left[u_k+\mathbb P\left\{\hat g(\bm X)\le-\xi_{(k)}^{\mathrm{te}}\right\}\right]\right\},\\
    \underline p_{\hat f}
      &=\max\left\{\mathbb P\left\{\hat g(\bm X)\le-b\right\},
        \max_{k=1,\ldots,K}\left[l_k+\mathbb P\left\{\hat g(\bm X)\le-\xi_{(k)}^{\mathrm{te}}\right\}-1\right]\right\}.
  \end{align*}
  Represent the support restriction on $\mathbb R$ by $\mathbb Q\{\xi<a\}\le0$ and $\mathbb Q\{\xi\le b\}\ge1$, and add $a,b$ to the ordered threshold list. At a newly added $a$, use upper and lower bounds zero; at a newly added $b$, use upper and lower bounds one. At coincident observed thresholds, tighten only the upper bound at $a$ and the lower bound at $b$. The enlarged sequences satisfy the compatibility conditions in Proposition~\ref{prop-finite-rank-constraints}, so the same finite-aggregation and ordered-cut arguments apply. The effective additional upper candidate is $\mathbb P\left\{\hat g(\bm X)\le-a\right\}$, and the effective additional lower candidate is $\mathbb P\left\{\hat g(\bm X)\le-b\right\}$; the other added candidates are redundant. If $a=b$, both bounds equal $\mathbb P\left\{\hat g(\bm X)\le-a\right\}$. For the partition on $\Xi=[a,b]$, intersect the residual cells from the augmented threshold list with $\Xi$ and omit empty cells. Including the support endpoints retains their atoms and the corresponding changes in failure status. For a response threshold $c$, shift each surrogate-response threshold in the displayed event probabilities by $c$.
\end{remark}

Remark~\ref{rem-interval-support} identifies inputs that must fail, $\hat g(\bm X)\le-b$, and those that are always safe, $\hat g(\bm X)>-a$. Thus $\mathbb P\left\{\hat g(\bm X)\le-a\right\}=0$ rules out failure.

% After zero failures in $N^{\mathrm{te}}$ independent tests, the one-sided exact binomial upper limit is $1-\delta^{1/N^{\mathrm{te}}}\approx\log(1/\delta)/N^{\mathrm{te}}$ \citep{C.P34}. At $c=0$, the simple DW case $K=1$ and $\mathbb P(\mathcal X_1)=0$ gives an upper bound $\beta_0(\theta)$ determined by the simultaneous-band calibration.

% Evaluation requires scalar DW inversion, sorting and counting the test residuals, and computing the $K$ input event probabilities $\mathbb P\left\{\hat g(\bm X)\le c-\xi_{(k)}^{\mathrm{te}}\right\}$ at each response threshold $c$. The endpoint table depends only on $N^{\mathrm{te}}$ and $\theta$ and can be precomputed. Given this table and the input event probabilities, the two scans cost $O(K)$ at each response threshold.

% For $\theta\ge\theta_{N^{\mathrm{te}}}(\delta)$, the exact expressions retain the coverage in Theorem~\ref{theorem-01} and Corollary~\ref{cor-simultaneous-threshold-bounds}. Independent auxiliary simulation of $\bm X\sim\mathbb P$ can supply simultaneous upper and lower confidence bounds on the input event probabilities. Substituting these bounds in the respective upper and lower expressions preserves conservative failure probability bounds. For one fixed $c$, simultaneous bounds at its $K$ evaluation thresholds suffice. To retain coverage of the entire response-threshold curve, the auxiliary bounds must hold for the surrogate-response CDF at every real threshold. They must hold conditionally on the test data, with their error probability included in the overall confidence budget.

\section{Implementation Procedure}
\label{sec-implementation}

Sections~\ref{sec-problem}--\ref{sec-reformulation} provide the joint model, residual calibration, and closed-form bounds. These results fit into an existing surrogate modeling workflow in three stages. Stage~I is conventional model fitting and selection; Stages~II and III add residual uncertainty calibration and reliability assessment. For an already fitted and selected surrogate, the procedure starts at Stage~II. The workflow's existing independent test data supply the labeled residual sample.

The following algorithm summarizes the three stages for $\Xi=\mathbb R$.

\begin{algorithm}[H]
  \caption{Residual calibration and reliability assessment within a surrogate workflow}
  \label{alg-implementation}
  \small
  \algrenewcommand{\algorithmicrequire}{\textbf{Input:}}
  \algrenewcommand{\algorithmicensure}{\textbf{Output:}}
  \begin{algorithmic}[1]
    \Require Training data $\mathcal D^{\mathrm{tr}}$ for $g$; an existing independent test set $\mathcal D^{\mathrm{te}}$ with size $N^{\mathrm{te}}$ fixed in advance; input law $\mathbb P$; confidence parameter $\delta\in(0,1)$; requested response thresholds $c$.
    \Ensure Intervals $\left[\underline p_{\hat f}(c),\overline p_{\hat f}(c)\right]$ at the requested thresholds.
    \Statex \textbf{Stage I. Conventional surrogate modeling}
    \State Fit and select $\hat g$ using only $\mathcal D^{\mathrm{tr}}$, then keep it fixed.
    \Statex \textbf{Stage II. Residual uncertainty calibration}
    \State Obtain $\theta\ge\theta_{N^{\mathrm{te}}}(\delta)$ from the uniform-reference quantile in \eqref{eq-dw-calibrated-threshold}.
    \State Compute $\xi_j^{\mathrm{te}}=g(\bm x_j^{\mathrm{te}})-\hat g(\bm x_j^{\mathrm{te}})$ for $j=1,\ldots,N^{\mathrm{te}}$ using the existing test responses.
    \State Sort the distinct residuals $\xi_{(1)}^{\mathrm{te}}<\cdots<\xi_{(K)}^{\mathrm{te}}$ and compute $c_0,\ldots,c_K$ as in \eqref{eq-dw-counts-and-levels}.
    \State Invert \eqref{eq-dw-rank-acceptance} to obtain $l_k=\alpha_{c_k}(\theta)$ and $u_k=\beta_{c_{k-1}}(\theta)$ for $k=1,\ldots,K$.
    \Statex \textbf{Stage III. Reliability assessment}
    \For{each requested response threshold $c$}
      \State Evaluate the surrogate-response CDF through $\mathbb P\left\{\hat g(\bm X)\le c-\xi_{(k)}^{\mathrm{te}}\right\}$ for $k=1,\ldots,K$.
      \State Compute $\overline p_{\hat f}(c)$ and $\underline p_{\hat f}(c)$ from \eqref{eq-threshold-closed-upper} and \eqref{eq-threshold-closed-lower}.
    \EndFor
  \end{algorithmic}
\end{algorithm}

\emph{Surrogate-model agnostic integration.}
Algorithm~\ref{alg-implementation} uses only surrogate predictions and signed test residuals. Gaussian processes, polynomial chaos expansions, neural networks, and other measurable predictors can therefore use the same calibration and evaluation steps. The fitting method and model architecture remain those of the existing workflow.

\emph{Post-training calibration and data reuse.}
The existing test set can serve both routine prediction assessment and residual calibration, provided it is independent of model fitting and selection. Tuning and early stopping use only $\mathcal D^{\mathrm{tr}}$, including any internal validation splits, and $\hat g$ remains fixed after calibration. When such test data are already available, Stage~II requires no additional evaluations of $g$.

\emph{Calibration numerics.}
For fixed $\phi$ and $\psi$, the reference quantile depends only on $N^{\mathrm{te}}$ and $\delta$ and can be tabulated offline. \citet[Supplement, Section~S.7]{L.J23} compute finite-sample critical values through uniform order-statistic boundary probabilities. For a chosen $\theta$, the scalar endpoints follow by bisection using Proposition~\ref{prop-finite-rank-constraints}; only the observed count levels are needed. Round $\theta$ upward and the acceptance intervals outward to preserve conservatism. If simulation estimates the reference quantile, include its estimation error in the confidence budget.

\emph{Optimization-free evaluation.}
Stage~III uses the scalar rank bounds and surrogate event probabilities directly. It requires no explicit joint partition, cell-mass polytope, or primal or dual LP. The finite reformulations and strong duality establish the sharpness of the closed-form values used by the algorithm. Sorting costs $O\left(N^{\mathrm{te}}\log N^{\mathrm{te}}\right)$; once the event probabilities are available, the two scans cost $O(K)$ per response threshold. Known interval support is incorporated through the endpoint candidates in Remark~\ref{rem-interval-support}.

\emph{Reusable calibration and probability evaluation.}
For fixed $g$, $\hat g$, and $\mathbb P$, Stage~II is reused across the entire response-threshold curve. With exact input event probabilities, Corollary~\ref{cor-simultaneous-threshold-bounds} gives confidence at least $1-\delta$ simultaneously for all $c$, including thresholds selected after calibration. These probabilities can be obtained analytically or estimated by evaluating only $\hat g$ on auxiliary draws from $\mathbb P$, independent of both datasets. For simulation, allocate the total error probability between residual calibration and a simultaneous CDF confidence band. Substitute the band's upper probabilities into \eqref{eq-threshold-closed-upper} and its lower probabilities into \eqref{eq-threshold-closed-lower}. Coverage of the entire curve requires the auxiliary band to hold at every real threshold.



% BEGIN OPTIONAL APPENDIX: INFORMATION CONTENT AND STATISTICAL LIMITS (COMMENTED OUT)
% This appendix is intentionally excluded from the compiled manuscript.
% To enable it, uncomment the block below, including the appendix command.
% It preserves the former Section 5, including all statements, proofs, and optional refinements.
%\appendix
%\section{Information Content and Statistical Limits}
%\label{sec-information}
%
%The closed-form expressions identify the information needed to certify a small failure probability. For $\Xi=\mathbb R$ and a prescribed target $p_{\mathrm{tar}}\in(0,1)$, Corollary~\ref{thm-closed-form-bounds} gives
%\begin{equation}
%  \overline p_{\hat f}\le p_{\mathrm{tar}}
%  \quad\Longleftrightarrow\quad
%  \text{there exists }k\in\{1,\ldots,K\}\text{ such that }
%  u_k+H\left(-\xi_{(k)}^{\mathrm{te}}\right)\le p_{\mathrm{tar}}.
%  \label{eq-target-certification}
%\end{equation}
%At that residual level, both the allowed lower-tail residual mass and the probability that the surrogate reaches the corresponding failure threshold must be small. A small mean squared residual alone does not ensure this condition. More generally, the upper bound is below one exactly when some $u_k+H\left(-\xi_{(k)}^{\mathrm{te}}\right)<1$, and the lower bound is positive exactly when some $l_k+H\left(-\xi_{(k)}^{\mathrm{te}}\right)>1$.
%
%Finite evaluation budgets impose a separate statistical limit. A function can agree with a safe surrogate at every evaluated input and fail on a region that the evaluations miss. The following theorem states this obstruction for arbitrary sampling strategies, including surrogate construction and adaptive evaluation locations. Its budget $N^{\mathrm{tot}}$ counts all calls to $g$, including training and calibration, whereas $N^{\mathrm{te}}$ denotes the independent test sample used in the DW construction.
%
%\begin{theorem}[Information limit for unrestricted limit-state functions]
%  \label{thm-information-limit}
%  Let $\delta\in(0,1)$, and suppose $\mathbb P$ is nonatomic. A procedure accesses an unknown measurable $g$ only through at most $N^{\mathrm{tot}}\ge1$ point evaluations and returns an upper bound $\widehat p_{\mathrm U}\in[0,1]$. It may randomize, choose inputs adaptively, and fit surrogates. Assume that, for every measurable $g$,
%  \[
%    \Pr_g\left\{p_f(g)\le\widehat p_{\mathrm U}\right\}\ge1-\delta,
%    \qquad p_f(g):=\mathbb P\{g(\bm X)\le0\},
%  \]
%  where $\Pr_g$ includes all sampling and internal randomization. On the safe function $g_0(\bm x)\equiv1$, every $\delta_0\in(\delta,1)$ satisfies
%  \begin{equation}
%    \Pr_{g_0}\left\{\widehat p_{\mathrm U}\ge
%      1-\left(\frac{\delta}{\delta_0}\right)^{1/N^{\mathrm{tot}}}\right\}
%      \ge1-\delta_0.
%    \label{eq-information-limit}
%  \end{equation}
%  For fixed $\delta$ and $\delta_0$, the threshold is
%  $\log(\delta_0/\delta)/N^{\mathrm{tot}}+O\left(\left(N^{\mathrm{tot}}\right)^{-2}\right)$.
%  If the output on $g_0$ is constant, that constant is at least $1-\delta^{1/N^{\mathrm{tot}}}$.
%\end{theorem}
%\begin{proof}
%  Fix $\varepsilon\in(0,1)$. Nonatomicity permits a partition of $\mathcal X$ into $L$ Borel cells of probability $1/L$. Choose uniformly a union $\mathcal V_L$ of $m_L:=\lceil\varepsilon L\rceil$ cells, and define $g_{\mathcal V_L}$ to be $-1$ on $\mathcal V_L$ and $1$ elsewhere. Every such function has failure probability $m_L/L\ge\varepsilon$.
%
%  Fix the procedure's random seed and follow its evaluation path under $g_0$. This path visits at most $N^{\mathrm{tot}}$ distinct cells. For sufficiently large $L$, the probability over $\mathcal V_L$ of missing every visited cell is at least
%  \[
%    \frac{\binom{L-N^{\mathrm{tot}}}{m_L}}{\binom{L}{m_L}}
%      \longrightarrow(1-\varepsilon)^{N^{\mathrm{tot}}}
%      \qquad\text{as }L\to\infty.
%  \]
%  On this event, all evaluations, subsequent adaptive choices, and the output are identical under $g_0$ and $g_{\mathcal V_L}$. If the baseline output is below $\varepsilon$, it therefore fails to cover $p_f(g_{\mathcal V_L})$. Averaging over the uniformly chosen regions and the procedure's randomness, the assumed coverage for every $g_{\mathcal V_L}$ gives
%  \[
%    \delta\ge
%    \Pr_{g_0}\left\{\widehat p_{\mathrm U}<\varepsilon\right\}
%    \frac{\binom{L-N^{\mathrm{tot}}}{m_L}}{\binom{L}{m_L}}.
%  \]
%  Letting $L\to\infty$ yields
%  $\Pr_{g_0}\left\{\widehat p_{\mathrm U}<\varepsilon\right\}
%    \le\delta/(1-\varepsilon)^{N^{\mathrm{tot}}}$.
%  Substituting $\varepsilon=1-(\delta/\delta_0)^{1/N^{\mathrm{tot}}}$ proves \eqref{eq-information-limit}. The expansion follows from the exponential series. For a constant baseline output below $1-\delta^{1/N^{\mathrm{tot}}}$, choosing $\varepsilon$ between these values makes the last inequality impossible.
%\end{proof}
%
%The theorem gives a probability statement about any uniformly valid procedure. It rules out a certificate that is $o\left(1/N^{\mathrm{tot}}\right)$ in probability even on the safe baseline. A familiar exact benchmark arises from independent testing. If a deterministic binomial upper confidence limit uses only the number of failures among $N^{\mathrm{te}}$ independent inputs, its zero-failure value must satisfy
%\begin{equation}
%  \widehat p_{\mathrm U}(0)\ge1-\delta^{1/N^{\mathrm{te}}}
%    =\frac{\log(1/\delta)}{N^{\mathrm{te}}}
%      +O\left(\left(N^{\mathrm{te}}\right)^{-2}\right).
%  \label{eq-zero-failure-benchmark}
%\end{equation}
%Indeed, if the reported value were smaller, a failure probability between it and $1-\delta^{1/N^{\mathrm{te}}}$ would be missed with probability $(1-p_f)^{N^{\mathrm{te}}}>\delta$. The one-sided exact binomial limit attains equality \citep{C.P34}. The coefficient depends on the confidence level and the calibration procedure. For example, when $K=1$ and $H\left(-\xi_{(1)}^{\mathrm{te}}\right)=0$, the DW upper bound equals $u_1=\beta_0(\theta)$, whose size also reflects the multiscale correction and the calibrated threshold. The information limit alone does not determine this cutoff's coefficient or sample-size dependence.
%
%Additional certified information changes the class of possible hidden failures. With known residual support $[a,b]$, Remark~\ref{rem-interval-support} gives $\overline p_{\hat f}\le H(-a)$. Thus $H(-a)\le p_{\mathrm{tar}}$ can certify the target directly, and $H(-a)=0$ certifies zero failure probability. Other verified restrictions on the residual function can exclude hidden alternatives as well. A known support interval helps when its location relative to the surrogate response rules out failures. On a finite input support, exhaustive evaluation also determines $p_f$ exactly; nonatomicity is essential to the oracle theorem.
%
%Dependence uncertainty has a different source. If the surrogate response takes values $0$ and $2$ with equal probabilities, and the known residual marginal assigns probability $1/2$ to each of $-1$ and $1$, failure occurs exactly when the response is $0$ and the residual is $-1$. Couplings with these marginals allow failure probabilities from $0$ to $1/2$. Independence gives $1/4$, while different deterministic residual functions can attain either endpoint. This gap remains under a known residual marginal because the marginal model leaves the locations of errors unrestricted.
%
%The paired test observations contain additional evidence through the true failure indicators. Define their count by
%\[
%  C_{\mathrm f}^{\mathrm{te}}
%    :=\sum_{j=1}^{N^{\mathrm{te}}}\mathbf1_{\{g(\bm x_j^{\mathrm{te}})\le0\}}.
%\]
%Conditional on the training data, this count is binomial with parameters $N^{\mathrm{te}}$ and $p_f$. An exact binomial interval \citep{C.P34} therefore constrains the joint failure probability directly. The following combination retains a closed-form evaluation.
%
%\begin{proposition}[Combining marginal and paired evidence]
%  \label{prop-paired-evidence}
%  Under the sampling conditions of Proposition~\ref{proposition-01}, choose $\delta_{\mathrm{DW}},\delta_{\mathrm{bin}}>0$ with $\delta_{\mathrm{DW}}+\delta_{\mathrm{bin}}\le\delta<1$, and calibrate $\theta\ge\theta_{N^{\mathrm{te}}}(\delta_{\mathrm{DW}})$. Let $[\underline p^{\mathrm{bin}},\overline p^{\mathrm{bin}}]$ be a binomial interval based on $C_{\mathrm f}^{\mathrm{te}}$ with conditional coverage at least $1-\delta_{\mathrm{bin}}$. When the intervals overlap, define
%  \[
%    \underline p_{\hat f}^{\mathrm{comb}}
%      :=\max\left\{\underline p_{\hat f},\underline p^{\mathrm{bin}}\right\},
%    \qquad
%    \overline p_{\hat f}^{\mathrm{comb}}
%      :=\min\left\{\overline p_{\hat f},\overline p^{\mathrm{bin}}\right\}.
%  \]
%  If they do not overlap, set the combined interval to $[0,1]$. Then
%  \begin{equation}
%    \Pr\left\{\underline p_{\hat f}^{\mathrm{comb}}\le p_f
%      \le\overline p_{\hat f}^{\mathrm{comb}}
%      \,\middle|\,\mathcal D^{\mathrm{tr}}\right\}
%      \ge1-\delta_{\mathrm{DW}}-\delta_{\mathrm{bin}}\ge1-\delta.
%    \label{eq-paired-evidence-coverage}
%  \end{equation}
%  Under the conditions of Corollary~\ref{thm-closed-form-bounds}, when the two component intervals overlap, the combined interval gives the exact lower and upper failure probabilities over
%  \[
%    \left\{\Pi\in\mathcal P\left(\mathcal Q^{\mathrm{DW}}(\theta)\right):
%      \underline p^{\mathrm{bin}}\le p_{\hat f}(\Pi)
%      \le\overline p^{\mathrm{bin}}\right\}.
%  \]
%\end{proposition}
%\begin{proof}
%  Theorem~\ref{theorem-01} gives conditional coverage $1-\delta_{\mathrm{DW}}$ for the DW interval. The binomial interval has conditional coverage $1-\delta_{\mathrm{bin}}$. A union bound gives simultaneous coverage at least $1-\delta_{\mathrm{DW}}-\delta_{\mathrm{bin}}$, even though both intervals use the same test data. On this event their intersection is nonempty and contains $p_f$, proving \eqref{eq-paired-evidence-coverage}.
%
%  For exactness, Proposition~\ref{prop-finite-rank-constraints} makes the joint ambiguity set convex, and Theorem~\ref{thm-finite-dual-reformulations} supplies joint laws attaining both original endpoints. Mixing these laws attains every intermediate failure probability because $p_{\hat f}(\Pi)$ is linear in $\Pi$. Adding the two binomial inequalities therefore intersects the attainable interval with $[\underline p^{\mathrm{bin}},\overline p^{\mathrm{bin}}]$, giving the stated endpoints.
%\end{proof}
%
%The combined interval is contained in each component interval whenever they overlap, with both components evaluated at their allocated error probabilities. The added joint constraints are linear, so the refinement remains within the measure LP framework. In this way, paired evidence can reduce a marginal coupling gap without assuming residual independence or a parametric error model.
%
%More detailed paired information can be retained by fixing input strata before calibration and constructing a DW band for the residual law within each stratum. Applying the results to the conditional input laws and summing the bounds with the stratum probabilities gives simultaneous bounds when the stratum error budgets sum to at most $\delta$. Conditional on the stratum membership labels, the observations within each stratum are independent draws from its conditional input law, so the DW calibration remains valid for every positive stratum sample size. A stratum with no observations contributes its full probability to the upper bound and zero to the lower bound. The partition must be fixed from training data or other information independent of the calibration sample. In particular, the test-dependent rank cells in \eqref{eq-input-cells} require fresh calibration data or a calibration that accounts for their selection if used as statistical strata.
%
%
% END OPTIONAL APPENDIX: INFORMATION CONTENT AND STATISTICAL LIMITS

\bibliographystyle{elsarticle-harv}
\bibliography{reference}

\end{document}

% BEGIN PRESERVED FUNCTIONAL-DUAL LEMMA AND PROOF
%Write $\bm u=(u_1,\ldots,u_K)^\top$ and $\bm l=(l_1,\ldots,l_K)^\top$. For nonnegative multipliers $\bm\lambda,\bm\mu\in\mathbb R_+^K$, define the residual step function
%\begin{equation}
%  \eta_{\bm\lambda,\bm\mu}(\xi)
%  :=\sum_{k=1}^K\lambda_k\mathbf1_{\{\xi<\xi_{(k)}^{\mathrm{te}}\}}
%    -\sum_{k=1}^K\mu_k\mathbf1_{\{\xi\le \xi_{(k)}^{\mathrm{te}}\}}.
%  \label{eq-dual-step-function}
%\end{equation}
%The function $\varphi$ below is the multiplier for the input marginal equality. Together, $\varphi$ and $\eta_{\bm\lambda,\bm\mu}$ bound the failure indicator at every input--residual pair.
%
%\begin{lemma}[Lagrangian dual problems and weak duality]
%  \label{lem-weak-duality}
%  The full upper and lower Lagrangian duals are
%  \begin{equation}
%    \begin{aligned}
%      \overline d_\theta
%      :=\inf_{\bm\lambda,\bm\mu\ge\bm 0,\ \varphi}\quad&
%        \mathbb E_{\mathbb P}[\varphi(\bm X)]
%        +\bm\lambda^\top\bm u-\bm\mu^\top\bm l\\
%      \text{s.t.}\quad&
%        \varphi(\bm x)+\eta_{\bm\lambda,\bm\mu}(\xi)
%        \ge\mathbf1_{\{\xi\le -\hat g(\bm x)\}},
%        \quad\forall(\bm x,\xi)\in\mathcal X\times\Xi,
%    \end{aligned}
%    \label{eq-upper-measure-dual}
%  \end{equation}
%  and
%  \begin{equation}
%    \begin{aligned}
%      \underline d_\theta
%      :=\sup_{\bm\lambda,\bm\mu\ge\bm 0,\ \varphi}\quad&
%        \mathbb E_{\mathbb P}[\varphi(\bm X)]
%        -\bm\lambda^\top\bm u+\bm\mu^\top\bm l\\
%      \text{s.t.}\quad&
%        \varphi(\bm x)-\eta_{\bm\lambda,\bm\mu}(\xi)
%        \le\mathbf1_{\{\xi\le -\hat g(\bm x)\}},
%        \quad\forall(\bm x,\xi)\in\mathcal X\times\Xi,
%    \end{aligned}
%    \label{eq-lower-measure-dual}
%  \end{equation}
%  where $\varphi$ ranges over bounded measurable functions on $\mathcal X$. Their optimal values satisfy
%  \begin{equation}
%    \underline d_\theta
%    \le\underline p_{\hat f}\le\overline p_{\hat f}
%    \le\overline d_\theta.
%    \label{eq-weak-duality}
%  \end{equation}
%\end{lemma}
%\begin{proof}
%  For the upper maximization problem, add the marginal term
%  $\mathbb E_{\mathbb P}[\varphi(\bm X)]-\int\varphi(\bm x)\,\mathrm d\Pi$
%  and the inequality penalties
%  $\lambda_k(u_k-\int\mathbf1_{\{\xi<\xi_{(k)}^{\mathrm{te}}\}}\,\mathrm d\Pi)$ and
%  $\mu_k(\int\mathbf1_{\{\xi\le \xi_{(k)}^{\mathrm{te}}\}}\,\mathrm d\Pi-l_k)$.
%  The resulting Lagrangian is \citep[Section~5.1]{B.V04}
%  \begin{equation}
%    \begin{aligned}
%      \mathcal L^+
%      ={}&\mathbb E_{\mathbb P}[\varphi(\bm X)]
%        +\bm\lambda^\top\bm u-\bm\mu^\top\bm l\\
%        &+\int_{\mathcal X\times\Xi}
%        [\mathbf1_{\{\xi\le -\hat g(\bm x)\}}
%        -\varphi(\bm x)-\eta_{\bm\lambda,\bm\mu}(\xi)]\,\mathrm d\Pi.
%    \end{aligned}
%    \label{eq-upper-lagrangian}
%  \end{equation}
%  Its supremum over finite nonnegative measures is finite exactly when the integrand is nonpositive everywhere. Otherwise, an arbitrarily large point mass at a point with positive integrand makes the supremum infinite. Under the pointwise inequality, the supremum integral is zero, attained by the zero measure. Minimizing the remaining constant term gives \eqref{eq-upper-measure-dual}.
%
%  For the lower minimization problem, reverse both inequality penalties:
%  \begin{equation}
%    \begin{aligned}
%      \mathcal L^-
%      ={}&\mathbb E_{\mathbb P}[\varphi(\bm X)]
%        -\bm\lambda^\top\bm u+\bm\mu^\top\bm l\\
%        &+\int_{\mathcal X\times\Xi}
%        [\mathbf1_{\{\xi\le -\hat g(\bm x)\}}
%        -\varphi(\bm x)+\eta_{\bm\lambda,\bm\mu}(\xi)]\,\mathrm d\Pi.
%    \end{aligned}
%    \label{eq-lower-lagrangian}
%  \end{equation}
%  Its infimum over finite nonnegative measures is finite exactly when the integrand is nonnegative everywhere, giving \eqref{eq-lower-measure-dual}. For primal feasible $\Pi$, the penalty signs imply
%  $\mathcal L^-\le p_{\hat f}(\Pi)\le\mathcal L^+$.
%  The feasible dual objectives therefore bound the primal optima, proving \eqref{eq-weak-duality}.
%\end{proof}
%
% END PRESERVED FUNCTIONAL-DUAL LEMMA AND PROOF

% BEGIN PRESERVED SECTION 4 BEFORE CELL-STRUCTURE REVISION (2026-11-24)
%\section{Tractable Reformulation}
%\label{sec-reformulation}
%
%The DW ambiguity set imposes a nonlinear constraint on CDF values. Its scalar inversion converts this constraint into finitely many linear restrictions on residual probabilities. We first use these restrictions to identify threshold-based failure probability bounds and derive the full Lagrangian duals. We then establish an exact finite-dimensional reformulation and exploit the ordering of the thresholds to obtain closed-form bounds. Throughout this section, the training data, the test residuals, and $\theta\ge0$ are fixed. We use the default residual support $\mathcal E=\mathbb R$ for the main results and give the extension to known interval support at the end.
%
%For each empirical probability level $q_j:=j/N^{\mathrm{te}}$, define the scalar acceptance interval
%\begin{equation}
%  [\alpha_j(\theta),\beta_j(\theta)]
%  :=\left\{p\in[0,1]:
%  N^{\mathrm{te}}\phi(q_j,p)-\psi(q_j,p)\le\theta\right\},
%  \qquad j=0,\ldots,N^{\mathrm{te}}.
%  \label{eq-dw-rank-acceptance}
%\end{equation}
%Let $\xi_{(1)}^{\mathrm{te}}<\cdots<\xi_{(K)}^{\mathrm{te}}$ denote the distinct ordered test residual values, with $k$ indexing distinct values, and set
%\begin{equation} 
%    c_0 :=0, \quad
%    c_k :=\sum_{j=1}^{N^{\mathrm{te}}}\mathbf1_{\{\xi_j^{\mathrm{te}}\le\xi_{(k)}^{\mathrm{te}}\}}, \quad
%    l_k :=\alpha_{c_k}(\theta),  \quad
%    u_k :=\beta_{c_{k-1}}(\theta),\qquad k=1,\ldots,K. 
%  \label{eq-dw-counts-and-levels}
%\end{equation}
%The cumulative counts retain all multiplicities, so $c_K=N^{\mathrm{te}}$ even when test residuals coincide.
%
%\begin{proposition}[CDF band and finite rank constraints]
%  \label{prop-finite-rank-constraints}
%  Each acceptance set in \eqref{eq-dw-rank-acceptance} is a nonempty closed interval containing $q_j$. Define
%  \[
%    L_\theta(t):=\alpha_{N^{\mathrm{te}}\widehat F^{\mathrm{te}}(t)}(\theta),
%    \qquad
%    U_\theta(t):=\beta_{N^{\mathrm{te}}\widehat F^{\mathrm{te}}(t)}(\theta).
%  \]
%  The DW ambiguity set has the equivalent representations
%  \begin{equation}
%    \mathcal Q^{\mathrm{DW}}(\theta) =\left\{\mathbb Q\in\mathcal M_1(\mathcal E): L_\theta(t)\le F_{\mathbb Q}(t)\le U_\theta(t), \, \forall t\in\mathbb R\right\}
%    \label{eq-dw-band-representation}
%  \end{equation}
%  and
%  \begin{equation}
%    \mathcal Q^{\mathrm{DW}}(\theta)
%    =\left\{\mathbb Q\in\mathcal M_1(\mathcal E): \mathbb Q\{\xi<\xi_{(k)}^{\mathrm{te}}\}\le u_k, \, \mathbb Q\{\xi\le \xi_{(k)}^{\mathrm{te}}\}\ge l_k, \, k=1,\ldots,K \right\}.
%    \label{eq-finite-rank-constraints}
%  \end{equation}
%  Both endpoint sequences are nondecreasing in the empirical rank. Consequently, $l_k$ and $u_k$ are nondecreasing and satisfy
%  \begin{equation}
%    0\le l_k,u_k\le1,\quad
%    l_{k-1}\le u_k,\quad
%    l_k\le u_{k+1},
%    \qquad k=1,\ldots,K,
%    \label{eq-rank-compatibility}
%  \end{equation}
%  where $l_0:=0$ and $u_{K+1}:=1$.
%\end{proposition}
%\begin{proof}
%  For fixed $q$, the Bernoulli divergence increases as $p$ moves away from $q$, while $\psi(q,p)$ cannot increase as the interval between $q$ and $p$ expands. Thus the corrected discrepancy decreases toward $q$ from below and increases above $q$. It is continuous for $p\in(0,1)$ and has value $-\psi_0(q)\le0$ at $p=q$. The boundary conventions in Section~\ref{sec-ambiguity-set} give a nonempty closed acceptance interval also when $q\in\{0,1\}$. Applying this scalar inversion at each empirical CDF level gives \eqref{eq-dw-band-representation}, as in the DW confidence band construction \citep[Supplement, Section~S.6]{L.J23}.
%
%  On $[\xi_{(k)}^{\mathrm{te}},\xi_{(k+1)}^{\mathrm{te}})$, the empirical CDF equals $c_k/N^{\mathrm{te}}$. Monotonicity of a candidate CDF makes its lower band constraint equivalent to $F_{\mathbb Q}\left(\xi_{(k)}^{\mathrm{te}}\right)\ge\alpha_{c_k}(\theta)$ and its upper band constraint equivalent to $F_{\mathbb Q}\left(\xi_{(k+1)}^{\mathrm{te}}-\right)\le\beta_{c_k}(\theta)$. The left exterior interval contributes $F_{\mathbb Q}\left(\xi_{(1)}^{\mathrm{te}}-\right)\le\beta_0(\theta)$, and the right exterior interval contributes $F_{\mathbb Q}\left(\xi_{(K)}^{\mathrm{te}}\right)\ge\alpha_{N^{\mathrm{te}}}(\theta)$. The other exterior constraints are automatic because $\alpha_0(\theta)=0$ and $\beta_{N^{\mathrm{te}}}(\theta)=1$. Using $F_{\mathbb Q}\left(\xi_{(k)}^{\mathrm{te}}-\right)=\mathbb Q\{\xi<\xi_{(k)}^{\mathrm{te}}\}$ and $F_{\mathbb Q}\left(\xi_{(k)}^{\mathrm{te}}\right)=\mathbb Q\{\xi\le \xi_{(k)}^{\mathrm{te}}\}$ proves \eqref{eq-finite-rank-constraints}.
%
%  With $p$ fixed, the corrected discrepancy also decreases as $q$ approaches $p$ from below and increases above $p$. Together with the fact that each acceptance interval contains $q$, this implies that both endpoints are nondecreasing in $q$. In particular,
%  $\alpha_{c_k}(\theta)\le c_k/N^{\mathrm{te}}\le\beta_{c_k}(\theta)$, which gives \eqref{eq-rank-compatibility}.
%\end{proof}
%
%The upper rank constraint limits residual mass strictly below $\xi_{(k)}^{\mathrm{te}}$, whereas the lower constraint guarantees mass at or below $\xi_{(k)}^{\mathrm{te}}$. Distinguishing these two half-lines retains possible atoms at the observed residuals. The scalar endpoints can be computed by one-dimensional bisection; the remaining constraints are linear in the residual law.
%
%To interpret the resulting probability bounds, define the CDF of the surrogate response by $H(t):=\mathbb P\left\{\hat g(\bm X)\le t\right\}$ for all $t\in\mathbb R$. For every rank $k$, the failure indicator satisfies 
%\begin{align}
%  \mathbf1_{\{\xi\le -\hat g(\bm x)\}}
%  &\le
%  \mathbf1_{\{\xi<\xi_{(k)}^{\mathrm{te}}\}}+\mathbf1_{\{\hat g(\bm x)\le-\xi_{(k)}^{\mathrm{te}}\}},
%  \label{eq-upper-threshold-certificate}\\
%  \mathbf1_{\{\xi\le -\hat g(\bm x)\}}
%  &\ge
%  \mathbf1_{\{\xi\le \xi_{(k)}^{\mathrm{te}}\}}+\mathbf1_{\{\hat g(\bm x)\le-\xi_{(k)}^{\mathrm{te}}\}}-1.
%  \label{eq-lower-threshold-certificate}
%\end{align}
%The upper relation follows from a union of events: failure requires either a residual below $\xi_{(k)}^{\mathrm{te}}$ or an input threshold at least $\xi_{(k)}^{\mathrm{te}}$. The lower relation follows from an intersection: a residual at or below $\xi_{(k)}^{\mathrm{te}}$ causes failure whenever the input threshold is at least $\xi_{(k)}^{\mathrm{te}}$. Hence every feasible joint law $\Pi$ satisfies
%\begin{equation}
%  l_k+H\left(-\xi_{(k)}^{\mathrm{te}}\right)-1\le p_{\hat f}(\Pi)\le u_k+H\left(-\xi_{(k)}^{\mathrm{te}}\right),
%  \qquad k=1,\ldots,K.
%  \label{eq-rankwise-failure-bounds}
%\end{equation}
%These bounds require no independence assumption. The question is whether choosing the best threshold merely gives a valid bound or recovers the exact optimum over all admissible joint laws. The following duality and finite-aggregation arguments establish exactness.
%
%Substituting \eqref{eq-finite-rank-constraints} into the upper failure probability problem gives the following linear program over $\mathcal M_+(\mathcal X\times\mathcal E)$:
%\begin{equation}
%  \begin{aligned}
%    \overline p_{\hat f}
%    =\sup_{\Pi\in\mathcal M_+(\mathcal X\times\mathcal E)}\quad&
%      \int_{\mathcal X\times\mathcal E}
%      \mathbf1_{\{\xi\le -\hat g(\bm x)\}}\,\Pi(\mathrm d\bm x,\mathrm d\xi)\\
%    \text{s.t.}\quad&
%      \int_{\mathcal X\times\mathcal E}
%      \mathbf1_{\mathcal U}(\bm x)\,\Pi(\mathrm d\bm x,\mathrm d\xi)
%      =\mathbb P(\bm X\in\mathcal U),
%      &&\forall \mathcal U\in\mathcal B(\mathcal X),\\
%    &\int_{\mathcal X\times\mathcal E}
%      \mathbf1_{\{\xi<\xi_{(k)}^{\mathrm{te}}\}}\,\Pi(\mathrm d\bm x,\mathrm d\xi)\le u_k,
%      &&k=1,\ldots,K,\\
%    &\int_{\mathcal X\times\mathcal E}
%      \mathbf1_{\{\xi\le \xi_{(k)}^{\mathrm{te}}\}}\,\Pi(\mathrm d\bm x,\mathrm d\xi)\ge l_k,
%      &&k=1,\ldots,K.
%  \end{aligned}
%  \label{eq-measure-upper}
%\end{equation}
%The lower problem minimizes the same integral under the same constraints:
%\begin{equation}
%  \underline p_{\hat f}
%  =\inf_{\Pi\in\mathcal P(\mathcal Q^{\mathrm{DW}}(\theta))}
%    \int_{\mathcal X\times\mathcal E}
%    \mathbf1_{\{\xi\le -\hat g(\bm x)\}}\,\Pi(\mathrm d\bm x,\mathrm d\xi).
%  \label{eq-measure-lower}
%\end{equation}
%The marginal equality fixes the input law and total mass one. The other constraints restrict the residual marginal while allowing its dependence on the input to vary. For $\mathcal E=\mathbb R$, both problems are feasible for every $\theta\ge0$: the product of $\mathbb P$ and the empirical residual law is feasible.
%
%\begin{lemma}[Lagrangian dual problems and weak duality]
%  \label{lem-weak-duality}
%  Consider the feasible measure programs \eqref{eq-measure-upper} and \eqref{eq-measure-lower}. Write $\bm u=(u_1,\ldots,u_K)^\top$ and $\bm l=(l_1,\ldots,l_K)^\top$. Let $\bm\lambda,\bm\gamma\in\mathbb R_+^K$ denote the Lagrange multipliers associated with the upper and lower rank constraints, respectively. The corresponding Lagrangian dual problems are given by
%  \begin{align}
%    \overline d_\theta
%    &:=\inf_{\bm\lambda,\bm\gamma\in\mathbb R_+^K}
%    \left\{\bm\lambda^\top\bm u-\bm\gamma^\top\bm l
%    +\mathbb E_{\mathbb P}\left[\sup_{\xi\in\mathcal E}\left\{
%    \mathbf1_{\{\xi\le-\hat g(\bm X)\}}
%    -\sum_{k=1}^K\lambda_k\mathbf1_{\{\xi<\xi_{(k)}^{\mathrm{te}}\}}
%    +\sum_{k=1}^K\gamma_k\mathbf1_{\{\xi\le\xi_{(k)}^{\mathrm{te}}\}}
%    \right\}\right]\right\},
%    \label{eq-upper-measure-dual}\\
%    \underline d_\theta
%    &:=\sup_{\bm\lambda,\bm\gamma\in\mathbb R_+^K}
%    \left\{-\bm\lambda^\top\bm u+\bm\gamma^\top\bm l
%    +\mathbb E_{\mathbb P}\left[\inf_{\xi\in\mathcal E}\left\{
%    \mathbf1_{\{\xi\le-\hat g(\bm X)\}}
%    +\sum_{k=1}^K\lambda_k\mathbf1_{\{\xi<\xi_{(k)}^{\mathrm{te}}\}}
%    -\sum_{k=1}^K\gamma_k\mathbf1_{\{\xi\le\xi_{(k)}^{\mathrm{te}}\}}
%    \right\}\right]\right\}.
%    \label{eq-lower-measure-dual}
%  \end{align}
%  Their optimal values satisfy
%  \begin{equation}
%    \underline d_\theta\le\underline p_{\hat f}
%    \le\overline p_{\hat f}\le\overline d_\theta.
%    \label{eq-weak-duality}
%  \end{equation}
%\end{lemma}
%\begin{proof}
%  Let $\eta\in\mathcal F_{\mathrm b}(\mathcal X)$ be the unrestricted Lagrange multiplier associated with the input marginal equality, where $\mathcal F_{\mathrm b}(\mathcal X)$ denotes the space of bounded Borel measurable real-valued functions on $\mathcal X$. Recall that the input marginal equality is equivalent to
%  \begin{equation*}
%    \int_{\mathcal X\times\mathcal E}\eta(\bm x)\,\Pi(\mathrm d\bm x,\mathrm d\xi)
%    =\mathbb E_{\mathbb P}\left[\eta(\bm X)\right],
%    \qquad\forall\eta\in\mathcal F_{\mathrm b}(\mathcal X).
%  \end{equation*}
%  We first consider the upper measure program \eqref{eq-measure-upper}. Introducing the nonnegative multipliers $\bm\lambda,\bm\gamma\in\mathbb R_+^K$ for the upper and lower rank constraints, respectively, we obtain the Lagrangian \citep[Section~5.1]{B.V04}
%  \begin{equation}
%    \begin{aligned}
%      \mathcal L^+(\Pi,\eta,\bm\lambda,\bm\gamma) = & \int_{\mathcal X\times\mathcal E}\mathbf1_{\{\xi\le-\hat g(\bm x)\}}\,\Pi(\mathrm d\bm x,\mathrm d\xi) +\mathbb E_{\mathbb P}\left[\eta(\bm X)\right]-\int_{\mathcal X\times\mathcal E}\eta(\bm x)\,\Pi(\mathrm d\bm x,\mathrm d\xi)\\
%      &+\sum_{k=1}^K\lambda_k\left(u_k-\int_{\mathcal X\times\mathcal E}\mathbf1_{\{\xi<\xi_{(k)}^{\mathrm{te}}\}}\,\Pi(\mathrm d\bm x,\mathrm d\xi)\right) +\sum_{k=1}^K\gamma_k\left(\int_{\mathcal X\times\mathcal E}\mathbf1_{\{\xi\le\xi_{(k)}^{\mathrm{te}}\}}\,\Pi(\mathrm d\bm x,\mathrm d\xi)-l_k\right).
%    \end{aligned}
%    \label{eq-upper-lagrangian-expanded}
%  \end{equation}
%  Collecting the terms involving $\Pi$ gives
%  \begin{equation}
%    \begin{aligned}
%      \mathcal L^+(\Pi,\eta,\bm\lambda,\bm\gamma)
%      ={}&\bm\lambda^\top\bm u-\bm\gamma^\top\bm l+\mathbb E_{\mathbb P}\left[\eta(\bm X)\right]\\
%      &+\int_{\mathcal X\times\mathcal E}\left[ \mathbf1_{\{\xi\le-\hat g(\bm x)\}}-\eta(\bm x) -\sum_{k=1}^K\lambda_k\mathbf1_{\{\xi<\xi_{(k)}^{\mathrm{te}}\}} +\sum_{k=1}^K\gamma_k\mathbf1_{\{\xi\le\xi_{(k)}^{\mathrm{te}}\}} 
%      \right]\,\Pi(\mathrm d\bm x,\mathrm d\xi).
%    \end{aligned}
%    \label{eq-upper-lagrangian}
%  \end{equation}
%  The Lagrangian dual function is defined as
%  \begin{equation}
%    D^+(\eta,\bm\lambda,\bm\gamma)
%    :=\sup_{\Pi\in\mathcal M_+(\mathcal X\times\mathcal E)}\mathcal L^+(\Pi,\eta,\bm\lambda,\bm\gamma),
%    \label{eq-upper-dual-function}
%  \end{equation}
%  Since $\Pi$ is no longer subject to a prescribed total mass in the Lagrangian relaxation, if the integrand in \eqref{eq-upper-lagrangian} is strictly positive at any point $(\bm x,\xi)$, placing mass $M$ at that point and letting $M\to\infty$ yields $D^+=+\infty$. Otherwise, the supremum of the integral term is zero, attained by the zero measure. Consequently, the Lagrangian dual problem is
%  \begin{equation}
%    \begin{aligned}
%      \overline d_\theta
%      =\inf_{\substack{\eta\in\mathcal F_{\mathrm b}(\mathcal X)\\
%      \bm\lambda,\bm\gamma\in\mathbb R_+^K}}
%      \quad &\bm\lambda^\top\bm u-\bm\gamma^\top\bm l+\mathbb E_{\mathbb P}\left[\eta(\bm X)\right]\\
%      \text{s.t.}\quad & \eta(\bm x)\ge
%      \mathbf1_{\{\xi\le-\hat g(\bm x)\}}
%      -\sum_{k=1}^K\lambda_k\mathbf1_{\{\xi<\xi_{(k)}^{\mathrm{te}}\}} 
%      +\sum_{k=1}^K\gamma_k\mathbf1_{\{\xi\le\xi_{(k)}^{\mathrm{te}}\}},
%      \qquad\forall(\bm x,\xi)\in\mathcal X\times\mathcal E.
%    \end{aligned}
%    \label{eq-upper-functional-dual}
%  \end{equation}
%  For fixed $\bm\lambda,\bm\gamma\in\mathbb R_+^K$, the pointwise smallest function satisfying the constraint in \eqref{eq-upper-functional-dual} is
%  \begin{equation}
%    \eta^+_{\bm\lambda,\bm\gamma}(\bm x)
%    :=\sup_{\xi\in\mathcal E}\left\{
%      \mathbf1_{\{\xi\le-\hat g(\bm x)\}}
%      -\sum_{k=1}^K\lambda_k\mathbf1_{\{\xi<\xi_{(k)}^{\mathrm{te}}\}}
%      +\sum_{k=1}^K\gamma_k\mathbf1_{\{\xi\le\xi_{(k)}^{\mathrm{te}}\}}
%    \right\}.
%    \label{eq-upper-marginal-envelope}
%  \end{equation}
%  This function is bounded in absolute value by $1+\|\bm\lambda\|_1+\|\bm\gamma\|_1$. Moreover, replacing $-\hat g(\bm x)$ by a real threshold $t$ in \eqref{eq-upper-marginal-envelope} yields a nondecreasing function of $t$, which is therefore Borel measurable. Since $\hat g$ is Borel measurable, it follows that $\eta^+_{\bm\lambda,\bm\gamma}\in\mathcal F_{\mathrm b}(\mathcal X)$. By monotonicity of the expectation, the optimal choice of $\eta$ in \eqref{eq-upper-functional-dual} for fixed $\bm\lambda,\bm\gamma$ is $\eta=\eta^+_{\bm\lambda,\bm\gamma}$. Substituting this choice into \eqref{eq-upper-functional-dual} yields \eqref{eq-upper-measure-dual}.
%
%  The derivation for the lower measure program \eqref{eq-measure-lower} is analogous, with the signs of the rank penalties reversed and the optimization directions interchanged. In this case, the optimal marginal multiplier for fixed $\bm\lambda,\bm\gamma$ is given by
%  \begin{equation}
%    \eta^-_{\bm\lambda,\bm\gamma}(\bm x)
%    :=\inf_{\xi\in\mathcal E}\left\{
%      \mathbf1_{\{\xi\le-\hat g(\bm x)\}}
%      +\sum_{k=1}^K\lambda_k\mathbf1_{\{\xi<\xi_{(k)}^{\mathrm{te}}\}}
%      -\sum_{k=1}^K\gamma_k\mathbf1_{\{\xi\le\xi_{(k)}^{\mathrm{te}}\}}
%    \right\},
%    \label{eq-lower-marginal-envelope}
%  \end{equation}
%  which also belongs to $\mathcal F_{\mathrm b}(\mathcal X)$ by the same argument. Substituting this choice into the corresponding dual problem yields \eqref{eq-lower-measure-dual}.
%
%  Finally, by weak duality, $\underline d_\theta\le\underline p_{\hat f}$ and $\overline p_{\hat f}\le\overline d_\theta$. Since the two primal programs share the same nonempty feasible set, we have $\underline p_{\hat f}\le\overline p_{\hat f}$. Combining these inequalities establishes \eqref{eq-weak-duality}.
%\end{proof}
%
%The functional multiplier has now been eliminated, but the dual objectives still involve an expectation of a residual supremum or infimum. To evaluate these quantities through the residual ranks, collect the rank penalties in the step function
%\begin{equation}
%  \zeta_{\bm\lambda,\bm\gamma}(\xi)
%  :=\sum_{k=1}^K\lambda_k\mathbf1_{\{\xi<\xi_{(k)}^{\mathrm{te}}\}}
%    -\sum_{k=1}^K\gamma_k\mathbf1_{\{\xi\le\xi_{(k)}^{\mathrm{te}}\}}.
%  \label{eq-dual-step-function}
%\end{equation}
%The expressions inside the supremum in \eqref{eq-upper-marginal-envelope} and the infimum in \eqref{eq-lower-marginal-envelope} are now $\mathbf1_{\{\xi\le-\hat g(\bm x)\}}-\zeta_{\bm\lambda,\bm\gamma}(\xi)$ and $\mathbf1_{\{\xi\le-\hat g(\bm x)\}}+\zeta_{\bm\lambda,\bm\gamma}(\xi)$, respectively. The following finite aggregation makes the ordered rank structure explicit.
%
%To obtain explicit LPs for these dual values, partition the input space according to the location of $-\hat g(\bm x)$ relative to the observed residuals. For boundary notation, set $\xi_{(0)}^{\mathrm{te}}:=-\infty$ and $\xi_{(K+1)}^{\mathrm{te}}:=+\infty$, with $H(-\infty):=0$ and $H(+\infty):=1$, and define
%\begin{equation}
%  \begin{aligned}
%    \mathcal A_i&:=\left\{\bm x:\xi_{(i)}^{\mathrm{te}}\le-\hat g(\bm x)<\xi_{(i+1)}^{\mathrm{te}}\right\},\\
%    w_i&:=\mathbb P(\mathcal A_i)
%    =H\left(-\xi_{(i)}^{\mathrm{te}}\right)-H\left(-\xi_{(i+1)}^{\mathrm{te}}\right),
%    \qquad i=0,\ldots,K.
%  \end{aligned}
%  \label{eq-input-cells}
%\end{equation}
%These groups retain threshold equalities and require no continuity of the surrogate response distribution. Partition the residual line into $J:=2K+1$ categories:
%\begin{equation}
%  \mathcal I_{2k}:=\{\xi_{(k)}^{\mathrm{te}}\}\quad(k=1,\ldots,K),\qquad
%  \mathcal I_{2k+1}:=(\xi_{(k)}^{\mathrm{te}},\xi_{(k+1)}^{\mathrm{te}})\quad(k=0,\ldots,K).
%  \label{eq-residual-cells}
%\end{equation}
%For $\bm x\in \mathcal A_i$ and $\xi\in \mathcal I_j$,
%\begin{equation}
%  \underline h_{ij}:=\mathbf1_{\{j\le2i\}}
%  \le\mathbf1_{\{\xi\le -\hat g(\bm x)\}}
%  \le\overline h_{ij}:=\mathbf1_{\{j\le2i+1\}}.
%  \label{eq-cell-costs}
%\end{equation}
%Categories $j\le2i$ always cause failure, and categories $j\ge2i+2$ are always safe. Only the overlapping interval $\mathcal I_{2i+1}$ distinguishes the two coefficients. At $-\hat g(\bm x)=\xi_{(i)}^{\mathrm{te}}$, this open interval is safe; attaining the upper coefficient then requires moving its assigned mass to $\xi_{(i)}^{\mathrm{te}}$. The proof below verifies that this move preserves feasibility. Thus exactness follows from a feasible probability allocation, rather than from assuming that the failure indicator is constant within every group.
%
%The residual step function is constant on each $\mathcal I_j$, with value
%\begin{equation}
%  \zeta_j:=\sum_{k=1}^K\lambda_k\mathbf1_{\{j<2k\}}
%    -\sum_{k=1}^K\gamma_k\mathbf1_{\{j\le2k\}}.
%  \label{eq-dual-cell-values}
%\end{equation}
%This notation makes the finite duals ordinary linear programs with explicit probability coefficients.
%
%\begin{theorem}[Finite-dimensional duals and strong duality]
%  \label{thm-finite-dual-reformulations}
%  For $\mathcal E=\mathbb R$, $\mathcal Q=\mathcal Q^{\mathrm{DW}}(\theta)$, and every $\theta\ge0$, the full dual values equal the optimal values of
%  \begin{equation}
%    \begin{aligned}
%      \overline d_\theta
%      =\min_{\bm\lambda,\bm\gamma\ge\bm 0,\ \bm z\in\mathbb R^{K+1}}\quad&
%        \bm\lambda^\top\bm u-\bm\gamma^\top\bm l+\sum_{i=0}^K w_i z_i\\
%      \text{s.t.}\quad&
%        z_i+\zeta_j\ge\overline h_{ij},
%        \qquad i=0,\ldots,K,\quad j=1,\ldots,J,
%    \end{aligned}
%    \label{eq-finite-upper-dual}
%  \end{equation}
%  and
%  \begin{equation}
%    \begin{aligned}
%      \underline d_\theta
%      =\max_{\bm\lambda,\bm\gamma\ge\bm 0,\ \bm z\in\mathbb R^{K+1}}\quad&
%        -\bm\lambda^\top\bm u+\bm\gamma^\top\bm l+\sum_{i=0}^K w_i z_i\\
%      \text{s.t.}\quad&
%        z_i-\zeta_j\le\underline h_{ij},
%        \qquad i=0,\ldots,K,\quad j=1,\ldots,J.
%    \end{aligned}
%    \label{eq-finite-lower-dual}
%  \end{equation}
%  Moreover,
%  \begin{equation}
%    \overline p_{\hat f}=\overline d_\theta,\qquad
%    \underline p_{\hat f}=\underline d_\theta.
%    \label{eq-strong-duality}
%  \end{equation}
%  Each finite LP has $3K+1$ scalar variables and $(K+1)(2K+1)$ pairwise inequalities, in addition to multiplier nonnegativity. The finite LP optima and the corresponding measure primal and dual optima are attained.
%\end{theorem}
%\begin{proof}
%  Every feasible joint law induces a probability matrix $\bm\pi=(\pi_{ij})$, with $\pi_{ij}:=\Pi(\mathcal A_i\times \mathcal I_j)$, in
%  \begin{equation}
%    \mathcal C_\theta:=
%    \left\{\bm\pi\in\mathbb R_+^{(K+1)\times J}:
%    \begin{array}{ll}
%      \displaystyle\sum_{j=1}^J\pi_{ij}=w_i,
%        &i=0,\ldots,K,\\[1mm]
%      \displaystyle\sum_{i=0}^K\sum_{j<2k}\pi_{ij}\le u_k,
%        &k=1,\ldots,K,\\[1mm]
%      \displaystyle\sum_{i=0}^K\sum_{j\le2k}\pi_{ij}\ge l_k,
%        &k=1,\ldots,K
%    \end{array}\right\}.
%    \label{eq-cell-mass-polytope}
%  \end{equation}
%  By \eqref{eq-cell-costs}, maximizing $\sum_{i,j}\overline h_{ij}\pi_{ij}$ over this polytope bounds the upper measure problem from above, while minimizing $\sum_{i,j}\underline h_{ij}\pi_{ij}$ bounds the lower problem from below.
%
%  For the reverse inequalities, fix $\bm\pi\in\mathcal C_\theta$. Draw $\bm x$ from $\mathbb P$ and, on $\mathcal A_i$ with $w_i>0$, select a provisional residual category $j$ with probability $\pi_{ij}/w_i$. For $j\ne2i+1$, choose a fixed representative of $\mathcal I_j$; both coefficients in \eqref{eq-cell-costs} then equal the actual failure indicator.
%
%  In the overlapping category, choose $\xi=-\hat g(\bm x)$ for the upper problem. This attains failure and remains in $\mathcal I_{2i+1}$ whenever $-\hat g(\bm x)>\xi_{(i)}^{\mathrm{te}}$. If $-\hat g(\bm x)=\xi_{(i)}^{\mathrm{te}}$, it moves the assigned mass to the left endpoint. For $i\ge1$, any $\widetilde\xi\in(\xi_{(i)}^{\mathrm{te}},\xi_{(i+1)}^{\mathrm{te}})$, and every observed threshold $\xi_{(k)}^{\mathrm{te}}$,
%  \begin{equation}
%    \mathbf1_{\{\xi_{(i)}^{\mathrm{te}}<\xi_{(k)}^{\mathrm{te}}\}}=\mathbf1_{\{\widetilde\xi<\xi_{(k)}^{\mathrm{te}}\}},\qquad
%    \mathbf1_{\{\xi_{(i)}^{\mathrm{te}}\le \xi_{(k)}^{\mathrm{te}}\}}\ge\mathbf1_{\{\widetilde\xi\le \xi_{(k)}^{\mathrm{te}}\}}.
%    \label{eq-endpoint-move}
%  \end{equation}
%  Thus the move leaves all open upper probabilities unchanged and can only increase the closed lower probabilities. It preserves every rank constraint and attains the assigned upper coefficient. For $i=0$, the finite threshold $-\hat g(\bm x)$ cannot equal $\xi_{(0)}^{\mathrm{te}}=-\infty$.
%
%  For the lower problem, choose
%  $\xi=\left(\xi_{(i+1)}^{\mathrm{te}}-\hat g(\bm x)\right)/2$ in the overlapping category when $i<K$, and $\xi=-\hat g(\bm x)+1$ when $i=K$. These choices remain in $\mathcal I_{2i+1}$ and exceed $-\hat g(\bm x)$, attaining safety even at $-\hat g(\bm x)=\xi_{(i)}^{\mathrm{te}}$. All selections are measurable and preserve the input marginal. A group with $w_i=0$ has a zero row and contributes no mass. We have therefore proved
%  \begin{equation}
%    \overline p_{\hat f}
%    =\max_{\bm\pi\in\mathcal C_\theta}\sum_{i=0}^K\sum_{j=1}^J
%      \overline h_{ij}\pi_{ij},\qquad
%    \underline p_{\hat f}
%    =\min_{\bm\pi\in\mathcal C_\theta}\sum_{i=0}^K\sum_{j=1}^J
%      \underline h_{ij}\pi_{ij}.
%    \label{eq-finite-primal-values}
%  \end{equation}
%  The upper construction may redistribute mass between adjacent categories; it preserves feasibility and the objective value needed for this equality.
%
%  The polytope is nonempty by empirical-product feasibility and compact because its nonnegative entries sum to one. Both finite primal optima are attained, and the constructions above lift them to optimal joint measures. Finite LP strong duality applies \citep[Section~5.2.3]{B.V04}. The row equalities have free multipliers $z_i$, and the rank inequalities have nonnegative multipliers $\lambda_k,\gamma_k$. The resulting upper constraints are $z_i+\zeta_j\ge\overline h_{ij}$, and the lower constraints are $z_i-\zeta_j\le\underline h_{ij}$, giving \eqref{eq-finite-upper-dual} and \eqref{eq-finite-lower-dual}.
%
%  Finally, an upper finite dual solution defines a bounded Borel function $\eta(\bm x)=z_i$ on $\mathcal A_i$ satisfying the functional upper-dual constraint, by \eqref{eq-cell-costs}. Thus $\eta^+_{\bm\lambda,\bm\gamma}(\bm x)\le z_i$ on $\mathcal A_i$. For a lower finite dual solution, the same argument gives $\eta^-_{\bm\lambda,\bm\gamma}(\bm x)\ge z_i$. At fixed rank multipliers, the reduced upper objective in \eqref{eq-upper-measure-dual} is therefore no greater than the finite upper-dual objective, and the reduced lower objective is no smaller than the finite lower-dual objective. Finite LP equality and Lemma~\ref{lem-weak-duality} force \eqref{eq-strong-duality}. Optimal finite dual solutions consequently provide optimal rank multipliers for the reduced duals and optimal marginal functions for the functional duals. The argument uses feasibility and boundedness of finite LPs and therefore includes $\theta=0$ and arbitrary atoms.
%\end{proof}
%
%The theorem establishes an exact finite reformulation. The nested residual categories yield a further simplification: the best threshold bounds in \eqref{eq-rankwise-failure-bounds} are already optimal.
%
%\begin{corollary}[Closed-form failure probability bounds]
%  \label{thm-closed-form-bounds}
%  Under the conditions of Theorem~\ref{thm-finite-dual-reformulations},
%  \begin{align}
%    \overline p_{\hat f}
%    &=\min\left\{1,\min_{k=1,\ldots,K}[u_k+H\left(-\xi_{(k)}^{\mathrm{te}}\right)]\right\},
%    \label{eq-dw-closed-upper}\\
%    \underline p_{\hat f}
%    &=\max\left\{0,\max_{k=1,\ldots,K}[l_k+H\left(-\xi_{(k)}^{\mathrm{te}}\right)-1]\right\}.
%    \label{eq-dw-closed-lower}
%  \end{align}
%\end{corollary}
%\begin{proof}
%  Fix a nonnegative column-sum vector $\bm y=(y_1,\ldots,y_J)^\top$ with $\sum_{j=1}^J y_j=1$, and let $Y_j:=\sum_{r=1}^j y_r$. Write $\mathcal T(\bm w,\bm y)$ for the transportation matrices with row-sum vector $\bm w=(w_0,\ldots,w_K)^\top$ and column-sum vector $\bm y$. For the upper problem, maximizing mass on pairs $j\le2i+1$ is a maximum-flow problem: connect the source to row $i$ with capacity $w_i$, each row to its permitted columns with capacity two, and column $j$ to the sink with capacity $y_j$. A partial flow can be completed to a transportation matrix using the remaining row and column masses, so the maximum flow equals the upper transport optimum.
%
%  There is a cut of capacity one. Hence a minimum cut crosses no middle edge of capacity two. The row neighborhoods are nested prefixes. If the largest source-side row is $m<K$, adding all smaller rows and retaining exactly the neighboring columns cannot increase the cut capacity. Writing $k=m+1$, the resulting cut has capacity
%  \[
%    \sum_{i=k}^K w_i+\sum_{j=1}^{2k-1}y_j
%    =H\left(-\xi_{(k)}^{\mathrm{te}}\right)+Y_{2k-1}.
%  \]
%  The empty row set and the case $m=K$ give capacity one. By the max-flow/min-cut theorem \citep{F.F56},
%  \begin{equation}
%    \max_{\bm\pi\in\mathcal T(\bm w,\bm y)}
%      \sum_{i,j}\overline h_{ij}\pi_{ij}
%    =\min\left\{1,\min_k[H\left(-\xi_{(k)}^{\mathrm{te}}\right)+Y_{2k-1}]\right\}.
%    \label{eq-fixed-column-upper}
%  \end{equation}
%
%  For the lower problem, maximize safe mass on pairs $j>2i$ and subtract from one. The neighborhoods are nested suffixes. If the smallest source-side row is $k\ge1$, the corresponding normalized cut has capacity
%  \[
%    \sum_{i=0}^{k-1}w_i+\sum_{j=2k+1}^Jy_j
%    =2-H\left(-\xi_{(k)}^{\mathrm{te}}\right)-Y_{2k}.
%  \]
%  The empty row set and the case $k=0$ give capacity one. The same flow argument yields
%  \begin{equation}
%    \min_{\bm\pi\in\mathcal T(\bm w,\bm y)}
%      \sum_{i,j}\underline h_{ij}\pi_{ij}
%    =\max\left\{0,\max_k[H\left(-\xi_{(k)}^{\mathrm{te}}\right)+Y_{2k}-1]\right\}.
%    \label{eq-fixed-column-lower}
%  \end{equation}
%
%  Feasible column totals satisfy $Y_{2k-1}\le u_k$ and $Y_{2k}\ge l_k$. To attain all upper prefix bounds simultaneously, choose
%  \begin{equation}
%    y_1^+:=u_1,\qquad
%    y_{2k}^+:=u_{k+1}-u_k,\qquad
%    y_{2k+1}^+:=0,\quad k=1,\ldots,K.
%    \label{eq-upper-column-allocation}
%  \end{equation}
%  These nonnegative masses sum to one and have
%  $Y_{2k-1}^+=u_k$ and $Y_{2k}^+=u_{k+1}\ge l_k$.
%  To attain all lower prefix bounds simultaneously, choose
%  \begin{equation}
%    y_{2k-1}^-:=0,\qquad
%    y_{2k}^-:=l_k-l_{k-1}\quad(k=1,\ldots,K),
%    \qquad y_J^-:=1-l_K.
%    \label{eq-lower-column-allocation}
%  \end{equation}
%  These masses have $Y_{2k}^-=l_k$ and
%  $Y_{2k-1}^-=l_{k-1}\le u_k$.
%  Monotonicity and \eqref{eq-rank-compatibility} make both choices feasible. Substituting them into \eqref{eq-fixed-column-upper} and \eqref{eq-fixed-column-lower}, and using \eqref{eq-finite-primal-values}, proves the two closed forms.
%\end{proof}
%
%The closed forms also identify optimal rank multipliers and their corresponding marginal functions. For a rank $k$, the upper choice
%\[
%  \bm\lambda=\bm e_k,\qquad \bm\gamma=\bm 0,\qquad
%  \eta^+_{\bm e_k,\bm 0}(\bm x)=\mathbf1_{\{\hat g(\bm x)\le-\xi_{(k)}^{\mathrm{te}}\}}
%\]
%is feasible by \eqref{eq-upper-threshold-certificate} and has objective $u_k+H\left(-\xi_{(k)}^{\mathrm{te}}\right)$. The lower choice
%\[
%  \bm\lambda=\bm 0,\qquad \bm\gamma=\bm e_k,\qquad
%  \eta^-_{\bm 0,\bm e_k}(\bm x)=\mathbf1_{\{\hat g(\bm x)\le-\xi_{(k)}^{\mathrm{te}}\}}-1
%\]
%is feasible by \eqref{eq-lower-threshold-certificate} and has objective $l_k+H\left(-\xi_{(k)}^{\mathrm{te}}\right)-1$. Here $\bm e_k$ is the $k$th unit vector. With zero rank multipliers, the upper and lower envelopes are the constant functions one and zero, giving the outer candidates.
%
%Selecting a minimizing upper candidate or a maximizing lower candidate therefore gives an optimal full dual solution. In particular, an optimum needs at most one upper rank multiplier for the upper bound and at most one lower rank multiplier for the lower bound. Both families of constraints still define the feasible residual laws; their ordering allows these sparse optimal certificates. The upper bound combines residual mass below a threshold with the probability that the surrogate threshold is at least that value. The lower bound quantifies overlap that no admissible dependence structure can avoid.
%
%\begin{remark}[Known interval support]
%  \label{rem-interval-support}
%  Suppose $\mathcal E=[a,b]$ with $-\infty<a\le b<\infty$ and the observed residuals lie in this interval. Then
%  \begin{align}
%    \overline p_{\hat f}
%    &=\min\left\{H(-a),\min_{k=1,\ldots,K}[u_k+H\left(-\xi_{(k)}^{\mathrm{te}}\right)]\right\},
%    \label{eq-dw-interval-upper}\\
%    \underline p_{\hat f}
%    &=\max\left\{H(-b),\max_{k=1,\ldots,K}[l_k+H\left(-\xi_{(k)}^{\mathrm{te}}\right)-1]\right\}.
%    \label{eq-dw-interval-lower}
%  \end{align}
%  To obtain these expressions, represent the support restriction on $\mathbb R$ by
%  $\mathbb Q\{\xi<a\}\le0$ and $\mathbb Q\{\xi\le b\}\ge1$.
%  Add $a,b$ to the ordered threshold list. At a newly added $a$, use upper and lower bounds zero; at a newly added $b$, use upper and lower bounds one. At coincident observed thresholds, tighten only the upper bound at $a$ and the lower bound at $b$. The enlarged sequences satisfy \eqref{eq-rank-compatibility}, so the same finite-aggregation and ordered-cut arguments apply. The effective additional upper candidate is $H(-a)$, and the effective additional lower candidate is $H(-b)$; the other added candidates are redundant. If $a=b$, both bounds equal $H(-a)$.
%\end{remark}
%
%For a fixed $\theta$, evaluation consists of scalar DW inversion, sorting and counting the test residuals, evaluating $H\left(-\xi_{(k)}^{\mathrm{te}}\right)$ under $\mathbb P$, and scanning the $K$ upper and lower candidates. The endpoint table depends only on $N^{\mathrm{te}}$ and $\theta$ and can be precomputed. Once this table and the surrogate CDF values are available, the final scan costs $O(K)$; solving the finite LPs is unnecessary.
%
%For $\theta\ge\theta_{N^{\mathrm{te}}}(\delta)$, these exact bounds inherit the finite-sample guarantee of Theorem~\ref{theorem-01}. If the surrogate CDF values are estimated by simulation, simultaneous upper CDF bounds can replace $H\left(-\xi_{(k)}^{\mathrm{te}}\right)$ in the upper expression, and simultaneous lower CDF bounds can replace $H\left(-\xi_{(k)}^{\mathrm{te}}\right)$ in the lower expression. An independent auxiliary simulation can provide these simultaneous bounds conditionally on the test data. Its error probability must be included in the overall confidence budget.
%
%
%
%
% END PRESERVED SECTION 4 BEFORE CELL-STRUCTURE REVISION

% BEGIN PRESERVED SECTION 4 BEFORE NUMBERING AND PROSE REVISION (2026-11-24)
%\section{Tractable Reformulation}
%\label{sec-reformulation}
%
%The DW ambiguity set imposes a nonlinear constraint on CDF values. Its scalar inversion converts this constraint into finitely many linear restrictions on residual probabilities. We first use these restrictions to identify threshold-based failure probability bounds and derive the full Lagrangian duals. We then establish an exact finite-dimensional reformulation and exploit the ordering of the thresholds to obtain closed-form bounds. Throughout this section, the training data, the test residuals, and $\theta\ge0$ are fixed. We use the default residual support $\Xi=\mathbb R$ for the main results and give the extension to known interval support at the end.
%
%For each empirical probability level $q_j:=j/N^{\mathrm{te}}$, define the scalar acceptance interval
%\begin{equation}
%  [\alpha_j(\theta),\beta_j(\theta)]
%  :=\left\{p\in[0,1]:
%  N^{\mathrm{te}}\phi(q_j,p)-\psi(q_j,p)\le\theta\right\},
%  \qquad j=0,\ldots,N^{\mathrm{te}}.
%  \label{eq-dw-rank-acceptance}
%\end{equation}
%Let $\xi_{(1)}^{\mathrm{te}}<\cdots<\xi_{(K)}^{\mathrm{te}}$ denote the distinct ordered test residual values, with $k$ indexing distinct values, and set
%\begin{equation} 
%    c_0 :=0, \quad
%    c_k :=\sum_{j=1}^{N^{\mathrm{te}}}\mathbf1_{\{\xi_j^{\mathrm{te}}\le\xi_{(k)}^{\mathrm{te}}\}}, \quad
%    l_k :=\alpha_{c_k}(\theta),  \quad
%    u_k :=\beta_{c_{k-1}}(\theta),\qquad k=1,\ldots,K. 
%  \label{eq-dw-counts-and-levels}
%\end{equation}
%The cumulative counts retain all multiplicities, so $c_K=N^{\mathrm{te}}$ even when test residuals coincide.
%
%\begin{proposition}[CDF band and finite rank constraints]
%  \label{prop-finite-rank-constraints}
%  Each acceptance set in \eqref{eq-dw-rank-acceptance} is a nonempty closed interval containing $q_j$. Define
%  \[
%    L_\theta(t):=\alpha_{N^{\mathrm{te}}\widehat F^{\mathrm{te}}(t)}(\theta),
%    \qquad
%    U_\theta(t):=\beta_{N^{\mathrm{te}}\widehat F^{\mathrm{te}}(t)}(\theta).
%  \]
%  The DW ambiguity set has the equivalent representations
%  \begin{equation}
%    \mathcal Q^{\mathrm{DW}}(\theta) =\left\{\mathbb Q\in\mathcal M_1(\Xi): L_\theta(t)\le F_{\mathbb Q}(t)\le U_\theta(t), \, \forall t\in\mathbb R\right\}
%    \label{eq-dw-band-representation}
%  \end{equation}
%  and
%  \begin{equation}
%    \mathcal Q^{\mathrm{DW}}(\theta)
%    =\left\{\mathbb Q\in\mathcal M_1(\Xi): \mathbb Q\{\xi<\xi_{(k)}^{\mathrm{te}}\}\le u_k, \, \mathbb Q\{\xi\le \xi_{(k)}^{\mathrm{te}}\}\ge l_k, \, k=1,\ldots,K \right\}.
%    \label{eq-finite-rank-constraints}
%  \end{equation}
%  Both endpoint sequences are nondecreasing in the empirical rank. Consequently, $l_k$ and $u_k$ are nondecreasing and satisfy
%  \begin{equation}
%    0\le l_k,u_k\le1,\quad
%    l_{k-1}\le u_k,\quad
%    l_k\le u_{k+1},
%    \qquad k=1,\ldots,K,
%    \label{eq-rank-compatibility}
%  \end{equation}
%  where $l_0:=0$ and $u_{K+1}:=1$.
%\end{proposition}
%\begin{proof}
%  For fixed $q$, the Bernoulli divergence increases as $p$ moves away from $q$, while $\psi(q,p)$ cannot increase as the interval between $q$ and $p$ expands. Thus the corrected discrepancy decreases toward $q$ from below and increases above $q$. It is continuous for $p\in(0,1)$ and has value $-\psi_0(q)\le0$ at $p=q$. The boundary conventions in Section~\ref{sec-ambiguity-set} give a nonempty closed acceptance interval also when $q\in\{0,1\}$. Applying this scalar inversion at each empirical CDF level gives \eqref{eq-dw-band-representation}, as in the DW confidence band construction \citep[Supplement, Section~S.6]{L.J23}.
%
%  On $[\xi_{(k)}^{\mathrm{te}},\xi_{(k+1)}^{\mathrm{te}})$, the empirical CDF equals $c_k/N^{\mathrm{te}}$. Monotonicity of a candidate CDF makes its lower band constraint equivalent to $F_{\mathbb Q}\left(\xi_{(k)}^{\mathrm{te}}\right)\ge\alpha_{c_k}(\theta)$ and its upper band constraint equivalent to $F_{\mathbb Q}\left(\xi_{(k+1)}^{\mathrm{te}}-\right)\le\beta_{c_k}(\theta)$. The left exterior interval contributes $F_{\mathbb Q}\left(\xi_{(1)}^{\mathrm{te}}-\right)\le\beta_0(\theta)$, and the right exterior interval contributes $F_{\mathbb Q}\left(\xi_{(K)}^{\mathrm{te}}\right)\ge\alpha_{N^{\mathrm{te}}}(\theta)$. The other exterior constraints are automatic because $\alpha_0(\theta)=0$ and $\beta_{N^{\mathrm{te}}}(\theta)=1$. Using $F_{\mathbb Q}\left(\xi_{(k)}^{\mathrm{te}}-\right)=\mathbb Q\{\xi<\xi_{(k)}^{\mathrm{te}}\}$ and $F_{\mathbb Q}\left(\xi_{(k)}^{\mathrm{te}}\right)=\mathbb Q\{\xi\le \xi_{(k)}^{\mathrm{te}}\}$ proves \eqref{eq-finite-rank-constraints}.
%
%  With $p$ fixed, the corrected discrepancy also decreases as $q$ approaches $p$ from below and increases above $p$. Together with the fact that each acceptance interval contains $q$, this implies that both endpoints are nondecreasing in $q$. In particular,
%  $\alpha_{c_k}(\theta)\le c_k/N^{\mathrm{te}}\le\beta_{c_k}(\theta)$, which gives \eqref{eq-rank-compatibility}.
%\end{proof}
%
%The upper rank constraint limits residual mass strictly below $\xi_{(k)}^{\mathrm{te}}$, whereas the lower constraint guarantees mass at or below $\xi_{(k)}^{\mathrm{te}}$. Distinguishing these two half-lines retains possible atoms at the observed residuals. The scalar endpoints can be computed by one-dimensional bisection; the remaining constraints are linear in the residual law.
%
%To interpret the resulting probability bounds, define the CDF of the surrogate response by $H(t):=\mathbb P\left\{\hat g(\bm X)\le t\right\}$ for all $t\in\mathbb R$. The ordered residual values classify inputs according to the location of the failure threshold $-\hat g(\bm x)$. Set $\xi_{(0)}^{\mathrm{te}}:=-\infty$ and $\xi_{(K+1)}^{\mathrm{te}}:=+\infty$, with $H(-\infty):=0$ and $H(+\infty):=1$, and define the input cells by 
%\begin{equation}
%  \mathcal X_i :=\left\{\bm x\in\mathcal X: \xi_{(i)}^{\mathrm{te}}\le-\hat g(\bm x)<\xi_{(i+1)}^{\mathrm{te}}\right\}, \qquad i=0,\ldots,K,
%\end{equation}
%with the corresponding probability mass
%\begin{equation}
%  w_i :=\mathbb P(\mathcal X_i) =H\left(-\xi_{(i)}^{\mathrm{te}}\right)-H\left(-\xi_{(i+1)}^{\mathrm{te}}\right)K. 
%\end{equation} 
%These $K+1$ cells form a disjoint partition of $\mathcal X$. Each nested threshold event is a union of consecutive input cells, so
%\begin{equation}
%  \mathbf1_{\{\hat g(\bm x)\le-\xi_{(k)}^{\mathrm{te}}\}}
%    =\sum_{i=k}^K\mathbf1_{\mathcal X_i}(\bm x),
%  \qquad
%  H\left(-\xi_{(k)}^{\mathrm{te}}\right)=\sum_{i=k}^K w_i,
%  \qquad k=1,\ldots,K.
%  \label{eq-input-threshold-masses}
%\end{equation}
%The half-open convention assigns an input with $-\hat g(\bm x)=\xi_{(k)}^{\mathrm{te}}$ to $\mathcal X_k$. Thus these identities retain threshold equalities even when $H$ has atoms.
%
%For every rank $k$, the failure indicator satisfies 
%\begin{align}
%  \mathbf1_{\{\xi\le -\hat g(\bm x)\}}
%  &\le
%  \mathbf1_{\{\xi<\xi_{(k)}^{\mathrm{te}}\}}+\mathbf1_{\{\hat g(\bm x)\le-\xi_{(k)}^{\mathrm{te}}\}},
%  \label{eq-upper-threshold-certificate}\\
%  \mathbf1_{\{\xi\le -\hat g(\bm x)\}}
%  &\ge
%  \mathbf1_{\{\xi\le \xi_{(k)}^{\mathrm{te}}\}}+\mathbf1_{\{\hat g(\bm x)\le-\xi_{(k)}^{\mathrm{te}}\}}-1.
%  \label{eq-lower-threshold-certificate}
%\end{align}
%The upper relation follows from a union of events: failure requires either a residual below $\xi_{(k)}^{\mathrm{te}}$ or an input threshold at least $\xi_{(k)}^{\mathrm{te}}$. The lower relation follows from an intersection: a residual at or below $\xi_{(k)}^{\mathrm{te}}$ causes failure whenever the input threshold is at least $\xi_{(k)}^{\mathrm{te}}$. Hence every feasible joint law $\Pi$ satisfies
%\begin{equation}
%  l_k+H\left(-\xi_{(k)}^{\mathrm{te}}\right)-1\le p_{\hat f}(\Pi)\le u_k+H\left(-\xi_{(k)}^{\mathrm{te}}\right),
%  \qquad k=1,\ldots,K.
%  \label{eq-rankwise-failure-bounds}
%\end{equation}
%These bounds require no independence assumption. The question is whether choosing the best threshold merely gives a valid bound or recovers the exact optimum over all admissible joint laws. The following duality and finite-aggregation arguments establish exactness.
%
%Substituting \eqref{eq-finite-rank-constraints} into the upper failure probability problem gives the following linear program over $\mathcal M_+(\mathcal X\times\Xi)$:
%\begin{equation}
%  \begin{aligned}
%    \overline p_{\hat f}
%    =\sup_{\Pi\in\mathcal M_+(\mathcal X\times\Xi)}\quad&
%      \int_{\mathcal X\times\Xi}
%      \mathbf1_{\{\xi\le -\hat g(\bm x)\}}\,\Pi(\mathrm d\bm x,\mathrm d\xi)\\
%    \text{s.t.}\quad&
%      \int_{\mathcal X\times\Xi}
%      \mathbf1_{\mathcal U}(\bm x)\,\Pi(\mathrm d\bm x,\mathrm d\xi)
%      =\mathbb P(\bm X\in\mathcal U),
%      &&\forall \mathcal U\in\mathcal B(\mathcal X),\\
%    &\int_{\mathcal X\times\Xi}
%      \mathbf1_{\{\xi<\xi_{(k)}^{\mathrm{te}}\}}\,\Pi(\mathrm d\bm x,\mathrm d\xi)\le u_k,
%      &&k=1,\ldots,K,\\
%    &\int_{\mathcal X\times\Xi}
%      \mathbf1_{\{\xi\le \xi_{(k)}^{\mathrm{te}}\}}\,\Pi(\mathrm d\bm x,\mathrm d\xi)\ge l_k,
%      &&k=1,\ldots,K.
%  \end{aligned}
%  \label{eq-measure-upper}
%\end{equation}
%The lower problem minimizes the same integral under the same constraints:
%\begin{equation}
%  \underline p_{\hat f}
%  =\inf_{\Pi\in\mathcal P(\mathcal Q^{\mathrm{DW}}(\theta))}
%    \int_{\mathcal X\times\Xi}
%    \mathbf1_{\{\xi\le -\hat g(\bm x)\}}\,\Pi(\mathrm d\bm x,\mathrm d\xi).
%  \label{eq-measure-lower}
%\end{equation}
%The marginal equality fixes the input law and total mass one. The other constraints restrict the residual marginal while allowing its dependence on the input to vary. For $\Xi=\mathbb R$, both problems are feasible for every $\theta\ge0$: the product of $\mathbb P$ and the empirical residual law is feasible.
%
%For $\bm\lambda,\bm\gamma\in\mathbb R_+^K$, collect the rank penalties in the residual step function
%\begin{equation}
%  \zeta_{\bm\lambda,\bm\gamma}(\xi)
%  :=\sum_{k=1}^K\lambda_k\mathbf1_{\{\xi<\xi_{(k)}^{\mathrm{te}}\}}
%    -\sum_{k=1}^K\gamma_k\mathbf1_{\{\xi\le\xi_{(k)}^{\mathrm{te}}\}}.
%  \label{eq-dual-step-function}
%\end{equation}
%
%\begin{lemma}[Lagrangian dual problems and weak duality]
%  \label{lem-weak-duality}
%  Consider the feasible measure programs \eqref{eq-measure-upper} and \eqref{eq-measure-lower}. Write $\bm u=(u_1,\ldots,u_K)^\top$ and $\bm l=(l_1,\ldots,l_K)^\top$. Let $\bm\lambda,\bm\gamma\in\mathbb R_+^K$ denote the Lagrange multipliers associated with the upper and lower rank constraints, respectively. The unrestricted multiplier for the input marginal equality is $\eta\in\mathcal F_{\mathrm b}(\mathcal X)$, where $\mathcal F_{\mathrm b}(\mathcal X)$ is the space of bounded Borel measurable real-valued functions. Eliminating $\eta$ pointwise gives the Lagrangian dual problems
%  \begin{align}
%    \overline d_\theta
%    &:=\inf_{\bm\lambda,\bm\gamma\in\mathbb R_+^K}
%    \left\{\bm\lambda^\top\bm u-\bm\gamma^\top\bm l
%    +\mathbb E_{\mathbb P}\left[\sup_{\xi\in\Xi}\left\{
%    \mathbf1_{\{\xi\le-\hat g(\bm X)\}}
%    -\zeta_{\bm\lambda,\bm\gamma}(\xi)
%    \right\}\right]\right\},
%    \label{eq-upper-measure-dual}\\
%    \underline d_\theta
%    &:=\sup_{\bm\lambda,\bm\gamma\in\mathbb R_+^K}
%    \left\{-\bm\lambda^\top\bm u+\bm\gamma^\top\bm l
%    +\mathbb E_{\mathbb P}\left[\inf_{\xi\in\Xi}\left\{
%    \mathbf1_{\{\xi\le-\hat g(\bm X)\}}
%    +\zeta_{\bm\lambda,\bm\gamma}(\xi)
%    \right\}\right]\right\}.
%    \label{eq-lower-measure-dual}
%  \end{align}
%  Their optimal values satisfy
%  \begin{equation}
%    \underline d_\theta\le\underline p_{\hat f}
%    \le\overline p_{\hat f}\le\overline d_\theta.
%    \label{eq-weak-duality}
%  \end{equation}
%\end{lemma}
%\begin{proof}
%  The input marginal equality is equivalent to
%  \begin{equation*}
%    \int_{\mathcal X\times\Xi}\eta(\bm x)\,\Pi(\mathrm d\bm x,\mathrm d\xi)
%    =\mathbb E_{\mathbb P}\left[\eta(\bm X)\right],
%    \qquad\forall\eta\in\mathcal F_{\mathrm b}(\mathcal X).
%  \end{equation*}
%  We first consider the upper measure program \eqref{eq-measure-upper}. Introducing the nonnegative multipliers $\bm\lambda,\bm\gamma\in\mathbb R_+^K$ for the upper and lower rank constraints, respectively, we obtain the Lagrangian \citep[Section~5.1]{B.V04}
%  \begin{equation}
%    \begin{aligned}
%      \mathcal L^+(\Pi,\eta,\bm\lambda,\bm\gamma) = & \int_{\mathcal X\times\Xi}\mathbf1_{\{\xi\le-\hat g(\bm x)\}}\,\Pi(\mathrm d\bm x,\mathrm d\xi) +\mathbb E_{\mathbb P}\left[\eta(\bm X)\right]-\int_{\mathcal X\times\Xi}\eta(\bm x)\,\Pi(\mathrm d\bm x,\mathrm d\xi)\\
%      &+\sum_{k=1}^K\lambda_k\left(u_k-\int_{\mathcal X\times\Xi}\mathbf1_{\{\xi<\xi_{(k)}^{\mathrm{te}}\}}\,\Pi(\mathrm d\bm x,\mathrm d\xi)\right) +\sum_{k=1}^K\gamma_k\left(\int_{\mathcal X\times\Xi}\mathbf1_{\{\xi\le\xi_{(k)}^{\mathrm{te}}\}}\,\Pi(\mathrm d\bm x,\mathrm d\xi)-l_k\right).
%    \end{aligned}
%    \label{eq-upper-lagrangian-expanded}
%  \end{equation}
%  Collecting the terms involving $\Pi$ gives
%  \begin{equation}
%    \begin{aligned}
%      \mathcal L^+(\Pi,\eta,\bm\lambda,\bm\gamma)
%      ={}&\bm\lambda^\top\bm u-\bm\gamma^\top\bm l+\mathbb E_{\mathbb P}\left[\eta(\bm X)\right]\\
%      &+\int_{\mathcal X\times\Xi}\left[ \mathbf1_{\{\xi\le-\hat g(\bm x)\}}-\eta(\bm x) -\zeta_{\bm\lambda,\bm\gamma}(\xi) 
%      \right]\,\Pi(\mathrm d\bm x,\mathrm d\xi).
%    \end{aligned}
%    \label{eq-upper-lagrangian}
%  \end{equation}
%  The Lagrangian dual function is defined as
%  \begin{equation}
%    D^+(\eta,\bm\lambda,\bm\gamma)
%    :=\sup_{\Pi\in\mathcal M_+(\mathcal X\times\Xi)}\mathcal L^+(\Pi,\eta,\bm\lambda,\bm\gamma),
%    \label{eq-upper-dual-function}
%  \end{equation}
%  Since $\Pi$ is no longer subject to a prescribed total mass in the Lagrangian relaxation, if the integrand in \eqref{eq-upper-lagrangian} is strictly positive at any point $(\bm x,\xi)$, placing mass $M$ at that point and letting $M\to\infty$ yields $D^+=+\infty$. Otherwise, the supremum of the integral term is zero, attained by the zero measure. Consequently, the Lagrangian dual problem is
%  \begin{equation}
%    \begin{aligned}
%      \overline d_\theta
%      =\inf_{\substack{\eta\in\mathcal F_{\mathrm b}(\mathcal X)\\
%      \bm\lambda,\bm\gamma\in\mathbb R_+^K}}
%      \quad &\bm\lambda^\top\bm u-\bm\gamma^\top\bm l+\mathbb E_{\mathbb P}\left[\eta(\bm X)\right]\\
%      \text{s.t.}\quad & \eta(\bm x)\ge
%      \mathbf1_{\{\xi\le-\hat g(\bm x)\}}
%      -\zeta_{\bm\lambda,\bm\gamma}(\xi),
%      \qquad\forall(\bm x,\xi)\in\mathcal X\times\Xi.
%    \end{aligned}
%    \label{eq-upper-functional-dual}
%  \end{equation}
%  For fixed $\bm\lambda,\bm\gamma\in\mathbb R_+^K$, the pointwise smallest function satisfying the constraint in \eqref{eq-upper-functional-dual} is
%  \begin{equation}
%    \eta^+_{\bm\lambda,\bm\gamma}(\bm x)
%    :=\sup_{\xi\in\Xi}\left\{
%      \mathbf1_{\{\xi\le-\hat g(\bm x)\}}
%      -\zeta_{\bm\lambda,\bm\gamma}(\xi)
%    \right\}.
%    \label{eq-upper-marginal-envelope}
%  \end{equation}
%  This function is bounded in absolute value by $1+\|\bm\lambda\|_1+\|\bm\gamma\|_1$. Moreover, replacing $-\hat g(\bm x)$ by a real threshold $t$ in \eqref{eq-upper-marginal-envelope} yields a nondecreasing function of $t$, which is therefore Borel measurable. Since $\hat g$ is Borel measurable, it follows that $\eta^+_{\bm\lambda,\bm\gamma}\in\mathcal F_{\mathrm b}(\mathcal X)$. By monotonicity of the expectation, the optimal choice of $\eta$ in \eqref{eq-upper-functional-dual} for fixed $\bm\lambda,\bm\gamma$ is $\eta=\eta^+_{\bm\lambda,\bm\gamma}$. Substituting this choice into \eqref{eq-upper-functional-dual} yields \eqref{eq-upper-measure-dual}.
%
%  The derivation for the lower measure program \eqref{eq-measure-lower} is analogous, with the signs of the rank penalties reversed and the optimization directions interchanged. In this case, the optimal marginal multiplier for fixed $\bm\lambda,\bm\gamma$ is given by
%  \begin{equation}
%    \eta^-_{\bm\lambda,\bm\gamma}(\bm x)
%    :=\inf_{\xi\in\Xi}\left\{
%      \mathbf1_{\{\xi\le-\hat g(\bm x)\}}
%      +\zeta_{\bm\lambda,\bm\gamma}(\xi)
%    \right\},
%    \label{eq-lower-marginal-envelope}
%  \end{equation}
%  which also belongs to $\mathcal F_{\mathrm b}(\mathcal X)$ by the same argument. Substituting this choice into the corresponding dual problem yields \eqref{eq-lower-measure-dual}.
%
%  Finally, by weak duality, $\underline d_\theta\le\underline p_{\hat f}$ and $\overline p_{\hat f}\le\overline d_\theta$. Since the two primal programs share the same nonempty feasible set, we have $\underline p_{\hat f}\le\overline p_{\hat f}$. Combining these inequalities establishes \eqref{eq-weak-duality}.
%\end{proof}
%
%The residual supremum and infimum in Lemma~\ref{lem-weak-duality} combine the rank penalty with the failure indicator. We first make the penalty values explicit, then identify the failure coefficients associated with the input cells in \eqref{eq-input-cells}. For the default support $\Xi=\mathbb R$, partition the residual space into $J:=2K+1$ cells:
%\begin{equation}
%  \begin{aligned}
%    \Xi_{2k}&:=\left\{\xi_{(k)}^{\mathrm{te}}\right\},
%      &&k=1,\ldots,K,\\
%    \Xi_{2k+1}&:=\left(\xi_{(k)}^{\mathrm{te}},\xi_{(k+1)}^{\mathrm{te}}\right),
%      &&k=0,\ldots,K.
%  \end{aligned}
%  \label{eq-residual-cells}
%\end{equation}
%The two partitions satisfy
%\begin{equation}
%  \mathcal X=\bigsqcup_{i=0}^K\mathcal X_i,
%  \qquad
%  \Xi=\bigsqcup_{j=1}^J\Xi_j.
%  \label{eq-space-partitions}
%\end{equation}
%The singleton residual cells retain atoms at the observed residuals and distinguish the open upper rank constraints from the closed lower rank constraints. The step function is constant on each $\Xi_j$, with value
%\begin{equation}
%  \zeta_j:=\sum_{k=1}^K\lambda_k\mathbf1_{\{j<2k\}}
%    -\sum_{k=1}^K\gamma_k\mathbf1_{\{j\le2k\}},
%    \qquad j=1,\ldots,J.
%  \label{eq-dual-cell-values}
%\end{equation}
%For $(\bm x,\xi)\in\mathcal X_i\times\Xi_j$, the failure indicator satisfies
%\begin{equation}
%  \underline h_{ij}:=\mathbf1_{\{j\le2i\}}
%  \le\mathbf1_{\{\xi\le-\hat g(\bm x)\}}
%  \le\overline h_{ij}:=\mathbf1_{\{j\le2i+1\}}.
%  \label{eq-cell-costs}
%\end{equation}
%Residual cells with $j\le2i$ always cause failure for inputs in $\mathcal X_i$, and those with $j\ge2i+2$ are always safe. Only the overlapping interval $\Xi_{2i+1}$ distinguishes the two coefficients. At $-\hat g(\bm x)=\xi_{(i)}^{\mathrm{te}}$, this open interval is safe; attaining the upper coefficient then requires moving its assigned mass to the left endpoint. The proof below handles this case explicitly and constructs joint laws attaining the assigned objective values while preserving every rank constraint. The coefficients $w_i$, $\zeta_j$, and $\underline h_{ij},\overline h_{ij}$ now give finite linear programs for the dual values.
%
%\begin{theorem}[Finite-dimensional duals and strong duality]
%  \label{thm-finite-dual-reformulations}
%  For $\Xi=\mathbb R$, $\mathcal Q=\mathcal Q^{\mathrm{DW}}(\theta)$, and every $\theta\ge0$, the full dual values equal the optimal values of
%  \begin{equation}
%    \begin{aligned}
%      \overline d_\theta
%      =\min_{\bm\lambda,\bm\gamma\ge\bm 0,\ \bm z\in\mathbb R^{K+1}}\quad&
%        \bm\lambda^\top\bm u-\bm\gamma^\top\bm l+\sum_{i=0}^K w_i z_i\\
%      \text{s.t.}\quad&
%        z_i+\zeta_j\ge\overline h_{ij},
%        \qquad i=0,\ldots,K,\quad j=1,\ldots,J,
%    \end{aligned}
%    \label{eq-finite-upper-dual}
%  \end{equation}
%  and
%  \begin{equation}
%    \begin{aligned}
%      \underline d_\theta
%      =\max_{\bm\lambda,\bm\gamma\ge\bm 0,\ \bm z\in\mathbb R^{K+1}}\quad&
%        -\bm\lambda^\top\bm u+\bm\gamma^\top\bm l+\sum_{i=0}^K w_i z_i\\
%      \text{s.t.}\quad&
%        z_i-\zeta_j\le\underline h_{ij},
%        \qquad i=0,\ldots,K,\quad j=1,\ldots,J.
%    \end{aligned}
%    \label{eq-finite-lower-dual}
%  \end{equation}
%  Moreover,
%  \begin{equation}
%    \overline p_{\hat f}=\overline d_\theta,\qquad
%    \underline p_{\hat f}=\underline d_\theta.
%    \label{eq-strong-duality}
%  \end{equation}
%  Each finite LP has $3K+1$ scalar variables and $(K+1)(2K+1)$ pairwise inequalities, in addition to multiplier nonnegativity. The finite LP optima and the corresponding measure primal and dual optima are attained.
%\end{theorem}
%\begin{proof}
%  Using the partitions in \eqref{eq-space-partitions}, associate each feasible joint law with its matrix of joint cell probabilities $\bm\pi=(\pi_{ij})$, where
%  \begin{equation}
%    \pi_{ij}:=\Pi(\mathcal X_i\times\Xi_j),
%    \qquad i=0,\ldots,K,\quad j=1,\ldots,J.
%    \label{eq-joint-cell-masses}
%  \end{equation}
%  The row sums are the fixed input cell probabilities, and the rank constraints bound cumulative column sums. Hence $\bm\pi\in\mathcal C_\theta$, where
%  \begin{equation}
%    \mathcal C_\theta:=
%    \left\{\bm\pi\in\mathbb R_+^{(K+1)\times J}:
%    \begin{array}{ll}
%      \displaystyle\sum_{j=1}^J\pi_{ij}=w_i,
%        &i=0,\ldots,K,\\[1mm]
%      \displaystyle\sum_{i=0}^K\sum_{j<2k}\pi_{ij}\le u_k,
%        &k=1,\ldots,K,\\[1mm]
%      \displaystyle\sum_{i=0}^K\sum_{j\le2k}\pi_{ij}\ge l_k,
%        &k=1,\ldots,K
%    \end{array}\right\}.
%    \label{eq-cell-mass-polytope}
%  \end{equation}
%  By \eqref{eq-cell-costs}, maximizing $\sum_{i,j}\overline h_{ij}\pi_{ij}$ over this polytope bounds the upper measure problem from above, while minimizing $\sum_{i,j}\underline h_{ij}\pi_{ij}$ bounds the lower problem from below.
%
%  For the reverse inequalities, fix $\bm\pi\in\mathcal C_\theta$. Draw $\bm x$ from $\mathbb P$ and, on $\mathcal X_i$ with $w_i>0$, select a provisional residual cell $j$ with probability $\pi_{ij}/w_i$. For $j\ne2i+1$, choose a fixed representative of $\Xi_j$; both coefficients in \eqref{eq-cell-costs} then equal the actual failure indicator.
%
%  In the overlapping residual cell, choose $\xi=-\hat g(\bm x)$ for the upper problem. This attains failure and remains in $\Xi_{2i+1}$ whenever $-\hat g(\bm x)>\xi_{(i)}^{\mathrm{te}}$. If $-\hat g(\bm x)=\xi_{(i)}^{\mathrm{te}}$, it moves the assigned mass to the left endpoint. For $i\ge1$, any $\widetilde\xi\in(\xi_{(i)}^{\mathrm{te}},\xi_{(i+1)}^{\mathrm{te}})$, and every observed threshold $\xi_{(k)}^{\mathrm{te}}$,
%  \begin{equation}
%    \mathbf1_{\{\xi_{(i)}^{\mathrm{te}}<\xi_{(k)}^{\mathrm{te}}\}}=\mathbf1_{\{\widetilde\xi<\xi_{(k)}^{\mathrm{te}}\}},\qquad
%    \mathbf1_{\{\xi_{(i)}^{\mathrm{te}}\le \xi_{(k)}^{\mathrm{te}}\}}\ge\mathbf1_{\{\widetilde\xi\le \xi_{(k)}^{\mathrm{te}}\}}.
%    \label{eq-endpoint-move}
%  \end{equation}
%  Thus the move leaves all open upper probabilities unchanged and can only increase the closed lower probabilities. It preserves every rank constraint and attains the assigned upper coefficient. For $i=0$, the finite threshold $-\hat g(\bm x)$ cannot equal $\xi_{(0)}^{\mathrm{te}}=-\infty$.
%
%  For the lower problem, choose
%  $\xi=\left(\xi_{(i+1)}^{\mathrm{te}}-\hat g(\bm x)\right)/2$ in the overlapping residual cell when $i<K$, and $\xi=-\hat g(\bm x)+1$ when $i=K$. These choices remain in $\Xi_{2i+1}$ and exceed $-\hat g(\bm x)$, attaining safety even at $-\hat g(\bm x)=\xi_{(i)}^{\mathrm{te}}$. All selections are measurable and preserve the input marginal. An input cell with $w_i=0$ has a zero row and contributes no mass. We have therefore proved
%  \begin{equation}
%    \overline p_{\hat f}
%    =\max_{\bm\pi\in\mathcal C_\theta}\sum_{i=0}^K\sum_{j=1}^J
%      \overline h_{ij}\pi_{ij},\qquad
%    \underline p_{\hat f}
%    =\min_{\bm\pi\in\mathcal C_\theta}\sum_{i=0}^K\sum_{j=1}^J
%      \underline h_{ij}\pi_{ij}.
%    \label{eq-finite-primal-values}
%  \end{equation}
%  The upper construction may redistribute mass between adjacent residual cells; it preserves feasibility and the objective value needed for this equality.
%
%  The polytope is nonempty by empirical-product feasibility and compact because its nonnegative entries sum to one. Both finite primal optima are attained, and the constructions above lift them to optimal joint measures. Finite LP strong duality applies \citep[Section~5.2.3]{B.V04}. The row equalities have free multipliers $z_i$, and the rank inequalities have nonnegative multipliers $\lambda_k,\gamma_k$. The resulting upper constraints are $z_i+\zeta_j\ge\overline h_{ij}$, and the lower constraints are $z_i-\zeta_j\le\underline h_{ij}$, giving \eqref{eq-finite-upper-dual} and \eqref{eq-finite-lower-dual}.
%
%  Finally, an upper finite dual solution defines a bounded Borel function $\eta(\bm x)=z_i$ on $\mathcal X_i$ satisfying the functional upper-dual constraint, by \eqref{eq-cell-costs}. Thus $\eta^+_{\bm\lambda,\bm\gamma}(\bm x)\le z_i$ on $\mathcal X_i$. For a lower finite dual solution, the same argument gives $\eta^-_{\bm\lambda,\bm\gamma}(\bm x)\ge z_i$. At fixed rank multipliers, the reduced upper objective in \eqref{eq-upper-measure-dual} is therefore no greater than the finite upper-dual objective, and the reduced lower objective is no smaller than the finite lower-dual objective. Finite LP equality and Lemma~\ref{lem-weak-duality} force \eqref{eq-strong-duality}. Optimal finite dual solutions consequently provide optimal rank multipliers for the reduced duals and optimal marginal functions for the functional duals. The argument uses feasibility and boundedness of finite LPs and therefore includes $\theta=0$ and arbitrary atoms.
%\end{proof}
%
%The theorem establishes an exact finite reformulation. The ordered residual cells yield a further simplification: the best threshold bounds in \eqref{eq-rankwise-failure-bounds} are already optimal.
%
%\begin{corollary}[Closed-form failure probability bounds]
%  \label{thm-closed-form-bounds}
%  Under the conditions of Theorem~\ref{thm-finite-dual-reformulations},
%  \begin{align}
%    \overline p_{\hat f}
%    &=\min\left\{1,\min_{k=1,\ldots,K}[u_k+H\left(-\xi_{(k)}^{\mathrm{te}}\right)]\right\},
%    \label{eq-dw-closed-upper}\\
%    \underline p_{\hat f}
%    &=\max\left\{0,\max_{k=1,\ldots,K}[l_k+H\left(-\xi_{(k)}^{\mathrm{te}}\right)-1]\right\}.
%    \label{eq-dw-closed-lower}
%  \end{align}
%\end{corollary}
%\begin{proof}
%  Fix a nonnegative column-sum vector $\bm y=(y_1,\ldots,y_J)^\top$ with $\sum_{j=1}^J y_j=1$, and let $Y_j:=\sum_{r=1}^j y_r$. Write $\mathcal T(\bm w,\bm y)$ for the transportation matrices with row-sum vector $\bm w=(w_0,\ldots,w_K)^\top$ and column-sum vector $\bm y$. For the upper problem, maximizing mass on pairs $j\le2i+1$ is a maximum-flow problem: connect the source to row $i$ with capacity $w_i$, each row to its permitted columns with capacity two, and column $j$ to the sink with capacity $y_j$. A partial flow can be completed to a transportation matrix using the remaining row and column masses, so the maximum flow equals the upper transport optimum.
%
%  There is a cut of capacity one. Hence a minimum cut crosses no middle edge of capacity two. The row neighborhoods are nested prefixes. If the largest source-side row is $m<K$, adding all smaller rows and retaining exactly the neighboring columns cannot increase the cut capacity. Writing $k=m+1$, the resulting cut has capacity
%  \[
%    \sum_{i=k}^K w_i+\sum_{j=1}^{2k-1}y_j
%    =H\left(-\xi_{(k)}^{\mathrm{te}}\right)+Y_{2k-1}.
%  \]
%  The empty row set and the case $m=K$ give capacity one. By the max-flow/min-cut theorem \citep{F.F56},
%  \begin{equation}
%    \max_{\bm\pi\in\mathcal T(\bm w,\bm y)}
%      \sum_{i,j}\overline h_{ij}\pi_{ij}
%    =\min\left\{1,\min_k[H\left(-\xi_{(k)}^{\mathrm{te}}\right)+Y_{2k-1}]\right\}.
%    \label{eq-fixed-column-upper}
%  \end{equation}
%
%  For the lower problem, maximize safe mass on pairs $j>2i$ and subtract from one. The neighborhoods are nested suffixes. If the smallest source-side row is $k\ge1$, the corresponding normalized cut has capacity
%  \[
%    \sum_{i=0}^{k-1}w_i+\sum_{j=2k+1}^Jy_j
%    =2-H\left(-\xi_{(k)}^{\mathrm{te}}\right)-Y_{2k}.
%  \]
%  The empty row set and the case $k=0$ give capacity one. The same flow argument yields
%  \begin{equation}
%    \min_{\bm\pi\in\mathcal T(\bm w,\bm y)}
%      \sum_{i,j}\underline h_{ij}\pi_{ij}
%    =\max\left\{0,\max_k[H\left(-\xi_{(k)}^{\mathrm{te}}\right)+Y_{2k}-1]\right\}.
%    \label{eq-fixed-column-lower}
%  \end{equation}
%
%  Feasible column totals satisfy $Y_{2k-1}\le u_k$ and $Y_{2k}\ge l_k$. To attain all upper prefix bounds simultaneously, choose
%  \begin{equation}
%    y_1^+:=u_1,\qquad
%    y_{2k}^+:=u_{k+1}-u_k,\qquad
%    y_{2k+1}^+:=0,\quad k=1,\ldots,K.
%    \label{eq-upper-column-allocation}
%  \end{equation}
%  These nonnegative masses sum to one and have
%  $Y_{2k-1}^+=u_k$ and $Y_{2k}^+=u_{k+1}\ge l_k$.
%  To attain all lower prefix bounds simultaneously, choose
%  \begin{equation}
%    y_{2k-1}^-:=0,\qquad
%    y_{2k}^-:=l_k-l_{k-1}\quad(k=1,\ldots,K),
%    \qquad y_J^-:=1-l_K.
%    \label{eq-lower-column-allocation}
%  \end{equation}
%  These masses have $Y_{2k}^-=l_k$ and
%  $Y_{2k-1}^-=l_{k-1}\le u_k$.
%  Monotonicity and \eqref{eq-rank-compatibility} make both choices feasible. Substituting them into \eqref{eq-fixed-column-upper} and \eqref{eq-fixed-column-lower}, and using \eqref{eq-finite-primal-values}, proves the two closed forms.
%\end{proof}
%
%The closed forms also identify optimal rank multipliers and their corresponding marginal functions. For a rank $k$, the upper choice
%\[
%  \bm\lambda=\bm e_k,\qquad \bm\gamma=\bm 0,\qquad
%  \eta^+_{\bm e_k,\bm 0}(\bm x)=\mathbf1_{\{\hat g(\bm x)\le-\xi_{(k)}^{\mathrm{te}}\}}
%\]
%is feasible by \eqref{eq-upper-threshold-certificate} and has objective $u_k+H\left(-\xi_{(k)}^{\mathrm{te}}\right)$. The lower choice
%\[
%  \bm\lambda=\bm 0,\qquad \bm\gamma=\bm e_k,\qquad
%  \eta^-_{\bm 0,\bm e_k}(\bm x)=\mathbf1_{\{\hat g(\bm x)\le-\xi_{(k)}^{\mathrm{te}}\}}-1
%\]
%is feasible by \eqref{eq-lower-threshold-certificate} and has objective $l_k+H\left(-\xi_{(k)}^{\mathrm{te}}\right)-1$. Here $\bm e_k$ is the $k$th unit vector. With zero rank multipliers, the upper and lower envelopes are the constant functions one and zero, giving the outer candidates.
%
%Selecting a minimizing upper candidate or a maximizing lower candidate therefore gives an optimal full dual solution. In particular, an optimum needs at most one upper rank multiplier for the upper bound and at most one lower rank multiplier for the lower bound. Both families of constraints still define the feasible residual laws; their ordering allows these sparse optimal certificates. The upper bound combines residual mass below a threshold with the probability that the surrogate threshold is at least that value. The lower bound quantifies overlap that no admissible dependence structure can avoid.
%
%\begin{remark}[Known interval support]
%  \label{rem-interval-support}
%  Suppose $\Xi=[a,b]$ with $-\infty<a\le b<\infty$ and the observed residuals lie in this interval. Then
%  \begin{align}
%    \overline p_{\hat f}
%    &=\min\left\{H(-a),\min_{k=1,\ldots,K}[u_k+H\left(-\xi_{(k)}^{\mathrm{te}}\right)]\right\},
%    \label{eq-dw-interval-upper}\\
%    \underline p_{\hat f}
%    &=\max\left\{H(-b),\max_{k=1,\ldots,K}[l_k+H\left(-\xi_{(k)}^{\mathrm{te}}\right)-1]\right\}.
%    \label{eq-dw-interval-lower}
%  \end{align}
%  To obtain these expressions, represent the support restriction on $\mathbb R$ by
%  $\mathbb Q\{\xi<a\}\le0$ and $\mathbb Q\{\xi\le b\}\ge1$.
%  Add $a,b$ to the ordered threshold list. At a newly added $a$, use upper and lower bounds zero; at a newly added $b$, use upper and lower bounds one. At coincident observed thresholds, tighten only the upper bound at $a$ and the lower bound at $b$. The enlarged sequences satisfy \eqref{eq-rank-compatibility}, so the same finite-aggregation and ordered-cut arguments apply. The effective additional upper candidate is $H(-a)$, and the effective additional lower candidate is $H(-b)$; the other added candidates are redundant. If $a=b$, both bounds equal $H(-a)$. For the resulting partition on $\Xi=[a,b]$, intersect the singleton and open-interval residual cells from the augmented threshold list with $\Xi$ and omit empty cells. Including the support endpoints in this list retains their atoms and the corresponding changes in failure status.
%\end{remark}
%
%For a fixed $\theta$, evaluation consists of scalar DW inversion, sorting and counting the test residuals, evaluating $H\left(-\xi_{(k)}^{\mathrm{te}}\right)$ under $\mathbb P$, and scanning the $K$ upper and lower candidates. The endpoint table depends only on $N^{\mathrm{te}}$ and $\theta$ and can be precomputed. Once this table and the surrogate CDF values are available, the final scan costs $O(K)$; solving the finite LPs is unnecessary.
%
%For $\theta\ge\theta_{N^{\mathrm{te}}}(\delta)$, these exact bounds inherit the finite-sample guarantee of Theorem~\ref{theorem-01}. If the surrogate CDF values are estimated by simulation, simultaneous upper CDF bounds can replace $H\left(-\xi_{(k)}^{\mathrm{te}}\right)$ in the upper expression, and simultaneous lower CDF bounds can replace $H\left(-\xi_{(k)}^{\mathrm{te}}\right)$ in the lower expression. An independent auxiliary simulation can provide these simultaneous bounds conditionally on the test data. Its error probability must be included in the overall confidence budget.
%
%
%
%
% END PRESERVED SECTION 4 BEFORE NUMBERING AND PROSE REVISION

% BEGIN PRESERVED SECTION 4 BEFORE RIGOR AUDIT (2026-11-24)
%\section{Tractable Reformulation}
%\label{sec-reformulation}
%
%Scalar inversion turns the DW constraint into linear restrictions on residual probabilities. These restrictions lead to Lagrangian duals, exact finite linear programs, and closed-form failure probability bounds. Throughout this section, the training data, the test residuals, and $\theta\ge0$ are fixed. The main results use $\Xi=\mathbb R$; the final remark treats known interval support.
%
%For each empirical probability level $q_j:=j/N^{\mathrm{te}}$, define the scalar acceptance interval
%\begin{equation}
%  [\alpha_j(\theta),\beta_j(\theta)]
%  :=\left\{p\in[0,1]:
%  N^{\mathrm{te}}\phi(q_j,p)-\psi(q_j,p)\le\theta\right\},
%  \qquad j=0,\ldots,N^{\mathrm{te}}.
%  \label{eq-dw-rank-acceptance}
%\end{equation}
%Let $\xi_{(1)}^{\mathrm{te}}<\cdots<\xi_{(K)}^{\mathrm{te}}$ be the distinct ordered test residuals. Define
%\begin{equation} 
%    c_0 :=0, \quad
%    c_k :=\sum_{j=1}^{N^{\mathrm{te}}}\mathbf1_{\{\xi_j^{\mathrm{te}}\le\xi_{(k)}^{\mathrm{te}}\}}, \quad
%    l_k :=\alpha_{c_k}(\theta),  \quad
%    u_k :=\beta_{c_{k-1}}(\theta),\qquad k=1,\ldots,K.
%  \label{eq-dw-counts-and-levels}
%\end{equation}
%The cumulative counts retain all multiplicities, so $c_K=N^{\mathrm{te}}$ even when test residuals coincide.
%
%\begin{proposition}[CDF band and finite rank constraints]
%  \label{prop-finite-rank-constraints}
%  Each acceptance set in \eqref{eq-dw-rank-acceptance} is a nonempty closed interval containing $q_j$. Define
%  \[
%    L_\theta(t):=\alpha_{N^{\mathrm{te}}\widehat F^{\mathrm{te}}(t)}(\theta),
%    \qquad
%    U_\theta(t):=\beta_{N^{\mathrm{te}}\widehat F^{\mathrm{te}}(t)}(\theta).
%  \]
%  The DW ambiguity set has the equivalent representations
%  \begin{equation*}
%    \mathcal Q^{\mathrm{DW}}(\theta) =\left\{\mathbb Q\in\mathcal M_1(\Xi): L_\theta(t)\le F_{\mathbb Q}(t)\le U_\theta(t), \, \forall t\in\mathbb R\right\}
%  \end{equation*}
%  and
%  \begin{equation}
%    \mathcal Q^{\mathrm{DW}}(\theta)
%    =\left\{\mathbb Q\in\mathcal M_1(\Xi): \mathbb Q\{\xi<\xi_{(k)}^{\mathrm{te}}\}\le u_k, \, \mathbb Q\{\xi\le \xi_{(k)}^{\mathrm{te}}\}\ge l_k, \, k=1,\ldots,K \right\}.
%    \label{eq-finite-rank-constraints}
%  \end{equation}
%  Both endpoint sequences are nondecreasing in the empirical rank. Consequently, $l_k$ and $u_k$ are nondecreasing and satisfy
%  \begin{equation*}
%    0\le l_k,u_k\le1,\quad
%    l_{k-1}\le u_k,\quad
%    l_k\le u_{k+1},
%    \qquad k=1,\ldots,K,
%  \end{equation*}
%  where $l_0:=0$ and $u_{K+1}:=1$.
%\end{proposition}
%\begin{proof}
%  For fixed $q$, the Bernoulli divergence increases as $p$ moves away from $q$, while $\psi(q,p)$ cannot increase as the interval between $q$ and $p$ expands. The corrected discrepancy therefore decreases toward $q$ from below and increases above $q$. It is continuous for $p\in(0,1)$ and equals $-\psi_0(q)\le0$ at $p=q$. The boundary conventions in Section~\ref{sec-ambiguity-set} give a nonempty closed acceptance interval also for $q\in\{0,1\}$. Scalar inversion at each empirical CDF level yields the stated CDF band, following \citet[Supplement, Section~S.6]{L.J23}.
%
%  On $[\xi_{(k)}^{\mathrm{te}},\xi_{(k+1)}^{\mathrm{te}})$, the empirical CDF equals $c_k/N^{\mathrm{te}}$. Monotonicity of a candidate CDF makes its lower band constraint equivalent to $F_{\mathbb Q}\left(\xi_{(k)}^{\mathrm{te}}\right)\ge\alpha_{c_k}(\theta)$ and its upper band constraint equivalent to $F_{\mathbb Q}\left(\xi_{(k+1)}^{\mathrm{te}}-\right)\le\beta_{c_k}(\theta)$. The left exterior interval contributes $F_{\mathbb Q}\left(\xi_{(1)}^{\mathrm{te}}-\right)\le\beta_0(\theta)$, and the right exterior interval contributes $F_{\mathbb Q}\left(\xi_{(K)}^{\mathrm{te}}\right)\ge\alpha_{N^{\mathrm{te}}}(\theta)$. The other exterior constraints are automatic because $\alpha_0(\theta)=0$ and $\beta_{N^{\mathrm{te}}}(\theta)=1$. Using $F_{\mathbb Q}\left(\xi_{(k)}^{\mathrm{te}}-\right)=\mathbb Q\{\xi<\xi_{(k)}^{\mathrm{te}}\}$ and $F_{\mathbb Q}\left(\xi_{(k)}^{\mathrm{te}}\right)=\mathbb Q\{\xi\le \xi_{(k)}^{\mathrm{te}}\}$ proves \eqref{eq-finite-rank-constraints}.
%
%  With $p$ fixed, the corrected discrepancy decreases as $q$ approaches $p$ from below and increases above $p$. Since each acceptance interval contains $q$, both endpoints are nondecreasing in $q$. In particular, $\alpha_{c_k}(\theta)\le c_k/N^{\mathrm{te}}\le\beta_{c_k}(\theta)$, which gives the stated compatibility inequalities.
%\end{proof}
%
%The upper rank constraints limit mass strictly below each observed residual, while the lower constraints guarantee mass at or below it. The two half-lines retain possible atoms at the observed values. The scalar endpoints can be found by one-dimensional bisection, after which the ambiguity set is described by linear constraints on the residual law.
%
%Define the CDF of the surrogate response by $H(t):=\mathbb P\left\{\hat g(\bm X)\le t\right\}$ for $t\in\mathbb R$. The observed residuals classify inputs by the location of the failure threshold $-\hat g(\bm x)$. Set $\xi_{(0)}^{\mathrm{te}}:=-\infty$ and $\xi_{(K+1)}^{\mathrm{te}}:=+\infty$, with $H(-\infty):=0$ and $H(+\infty):=1$. For $i=0,\ldots,K$, the input cells and their probabilities are
%\begin{equation}
%  \mathcal X_i:=\left\{\bm x\in\mathcal X:
%    \xi_{(i)}^{\mathrm{te}}\le-\hat g(\bm x)<\xi_{(i+1)}^{\mathrm{te}}\right\},
%  \qquad
%  w_i:=\mathbb P(\mathcal X_i)
%    =H\left(-\xi_{(i)}^{\mathrm{te}}\right)-H\left(-\xi_{(i+1)}^{\mathrm{te}}\right).
%  \label{eq-input-cells}
%\end{equation}
%These $K+1$ cells partition $\mathcal X$, and each threshold event is a union of consecutive cells:
%\[
%  \mathbf1_{\{\hat g(\bm x)\le-\xi_{(k)}^{\mathrm{te}}\}}
%    =\sum_{i=k}^K\mathbf1_{\mathcal X_i}(\bm x),
%  \qquad
%  H\left(-\xi_{(k)}^{\mathrm{te}}\right)=\sum_{i=k}^K w_i,
%  \qquad k=1,\ldots,K.
%\]
%The half-open convention assigns an input with $-\hat g(\bm x)=\xi_{(k)}^{\mathrm{te}}$ to $\mathcal X_k$, preserving threshold equalities even when $H$ has atoms.
%
%For every rank $k$, the failure indicator satisfies 
%\[
%  \mathbf1_{\{\xi\le\xi_{(k)}^{\mathrm{te}}\}}
%    +\mathbf1_{\{\hat g(\bm x)\le-\xi_{(k)}^{\mathrm{te}}\}}-1
%  \le\mathbf1_{\{\xi\le-\hat g(\bm x)\}}
%  \le\mathbf1_{\{\xi<\xi_{(k)}^{\mathrm{te}}\}}
%    +\mathbf1_{\{\hat g(\bm x)\le-\xi_{(k)}^{\mathrm{te}}\}}.
%\]
%The upper inequality covers failure by a union of two events. The lower inequality guarantees failure on their intersection. Integrating under any feasible joint law gives
%\begin{equation*}
%  l_k+H\left(-\xi_{(k)}^{\mathrm{te}}\right)-1\le p_{\hat f}(\Pi)\le u_k+H\left(-\xi_{(k)}^{\mathrm{te}}\right),
%  \qquad k=1,\ldots,K.
%\end{equation*}
%These candidate bounds hold for every admissible dependence structure. The duality argument below establishes that the best candidate in each direction is exact.
%
%Substituting \eqref{eq-finite-rank-constraints} into the upper failure probability problem gives the following linear program over $\mathcal M_+(\mathcal X\times\Xi)$:
%\begin{equation}
%  \begin{aligned}
%    \overline p_{\hat f}
%    =\sup_{\Pi\in\mathcal M_+(\mathcal X\times\Xi)}\quad&
%      \int_{\mathcal X\times\Xi}
%      \mathbf1_{\{\xi\le -\hat g(\bm x)\}}\,\Pi(\mathrm d\bm x,\mathrm d\xi)\\
%    \text{s.t.}\quad&
%      \int_{\mathcal X\times\Xi}
%      \mathbf1_{\mathcal U}(\bm x)\,\Pi(\mathrm d\bm x,\mathrm d\xi)
%      =\mathbb P(\bm X\in\mathcal U),
%      &&\forall \mathcal U\in\mathcal B(\mathcal X),\\
%    &\int_{\mathcal X\times\Xi}
%      \mathbf1_{\{\xi<\xi_{(k)}^{\mathrm{te}}\}}\,\Pi(\mathrm d\bm x,\mathrm d\xi)\le u_k,
%      &&k=1,\ldots,K,\\
%    &\int_{\mathcal X\times\Xi}
%      \mathbf1_{\{\xi\le \xi_{(k)}^{\mathrm{te}}\}}\,\Pi(\mathrm d\bm x,\mathrm d\xi)\ge l_k,
%      &&k=1,\ldots,K.
%  \end{aligned}
%  \label{eq-measure-upper}
%\end{equation}
%The lower problem minimizes the same integral subject to the same constraints:
%\begin{equation*}
%  \underline p_{\hat f}
%  =\inf_{\Pi\in\mathcal P(\mathcal Q^{\mathrm{DW}}(\theta))}
%    \int_{\mathcal X\times\Xi}
%    \mathbf1_{\{\xi\le -\hat g(\bm x)\}}\,\Pi(\mathrm d\bm x,\mathrm d\xi).
%\end{equation*}
%The marginal equality fixes the input law and total mass one. The rank inequalities restrict the residual marginal while allowing its dependence on the input to vary. For every $\theta\ge0$, the product of $\mathbb P$ and the empirical residual law is feasible when $\Xi=\mathbb R$.
%
%For $\bm\lambda,\bm\gamma\in\mathbb R_+^K$, collect the rank penalties in the residual step function
%\begin{equation}
%  \zeta_{\bm\lambda,\bm\gamma}(\xi)
%  :=\sum_{k=1}^K\lambda_k\mathbf1_{\{\xi<\xi_{(k)}^{\mathrm{te}}\}}
%    -\sum_{k=1}^K\gamma_k\mathbf1_{\{\xi\le\xi_{(k)}^{\mathrm{te}}\}}.
%  \label{eq-dual-step-function}
%\end{equation}
%
%\begin{lemma}[Lagrangian dual problems and weak duality]
%  \label{lem-weak-duality}
%  Consider the upper measure program \eqref{eq-measure-upper} and its lower minimization counterpart. Write $\bm u=(u_1,\ldots,u_K)^\top$ and $\bm l=(l_1,\ldots,l_K)^\top$. The nonnegative multipliers $\bm\lambda,\bm\gamma\in\mathbb R_+^K$ correspond to the upper and lower rank constraints. The input marginal equality has an unrestricted multiplier $\eta\in\mathcal F_{\mathrm b}(\mathcal X)$, the space of bounded Borel measurable real-valued functions on $\mathcal X$. Eliminating $\eta$ pointwise gives the dual problems
%  \begin{align}
%    \overline d_\theta
%    &:=\inf_{\bm\lambda,\bm\gamma\in\mathbb R_+^K}
%    \left\{\bm\lambda^\top\bm u-\bm\gamma^\top\bm l
%    +\mathbb E_{\mathbb P}\left[\sup_{\xi\in\Xi}\left\{
%    \mathbf1_{\{\xi\le-\hat g(\bm X)\}}
%    -\zeta_{\bm\lambda,\bm\gamma}(\xi)
%    \right\}\right]\right\},
%    \label{eq-upper-measure-dual}\\
%    \underline d_\theta
%    &:=\sup_{\bm\lambda,\bm\gamma\in\mathbb R_+^K}
%    \left\{-\bm\lambda^\top\bm u+\bm\gamma^\top\bm l
%    +\mathbb E_{\mathbb P}\left[\inf_{\xi\in\Xi}\left\{
%    \mathbf1_{\{\xi\le-\hat g(\bm X)\}}
%    +\zeta_{\bm\lambda,\bm\gamma}(\xi)
%    \right\}\right]\right\}.
%    \label{eq-lower-measure-dual}
%  \end{align}
%  Their optimal values satisfy
%  \begin{equation}
%    \underline d_\theta\le\underline p_{\hat f}
%    \le\overline p_{\hat f}\le\overline d_\theta.
%    \label{eq-weak-duality}
%  \end{equation}
%\end{lemma}
%\begin{proof}
%  The input marginal equality is equivalent to
%  \begin{equation*}
%    \int_{\mathcal X\times\Xi}\eta(\bm x)\,\Pi(\mathrm d\bm x,\mathrm d\xi)
%    =\mathbb E_{\mathbb P}\left[\eta(\bm X)\right],
%    \qquad\forall\eta\in\mathcal F_{\mathrm b}(\mathcal X).
%  \end{equation*}
%  Relaxing the marginal equality and the rank inequalities gives the upper Lagrangian \citep[Section~5.1]{B.V04}
%  \begin{equation*}
%    \begin{aligned}
%      \mathcal L^+(\Pi,\eta,\bm\lambda,\bm\gamma) = & \int_{\mathcal X\times\Xi}\mathbf1_{\{\xi\le-\hat g(\bm x)\}}\,\Pi(\mathrm d\bm x,\mathrm d\xi) +\mathbb E_{\mathbb P}\left[\eta(\bm X)\right]-\int_{\mathcal X\times\Xi}\eta(\bm x)\,\Pi(\mathrm d\bm x,\mathrm d\xi)\\
%      &+\sum_{k=1}^K\lambda_k\left(u_k-\int_{\mathcal X\times\Xi}\mathbf1_{\{\xi<\xi_{(k)}^{\mathrm{te}}\}}\,\Pi(\mathrm d\bm x,\mathrm d\xi)\right) +\sum_{k=1}^K\gamma_k\left(\int_{\mathcal X\times\Xi}\mathbf1_{\{\xi\le\xi_{(k)}^{\mathrm{te}}\}}\,\Pi(\mathrm d\bm x,\mathrm d\xi)-l_k\right).
%    \end{aligned}
%  \end{equation*}
%  Collecting the terms involving $\Pi$ gives
%  \begin{equation*}
%    \mathcal L^+
%    =\bm\lambda^\top\bm u-\bm\gamma^\top\bm l
%      +\mathbb E_{\mathbb P}\left[\eta(\bm X)\right]
%      +\int_{\mathcal X\times\Xi}\left[
%        \mathbf1_{\{\xi\le-\hat g(\bm x)\}}-\eta(\bm x)-\zeta_{\bm\lambda,\bm\gamma}(\xi)
%      \right]\,\Pi(\mathrm d\bm x,\mathrm d\xi).
%  \end{equation*}
%  The Lagrangian dual function is defined as
%  \begin{equation*}
%    D^+(\eta,\bm\lambda,\bm\gamma)
%    :=\sup_{\Pi\in\mathcal M_+(\mathcal X\times\Xi)}\mathcal L^+(\Pi,\eta,\bm\lambda,\bm\gamma),
%  \end{equation*}
%  The relaxed measure has arbitrary finite nonnegative mass. If the integrand in the collected Lagrangian is positive at any point, scaling a point mass there by $M\to\infty$ gives $D^+=+\infty$. If the integrand is nonpositive everywhere, its supremum integral is zero, attained by the zero measure. The resulting functional dual is
%  \begin{equation*}
%    \begin{aligned}
%      \overline d_\theta
%      =\inf_{\substack{\eta\in\mathcal F_{\mathrm b}(\mathcal X)\\
%      \bm\lambda,\bm\gamma\in\mathbb R_+^K}}
%      \quad &\bm\lambda^\top\bm u-\bm\gamma^\top\bm l+\mathbb E_{\mathbb P}\left[\eta(\bm X)\right]\\
%      \text{s.t.}\quad & \eta(\bm x)\ge
%      \mathbf1_{\{\xi\le-\hat g(\bm x)\}}
%      -\zeta_{\bm\lambda,\bm\gamma}(\xi),
%      \qquad\forall(\bm x,\xi)\in\mathcal X\times\Xi.
%    \end{aligned}
%  \end{equation*}
%  For fixed rank multipliers, the pointwise smallest feasible marginal function is
%  \begin{equation*}
%    \eta^+_{\bm\lambda,\bm\gamma}(\bm x)
%    :=\sup_{\xi\in\Xi}\left\{
%      \mathbf1_{\{\xi\le-\hat g(\bm x)\}}
%      -\zeta_{\bm\lambda,\bm\gamma}(\xi)
%    \right\}.
%  \end{equation*}
%  The envelope is bounded in absolute value by $1+\|\bm\lambda\|_1+\|\bm\gamma\|_1$. As a function of the real threshold replacing $-\hat g(\bm x)$, it is nondecreasing and hence Borel measurable. Composing it with the measurable surrogate gives $\eta^+_{\bm\lambda,\bm\gamma}\in\mathcal F_{\mathrm b}(\mathcal X)$. This smallest feasible function minimizes the expectation and yields the upper dual \eqref{eq-upper-measure-dual}.
%
%  Reversing the inequality penalties for the lower problem gives the pointwise largest feasible marginal function
%  \begin{equation*}
%    \eta^-_{\bm\lambda,\bm\gamma}(\bm x)
%    :=\inf_{\xi\in\Xi}\left\{
%      \mathbf1_{\{\xi\le-\hat g(\bm x)\}}
%      +\zeta_{\bm\lambda,\bm\gamma}(\xi)
%    \right\},
%  \end{equation*}
%  which is bounded and Borel measurable by the same argument. Substitution gives the lower dual \eqref{eq-lower-measure-dual}.
%
%  For a primal feasible $\Pi$, the upper Lagrangian is at least $p_{\hat f}(\Pi)$ and the lower Lagrangian is at most $p_{\hat f}(\Pi)$. Every feasible upper dual solution therefore bounds the primal objective from above, and every feasible lower dual solution bounds it from below. Taking the respective optima and using the common nonempty primal feasible set gives \eqref{eq-weak-duality}.
%\end{proof}
%
%The remaining supremum and infimum combine the residual penalty with the failure indicator. For $\Xi=\mathbb R$, the penalty in \eqref{eq-dual-step-function} changes only at observed residuals. Separate these values from the intervening intervals to form $J:=2K+1$ residual cells:
%\begin{equation}
%  \Xi_{2k}:=\left\{\xi_{(k)}^{\mathrm{te}}\right\}
%    \quad\left(k=1,\ldots,K\right),
%  \qquad
%  \Xi_{2k+1}:=\left(\xi_{(k)}^{\mathrm{te}},\xi_{(k+1)}^{\mathrm{te}}\right)
%    \quad\left(k=0,\ldots,K\right).
%  \label{eq-residual-cells}
%\end{equation}
%The input cells in \eqref{eq-input-cells} and the residual cells in \eqref{eq-residual-cells} satisfy
%\begin{equation*}
%  \mathcal X=\bigsqcup_{i=0}^K\mathcal X_i,
%  \qquad
%  \Xi=\bigsqcup_{j=1}^J\Xi_j.
%\end{equation*}
%The singleton cells retain residual atoms and distinguish the open and closed half-lines in the rank constraints. The penalty is constant on each $\Xi_j$, with value
%\begin{equation}
%  \zeta_j:=\sum_{k=1}^K\lambda_k\mathbf1_{\{j<2k\}}
%    -\sum_{k=1}^K\gamma_k\mathbf1_{\{j\le2k\}},
%    \qquad j=1,\ldots,J.
%  \label{eq-dual-cell-values}
%\end{equation}
%For $(\bm x,\xi)\in\mathcal X_i\times\Xi_j$, the failure indicator satisfies
%\begin{equation}
%  \underline h_{ij}:=\mathbf1_{\{j\le2i\}}
%  \le\mathbf1_{\{\xi\le-\hat g(\bm x)\}}
%  \le\overline h_{ij}:=\mathbf1_{\{j\le2i+1\}}.
%  \label{eq-cell-costs}
%\end{equation}
%For inputs in $\mathcal X_i$, residual cells with $j\le2i$ always cause failure, and those with $j\ge2i+2$ are always safe. The overlapping interval $\Xi_{2i+1}$ determines the difference between the two coefficients. When $-\hat g(\bm x)=\xi_{(i)}^{\mathrm{te}}$, that interval is safe; the upper construction assigns its mass to the left endpoint. The proof below verifies this assignment. The weights $w_i$, penalty values in \eqref{eq-dual-cell-values}, and failure coefficients now give finite linear programs for the duals.
%
%\begin{theorem}[Finite-dimensional duals and strong duality]
%  \label{thm-finite-dual-reformulations}
%  For $\Xi=\mathbb R$, $\mathcal Q=\mathcal Q^{\mathrm{DW}}(\theta)$, and every $\theta\ge0$, the dual values in Lemma~\ref{lem-weak-duality} equal the optimal values of
%  \begin{equation}
%    \begin{aligned}
%      \overline d_\theta
%      =\min_{\bm\lambda,\bm\gamma\ge\bm 0,\ \bm z\in\mathbb R^{K+1}}\quad&
%        \bm\lambda^\top\bm u-\bm\gamma^\top\bm l+\sum_{i=0}^K w_i z_i\\
%      \text{s.t.}\quad&
%        z_i+\zeta_j\ge\overline h_{ij},
%        \qquad i=0,\ldots,K,\quad j=1,\ldots,J,
%    \end{aligned}
%    \label{eq-finite-upper-dual}
%  \end{equation}
%  and
%  \begin{equation}
%    \begin{aligned}
%      \underline d_\theta
%      =\max_{\bm\lambda,\bm\gamma\ge\bm 0,\ \bm z\in\mathbb R^{K+1}}\quad&
%        -\bm\lambda^\top\bm u+\bm\gamma^\top\bm l+\sum_{i=0}^K w_i z_i\\
%      \text{s.t.}\quad&
%        z_i-\zeta_j\le\underline h_{ij},
%        \qquad i=0,\ldots,K,\quad j=1,\ldots,J.
%    \end{aligned}
%    \label{eq-finite-lower-dual}
%  \end{equation}
%  Moreover,
%  \begin{equation}
%    \overline p_{\hat f}=\overline d_\theta,\qquad
%    \underline p_{\hat f}=\underline d_\theta.
%    \label{eq-strong-duality}
%  \end{equation}
%  Each LP has $3K+1$ scalar variables and $(K+1)(2K+1)$ pairwise inequalities, together with multiplier nonnegativity. The free variables $z_i$ correspond to the input cell probability equalities. The finite LP optima and the corresponding measure primal and functional dual optima are attained.
%\end{theorem}
%\begin{proof}
%  For a feasible joint law, let $\bm\pi=(\pi_{ij})$ be its matrix of joint cell probabilities,
%  \begin{equation*}
%    \pi_{ij}:=\Pi(\mathcal X_i\times\Xi_j),
%    \qquad i=0,\ldots,K,\quad j=1,\ldots,J.
%  \end{equation*}
%  Its row sums are the fixed input cell probabilities, and its cumulative column sums satisfy the rank constraints. Thus $\bm\pi\in\mathcal C_\theta$, where
%  \begin{equation}
%    \mathcal C_\theta:=
%    \left\{\bm\pi\in\mathbb R_+^{(K+1)\times J}:
%    \begin{array}{ll}
%      \displaystyle\sum_{j=1}^J\pi_{ij}=w_i,
%        &i=0,\ldots,K,\\[1mm]
%      \displaystyle\sum_{i=0}^K\sum_{j<2k}\pi_{ij}\le u_k,
%        &k=1,\ldots,K,\\[1mm]
%      \displaystyle\sum_{i=0}^K\sum_{j\le2k}\pi_{ij}\ge l_k,
%        &k=1,\ldots,K
%    \end{array}\right\}.
%    \label{eq-cell-mass-polytope}
%  \end{equation}
%  By \eqref{eq-cell-costs}, maximizing $\sum_{i,j}\overline h_{ij}\pi_{ij}$ over this polytope bounds the upper measure problem from above, while minimizing $\sum_{i,j}\underline h_{ij}\pi_{ij}$ bounds the lower problem from below.
%
%  To establish the reverse inequalities, fix $\bm\pi\in\mathcal C_\theta$. Draw $\bm x$ from $\mathbb P$ and, on $\mathcal X_i$ with $w_i>0$, select a provisional residual cell $j$ with probability $\pi_{ij}/w_i$. For $j\ne2i+1$, choose a fixed representative of $\Xi_j$. Both coefficients in \eqref{eq-cell-costs} then equal the actual failure indicator.
%
%  For the upper problem, assign $\xi=-\hat g(\bm x)$ in the overlapping cell. This causes failure and lies in $\Xi_{2i+1}$ when $-\hat g(\bm x)>\xi_{(i)}^{\mathrm{te}}$. At equality, the assignment moves the provisional mass to the left endpoint. For $i\ge1$, any $\widetilde\xi\in\left(\xi_{(i)}^{\mathrm{te}},\xi_{(i+1)}^{\mathrm{te}}\right)$, and each observed threshold,
%  \begin{equation*}
%    \mathbf1_{\{\xi_{(i)}^{\mathrm{te}}<\xi_{(k)}^{\mathrm{te}}\}}=\mathbf1_{\{\widetilde\xi<\xi_{(k)}^{\mathrm{te}}\}},\qquad
%    \mathbf1_{\{\xi_{(i)}^{\mathrm{te}}\le \xi_{(k)}^{\mathrm{te}}\}}\ge\mathbf1_{\{\widetilde\xi\le \xi_{(k)}^{\mathrm{te}}\}}.
%  \end{equation*}
%  The assignment leaves every open upper probability unchanged and can only increase the closed lower probabilities. It therefore preserves all rank constraints and attains the upper coefficient. For $i=0$, the finite threshold $-\hat g(\bm x)$ is strictly above $\xi_{(0)}^{\mathrm{te}}=-\infty$.
%
%  For the lower problem, choose
%  $\xi=\left(\xi_{(i+1)}^{\mathrm{te}}-\hat g(\bm x)\right)/2$ in the overlapping residual cell when $i<K$, and $\xi=-\hat g(\bm x)+1$ when $i=K$. These choices remain in $\Xi_{2i+1}$ and exceed $-\hat g(\bm x)$, attaining safety even at $-\hat g(\bm x)=\xi_{(i)}^{\mathrm{te}}$. All selections are measurable and preserve the input marginal. An input cell with $w_i=0$ has a zero row and contributes no mass. We have therefore proved
%  \begin{equation}
%    \overline p_{\hat f}
%    =\max_{\bm\pi\in\mathcal C_\theta}\sum_{i=0}^K\sum_{j=1}^J
%      \overline h_{ij}\pi_{ij},\qquad
%    \underline p_{\hat f}
%    =\min_{\bm\pi\in\mathcal C_\theta}\sum_{i=0}^K\sum_{j=1}^J
%      \underline h_{ij}\pi_{ij}.
%    \label{eq-finite-primal-values}
%  \end{equation}
%  The upper endpoint assignment can move mass between adjacent residual cells while preserving every rank constraint and the assigned objective value.
%
%  The empirical product law gives a point in $\mathcal C_\theta$. The polytope is compact because its nonnegative entries sum to one. Both finite primal optima are attained, and the preceding constructions give optimal joint measures. Finite LP strong duality applies \citep[Section~5.2.3]{B.V04}. The row equalities have free multipliers $z_i$, and the rank inequalities have nonnegative multipliers $\lambda_k,\gamma_k$. Their dual constraints are $z_i+\zeta_j\ge\overline h_{ij}$ and $z_i-\zeta_j\le\underline h_{ij}$, giving \eqref{eq-finite-upper-dual} and \eqref{eq-finite-lower-dual}.
%
%  An upper finite dual solution defines a bounded Borel function $\eta(\bm x)=z_i$ on $\mathcal X_i$. The indicator bounds \eqref{eq-cell-costs} make it feasible for the functional upper dual, so $\eta^+_{\bm\lambda,\bm\gamma}(\bm x)\le z_i$. A lower finite dual solution similarly gives $\eta^-_{\bm\lambda,\bm\gamma}(\bm x)\ge z_i$. The reduced upper objective is therefore no greater than the finite upper dual objective, and the reduced lower objective is no smaller than the finite lower dual objective.
%
%  Finite LP equality and Lemma~\ref{lem-weak-duality} give \eqref{eq-strong-duality}. Optimal finite dual solutions supply optimal rank multipliers and marginal functions for the functional duals. This argument uses feasibility and boundedness, so it includes $\theta=0$ and distributions with atoms.
%\end{proof}
%
%The finite LPs establish exactness of the duals. Their ordered failure coefficients give a further simplification: the best single-threshold candidates derived above are already optimal.
%
%\begin{corollary}[Closed-form bounds and optimal certificates]
%  \label{thm-closed-form-bounds}
%  Under the conditions of Theorem~\ref{thm-finite-dual-reformulations},
%  \begin{align}
%    \overline p_{\hat f}
%    &=\min\left\{1,\min_{k=1,\ldots,K}[u_k+H\left(-\xi_{(k)}^{\mathrm{te}}\right)]\right\},
%    \label{eq-dw-closed-upper}\\
%    \underline p_{\hat f}
%    &=\max\left\{0,\max_{k=1,\ldots,K}[l_k+H\left(-\xi_{(k)}^{\mathrm{te}}\right)-1]\right\}.
%    \label{eq-dw-closed-lower}
%  \end{align}
%  Each dual optimum admits a certificate with at most one nonzero rank multiplier, of value one.
%\end{corollary}
%\begin{proof}
%  Fix a nonnegative column-sum vector $\bm y=(y_1,\ldots,y_J)^\top$ with $\sum_{j=1}^J y_j=1$, and let $Y_j:=\sum_{r=1}^j y_r$. Write $\mathcal T(\bm w,\bm y)$ for the transportation matrices with row-sum vector $\bm w=(w_0,\ldots,w_K)^\top$ and column-sum vector $\bm y$. For the upper problem, maximizing mass on pairs $j\le2i+1$ is a maximum-flow problem: connect the source to row $i$ with capacity $w_i$, each row to its permitted columns with capacity two, and column $j$ to the sink with capacity $y_j$. A partial flow can be completed to a transportation matrix using the remaining row and column masses, so the maximum flow equals the upper transport optimum.
%
%  A cut of capacity one is available, so a minimum cut crosses no middle edge of capacity two. The row neighborhoods are nested prefixes. If the largest source-side row is $m<K$, including all smaller rows and retaining exactly their neighboring columns cannot increase the cut capacity. With $k=m+1$, this cut has capacity
%  \[
%    \sum_{i=k}^K w_i+\sum_{j=1}^{2k-1}y_j
%    =H\left(-\xi_{(k)}^{\mathrm{te}}\right)+Y_{2k-1}.
%  \]
%  The empty row set and the case $m=K$ give capacity one. By the max-flow/min-cut theorem \citep{F.F56},
%  \begin{equation*}
%    \max_{\bm\pi\in\mathcal T(\bm w,\bm y)}
%      \sum_{i,j}\overline h_{ij}\pi_{ij}
%    =\min\left\{1,\min_k[H\left(-\xi_{(k)}^{\mathrm{te}}\right)+Y_{2k-1}]\right\}.
%  \end{equation*}
%
%  For the lower problem, maximize safe mass on pairs $j>2i$ and subtract from one. The neighborhoods are nested suffixes. If the smallest source-side row is $k\ge1$, the corresponding normalized cut has capacity
%  \[
%    \sum_{i=0}^{k-1}w_i+\sum_{j=2k+1}^Jy_j
%    =2-H\left(-\xi_{(k)}^{\mathrm{te}}\right)-Y_{2k}.
%  \]
%  The empty row set and the case $k=0$ give capacity one. The same flow argument yields
%  \begin{equation*}
%    \min_{\bm\pi\in\mathcal T(\bm w,\bm y)}
%      \sum_{i,j}\underline h_{ij}\pi_{ij}
%    =\max\left\{0,\max_k[H\left(-\xi_{(k)}^{\mathrm{te}}\right)+Y_{2k}-1]\right\}.
%  \end{equation*}
%
%  Feasible column totals satisfy $Y_{2k-1}\le u_k$ and $Y_{2k}\ge l_k$. To attain all upper prefix bounds simultaneously, choose
%  \begin{equation*}
%    y_1^+:=u_1,\qquad
%    y_{2k}^+:=u_{k+1}-u_k,\qquad
%    y_{2k+1}^+:=0,\quad k=1,\ldots,K.
%  \end{equation*}
%  These nonnegative masses sum to one and have
%  $Y_{2k-1}^+=u_k$ and $Y_{2k}^+=u_{k+1}\ge l_k$.
%  To attain all lower prefix bounds simultaneously, choose
%  \begin{equation*}
%    y_{2k-1}^-:=0,\qquad
%    y_{2k}^-:=l_k-l_{k-1}\quad(k=1,\ldots,K),
%    \qquad y_J^-:=1-l_K.
%  \end{equation*}
%  These masses have $Y_{2k}^-=l_k$ and
%  $Y_{2k-1}^-=l_{k-1}\le u_k$.
%  Monotonicity and the compatibility conditions in Proposition~\ref{prop-finite-rank-constraints} make both allocations feasible. Substituting them into the upper and lower transport values above, then using \eqref{eq-finite-primal-values}, gives \eqref{eq-dw-closed-upper} and \eqref{eq-dw-closed-lower}.
%
%  Let $\bm e_k$ be the $k$th unit vector. An upper dual certificate is
%\[
%  \bm\lambda=\bm e_k,\qquad \bm\gamma=\bm 0,\qquad
%  \eta^+_{\bm e_k,\bm 0}(\bm x)=\mathbf1_{\{\hat g(\bm x)\le-\xi_{(k)}^{\mathrm{te}}\}}
%\]
%with objective $u_k+H\left(-\xi_{(k)}^{\mathrm{te}}\right)$. A lower dual certificate is
%\[
%  \bm\lambda=\bm 0,\qquad \bm\gamma=\bm e_k,\qquad
%  \eta^-_{\bm 0,\bm e_k}(\bm x)=\mathbf1_{\{\hat g(\bm x)\le-\xi_{(k)}^{\mathrm{te}}\}}-1
%\]
%with objective $l_k+H\left(-\xi_{(k)}^{\mathrm{te}}\right)-1$. The pointwise threshold inequalities derived above ensure feasibility. With zero rank multipliers, the optimal marginal functions are the constants one and zero, giving the outer candidates in the two closed forms.
%
%Choosing a minimizing candidate in \eqref{eq-dw-closed-upper} or a maximizing candidate in \eqref{eq-dw-closed-lower} gives an optimal full dual solution. Both families of rank constraints enter feasibility; their ordering permits an optimal certificate using at most one rank multiplier.
%\end{proof}
%
%The upper bound combines residual mass below an observed value with the probability that the surrogate threshold reaches that value. The lower bound measures overlap that every admissible dependence structure must retain.
%
%\begin{remark}[Known interval support]
%  \label{rem-interval-support}
%  Suppose $\Xi=[a,b]$ with $-\infty<a\le b<\infty$ and the observed residuals lie in this interval. Then
%  \begin{align*}
%    \overline p_{\hat f}
%    &=\min\left\{H(-a),\min_{k=1,\ldots,K}[u_k+H\left(-\xi_{(k)}^{\mathrm{te}}\right)]\right\},\\
%    \underline p_{\hat f}
%    &=\max\left\{H(-b),\max_{k=1,\ldots,K}[l_k+H\left(-\xi_{(k)}^{\mathrm{te}}\right)-1]\right\}.
%  \end{align*}
%  To obtain these expressions, represent the support restriction on $\mathbb R$ by
%  $\mathbb Q\{\xi<a\}\le0$ and $\mathbb Q\{\xi\le b\}\ge1$.
%  Add $a,b$ to the ordered threshold list. At a newly added $a$, use upper and lower bounds zero; at a newly added $b$, use upper and lower bounds one. At coincident observed thresholds, tighten only the upper bound at $a$ and the lower bound at $b$. The enlarged sequences satisfy the compatibility conditions in Proposition~\ref{prop-finite-rank-constraints}, so the same finite-aggregation and ordered-cut arguments apply. The effective additional upper candidate is $H(-a)$, and the effective additional lower candidate is $H(-b)$; the other added candidates are redundant. If $a=b$, both bounds equal $H(-a)$. For the resulting partition on $\Xi=[a,b]$, intersect the singleton and open-interval residual cells from the augmented threshold list with $\Xi$ and omit empty cells. Including the support endpoints in this list retains their atoms and the corresponding changes in failure status.
%\end{remark}
%
%For fixed $\theta$, evaluation requires scalar DW inversion, sorting and counting the test residuals, and computing the $K$ values $H\left(-\xi_{(k)}^{\mathrm{te}}\right)$. The endpoint table depends only on $N^{\mathrm{te}}$ and $\theta$ and can be precomputed. Given this table and the surrogate CDF values, the final upper and lower scans cost $O(K)$.
%
%For $\theta\ge\theta_{N^{\mathrm{te}}}(\delta)$, these exact bounds inherit the finite-sample guarantee of Theorem~\ref{theorem-01}. If the surrogate CDF values are estimated by simulation, simultaneous upper CDF bounds can replace $H\left(-\xi_{(k)}^{\mathrm{te}}\right)$ in the upper expression, and simultaneous lower CDF bounds can replace $H\left(-\xi_{(k)}^{\mathrm{te}}\right)$ in the lower expression. An independent auxiliary simulation can provide these simultaneous bounds conditionally on the test data. Its error probability must be included in the overall confidence budget.
%
%
% END PRESERVED SECTION 4 BEFORE RIGOR AUDIT


% BEGIN PRESERVED CLOSED-FORM PROOF BEFORE STRUCTURAL REFINEMENT (2026-11-24)
%\begin{proof}
%  Fix a nonnegative column-sum vector $\bm y=(y_1,\ldots,y_J)^\top$ with $\sum_{j=1}^J y_j=1$, and let $Y_j:=\sum_{r=1}^j y_r$. Write $\mathcal T(\bm w,\bm y)$ for the transportation matrices with row-sum vector $\bm w=(w_0,\ldots,w_K)^\top$ and column-sum vector $\bm y$. This set is nonempty because it contains the product matrix $\pi_{ij}=w_i y_j$, and the rank probabilities depend only on $\bm y$. For the upper problem, maximizing mass on pairs $j\le2i+1$ is a maximum-flow problem: connect the source to row $i$ with capacity $w_i$, each row to its permitted columns with capacity two, and column $j$ to the sink with capacity $y_j$. A transportation matrix restricts to a feasible flow on these permitted pairs. Conversely, a partial flow can be completed to a transportation matrix because the remaining row and column masses are nonnegative and have equal totals. Completion can only add permitted mass, so these two constructions identify the maximum flow with the upper transport optimum.
%
%  A cut of capacity one is available, so a minimum cut crosses no middle edge of capacity two. The row neighborhoods are nested prefixes. If the largest source-side row is $m<K$, including all smaller rows and retaining exactly their neighboring columns cannot increase the cut capacity. With $k=m+1$, this cut has capacity
%  \[
%    \sum_{i=k}^K w_i+\sum_{j=1}^{2k-1}y_j
%    =H\left(-\xi_{(k)}^{\mathrm{te}}\right)+Y_{2k-1}.
%  \]
%  The empty row set and the case $m=K$ give capacity one. By the max-flow/min-cut theorem \citep{F.F56},
%  \begin{equation*}
%    \max_{\bm\pi\in\mathcal T(\bm w,\bm y)}
%      \sum_{i,j}\overline h_{ij}\pi_{ij}
%    =\min\left\{1,\min_k[H\left(-\xi_{(k)}^{\mathrm{te}}\right)+Y_{2k-1}]\right\}.
%  \end{equation*}
%
%  For the lower problem, maximize safe mass on pairs $j>2i$ and subtract from one. The neighborhoods are nested suffixes. A finite minimum cut with a nonempty source-side row set can be normalized by including every row with index $i\ge k$, where $k$ is the smallest source-side row index, and retaining exactly the neighboring columns. For $k\ge1$, the cut has capacity
%  \[
%    \sum_{i=0}^{k-1}w_i+\sum_{j=2k+1}^Jy_j
%    =2-H\left(-\xi_{(k)}^{\mathrm{te}}\right)-Y_{2k}.
%  \]
%  The empty row set and the case $k=0$ give capacity one. The same flow argument yields
%  \begin{equation*}
%    \min_{\bm\pi\in\mathcal T(\bm w,\bm y)}
%      \sum_{i,j}\underline h_{ij}\pi_{ij}
%    =\max\left\{0,\max_k[H\left(-\xi_{(k)}^{\mathrm{te}}\right)+Y_{2k}-1]\right\}.
%  \end{equation*}
%
%  The upper transport value is nondecreasing in the open prefixes $Y_{2k-1}$, and the lower transport value is nondecreasing in the closed prefixes $Y_{2k}$. Feasibility gives $Y_{2k-1}\le u_k$ and $Y_{2k}\ge l_k$. Thus maximizing the former and minimizing the latter give the claimed bounds if all relevant prefix limits can be attained simultaneously. For the upper prefixes, choose
%  \begin{equation*}
%    y_1^+:=u_1,\qquad
%    y_{2k}^+:=u_{k+1}-u_k,\qquad
%    y_{2k+1}^+:=0,\quad k=1,\ldots,K.
%  \end{equation*}
%  These nonnegative masses sum to one and have
%  $Y_{2k-1}^+=u_k$ and $Y_{2k}^+=u_{k+1}\ge l_k$.
%  To attain all lower prefix bounds simultaneously, choose
%  \begin{equation*}
%    y_{2k-1}^-:=0,\qquad
%    y_{2k}^-:=l_k-l_{k-1}\quad(k=1,\ldots,K),
%    \qquad y_J^-:=1-l_K.
%  \end{equation*}
%  These masses have $Y_{2k}^-=l_k$ and
%  $Y_{2k-1}^-=l_{k-1}\le u_k$.
%  Monotonicity and the compatibility conditions in Proposition~\ref{prop-finite-rank-constraints} make both allocations feasible. Substituting them into the upper and lower transport values above, then using \eqref{eq-finite-primal-values}, gives \eqref{eq-dw-closed-upper} and \eqref{eq-dw-closed-lower}.
%
%  Let $\bm e_k$ be the $k$th unit vector. An upper dual certificate is
%\[
%  \bm\lambda=\bm e_k,\qquad \bm\gamma=\bm 0,\qquad
%  \eta^+_{\bm e_k,\bm 0}(\bm x)=\mathbf1_{\{\hat g(\bm x)\le-\xi_{(k)}^{\mathrm{te}}\}}
%\]
%with objective $u_k+H\left(-\xi_{(k)}^{\mathrm{te}}\right)$. A lower dual certificate is
%\[
%  \bm\lambda=\bm 0,\qquad \bm\gamma=\bm e_k,\qquad
%  \eta^-_{\bm 0,\bm e_k}(\bm x)=\mathbf1_{\{\hat g(\bm x)\le-\xi_{(k)}^{\mathrm{te}}\}}-1
%\]
%with objective $l_k+H\left(-\xi_{(k)}^{\mathrm{te}}\right)-1$. The pointwise threshold inequalities give feasibility. For the upper envelope, when $-\hat g(\bm x)\ge\xi_{(k)}^{\mathrm{te}}$, choosing $\xi=\xi_{(k)}^{\mathrm{te}}$ attains value one. For the lower envelope, when $-\hat g(\bm x)<\xi_{(k)}^{\mathrm{te}}$, choosing $\xi\in\left(-\hat g(\bm x),\xi_{(k)}^{\mathrm{te}}\right]$ attains value minus one. In the complementary cases, a residual above both thresholds attains value zero. At threshold equality, the upper and lower envelopes are therefore one and zero, respectively. Inputs at equality are included in the closed CDF value $H\left(-\xi_{(k)}^{\mathrm{te}}\right)$. With zero rank multipliers, the marginal functions are the constants one and zero, giving the outer candidates.
%
%Choosing a minimizing candidate in \eqref{eq-dw-closed-upper} or a maximizing candidate in \eqref{eq-dw-closed-lower} gives an optimal full dual solution. Both families of rank constraints enter feasibility; their ordering permits an optimal certificate using at most one rank multiplier.
%\end{proof}
% END PRESERVED CLOSED-FORM PROOF BEFORE STRUCTURAL REFINEMENT


% BEGIN PRESERVED ACTIVE MANUSCRIPT BEFORE INFORMATION-LIMIT EXTENSION (2026-11-24)
%\documentclass[review,11pt,3p,authoryear,nopreprintline]{elsarticle}
%
%\usepackage[T1]{fontenc}
%\usepackage{amsmath,amssymb,amsthm}
%\usepackage{mathrsfs}
%\usepackage{bm}
%\usepackage{xcolor}
%\usepackage[colorlinks=true,citecolor=blue,linkcolor=blue,urlcolor=blue]{hyperref}
%
%\newtheorem{theorem}{Theorem}
%\newtheorem{proposition}{Proposition}
%\newtheorem{corollary}{Corollary}
%\newtheorem{lemma}{Lemma}
%\newdefinition{example}{Example}
%\newdefinition{definition}{Definition}
%\newdefinition{remark}{Remark}
%\newdefinition{property}{Property}
%\newdefinition{condition}{Condition}
%
%\begin{document}
%
%\begin{frontmatter}
%  \title{Failure Probability Bounds with Uncertain Surrogate Residual Distributions}
%
%  \author[inst1]{Renlong Wang}
%  \ead{rw106@rice.edu}
%  \affiliation[inst1]{organization={Department of Civil and Environmental Engineering, Rice University}, city={Houston}, postcode={TX 77005}, country={USA}}
%
%  \begin{abstract}
%  \end{abstract}
%
%  \begin{keyword}
%  \end{keyword}
%\end{frontmatter}
%
%\section{Introduction}
%
%\section{Problem Description}
%\label{sec-problem}
%
%Let $\bm X\sim\mathbb P$ be a random input taking values in $\mathcal X\subseteq\mathbb R^n$, and let $g:\mathcal X\to\mathbb R$ be a measurable limit-state function. The failure probability is defined as
%\[
%  p_f:=\Pr\{g(\bm X)\le0\}=\mathbb E_{\bm X\sim\mathbb P}\left[\mathbf 1_{\{g(\bm X)\le0\}}\right]=\int_{\mathcal X}\mathbf 1_{\{g(\bm x)\le0\}}\,\mathbb P(\mathrm d\bm x).
%\]
%For fixed $g$, this functional is linear in $\mathbb P$.
%
%Evaluating $g$ can require an expensive high-fidelity simulation. We use the training dataset
%\[
%  \mathcal D^{\mathrm{tr}}:=\left\{\left(\bm x_i^{\mathrm{tr}},g(\bm x_i^{\mathrm{tr}})\right)\right\}_{i=1}^{N^{\mathrm{tr}}}
%\]
%to fit a computationally efficient, measurable surrogate $\hat g:=\hat g_{\mathcal D^{\mathrm{tr}}}$. Model fitting and selection use only the training data and are completed before the test data are used. We condition on $\mathcal D^{\mathrm{tr}}$, treat $\hat g$ as fixed, and suppress this dependence in the notation.
%
%Define the signed residual function
%\[
%  \xi(\bm x):=g(\bm x)-\hat g(\bm x),\qquad \forall \bm x\in\mathcal X.
%\]
%For a new input $\bm X\sim\mathbb P$ independent of the training data, let $\mathbb Q^\star$ be the distribution of $\xi^\star:=\xi(\bm X)$. Let $\Xi\subseteq\mathbb R$ be a known Borel set with $\mathbb Q^\star(\Xi)=1$; the default choice is $\Xi=\mathbb R$. For any Borel set $\mathcal S$ in a Euclidean space, write $\mathcal B(\mathcal S)$ for its Borel sigma-algebra. Let $\mathcal M_+(\mathcal S)$ denote the finite nonnegative Borel measures concentrated on $\mathcal S$, and let $\mathcal M_1(\mathcal S)$ denote the probability measures.
%
%For a candidate residual law $\mathbb Q\in\mathcal M_1(\Xi)$, let $\Pi$ be a joint distribution on $\mathcal X\times\Xi$ with marginals $\Pi_{\bm X}=\mathbb P$ and $\Pi_\xi=\mathbb Q$. Its surrogate-based failure probability is
%\begin{equation}
%  p_{\hat f}(\Pi):=\Pr\{\hat g(\bm X)+\xi\le0\}=\int_{\mathcal X\times\Xi}\mathbf 1_{\{\hat g(\bm x)+\xi\le0\}}\,\Pi(\mathrm d\bm x,\mathrm d\xi).
%  \label{eq-surrogate-failure-probability}
%\end{equation}
%For fixed $\hat g$, this functional is linear in $\Pi$. In particular, the true joint distribution $\Pi^\star$ satisfies $p_{\hat f}(\Pi^\star)=p_f$.
%
%The failure probability depends on the residual law and its dependence on the input. Near the surrogate failure threshold, even a small residual can change the failure classification. We represent uncertainty in the residual law by an ambiguity set $\mathcal Q$ calibrated to contain $\mathbb Q^\star$ with high confidence. For a given $\mathcal Q$, define the joint ambiguity set
%\[
%  \mathcal P(\mathcal Q) :=\left\{\Pi\in\mathcal M_1(\mathcal X\times\Xi):\Pi_{\bm X}=\mathbb P,\ \Pi_\xi\in\mathcal Q\right\}.
%\]
%The joint ambiguity set allows all dependence structures consistent with the specified marginals. The upper and lower failure probability bounds are
%\begin{equation}
%  \overline p_{\hat f}:=\sup_{\Pi\in\mathcal P(\mathcal Q)}p_{\hat f}(\Pi),
%  \label{eq-upper-bound}
%\end{equation}
%and
%\begin{equation}
%  \underline p_{\hat f}:=\inf_{\Pi\in\mathcal P(\mathcal Q)}p_{\hat f}(\Pi),
%  \label{eq-lower-bound}
%\end{equation}
%respectively. If $\mathbb Q^\star\in\mathcal Q$, then $\Pi^\star\in\mathcal P(\mathcal Q)$ and hence $\underline p_{\hat f}\le p_f\le\overline p_{\hat f}$.
%
%The analysis has two tasks:
%\begin{enumerate}
%  \item \textbf{Ambiguity set design and statistical calibration.} Construct a residual ambiguity set and establish its coverage guarantee under explicit sampling conditions.
%  \item \textbf{Tractable reformulations.} Derive exact and computationally tractable reformulations of the infinite-dimensional optimization problems defining the failure probability bounds.
%\end{enumerate}
%
%
%\section{Ambiguity Set Design}
%\label{sec-ambiguity-set}
%
%Failure occurs when $\xi\le-\hat g(\bm X)$, so the residual cumulative distribution function (CDF) gives the marginal mass at or below each failure threshold. Accurate tail control is needed for small failure probabilities. We use the D\"umbgen--Wellner (DW) construction \citep{L.J23} to control empirical CDF deviations simultaneously over all thresholds, with tolerances that vary across probability levels.
%
%Let $q$ and $p$ denote an empirical probability and a candidate probability. Following \citet{B.J79} and the CDF confidence bands of \citet{O95}, we measure their local discrepancy by the Bernoulli Kullback--Leibler divergence
%\begin{equation}
%  \phi(q,p):=q\log\frac{q}{p}+(1-q)\log\frac{1-q}{1-p},\qquad \forall q\in[0,1],\, p\in(0,1),
%  \label{eq-bernoulli-kl}
%\end{equation}
%where logarithms are natural and a term with zero numerator is interpreted as zero. For $q$ close to an interior $p$, the approximation $\phi(q,p)\approx(q-p)^2/[2p(1-p)]$ scales the deviation by the Bernoulli variance. Thus, the same absolute deviation receives a larger weight near the probability boundaries. The full divergence also covers zero and full empirical counts through $\phi(0,p)=-\log(1-p)$ and $\phi(1,p)=-\log p$. At the boundaries of the candidate probability, we set $\phi(0,0)=\phi(1,1)=0$, $\phi(q,0)=+\infty$ for $q>0$, and $\phi(q,1)=+\infty$ for $q<1$.
%
%To compare deviations simultaneously across probability levels, we use the additive correction of \citet{L.J23}. For $v\in(0,1)$, define
%\begin{equation*}
%  \rho(v):=\log\log\frac{e}{4v(1-v)},\qquad \omega(v):=\log\!\left(1+\rho(v)^2\right).
%\end{equation*}
%Before sampling, we fix the divergence parameter $s=1$, which gives \eqref{eq-bernoulli-kl}, and the correction parameter $\nu=1$. Set $\psi_0(v):=\rho(v)+\omega(v)$, with $\psi_0(0)=\psi_0(1)=+\infty$, and define the bivariate correction
%\begin{equation*}
%  \psi(q,p):=\min_{v\in[\min\{q,p\},\,\max\{q,p\}]}\psi_0(v).
%\end{equation*}
%The functions $\rho$ and $\omega$ vanish at $v=1/2$ and increase toward the probability boundaries. The minimum over the interval between $q$ and $p$ is therefore large only when both probabilities lie in the same tail. This correction balances fluctuations across probability scales in the simultaneous CDF constraint.
%
%We construct the empirical CDF from a test dataset
%\[
%  \mathcal D^{\mathrm{te}}:=\left\{\left(\bm x_j^{\mathrm{te}},g(\bm x_j^{\mathrm{te}})\right)\right\}_{j=1}^{N^{\mathrm{te}}}.
%\]
%The test inputs are independent draws from $\mathbb P$ and are independent of $\mathcal D^{\mathrm{tr}}$. We fix $N^{\mathrm{te}}\ge1$ before sampling and use the test data to calibrate the residual ambiguity set, with the fitted surrogate held fixed.
%
%Define the signed test residuals
%\[
%  \xi_j^{\mathrm{te}}:=\xi(\bm x_j^{\mathrm{te}})=g(\bm x_j^{\mathrm{te}})-\hat g(\bm x_j^{\mathrm{te}}),\qquad j=1,\ldots,N^{\mathrm{te}}.
%\]
%Conditional on $\mathcal D^{\mathrm{tr}}$, these residuals are independent and identically distributed according to $\mathbb Q^\star$. Their empirical CDF is
%\begin{equation}
%  \widehat F^{\mathrm{te}}(t):=\frac{1}{N^{\mathrm{te}}}\sum_{j=1}^{N^{\mathrm{te}}}\mathbf 1_{\{\xi_j^{\mathrm{te}}\le t\}},\qquad \forall t\in\mathbb R.
%  \label{eq-empirical-cdf}
%\end{equation}
%For $\mathbb Q\in\mathcal M_1(\Xi)$, write $F_{\mathbb Q}(t):=\mathbb Q\left(({-\infty},t]\right)$. At a threshold $t$, $q=\widehat F^{\mathrm{te}}(t)$ is the mean of the Bernoulli indicators in \eqref{eq-empirical-cdf}, and $p=F_{\mathbb Q}(t)$ is the candidate probability. Thus $N^{\mathrm{te}}\phi(q,p)$ is the Bernoulli log-likelihood difference between the empirical and candidate probabilities. Combining it with the multiscale correction gives $N^{\mathrm{te}}\phi(q,p)-\psi(q,p)$.
%
%At the probability boundaries, the corrected discrepancy is $-\infty$ for $q=p\in\{0,1\}$ and $+\infty$ for $p=0<q$ or $p=1>q$. These conventions extend the constraint to every residual threshold.
%
%\begin{definition}[DW-based CDF ambiguity set]
%  \label{definition-dw-ambiguity-set}
%  Given the empirical CDF $\widehat F^{\mathrm{te}}$ and a threshold $\theta\ge0$, define
%  \begin{equation} 
%    \label{eq-cdf-ambiguity-set}
%      \mathcal Q^{\mathrm{DW}}(\theta):=\left\{\mathbb Q\in\mathcal M_1(\Xi):\ 
%      \sup_{t\in\mathbb R}\left[ N^{\mathrm{te}}\phi\!\left(\widehat F^{\mathrm{te}}(t),F_{\mathbb Q}(t)\right) 
%       {}-\psi\!\left(\widehat F^{\mathrm{te}}(t),F_{\mathbb Q}(t)\right)\right]\le\theta\right\}. 
%  \end{equation} 
%\end{definition}
%
%We calibrate $\theta$ against a distribution-free reference law. Following \citet{L.J23}, let $V_1,\ldots,V_{N^{\mathrm{te}}}$ be independent standard uniform variables, independent of the training and test data. Their empirical CDF is
%\begin{equation*}
%  G_{N^{\mathrm{te}}}(u):=\frac{1}{N^{\mathrm{te}}}\sum_{j=1}^{N^{\mathrm{te}}}\mathbf 1_{\{V_j\le u\}},\qquad \forall u\in[0,1].
%\end{equation*}
%Define the uniform reference statistic
%\begin{equation}
%  T_{N^{\mathrm{te}}}^{\mathrm U}:=\sup_{u\in(0,1)}\left[N^{\mathrm{te}}\phi\!\left(G_{N^{\mathrm{te}}}(u),u\right)-\psi\!\left(G_{N^{\mathrm{te}}}(u),u\right)\right].
%  \label{eq-dw-reference-statistic}
%\end{equation}
%
%A quantile of this reference statistic determines the calibration threshold.
%
%\begin{proposition}[Finite-sample threshold calibration]
%  \label{proposition-01}
%  Under the sampling conditions stated above, let $\delta\in(0,1)$ and define
%  \begin{equation}
%    \theta_{N^{\mathrm{te}}}(\delta):=\inf\left\{\tau\ge0:\Pr\left\{T_{N^{\mathrm{te}}}^{\mathrm U}\le\tau\right\}\ge1-\delta\right\}.
%    \label{eq-dw-calibrated-threshold}
%  \end{equation}
%
%  Then, for every $\theta\ge\theta_{N^{\mathrm{te}}}(\delta)$, the DW-based CDF ambiguity set satisfies
%  \begin{equation}
%    \Pr\left\{\mathbb Q^\star\in\mathcal Q^{\mathrm{DW}}(\theta)\,\middle|\,\mathcal D^{\mathrm{tr}}\right\}\ge1-\delta.
%    \label{eq-cdf-coverage}
%  \end{equation}
%\end{proposition}
%\begin{proof}
%  The reference statistic is finite almost surely, and its distribution function is right-continuous. Hence \eqref{eq-dw-calibrated-threshold} gives
%  \[
%    \Pr\left\{T_{N^{\mathrm{te}}}^{\mathrm U}>\theta_{N^{\mathrm{te}}}(\delta)\right\}\le\delta.
%  \]
%  Conditioning on $\mathcal D^{\mathrm{tr}}$ fixes $\hat g$ and $\mathbb Q^\star$. Define the generalized quantile function $F_{\mathbb Q^\star}^{-1}(u):=\inf\{t\in\mathbb R:F_{\mathbb Q^\star}(t)\ge u\}$ for all $u\in(0,1)$. The sample $F_{\mathbb Q^\star}^{-1}(V_j)$ for $j=1,\ldots,N^{\mathrm{te}}$, has the same conditional joint law as the test residuals. Thus, given $\mathcal D^{\mathrm{tr}}$,
%  \[
%    \left(\widehat F^{\mathrm{te}}(t)\right)_{t\in\mathbb R}\overset{\mathrm d}{=}\left(G_{N^{\mathrm{te}}}\!\left(F_{\mathbb Q^\star}(t)\right)\right)_{t\in\mathbb R}.
%  \]
%  For $0<F_{\mathbb Q^\star}(t)<1$, the discrepancy at $\mathbb Q^\star$ is evaluated at a probability level included in \eqref{eq-dw-reference-statistic}. When $F_{\mathbb Q^\star}(t)$ is zero or one, the empirical CDF equals that boundary almost surely, giving discrepancy $-\infty$. The supremum in \eqref{eq-cdf-ambiguity-set}, evaluated at $\mathbb Q^\star$, is therefore bounded by $T_{N^{\mathrm{te}}}^{\mathrm U}$ under this coupling. Since $\mathbb Q^\star(\Xi)=1$, we obtain
%  \[
%    \Pr\left\{\mathbb Q^\star\notin\mathcal Q^{\mathrm{DW}}\!\left(\theta_{N^{\mathrm{te}}}(\delta)\right)\,\middle|\,\mathcal D^{\mathrm{tr}}\right\}
%    \le\Pr\left\{T_{N^{\mathrm{te}}}^{\mathrm U}>\theta_{N^{\mathrm{te}}}(\delta)\right\}\le\delta.
%  \]
%  For every test sample, $\theta\ge\theta_{N^{\mathrm{te}}}(\delta)$ implies $\mathcal Q^{\mathrm{DW}}\!\left(\theta_{N^{\mathrm{te}}}(\delta)\right)\subseteq\mathcal Q^{\mathrm{DW}}(\theta)$. This inclusion proves \eqref{eq-cdf-coverage} for every such $\theta$.
%\end{proof}
%
%Proposition~\ref{proposition-01} applies to arbitrary residual distributions, including laws with atoms or infinite moments. With $\phi$ and $\psi$ fixed, the reference law depends only on $N^{\mathrm{te}}$. The threshold $\theta_{N^{\mathrm{te}}}(\delta)$ can therefore be calibrated offline for each confidence level $1-\delta$. \citet{L.J23} describe finite-sample evaluation using uniform order-statistic boundary probabilities. Any conservative upper bound on the quantile in \eqref{eq-dw-calibrated-threshold} preserves \eqref{eq-cdf-coverage}.
%% If calibration is performed by simulation, quantile estimation error must be accounted for in the confidence guarantee. A threshold chosen from the test data also preserves \eqref{eq-cdf-coverage}, provided $\theta\ge\theta_{N^{\mathrm{te}}}(\delta)$ almost surely, by the same nesting argument.
%
%Residual-law coverage gives simultaneous bounds on the failure probability under the joint model.
%
%\begin{theorem}[Finite-sample failure probability bounds]
%  \label{theorem-01}
%  Under the sampling conditions of Proposition~\ref{proposition-01}, let $\delta\in(0,1)$ and $\theta\ge\theta_{N^{\mathrm{te}}}(\delta)$, with the calibrated threshold defined in \eqref{eq-dw-calibrated-threshold}. Then the bounds defined in \eqref{eq-upper-bound} and \eqref{eq-lower-bound}, evaluated with $\mathcal Q=\mathcal Q^{\mathrm{DW}}(\theta)$, satisfy
%  \begin{equation}
%    \Pr\left\{\underline p_{\hat f}\le p_f\le\overline p_{\hat f}\,\middle|\,\mathcal D^{\mathrm{tr}}\right\}\ge1-\delta.
%    \label{eq-failure-probability-coverage}
%  \end{equation}
%\end{theorem}
%
%\begin{proof}
%  On the event $\mathbb Q^\star\in\mathcal Q^{\mathrm{DW}}(\theta)$, the true joint distribution $\Pi^\star$ belongs to $\mathcal P\left(\mathcal Q^{\mathrm{DW}}(\theta)\right)$. The definitions of the bounds and the identity $p_{\hat f}(\Pi^\star)=p_f$ give
%  \[
%    \underline p_{\hat f}
%    \le p_{\hat f}(\Pi^\star)
%    =p_f
%    \le\overline p_{\hat f}.
%  \]
%  Consequently,
%  \[
%    \Pr\left\{
%      \underline p_{\hat f}\le p_f\le\overline p_{\hat f}
%      \,\middle|\,\mathcal D^{\mathrm{tr}}
%    \right\}
%    \ge
%    \Pr\left\{
%      \mathbb Q^\star\in\mathcal Q^{\mathrm{DW}}(\theta)
%      \,\middle|\,\mathcal D^{\mathrm{tr}}
%    \right\}.
%  \]
%  Proposition~\ref{proposition-01} bounds the right-hand side below by $1-\delta$, proving \eqref{eq-failure-probability-coverage}.
%\end{proof}
%
%
%The event $\mathbb Q^\star\in\mathcal Q^{\mathrm{DW}}(\theta)$ gives both inequalities, so $1-\delta$ is the confidence level for the entire interval. The threshold $\theta$ controls the corrected likelihood discrepancy; the resulting CDF tolerances vary with the probability level and can be asymmetric about the empirical CDF. For fixed test data, $0\le\theta_1\le\theta_2$ implies $\mathcal Q^{\mathrm{DW}}(\theta_1)\subseteq\mathcal Q^{\mathrm{DW}}(\theta_2)$. The lower failure bound is therefore nonincreasing in $\theta$, and the upper bound is nondecreasing. Increasing the prescribed confidence level gives a nondecreasing calibrated threshold and hence a wider or unchanged interval.
%
%The interval accounts for finite-sample uncertainty in the residual law and uncertainty in its dependence on the input. Because the joint ambiguity set includes all couplings with the specified marginals, a gap can remain even when the residual marginal is known exactly. The same coverage argument applies to any residual ambiguity set containing $\mathbb Q^\star$ with conditional probability at least $1-\delta$ given $\mathcal D^{\mathrm{tr}}$.
%
%
%\section{Tractable Reformulation}
%\label{sec-reformulation}
%
%Scalar inversion turns the DW constraint into linear restrictions on residual probabilities. These restrictions lead to Lagrangian duals, exact finite linear programs, and closed-form failure probability bounds. Throughout this section, the training data, the test residuals, and $\theta\ge0$ are fixed. The main results use $\Xi=\mathbb R$; the final remark treats known interval support.
%
%For each empirical probability level $q_j:=j/N^{\mathrm{te}}$, define the scalar acceptance interval
%\begin{equation}
%  [\alpha_j(\theta),\beta_j(\theta)]
%  :=\left\{p\in[0,1]:
%  N^{\mathrm{te}}\phi(q_j,p)-\psi(q_j,p)\le\theta\right\},
%  \qquad j=0,\ldots,N^{\mathrm{te}}.
%  \label{eq-dw-rank-acceptance}
%\end{equation}
%Let $\xi_{(1)}^{\mathrm{te}}<\cdots<\xi_{(K)}^{\mathrm{te}}$ be the distinct ordered test residuals. Define
%\begin{equation} 
%    c_0 :=0, \quad
%    c_k :=\sum_{j=1}^{N^{\mathrm{te}}}\mathbf1_{\{\xi_j^{\mathrm{te}}\le\xi_{(k)}^{\mathrm{te}}\}}, \quad
%    l_k :=\alpha_{c_k}(\theta),  \quad
%    u_k :=\beta_{c_{k-1}}(\theta),\qquad k=1,\ldots,K.
%  \label{eq-dw-counts-and-levels}
%\end{equation}
%The cumulative counts retain all multiplicities, so $c_K=N^{\mathrm{te}}$ even when test residuals coincide.
%
%\begin{proposition}[CDF band and finite rank constraints]
%  \label{prop-finite-rank-constraints}
%  Each acceptance set in \eqref{eq-dw-rank-acceptance} is a nonempty closed interval containing $q_j$. Define
%  \[
%    L_\theta(t):=\alpha_{N^{\mathrm{te}}\widehat F^{\mathrm{te}}(t)}(\theta),
%    \qquad
%    U_\theta(t):=\beta_{N^{\mathrm{te}}\widehat F^{\mathrm{te}}(t)}(\theta).
%  \]
%  The DW ambiguity set has the equivalent representations
%  \begin{equation*}
%    \mathcal Q^{\mathrm{DW}}(\theta) =\left\{\mathbb Q\in\mathcal M_1(\Xi): L_\theta(t)\le F_{\mathbb Q}(t)\le U_\theta(t), \, \forall t\in\mathbb R\right\}
%  \end{equation*}
%  and
%  \begin{equation}
%    \mathcal Q^{\mathrm{DW}}(\theta)
%    =\left\{\mathbb Q\in\mathcal M_1(\Xi): \mathbb Q\{\xi<\xi_{(k)}^{\mathrm{te}}\}\le u_k, \, \mathbb Q\{\xi\le \xi_{(k)}^{\mathrm{te}}\}\ge l_k, \, k=1,\ldots,K \right\}.
%    \label{eq-finite-rank-constraints}
%  \end{equation}
%  Both endpoint sequences are nondecreasing in the empirical rank. Consequently, $l_k$ and $u_k$ are nondecreasing and satisfy
%  \begin{equation*}
%    0\le l_k,u_k\le1,\quad
%    l_{k-1}\le u_k,\quad
%    l_k\le u_{k+1},
%    \qquad k=1,\ldots,K,
%  \end{equation*}
%  where $l_0:=0$ and $u_{K+1}:=1$.
%\end{proposition}
%\begin{proof}
%  For fixed $q$, the Bernoulli divergence increases as $p$ moves away from $q$, while $\psi(q,p)$ cannot increase as the interval between $q$ and $p$ expands. The corrected discrepancy therefore decreases toward $q$ from below and increases above $q$. It is continuous for $p\in(0,1)$ and equals $-\psi_0(q)\le0$ at $p=q$. The boundary conventions in Section~\ref{sec-ambiguity-set} give a nonempty closed acceptance interval also for $q\in\{0,1\}$. Scalar inversion at each empirical CDF level yields the stated CDF band, following \citet[Supplement, Section~S.6]{L.J23}.
%
%  For $k=1,\ldots,K-1$, the empirical CDF on $\left[\xi_{(k)}^{\mathrm{te}},\xi_{(k+1)}^{\mathrm{te}}\right)$ equals $c_k/N^{\mathrm{te}}$. Monotonicity of a candidate CDF makes its lower band constraint equivalent to $F_{\mathbb Q}\left(\xi_{(k)}^{\mathrm{te}}\right)\ge\alpha_{c_k}(\theta)$ and its upper band constraint equivalent to $F_{\mathbb Q}\left(\xi_{(k+1)}^{\mathrm{te}}-\right)\le\beta_{c_k}(\theta)$. The left exterior interval contributes $F_{\mathbb Q}\left(\xi_{(1)}^{\mathrm{te}}-\right)\le\beta_0(\theta)$, and the right exterior interval contributes $F_{\mathbb Q}\left(\xi_{(K)}^{\mathrm{te}}\right)\ge\alpha_{N^{\mathrm{te}}}(\theta)$. The other exterior constraints are automatic because $\alpha_0(\theta)=0$ and $\beta_{N^{\mathrm{te}}}(\theta)=1$. Using $F_{\mathbb Q}\left(\xi_{(k)}^{\mathrm{te}}-\right)=\mathbb Q\{\xi<\xi_{(k)}^{\mathrm{te}}\}$ and $F_{\mathbb Q}\left(\xi_{(k)}^{\mathrm{te}}\right)=\mathbb Q\{\xi\le \xi_{(k)}^{\mathrm{te}}\}$ proves \eqref{eq-finite-rank-constraints}.
%
%  With $p$ fixed, the corrected discrepancy decreases as $q$ approaches $p$ from either side. For empirical levels $q<q'$, an accepted value $p<q$ at level $q'$ remains accepted at level $q$; an accepted value $p>q'$ at level $q$ remains accepted at level $q'$. Since each acceptance interval contains its empirical level, these implications make both endpoints nondecreasing. The inequalities $\alpha_{c_k}(\theta)\le c_k/N^{\mathrm{te}}\le\beta_{c_k}(\theta)$ then give the stated compatibility conditions. This argument permits singleton intervals: if $q_j=1/2$ and $\theta=0$, the acceptance interval is $\{1/2\}$ because $\psi(1/2,p)=0$.
%\end{proof}
%
%The upper rank constraints limit mass strictly below each observed residual, while the lower constraints guarantee mass at or below it. The two half-lines retain possible atoms at the observed values. The scalar endpoints can be found by one-dimensional bisection, after which the ambiguity set is described by linear constraints on the residual law.
%
%Define the CDF of the surrogate response by $H(t):=\mathbb P\left\{\hat g(\bm X)\le t\right\}$ for $t\in\mathbb R$. The observed residuals classify inputs by the location of the failure threshold $-\hat g(\bm x)$. Set $\xi_{(0)}^{\mathrm{te}}:=-\infty$ and $\xi_{(K+1)}^{\mathrm{te}}:=+\infty$, with $H(-\infty):=0$ and $H(+\infty):=1$. For $i=0,\ldots,K$, the input cells and their probabilities are
%\begin{equation}
%  \mathcal X_i:=\left\{\bm x\in\mathcal X:
%    \xi_{(i)}^{\mathrm{te}}\le-\hat g(\bm x)<\xi_{(i+1)}^{\mathrm{te}}\right\},
%  \qquad
%  w_i:=\mathbb P(\mathcal X_i)
%    =H\left(-\xi_{(i)}^{\mathrm{te}}\right)-H\left(-\xi_{(i+1)}^{\mathrm{te}}\right).
%  \label{eq-input-cells}
%\end{equation}
%These $K+1$ cells partition $\mathcal X$, and each threshold event is a union of consecutive cells:
%\[
%  \mathbf1_{\{\hat g(\bm x)\le-\xi_{(k)}^{\mathrm{te}}\}}
%    =\sum_{i=k}^K\mathbf1_{\mathcal X_i}(\bm x),
%  \qquad
%  H\left(-\xi_{(k)}^{\mathrm{te}}\right)=\sum_{i=k}^K w_i,
%  \qquad k=1,\ldots,K.
%\]
%The half-open convention assigns an input with $-\hat g(\bm x)=\xi_{(k)}^{\mathrm{te}}$ to $\mathcal X_k$, preserving threshold equalities even when $H$ has atoms.
%
%For every rank $k$, the failure indicator satisfies 
%\[
%  \mathbf1_{\{\xi\le\xi_{(k)}^{\mathrm{te}}\}}
%    +\mathbf1_{\{\hat g(\bm x)\le-\xi_{(k)}^{\mathrm{te}}\}}-1
%  \le\mathbf1_{\{\xi\le-\hat g(\bm x)\}}
%  \le\mathbf1_{\{\xi<\xi_{(k)}^{\mathrm{te}}\}}
%    +\mathbf1_{\{\hat g(\bm x)\le-\xi_{(k)}^{\mathrm{te}}\}}.
%\]
%The upper inequality covers failure by a union of two events. The lower inequality guarantees failure on their intersection. Integrating under any feasible joint law gives
%\begin{equation*}
%  l_k+H\left(-\xi_{(k)}^{\mathrm{te}}\right)-1\le p_{\hat f}(\Pi)\le u_k+H\left(-\xi_{(k)}^{\mathrm{te}}\right),
%  \qquad k=1,\ldots,K.
%\end{equation*}
%These candidate bounds hold for every admissible dependence structure. The duality argument below establishes that the best candidate in each direction is exact.
%
%Substituting \eqref{eq-finite-rank-constraints} into the upper failure probability problem gives the following linear program over $\mathcal M_+(\mathcal X\times\Xi)$:
%\begin{equation}
%  \begin{aligned}
%    \overline p_{\hat f}
%    =\sup_{\Pi\in\mathcal M_+(\mathcal X\times\Xi)}\quad&
%      \int_{\mathcal X\times\Xi}
%      \mathbf1_{\{\xi\le -\hat g(\bm x)\}}\,\Pi(\mathrm d\bm x,\mathrm d\xi)\\
%    \text{s.t.}\quad&
%      \int_{\mathcal X\times\Xi}
%      \mathbf1_{\mathcal U}(\bm x)\,\Pi(\mathrm d\bm x,\mathrm d\xi)
%      =\mathbb P(\bm X\in\mathcal U),
%      &&\forall \mathcal U\in\mathcal B(\mathcal X),\\
%    &\int_{\mathcal X\times\Xi}
%      \mathbf1_{\{\xi<\xi_{(k)}^{\mathrm{te}}\}}\,\Pi(\mathrm d\bm x,\mathrm d\xi)\le u_k,
%      &&k=1,\ldots,K,\\
%    &\int_{\mathcal X\times\Xi}
%      \mathbf1_{\{\xi\le \xi_{(k)}^{\mathrm{te}}\}}\,\Pi(\mathrm d\bm x,\mathrm d\xi)\ge l_k,
%      &&k=1,\ldots,K.
%  \end{aligned}
%  \label{eq-measure-upper}
%\end{equation}
%The lower problem minimizes the same integral subject to the same constraints:
%\begin{equation*}
%  \underline p_{\hat f}
%  =\inf_{\Pi\in\mathcal P(\mathcal Q^{\mathrm{DW}}(\theta))}
%    \int_{\mathcal X\times\Xi}
%    \mathbf1_{\{\xi\le -\hat g(\bm x)\}}\,\Pi(\mathrm d\bm x,\mathrm d\xi).
%\end{equation*}
%The marginal equality fixes the input law and total mass one. The rank inequalities restrict the residual marginal while allowing its dependence on the input to vary. For every $\theta\ge0$, the product of $\mathbb P$ and the empirical residual law is feasible when $\Xi=\mathbb R$.
%
%For $\bm\lambda,\bm\gamma\in\mathbb R_+^K$, collect the rank penalties in the residual step function
%\begin{equation}
%  \zeta_{\bm\lambda,\bm\gamma}(\xi)
%  :=\sum_{k=1}^K\lambda_k\mathbf1_{\{\xi<\xi_{(k)}^{\mathrm{te}}\}}
%    -\sum_{k=1}^K\gamma_k\mathbf1_{\{\xi\le\xi_{(k)}^{\mathrm{te}}\}}.
%  \label{eq-dual-step-function}
%\end{equation}
%
%\begin{lemma}[Lagrangian dual problems and weak duality]
%  \label{lem-weak-duality}
%  Consider the feasible upper measure program \eqref{eq-measure-upper} and its lower minimization counterpart. Write $\bm u=(u_1,\ldots,u_K)^\top$ and $\bm l=(l_1,\ldots,l_K)^\top$. The nonnegative multipliers $\bm\lambda,\bm\gamma\in\mathbb R_+^K$ correspond to the upper and lower rank constraints. The input marginal equality has an unrestricted multiplier $\eta\in\mathcal F_{\mathrm b}(\mathcal X)$, the space of bounded Borel measurable real-valued functions on $\mathcal X$. Eliminating $\eta$ pointwise gives the dual problems
%  \begin{align}
%    \overline d_\theta
%    &:=\inf_{\bm\lambda,\bm\gamma\in\mathbb R_+^K}
%    \left\{\bm\lambda^\top\bm u-\bm\gamma^\top\bm l
%    +\mathbb E_{\mathbb P}\left[\sup_{\xi\in\Xi}\left\{
%    \mathbf1_{\{\xi\le-\hat g(\bm X)\}}
%    -\zeta_{\bm\lambda,\bm\gamma}(\xi)
%    \right\}\right]\right\},
%    \label{eq-upper-measure-dual}\\
%    \underline d_\theta
%    &:=\sup_{\bm\lambda,\bm\gamma\in\mathbb R_+^K}
%    \left\{-\bm\lambda^\top\bm u+\bm\gamma^\top\bm l
%    +\mathbb E_{\mathbb P}\left[\inf_{\xi\in\Xi}\left\{
%    \mathbf1_{\{\xi\le-\hat g(\bm X)\}}
%    +\zeta_{\bm\lambda,\bm\gamma}(\xi)
%    \right\}\right]\right\}.
%    \label{eq-lower-measure-dual}
%  \end{align}
%  Their optimal values satisfy
%  \begin{equation}
%    \underline d_\theta\le\underline p_{\hat f}
%    \le\overline p_{\hat f}\le\overline d_\theta.
%    \label{eq-weak-duality}
%  \end{equation}
%\end{lemma}
%\begin{proof}
%  The input marginal equality is equivalent to
%  \begin{equation*}
%    \int_{\mathcal X\times\Xi}\eta(\bm x)\,\Pi(\mathrm d\bm x,\mathrm d\xi)
%    =\mathbb E_{\mathbb P}\left[\eta(\bm X)\right],
%    \qquad\forall\eta\in\mathcal F_{\mathrm b}(\mathcal X).
%  \end{equation*}
%  Relaxing the marginal equality and the rank inequalities gives the upper Lagrangian \citep[Section~5.1]{B.V04}
%  \begin{equation*}
%    \begin{aligned}
%      \mathcal L^+(\Pi,\eta,\bm\lambda,\bm\gamma) = & \int_{\mathcal X\times\Xi}\mathbf1_{\{\xi\le-\hat g(\bm x)\}}\,\Pi(\mathrm d\bm x,\mathrm d\xi) +\mathbb E_{\mathbb P}\left[\eta(\bm X)\right]-\int_{\mathcal X\times\Xi}\eta(\bm x)\,\Pi(\mathrm d\bm x,\mathrm d\xi)\\
%      &+\sum_{k=1}^K\lambda_k\left(u_k-\int_{\mathcal X\times\Xi}\mathbf1_{\{\xi<\xi_{(k)}^{\mathrm{te}}\}}\,\Pi(\mathrm d\bm x,\mathrm d\xi)\right) +\sum_{k=1}^K\gamma_k\left(\int_{\mathcal X\times\Xi}\mathbf1_{\{\xi\le\xi_{(k)}^{\mathrm{te}}\}}\,\Pi(\mathrm d\bm x,\mathrm d\xi)-l_k\right).
%    \end{aligned}
%  \end{equation*}
%  Collecting the terms involving $\Pi$ gives
%  \begin{equation*}
%    \mathcal L^+
%    =\bm\lambda^\top\bm u-\bm\gamma^\top\bm l
%      +\mathbb E_{\mathbb P}\left[\eta(\bm X)\right]
%      +\int_{\mathcal X\times\Xi}\left[
%        \mathbf1_{\{\xi\le-\hat g(\bm x)\}}-\eta(\bm x)-\zeta_{\bm\lambda,\bm\gamma}(\xi)
%      \right]\,\Pi(\mathrm d\bm x,\mathrm d\xi).
%  \end{equation*}
%  The Lagrangian dual function is defined as
%  \begin{equation*}
%    D^+(\eta,\bm\lambda,\bm\gamma)
%    :=\sup_{\Pi\in\mathcal M_+(\mathcal X\times\Xi)}\mathcal L^+(\Pi,\eta,\bm\lambda,\bm\gamma),
%  \end{equation*}
%  The relaxed measure has arbitrary finite nonnegative mass. If the integrand in the collected Lagrangian is positive at any point, scaling a point mass there by $M\to\infty$ gives $D^+=+\infty$. If the integrand is nonpositive everywhere, its supremum integral is zero, attained by the zero measure. The resulting functional dual is
%  \begin{equation*}
%    \begin{aligned}
%      \overline d_\theta
%      =\inf_{\substack{\eta\in\mathcal F_{\mathrm b}(\mathcal X)\\
%      \bm\lambda,\bm\gamma\in\mathbb R_+^K}}
%      \quad &\bm\lambda^\top\bm u-\bm\gamma^\top\bm l+\mathbb E_{\mathbb P}\left[\eta(\bm X)\right]\\
%      \text{s.t.}\quad & \eta(\bm x)\ge
%      \mathbf1_{\{\xi\le-\hat g(\bm x)\}}
%      -\zeta_{\bm\lambda,\bm\gamma}(\xi),
%      \qquad\forall(\bm x,\xi)\in\mathcal X\times\Xi.
%    \end{aligned}
%  \end{equation*}
%  For fixed rank multipliers, the pointwise smallest feasible marginal function is
%  \begin{equation*}
%    \eta^+_{\bm\lambda,\bm\gamma}(\bm x)
%    :=\sup_{\xi\in\Xi}\left\{
%      \mathbf1_{\{\xi\le-\hat g(\bm x)\}}
%      -\zeta_{\bm\lambda,\bm\gamma}(\xi)
%    \right\}.
%  \end{equation*}
%  The envelope is bounded in absolute value by $1+\|\bm\lambda\|_1+\|\bm\gamma\|_1$. As a function of the real threshold replacing $-\hat g(\bm x)$, it is nondecreasing and hence Borel measurable. Composing it with the measurable surrogate gives $\eta^+_{\bm\lambda,\bm\gamma}\in\mathcal F_{\mathrm b}(\mathcal X)$. This smallest feasible function minimizes the expectation and yields the upper dual \eqref{eq-upper-measure-dual}.
%
%  For the lower problem, reversing the inequality penalties gives
%  \begin{equation*}
%    \mathcal L^-
%    =-\bm\lambda^\top\bm u+\bm\gamma^\top\bm l
%      +\mathbb E_{\mathbb P}\left[\eta(\bm X)\right]
%      +\int_{\mathcal X\times\Xi}\left[
%        \mathbf1_{\{\xi\le-\hat g(\bm x)\}}-\eta(\bm x)+\zeta_{\bm\lambda,\bm\gamma}(\xi)
%      \right]\,\Pi(\mathrm d\bm x,\mathrm d\xi).
%  \end{equation*}
%  Its infimum over finite nonnegative measures is finite exactly when $\eta(\bm x)\le\mathbf1_{\{\xi\le-\hat g(\bm x)\}}+\zeta_{\bm\lambda,\bm\gamma}(\xi)$ for every $(\bm x,\xi)$. Under this condition, the infimum integral is zero. The pointwise largest feasible marginal function is therefore
%  \begin{equation*}
%    \eta^-_{\bm\lambda,\bm\gamma}(\bm x)
%    :=\inf_{\xi\in\Xi}\left\{
%      \mathbf1_{\{\xi\le-\hat g(\bm x)\}}
%      +\zeta_{\bm\lambda,\bm\gamma}(\xi)
%    \right\},
%  \end{equation*}
%  which is bounded and Borel measurable by the same argument. Substitution gives the lower dual \eqref{eq-lower-measure-dual}.
%
%  For a primal feasible $\Pi$, the upper Lagrangian is at least $p_{\hat f}(\Pi)$ and the lower Lagrangian is at most $p_{\hat f}(\Pi)$. Every feasible upper dual solution therefore bounds the primal objective from above, and every feasible lower dual solution bounds it from below. Taking the respective optima and using the common nonempty primal feasible set gives \eqref{eq-weak-duality}.
%\end{proof}
%
%The remaining supremum and infimum combine the residual penalty with the failure indicator. For $\Xi=\mathbb R$, the penalty in \eqref{eq-dual-step-function} changes only at observed residuals. Separate these values from the intervening intervals to form $J:=2K+1$ residual cells:
%\begin{equation}
%  \Xi_{2k}:=\left\{\xi_{(k)}^{\mathrm{te}}\right\}
%    \quad\left(k=1,\ldots,K\right),
%  \qquad
%  \Xi_{2k+1}:=\left(\xi_{(k)}^{\mathrm{te}},\xi_{(k+1)}^{\mathrm{te}}\right)
%    \quad\left(k=0,\ldots,K\right).
%  \label{eq-residual-cells}
%\end{equation}
%The input cells in \eqref{eq-input-cells} and the residual cells in \eqref{eq-residual-cells} satisfy
%\begin{equation*}
%  \mathcal X=\bigsqcup_{i=0}^K\mathcal X_i,
%  \qquad
%  \Xi=\bigsqcup_{j=1}^J\Xi_j.
%\end{equation*}
%The singleton cells retain residual atoms and distinguish the open and closed half-lines in the rank constraints. The penalty is constant on each $\Xi_j$, with value
%\begin{equation}
%  \zeta_j:=\sum_{k=1}^K\lambda_k\mathbf1_{\{j<2k\}}
%    -\sum_{k=1}^K\gamma_k\mathbf1_{\{j\le2k\}},
%    \qquad j=1,\ldots,J.
%  \label{eq-dual-cell-values}
%\end{equation}
%For $(\bm x,\xi)\in\mathcal X_i\times\Xi_j$, the failure indicator satisfies
%\begin{equation}
%  \underline h_{ij}:=\mathbf1_{\{j\le2i\}}
%  \le\mathbf1_{\{\xi\le-\hat g(\bm x)\}}
%  \le\overline h_{ij}:=\mathbf1_{\{j\le2i+1\}}.
%  \label{eq-cell-costs}
%\end{equation}
%For inputs in $\mathcal X_i$, residual cells with $j\le2i$ always cause failure, and those with $j\ge2i+2$ are always safe. The overlapping interval $\Xi_{2i+1}$ determines the difference between the two coefficients. When $-\hat g(\bm x)=\xi_{(i)}^{\mathrm{te}}$, that interval is safe; the upper construction assigns its mass to the left endpoint. The proof below verifies this assignment. The weights $w_i$, penalty values in \eqref{eq-dual-cell-values}, and failure coefficients now give finite linear programs for the duals.
%
%\begin{theorem}[Finite-dimensional duals and strong duality]
%  \label{thm-finite-dual-reformulations}
%  For $\Xi=\mathbb R$, $\mathcal Q=\mathcal Q^{\mathrm{DW}}(\theta)$, and every $\theta\ge0$, the dual values in Lemma~\ref{lem-weak-duality} equal the optimal values of
%  \begin{equation}
%    \begin{aligned}
%      \overline d_\theta
%      =\min_{\bm\lambda,\bm\gamma\ge\bm 0,\ \bm z\in\mathbb R^{K+1}}\quad&
%        \bm\lambda^\top\bm u-\bm\gamma^\top\bm l+\sum_{i=0}^K w_i z_i\\
%      \text{s.t.}\quad&
%        z_i+\zeta_j\ge\overline h_{ij},
%        \qquad i=0,\ldots,K,\quad j=1,\ldots,J,
%    \end{aligned}
%    \label{eq-finite-upper-dual}
%  \end{equation}
%  and
%  \begin{equation}
%    \begin{aligned}
%      \underline d_\theta
%      =\max_{\bm\lambda,\bm\gamma\ge\bm 0,\ \bm z\in\mathbb R^{K+1}}\quad&
%        -\bm\lambda^\top\bm u+\bm\gamma^\top\bm l+\sum_{i=0}^K w_i z_i\\
%      \text{s.t.}\quad&
%        z_i-\zeta_j\le\underline h_{ij},
%        \qquad i=0,\ldots,K,\quad j=1,\ldots,J.
%    \end{aligned}
%    \label{eq-finite-lower-dual}
%  \end{equation}
%  Moreover,
%  \begin{equation}
%    \overline p_{\hat f}=\overline d_\theta,\qquad
%    \underline p_{\hat f}=\underline d_\theta.
%    \label{eq-strong-duality}
%  \end{equation}
%  Each LP has $3K+1$ scalar variables and $(K+1)(2K+1)$ pairwise inequalities, together with multiplier nonnegativity. The free variables $z_i$ correspond to the input cell probability equalities. The finite LP optima and the corresponding measure primal and functional dual optima are attained.
%\end{theorem}
%\begin{proof}
%  For a feasible joint law, let $\bm\pi=(\pi_{ij})$ be its matrix of joint cell probabilities,
%  \begin{equation*}
%    \pi_{ij}:=\Pi(\mathcal X_i\times\Xi_j),
%    \qquad i=0,\ldots,K,\quad j=1,\ldots,J.
%  \end{equation*}
%  Its row sums are the fixed input cell probabilities, and its cumulative column sums satisfy the rank constraints. Thus $\bm\pi\in\mathcal C_\theta$, where
%  \begin{equation}
%    \mathcal C_\theta:=
%    \left\{\bm\pi\in\mathbb R_+^{(K+1)\times J}:
%    \begin{array}{ll}
%      \displaystyle\sum_{j=1}^J\pi_{ij}=w_i,
%        &i=0,\ldots,K,\\[1mm]
%      \displaystyle\sum_{i=0}^K\sum_{j<2k}\pi_{ij}\le u_k,
%        &k=1,\ldots,K,\\[1mm]
%      \displaystyle\sum_{i=0}^K\sum_{j\le2k}\pi_{ij}\ge l_k,
%        &k=1,\ldots,K
%    \end{array}\right\}.
%    \label{eq-cell-mass-polytope}
%  \end{equation}
%  By \eqref{eq-cell-costs}, maximizing $\sum_{i,j}\overline h_{ij}\pi_{ij}$ over this polytope bounds the upper measure problem from above, while minimizing $\sum_{i,j}\underline h_{ij}\pi_{ij}$ bounds the lower problem from below.
%
%  To establish the reverse inequalities, fix $\bm\pi\in\mathcal C_\theta$. Draw $\bm x$ from $\mathbb P$ and, on $\mathcal X_i$ with $w_i>0$, select a provisional residual cell $j$ with probability $\pi_{ij}/w_i$. These probabilities sum to one by the row equality. For $j\ne2i+1$, choose a fixed representative of $\Xi_j$; the two coefficients in \eqref{eq-cell-costs} then equal the actual failure indicator. The following assignments in the overlapping cell define the upper and lower joint laws $\Pi^+$ and $\Pi^-$, respectively.
%
%  For the upper problem, assign $\xi=-\hat g(\bm x)$ in the overlapping cell. This causes failure and lies in $\Xi_{2i+1}$ when $-\hat g(\bm x)>\xi_{(i)}^{\mathrm{te}}$. At equality, the assignment moves the provisional mass to the left endpoint. For $i\ge1$, any $\widetilde\xi\in\left(\xi_{(i)}^{\mathrm{te}},\xi_{(i+1)}^{\mathrm{te}}\right)$, and each observed threshold,
%  \begin{equation*}
%    \mathbf1_{\{\xi_{(i)}^{\mathrm{te}}<\xi_{(k)}^{\mathrm{te}}\}}=\mathbf1_{\{\widetilde\xi<\xi_{(k)}^{\mathrm{te}}\}},\qquad
%    \mathbf1_{\{\xi_{(i)}^{\mathrm{te}}\le \xi_{(k)}^{\mathrm{te}}\}}\ge\mathbf1_{\{\widetilde\xi\le \xi_{(k)}^{\mathrm{te}}\}}.
%  \end{equation*}
%  The assignment leaves every open upper probability unchanged and can only increase the closed lower probabilities. It therefore preserves all rank constraints and attains the upper coefficient. For $i=0$, the finite threshold $-\hat g(\bm x)$ is strictly above $\xi_{(0)}^{\mathrm{te}}=-\infty$.
%
%  For the lower problem, assign $\xi=\left(\xi_{(i+1)}^{\mathrm{te}}-\hat g(\bm x)\right)/2$ when $i<K$, and $\xi=-\hat g(\bm x)+1$ when $i=K$. Each value lies in $\Xi_{2i+1}$ and strictly exceeds $-\hat g(\bm x)$, including when the threshold equals $\xi_{(i)}^{\mathrm{te}}$. This attains the lower coefficient and preserves the provisional residual cell.
%
%  On an input cell with $w_i=0$, the row of $\bm\pi$ is zero; set $\xi=\xi_{(1)}^{\mathrm{te}}$ there. All residual assignments are Borel measurable functions of $\bm x$. Their finite mixtures define probability kernels, which, combined with $\mathbb P$, give joint probability laws with the prescribed input marginal. This construction also allows randomization at input atoms. Both joint laws satisfy the rank constraints and attain the assigned matrix objectives:
%  \begin{equation*}
%    p_{\hat f}(\Pi^+)=\sum_{i=0}^K\sum_{j=1}^J\overline h_{ij}\pi_{ij},
%    \qquad
%    p_{\hat f}(\Pi^-)=\sum_{i=0}^K\sum_{j=1}^J\underline h_{ij}\pi_{ij}.
%  \end{equation*}
%  Combining these identities with the initial bounds gives
%  \begin{equation}
%    \overline p_{\hat f}
%    =\max_{\bm\pi\in\mathcal C_\theta}\sum_{i=0}^K\sum_{j=1}^J
%      \overline h_{ij}\pi_{ij},\qquad
%    \underline p_{\hat f}
%    =\min_{\bm\pi\in\mathcal C_\theta}\sum_{i=0}^K\sum_{j=1}^J
%      \underline h_{ij}\pi_{ij}.
%    \label{eq-finite-primal-values}
%  \end{equation}
%  The lower assignment retains every joint cell mass $\pi_{ij}$. The upper assignment can transfer mass from an open interval to its left endpoint while preserving the rank inequalities and the assigned objective value; equality of the resulting cell matrix with $\bm\pi$ is not needed for \eqref{eq-finite-primal-values}.
%
%  For every $\theta\ge0$, a feasible matrix is induced by the empirical product law:
%  \begin{equation*}
%    \pi_{i,2k}^{\mathrm{emp}}:=w_i\frac{c_k-c_{k-1}}{N^{\mathrm{te}}},
%    \qquad i=0,\ldots,K,\quad k=1,\ldots,K,
%  \end{equation*}
%  with all odd columns zero. Its open and closed rank probabilities are $c_{k-1}/N^{\mathrm{te}}$ and $c_k/N^{\mathrm{te}}$, respectively, so \eqref{eq-dw-counts-and-levels} makes it feasible also at $\theta=0$. The set $\mathcal C_\theta$ is closed, and its nonnegative entries sum to one, so it is compact. Both matrix objectives lie in $[0,1]$ and attain their optima, which the preceding constructions realize by joint laws. Finite LP strong duality therefore applies \citep[Section~5.2.3]{B.V04}. The row equalities have free multipliers $z_i$, and the rank inequalities have nonnegative multipliers $\lambda_k,\gamma_k$. Their dual constraints are $z_i+\zeta_j\ge\overline h_{ij}$ and $z_i-\zeta_j\le\underline h_{ij}$, giving \eqref{eq-finite-upper-dual} and \eqref{eq-finite-lower-dual}.
%
%  An upper finite dual solution defines a bounded Borel function $\eta(\bm x)=z_i$ on $\mathcal X_i$. The indicator bounds \eqref{eq-cell-costs} make it feasible for the functional upper dual, so $\eta^+_{\bm\lambda,\bm\gamma}(\bm x)\le z_i$. A lower finite dual solution similarly gives $\eta^-_{\bm\lambda,\bm\gamma}(\bm x)\ge z_i$. The reduced upper objective is therefore no greater than the finite upper dual objective, and the reduced lower objective is no smaller than the finite lower dual objective.
%
%  Finite LP equality and Lemma~\ref{lem-weak-duality} give \eqref{eq-strong-duality}. Optimal finite dual solutions supply optimal rank multipliers and marginal functions for the functional duals. Feasibility and finite optimal values suffice for this LP argument, including at $\theta=0$ when some rank inequalities can be binding for every feasible law.
%\end{proof}
%
%The finite LPs establish exactness of the duals. Their ordered failure coefficients give a further simplification: the best single-threshold candidates derived above are already optimal.
%
%\begin{corollary}[Closed-form bounds and optimal certificates]
%  \label{thm-closed-form-bounds}
%  Under the conditions of Theorem~\ref{thm-finite-dual-reformulations},
%  \begin{align}
%    \overline p_{\hat f}
%    &=\min\left\{1,\min_{k=1,\ldots,K}[u_k+H\left(-\xi_{(k)}^{\mathrm{te}}\right)]\right\},
%    \label{eq-dw-closed-upper}\\
%    \underline p_{\hat f}
%    &=\max\left\{0,\max_{k=1,\ldots,K}[l_k+H\left(-\xi_{(k)}^{\mathrm{te}}\right)-1]\right\}.
%    \label{eq-dw-closed-lower}
%  \end{align}
%  Each dual optimum admits a certificate with at most one nonzero rank multiplier, of value one.
%\end{corollary}
%\begin{proof}
%  Fix nonnegative column masses $\bm y=(y_1,\ldots,y_J)^\top$ with $\sum_j y_j=1$, and write $Y_j:=\sum_{r=1}^j y_r$. Let $\mathcal T(\bm w,\bm y)$ be the set of transportation matrices with row sums $\bm w=(w_0,\ldots,w_K)^\top$ and column sums $\bm y$. It contains the product matrix $\pi_{ij}=w_i y_j$. The rank constraints depend only on $\bm y$, through $Y_{2k-1}\le u_k$ and $Y_{2k}\ge l_k$.
%
%  For fixed $\bm y$, the upper problem maximizes mass on failure pairs $j\le2i+1$; the lower problem maximizes mass on safe pairs $j>2i$ and subtracts it from one. Both are maximum-flow problems: source-to-row capacities are $w_i$, column-to-sink capacities are $y_j$, and a row is connected to each permitted column by an edge of capacity two. Restricting a transportation matrix to permitted pairs gives a feasible flow. Conversely, any flow can be completed to a transportation matrix because the remaining row and column masses are nonnegative and have equal totals. These two operations show that the maximum flow equals the maximum permitted mass.
%
%  A cut of capacity one is available, so a minimum cut crosses no edge of capacity two. Its source-side columns can be taken to be exactly the neighbors of its source-side rows. Because the neighborhoods are nested, adding any row whose neighbors are already included cannot increase the cut capacity. Thus, apart from empty or full row sets, upper-problem cuts reduce to the rows $i<k$ and columns $j<2k$, while safe-mass cuts reduce to the rows $i\ge k$ and columns $j>2k$. For $k=1,\ldots,K$, their capacities are, respectively,
%  \begin{align*}
%    \sum_{i=k}^K w_i+Y_{2k-1}
%      &=H\left(-\xi_{(k)}^{\mathrm{te}}\right)+Y_{2k-1},\\
%    \sum_{i=0}^{k-1}w_i+1-Y_{2k}
%      &=2-H\left(-\xi_{(k)}^{\mathrm{te}}\right)-Y_{2k}.
%  \end{align*}
%  Empty and full row sets give capacity one. Each remaining cut therefore tests a single separation between ordered cells. The max-flow/min-cut theorem \citep{F.F56} gives
%  \begin{align*}
%    \max_{\bm\pi\in\mathcal T(\bm w,\bm y)}\sum_{i,j}\overline h_{ij}\pi_{ij}
%      &=\min\left\{1,\min_k\left[H\left(-\xi_{(k)}^{\mathrm{te}}\right)+Y_{2k-1}\right]\right\},\\
%    \min_{\bm\pi\in\mathcal T(\bm w,\bm y)}\sum_{i,j}\underline h_{ij}\pi_{ij}
%      &=\max\left\{0,\max_k\left[H\left(-\xi_{(k)}^{\mathrm{te}}\right)+Y_{2k}-1\right]\right\}.
%  \end{align*}
%
%  The first value is nondecreasing in every open prefix $Y_{2k-1}$, and the second is nondecreasing in every closed prefix $Y_{2k}$. It remains to attain all upper prefix limits simultaneously for the maximization, and all lower prefix limits simultaneously for the minimization. With $u_{K+1}=1$ and $l_0=0$, choose
%  \begin{align*}
%    &y_1^+:=u_1,\qquad y_{2k}^+:=u_{k+1}-u_k,\qquad
%      y_{2k+1}^+:=0,\quad k=1,\ldots,K,\\
%    &y_{2k-1}^-:=0,\qquad y_{2k}^-:=l_k-l_{k-1},\quad k=1,\ldots,K,
%      \qquad y_J^-:=1-l_K.
%  \end{align*}
%  The monotonicity and compatibility conditions in Proposition~\ref{prop-finite-rank-constraints} make both vectors nonnegative with total mass one. Their prefixes satisfy
%  \[
%    Y_{2k-1}^+=u_k,\qquad Y_{2k}^+=u_{k+1}\ge l_k,
%    \qquad Y_{2k}^-=l_k,\qquad Y_{2k-1}^-=l_{k-1}\le u_k.
%  \]
%  Hence both allocations are feasible and attain the required prefix limits. Substituting them into the fixed-column values and using \eqref{eq-finite-primal-values} proves \eqref{eq-dw-closed-upper} and \eqref{eq-dw-closed-lower}. Zero masses and coincident prefix limits are allowed, so this argument includes $\theta=0$.
%
%  Finally, let $\bm e_k$ be the $k$th unit vector. The upper and lower dual certificates are, respectively,
%  \begin{align*}
%    &\bm\lambda=\bm e_k,\qquad \bm\gamma=\bm 0,\qquad
%      \eta^+_{\bm e_k,\bm 0}(\bm x)
%      =\mathbf1_{\{\hat g(\bm x)\le-\xi_{(k)}^{\mathrm{te}}\}},\\
%    &\bm\lambda=\bm 0,\qquad \bm\gamma=\bm e_k,\qquad
%      \eta^-_{\bm 0,\bm e_k}(\bm x)
%      =\mathbf1_{\{\hat g(\bm x)\le-\xi_{(k)}^{\mathrm{te}}\}}-1.
%  \end{align*}
%  The threshold indicator inequalities established before Lemma~\ref{lem-weak-duality} give feasibility. To verify the envelope values, take $\xi=\xi_{(k)}^{\mathrm{te}}$ in the upper problem when $-\hat g(\bm x)\ge\xi_{(k)}^{\mathrm{te}}$, and $\xi\in\left(-\hat g(\bm x),\xi_{(k)}^{\mathrm{te}}\right]$ in the lower problem when $-\hat g(\bm x)<\xi_{(k)}^{\mathrm{te}}$. These choices attain one and minus one, respectively; a residual above both thresholds attains zero in the complementary cases. At threshold equality, the upper and lower envelopes are therefore one and zero, and the input is included in $H\left(-\xi_{(k)}^{\mathrm{te}}\right)$. The certificate objectives are exactly the inner candidates in \eqref{eq-dw-closed-upper} and \eqref{eq-dw-closed-lower}. Zero rank multipliers give the constant marginal functions one and zero, attaining the outer candidates. Selecting the best candidate thus supplies an optimal full dual solution with at most one nonzero rank multiplier.
%\end{proof}
%
%The upper bound combines residual mass below an observed value with the probability that the surrogate threshold reaches that value. The lower bound measures overlap that every admissible dependence structure must retain.
%
%\begin{remark}[Known interval support]
%  \label{rem-interval-support}
%  Suppose $\Xi=[a,b]$ with $-\infty<a\le b<\infty$ and the observed residuals lie in this interval. Then
%  \begin{align*}
%    \overline p_{\hat f}
%    &=\min\left\{H(-a),\min_{k=1,\ldots,K}[u_k+H\left(-\xi_{(k)}^{\mathrm{te}}\right)]\right\},\\
%    \underline p_{\hat f}
%    &=\max\left\{H(-b),\max_{k=1,\ldots,K}[l_k+H\left(-\xi_{(k)}^{\mathrm{te}}\right)-1]\right\}.
%  \end{align*}
%  To obtain these expressions, represent the support restriction on $\mathbb R$ by
%  $\mathbb Q\{\xi<a\}\le0$ and $\mathbb Q\{\xi\le b\}\ge1$.
%  Add $a,b$ to the ordered threshold list. At a newly added $a$, use upper and lower bounds zero; at a newly added $b$, use upper and lower bounds one. At coincident observed thresholds, tighten only the upper bound at $a$ and the lower bound at $b$. The enlarged sequences satisfy the compatibility conditions in Proposition~\ref{prop-finite-rank-constraints}, so the same finite-aggregation and ordered-cut arguments apply. The effective additional upper candidate is $H(-a)$, and the effective additional lower candidate is $H(-b)$; the other added candidates are redundant. If $a=b$, both bounds equal $H(-a)$. For the resulting partition on $\Xi=[a,b]$, intersect the singleton and open-interval residual cells from the augmented threshold list with $\Xi$ and omit empty cells. Including the support endpoints in this list retains their atoms and the corresponding changes in failure status.
%\end{remark}
%
%For fixed $\theta$, evaluation requires scalar DW inversion, sorting and counting the test residuals, and computing the $K$ values $H\left(-\xi_{(k)}^{\mathrm{te}}\right)$. The endpoint table depends only on $N^{\mathrm{te}}$ and $\theta$ and can be precomputed. Given this table and the surrogate CDF values, the final upper and lower scans cost $O(K)$.
%
%For $\theta\ge\theta_{N^{\mathrm{te}}}(\delta)$, these exact bounds inherit the finite-sample guarantee of Theorem~\ref{theorem-01}. If the surrogate CDF values are estimated by simulation, simultaneous upper CDF bounds can replace $H\left(-\xi_{(k)}^{\mathrm{te}}\right)$ in the upper expression, and simultaneous lower CDF bounds can replace $H\left(-\xi_{(k)}^{\mathrm{te}}\right)$ in the lower expression. An independent auxiliary simulation can provide these simultaneous bounds conditionally on the test data. Its error probability must be included in the overall confidence budget.
%
%
%\bibliographystyle{elsarticle-harv}
%\bibliography{reference}
%
% END PRESERVED ACTIVE MANUSCRIPT BEFORE INFORMATION-LIMIT EXTENSION

% BEGIN PRESERVED SECTIONS 2--4 BEFORE NARRATIVE REVISION (2026-11-24)
%\section{Problem Description}
%\label{sec-problem}
%
%Let $\bm X\sim\mathbb P$ be a random input taking values in $\mathcal X\subseteq\mathbb R^n$, and let $g:\mathcal X\to\mathbb R$ be a measurable limit-state function. The failure probability is defined as
%\[
%  p_f:=\Pr\{g(\bm X)\le0\}=\mathbb E_{\bm X\sim\mathbb P}\left[\mathbf 1_{\{g(\bm X)\le0\}}\right]=\int_{\mathcal X}\mathbf 1_{\{g(\bm x)\le0\}}\,\mathbb P(\mathrm d\bm x).
%\]
%For fixed $g$, this functional is linear in $\mathbb P$.
%
%Evaluating $g$ can require an expensive high-fidelity simulation. We use the training dataset
%\[
%  \mathcal D^{\mathrm{tr}}:=\left\{\left(\bm x_i^{\mathrm{tr}},g(\bm x_i^{\mathrm{tr}})\right)\right\}_{i=1}^{N^{\mathrm{tr}}}
%\]
%to fit a computationally efficient, measurable surrogate $\hat g:=\hat g_{\mathcal D^{\mathrm{tr}}}$. Model fitting and selection use only the training data and are completed before the test data are used. We condition on $\mathcal D^{\mathrm{tr}}$, treat $\hat g$ as fixed, and suppress this dependence in the notation.
%
%Define the signed residual function
%\[
%  \xi(\bm x):=g(\bm x)-\hat g(\bm x),\qquad \forall \bm x\in\mathcal X.
%\]
%For a new input $\bm X\sim\mathbb P$ independent of the training data, let $\mathbb Q^\star$ be the distribution of $\xi^\star:=\xi(\bm X)$. Let $\Xi\subseteq\mathbb R$ be a known Borel set with $\mathbb Q^\star(\Xi)=1$; the default choice is $\Xi=\mathbb R$. For any Borel set $\mathcal S$ in a Euclidean space, write $\mathcal B(\mathcal S)$ for its Borel sigma-algebra. Let $\mathcal M_+(\mathcal S)$ denote the finite nonnegative Borel measures concentrated on $\mathcal S$, and let $\mathcal M_1(\mathcal S)$ denote the probability measures.
%
%For a candidate residual law $\mathbb Q\in\mathcal M_1(\Xi)$, let $\Pi$ be a joint distribution on $\mathcal X\times\Xi$ with marginals $\Pi_{\bm X}=\mathbb P$ and $\Pi_\xi=\mathbb Q$. Its surrogate-based failure probability is
%\begin{equation}
%  p_{\hat f}(\Pi):=\Pr\{\hat g(\bm X)+\xi\le0\}=\int_{\mathcal X\times\Xi}\mathbf 1_{\{\hat g(\bm x)+\xi\le0\}}\,\Pi(\mathrm d\bm x,\mathrm d\xi).
%  \label{eq-surrogate-failure-probability}
%\end{equation}
%For fixed $\hat g$, this functional is linear in $\Pi$. In particular, the true joint distribution $\Pi^\star$ satisfies $p_{\hat f}(\Pi^\star)=p_f$.
%
%The failure probability depends on the residual law and its dependence on the input. Near the surrogate failure threshold, even a small residual can change the failure classification. We represent uncertainty in the residual law by an ambiguity set $\mathcal Q$ calibrated to contain $\mathbb Q^\star$ with high confidence. For a given $\mathcal Q$, define the joint ambiguity set
%\[
%  \mathcal P(\mathcal Q) :=\left\{\Pi\in\mathcal M_1(\mathcal X\times\Xi):\Pi_{\bm X}=\mathbb P,\ \Pi_\xi\in\mathcal Q\right\}.
%\]
%The joint ambiguity set allows all dependence structures consistent with the specified marginals. The upper and lower failure probability bounds are
%\begin{equation}
%  \overline p_{\hat f}:=\sup_{\Pi\in\mathcal P(\mathcal Q)}p_{\hat f}(\Pi),
%  \label{eq-upper-bound}
%\end{equation}
%and
%\begin{equation}
%  \underline p_{\hat f}:=\inf_{\Pi\in\mathcal P(\mathcal Q)}p_{\hat f}(\Pi),
%  \label{eq-lower-bound}
%\end{equation}
%respectively. If $\mathbb Q^\star\in\mathcal Q$, then $\Pi^\star\in\mathcal P(\mathcal Q)$ and hence $\underline p_{\hat f}\le p_f\le\overline p_{\hat f}$.
%
%The true residual is a deterministic function of the input. The joint ambiguity set enlarges its unknown graph to all joint laws satisfying the marginal restrictions, including randomized conditional residual laws. This relaxation provides bounds from residual-marginal information without specifying an input--residual dependence model. Paired observations can refine the information model through additional joint constraints.
%
%The analysis has two tasks:
%\begin{enumerate}
%  \item \textbf{Ambiguity set design and statistical calibration.} Construct a residual ambiguity set and establish its coverage guarantee under explicit sampling conditions.
%  \item \textbf{Tractable reformulations.} Derive exact and computationally tractable reformulations of the infinite-dimensional optimization problems defining the failure probability bounds.
%\end{enumerate}
%
%
%\section{Ambiguity Set Design}
%\label{sec-ambiguity-set}
%
%Failure occurs when $\xi\le-\hat g(\bm X)$, so the residual cumulative distribution function (CDF) gives the marginal mass at or below each failure threshold. Accurate tail control is needed for small failure probabilities. We use the D\"umbgen--Wellner (DW) construction \citep{L.J23} to control empirical CDF deviations simultaneously over all thresholds, with tolerances that vary across probability levels.
%
%Let $q$ and $p$ denote an empirical probability and a candidate probability. Following \citet{B.J79} and the CDF confidence bands of \citet{O95}, we measure their local discrepancy by the Bernoulli Kullback--Leibler divergence
%\begin{equation}
%  \phi(q,p):=q\log\frac{q}{p}+(1-q)\log\frac{1-q}{1-p},\qquad \forall q\in[0,1],\, p\in(0,1),
%  \label{eq-bernoulli-kl}
%\end{equation}
%where logarithms are natural and a term with zero numerator is interpreted as zero. For $q$ close to an interior $p$, the approximation $\phi(q,p)\approx(q-p)^2/[2p(1-p)]$ scales the deviation by the Bernoulli variance. Thus, the same absolute deviation receives a larger weight near the probability boundaries. The full divergence also covers zero and full empirical counts through $\phi(0,p)=-\log(1-p)$ and $\phi(1,p)=-\log p$. At the boundaries of the candidate probability, we set $\phi(0,0)=\phi(1,1)=0$, $\phi(q,0)=+\infty$ for $q>0$, and $\phi(q,1)=+\infty$ for $q<1$.
%
%To compare deviations simultaneously across probability levels, we use the additive correction of \citet{L.J23}. For $v\in(0,1)$, define
%\begin{equation*}
%  \rho(v):=\log\log\frac{e}{4v(1-v)},\qquad \omega(v):=\log\!\left(1+\rho(v)^2\right).
%\end{equation*}
%Before sampling, we fix the divergence parameter $s=1$, which gives \eqref{eq-bernoulli-kl}, and the correction parameter $\nu=1$. Set $\psi_0(v):=\rho(v)+\omega(v)$, with $\psi_0(0)=\psi_0(1)=+\infty$, and define the bivariate correction
%\begin{equation*}
%  \psi(q,p):=\min_{v\in[\min\{q,p\},\,\max\{q,p\}]}\psi_0(v).
%\end{equation*}
%The functions $\rho$ and $\omega$ vanish at $v=1/2$ and increase toward the probability boundaries. The minimum over the interval between $q$ and $p$ is therefore large only when both probabilities lie in the same tail. This correction balances fluctuations across probability scales in the simultaneous CDF constraint.
%
%We construct the empirical CDF from a test dataset
%\[
%  \mathcal D^{\mathrm{te}}:=\left\{\left(\bm x_j^{\mathrm{te}},g(\bm x_j^{\mathrm{te}})\right)\right\}_{j=1}^{N^{\mathrm{te}}}.
%\]
%The test inputs are independent draws from $\mathbb P$ and are independent of $\mathcal D^{\mathrm{tr}}$. We fix $N^{\mathrm{te}}\ge1$ before sampling and use the test data to calibrate the residual ambiguity set, with the fitted surrogate held fixed.
%
%Define the signed test residuals
%\[
%  \xi_j^{\mathrm{te}}:=\xi(\bm x_j^{\mathrm{te}})=g(\bm x_j^{\mathrm{te}})-\hat g(\bm x_j^{\mathrm{te}}),\qquad j=1,\ldots,N^{\mathrm{te}}.
%\]
%Conditional on $\mathcal D^{\mathrm{tr}}$, these residuals are independent and identically distributed according to $\mathbb Q^\star$. Their empirical CDF is
%\begin{equation}
%  \widehat F^{\mathrm{te}}(t):=\frac{1}{N^{\mathrm{te}}}\sum_{j=1}^{N^{\mathrm{te}}}\mathbf 1_{\{\xi_j^{\mathrm{te}}\le t\}},\qquad \forall t\in\mathbb R.
%  \label{eq-empirical-cdf}
%\end{equation}
%For $\mathbb Q\in\mathcal M_1(\Xi)$, write $F_{\mathbb Q}(t):=\mathbb Q\left(({-\infty},t]\right)$. At a threshold $t$, $q=\widehat F^{\mathrm{te}}(t)$ is the mean of the Bernoulli indicators in \eqref{eq-empirical-cdf}, and $p=F_{\mathbb Q}(t)$ is the candidate probability. Thus $N^{\mathrm{te}}\phi(q,p)$ is the Bernoulli log-likelihood difference between the empirical and candidate probabilities. Combining it with the multiscale correction gives $N^{\mathrm{te}}\phi(q,p)-\psi(q,p)$.
%
%At the probability boundaries, the corrected discrepancy is $-\infty$ for $q=p\in\{0,1\}$ and $+\infty$ for $p=0<q$ or $p=1>q$. These conventions extend the constraint to every residual threshold.
%
%\begin{definition}[DW-based CDF ambiguity set]
%  \label{definition-dw-ambiguity-set}
%  Given the empirical CDF $\widehat F^{\mathrm{te}}$ and a threshold $\theta\ge0$, define
%  \begin{equation} 
%    \label{eq-cdf-ambiguity-set}
%      \mathcal Q^{\mathrm{DW}}(\theta):=\left\{\mathbb Q\in\mathcal M_1(\Xi):\ 
%      \sup_{t\in\mathbb R}\left[ N^{\mathrm{te}}\phi\!\left(\widehat F^{\mathrm{te}}(t),F_{\mathbb Q}(t)\right) 
%       {}-\psi\!\left(\widehat F^{\mathrm{te}}(t),F_{\mathbb Q}(t)\right)\right]\le\theta\right\}. 
%  \end{equation} 
%\end{definition}
%
%We calibrate $\theta$ against a distribution-free reference law. Following \citet{L.J23}, let $V_1,\ldots,V_{N^{\mathrm{te}}}$ be independent standard uniform variables, independent of the training and test data. Their empirical CDF is
%\begin{equation*}
%  G_{N^{\mathrm{te}}}(u):=\frac{1}{N^{\mathrm{te}}}\sum_{j=1}^{N^{\mathrm{te}}}\mathbf 1_{\{V_j\le u\}},\qquad \forall u\in[0,1].
%\end{equation*}
%Define the uniform reference statistic
%\begin{equation}
%  T_{N^{\mathrm{te}}}^{\mathrm U}:=\sup_{u\in(0,1)}\left[N^{\mathrm{te}}\phi\!\left(G_{N^{\mathrm{te}}}(u),u\right)-\psi\!\left(G_{N^{\mathrm{te}}}(u),u\right)\right].
%  \label{eq-dw-reference-statistic}
%\end{equation}
%
%A quantile of this reference statistic determines the calibration threshold.
%
%\begin{proposition}[Finite-sample threshold calibration]
%  \label{proposition-01}
%  Under the sampling conditions stated above, let $\delta\in(0,1)$ and define
%  \begin{equation}
%    \theta_{N^{\mathrm{te}}}(\delta):=\inf\left\{\tau\ge0:\Pr\left\{T_{N^{\mathrm{te}}}^{\mathrm U}\le\tau\right\}\ge1-\delta\right\}.
%    \label{eq-dw-calibrated-threshold}
%  \end{equation}
%
%  Then, for every $\theta\ge\theta_{N^{\mathrm{te}}}(\delta)$, the DW-based CDF ambiguity set satisfies
%  \begin{equation}
%    \Pr\left\{\mathbb Q^\star\in\mathcal Q^{\mathrm{DW}}(\theta)\,\middle|\,\mathcal D^{\mathrm{tr}}\right\}\ge1-\delta.
%    \label{eq-cdf-coverage}
%  \end{equation}
%\end{proposition}
%\begin{proof}
%  The reference statistic is finite almost surely, and its distribution function is right-continuous. Hence \eqref{eq-dw-calibrated-threshold} gives
%  \[
%    \Pr\left\{T_{N^{\mathrm{te}}}^{\mathrm U}>\theta_{N^{\mathrm{te}}}(\delta)\right\}\le\delta.
%  \]
%  Conditioning on $\mathcal D^{\mathrm{tr}}$ fixes $\hat g$ and $\mathbb Q^\star$. Define the generalized quantile function $F_{\mathbb Q^\star}^{-1}(u):=\inf\{t\in\mathbb R:F_{\mathbb Q^\star}(t)\ge u\}$ for all $u\in(0,1)$. The sample $F_{\mathbb Q^\star}^{-1}(V_j)$ for $j=1,\ldots,N^{\mathrm{te}}$, has the same conditional joint law as the test residuals. Thus, given $\mathcal D^{\mathrm{tr}}$,
%  \[
%    \left(\widehat F^{\mathrm{te}}(t)\right)_{t\in\mathbb R}\overset{\mathrm d}{=}\left(G_{N^{\mathrm{te}}}\!\left(F_{\mathbb Q^\star}(t)\right)\right)_{t\in\mathbb R}.
%  \]
%  For $0<F_{\mathbb Q^\star}(t)<1$, the discrepancy at $\mathbb Q^\star$ is evaluated at a probability level included in \eqref{eq-dw-reference-statistic}. When $F_{\mathbb Q^\star}(t)$ is zero or one, the empirical CDF equals that boundary almost surely, giving discrepancy $-\infty$. The supremum in \eqref{eq-cdf-ambiguity-set}, evaluated at $\mathbb Q^\star$, is therefore bounded by $T_{N^{\mathrm{te}}}^{\mathrm U}$ under this coupling. Since $\mathbb Q^\star(\Xi)=1$, we obtain
%  \[
%    \Pr\left\{\mathbb Q^\star\notin\mathcal Q^{\mathrm{DW}}\!\left(\theta_{N^{\mathrm{te}}}(\delta)\right)\,\middle|\,\mathcal D^{\mathrm{tr}}\right\}
%    \le\Pr\left\{T_{N^{\mathrm{te}}}^{\mathrm U}>\theta_{N^{\mathrm{te}}}(\delta)\right\}\le\delta.
%  \]
%  For every test sample, $\theta\ge\theta_{N^{\mathrm{te}}}(\delta)$ implies $\mathcal Q^{\mathrm{DW}}\!\left(\theta_{N^{\mathrm{te}}}(\delta)\right)\subseteq\mathcal Q^{\mathrm{DW}}(\theta)$. This inclusion proves \eqref{eq-cdf-coverage} for every such $\theta$.
%\end{proof}
%
%Proposition~\ref{proposition-01} applies to arbitrary residual distributions, including laws with atoms or infinite moments. With $\phi$ and $\psi$ fixed, the reference law depends only on $N^{\mathrm{te}}$. The threshold $\theta_{N^{\mathrm{te}}}(\delta)$ can therefore be calibrated offline for each confidence level $1-\delta$. \citet{L.J23} describe finite-sample evaluation using uniform order-statistic boundary probabilities. Any conservative upper bound on the quantile in \eqref{eq-dw-calibrated-threshold} preserves \eqref{eq-cdf-coverage}.
%% If calibration is performed by simulation, quantile estimation error must be accounted for in the confidence guarantee. A threshold chosen from the test data also preserves \eqref{eq-cdf-coverage}, provided $\theta\ge\theta_{N^{\mathrm{te}}}(\delta)$ almost surely, by the same nesting argument.
%
%Residual-law coverage gives simultaneous bounds on the failure probability under the joint model.
%
%\begin{theorem}[Finite-sample failure probability bounds]
%  \label{theorem-01}
%  Under the sampling conditions of Proposition~\ref{proposition-01}, let $\delta\in(0,1)$ and $\theta\ge\theta_{N^{\mathrm{te}}}(\delta)$, with the calibrated threshold defined in \eqref{eq-dw-calibrated-threshold}. Then the bounds defined in \eqref{eq-upper-bound} and \eqref{eq-lower-bound}, evaluated with $\mathcal Q=\mathcal Q^{\mathrm{DW}}(\theta)$, satisfy
%  \begin{equation}
%    \Pr\left\{\underline p_{\hat f}\le p_f\le\overline p_{\hat f}\,\middle|\,\mathcal D^{\mathrm{tr}}\right\}\ge1-\delta.
%    \label{eq-failure-probability-coverage}
%  \end{equation}
%\end{theorem}
%
%\begin{proof}
%  On the event $\mathbb Q^\star\in\mathcal Q^{\mathrm{DW}}(\theta)$, the true joint distribution $\Pi^\star$ belongs to $\mathcal P\left(\mathcal Q^{\mathrm{DW}}(\theta)\right)$. The definitions of the bounds and the identity $p_{\hat f}(\Pi^\star)=p_f$ give
%  \[
%    \underline p_{\hat f}
%    \le p_{\hat f}(\Pi^\star)
%    =p_f
%    \le\overline p_{\hat f}.
%  \]
%  Consequently,
%  \[
%    \Pr\left\{
%      \underline p_{\hat f}\le p_f\le\overline p_{\hat f}
%      \,\middle|\,\mathcal D^{\mathrm{tr}}
%    \right\}
%    \ge
%    \Pr\left\{
%      \mathbb Q^\star\in\mathcal Q^{\mathrm{DW}}(\theta)
%      \,\middle|\,\mathcal D^{\mathrm{tr}}
%    \right\}.
%  \]
%  Proposition~\ref{proposition-01} bounds the right-hand side below by $1-\delta$, proving \eqref{eq-failure-probability-coverage}.
%\end{proof}
%
%
%The event $\mathbb Q^\star\in\mathcal Q^{\mathrm{DW}}(\theta)$ gives both inequalities, so $1-\delta$ is the confidence level for the entire interval. The threshold $\theta$ controls the corrected likelihood discrepancy; the resulting CDF tolerances vary with the probability level and can be asymmetric about the empirical CDF. For fixed test data, $0\le\theta_1\le\theta_2$ implies $\mathcal Q^{\mathrm{DW}}(\theta_1)\subseteq\mathcal Q^{\mathrm{DW}}(\theta_2)$. The lower failure bound is therefore nonincreasing in $\theta$, and the upper bound is nondecreasing. Increasing the prescribed confidence level gives a nondecreasing calibrated threshold and hence a wider or unchanged interval.
%
%The interval accounts for finite-sample uncertainty in the residual law and uncertainty in its dependence on the input. Because the joint ambiguity set includes all couplings with the specified marginals, a gap can remain even when the residual marginal is known exactly. The same coverage argument applies to any residual ambiguity set containing $\mathbb Q^\star$ with conditional probability at least $1-\delta$ given $\mathcal D^{\mathrm{tr}}$.
%
%
%\section{Tractable Reformulation}
%\label{sec-reformulation}
%
%Scalar inversion turns the DW constraint into linear restrictions on residual probabilities. These restrictions lead to Lagrangian duals, exact finite linear programs, and closed-form failure probability bounds. Throughout this section, the training data, the test residuals, and $\theta\ge0$ are fixed. The main results use $\Xi=\mathbb R$; the final remark treats known interval support.
%
%For each empirical probability level $q_j:=j/N^{\mathrm{te}}$, define the scalar acceptance interval
%\begin{equation}
%  [\alpha_j(\theta),\beta_j(\theta)]
%  :=\left\{p\in[0,1]:
%  N^{\mathrm{te}}\phi(q_j,p)-\psi(q_j,p)\le\theta\right\},
%  \qquad j=0,\ldots,N^{\mathrm{te}}.
%  \label{eq-dw-rank-acceptance}
%\end{equation}
%Let $\xi_{(1)}^{\mathrm{te}}<\cdots<\xi_{(K)}^{\mathrm{te}}$ be the distinct ordered test residuals. Define
%\begin{equation} 
%    c_0 :=0, \quad
%    c_k :=\sum_{j=1}^{N^{\mathrm{te}}}\mathbf1_{\{\xi_j^{\mathrm{te}}\le\xi_{(k)}^{\mathrm{te}}\}}, \quad
%    l_k :=\alpha_{c_k}(\theta),  \quad
%    u_k :=\beta_{c_{k-1}}(\theta),\qquad k=1,\ldots,K.
%  \label{eq-dw-counts-and-levels}
%\end{equation}
%The cumulative counts retain all multiplicities, so $c_K=N^{\mathrm{te}}$ even when test residuals coincide.
%
%\begin{proposition}[CDF band and finite rank constraints]
%  \label{prop-finite-rank-constraints}
%  Each acceptance set in \eqref{eq-dw-rank-acceptance} is a nonempty closed interval containing $q_j$. Define
%  \[
%    L_\theta(t):=\alpha_{N^{\mathrm{te}}\widehat F^{\mathrm{te}}(t)}(\theta),
%    \qquad
%    U_\theta(t):=\beta_{N^{\mathrm{te}}\widehat F^{\mathrm{te}}(t)}(\theta).
%  \]
%  The DW ambiguity set has the equivalent representations
%  \begin{equation*}
%    \mathcal Q^{\mathrm{DW}}(\theta) =\left\{\mathbb Q\in\mathcal M_1(\Xi): L_\theta(t)\le F_{\mathbb Q}(t)\le U_\theta(t), \, \forall t\in\mathbb R\right\}
%  \end{equation*}
%  and
%  \begin{equation}
%    \mathcal Q^{\mathrm{DW}}(\theta)
%    =\left\{\mathbb Q\in\mathcal M_1(\Xi): \mathbb Q\{\xi<\xi_{(k)}^{\mathrm{te}}\}\le u_k, \, \mathbb Q\{\xi\le \xi_{(k)}^{\mathrm{te}}\}\ge l_k, \, k=1,\ldots,K \right\}.
%    \label{eq-finite-rank-constraints}
%  \end{equation}
%  Both endpoint sequences are nondecreasing in the empirical rank. Consequently, $l_k$ and $u_k$ are nondecreasing and satisfy
%  \begin{equation*}
%    0\le l_k,u_k\le1,\quad
%    l_{k-1}\le u_k,\quad
%    l_k\le u_{k+1},
%    \qquad k=1,\ldots,K,
%  \end{equation*}
%  where $l_0:=0$ and $u_{K+1}:=1$.
%\end{proposition}
%\begin{proof}
%  For fixed $q$, the Bernoulli divergence increases as $p$ moves away from $q$, while $\psi(q,p)$ cannot increase as the interval between $q$ and $p$ expands. The corrected discrepancy therefore decreases toward $q$ from below and increases above $q$. It is continuous for $p\in(0,1)$ and equals $-\psi_0(q)\le0$ at $p=q$. The boundary conventions in Section~\ref{sec-ambiguity-set} give a nonempty closed acceptance interval also for $q\in\{0,1\}$. Scalar inversion at each empirical CDF level yields the stated CDF band, following \citet[Supplement, Section~S.6]{L.J23}.
%
%  For $k=1,\ldots,K-1$, the empirical CDF on $\left[\xi_{(k)}^{\mathrm{te}},\xi_{(k+1)}^{\mathrm{te}}\right)$ equals $c_k/N^{\mathrm{te}}$. Monotonicity of a candidate CDF makes its lower band constraint equivalent to $F_{\mathbb Q}\left(\xi_{(k)}^{\mathrm{te}}\right)\ge\alpha_{c_k}(\theta)$ and its upper band constraint equivalent to $F_{\mathbb Q}\left(\xi_{(k+1)}^{\mathrm{te}}-\right)\le\beta_{c_k}(\theta)$. The left exterior interval contributes $F_{\mathbb Q}\left(\xi_{(1)}^{\mathrm{te}}-\right)\le\beta_0(\theta)$, and the right exterior interval contributes $F_{\mathbb Q}\left(\xi_{(K)}^{\mathrm{te}}\right)\ge\alpha_{N^{\mathrm{te}}}(\theta)$. The other exterior constraints are automatic because $\alpha_0(\theta)=0$ and $\beta_{N^{\mathrm{te}}}(\theta)=1$. Using $F_{\mathbb Q}\left(\xi_{(k)}^{\mathrm{te}}-\right)=\mathbb Q\{\xi<\xi_{(k)}^{\mathrm{te}}\}$ and $F_{\mathbb Q}\left(\xi_{(k)}^{\mathrm{te}}\right)=\mathbb Q\{\xi\le \xi_{(k)}^{\mathrm{te}}\}$ proves \eqref{eq-finite-rank-constraints}.
%
%  With $p$ fixed, the corrected discrepancy decreases as $q$ approaches $p$ from either side. For empirical levels $q<q'$, an accepted value $p<q$ at level $q'$ remains accepted at level $q$; an accepted value $p>q'$ at level $q$ remains accepted at level $q'$. Since each acceptance interval contains its empirical level, these implications make both endpoints nondecreasing. The inequalities $\alpha_{c_k}(\theta)\le c_k/N^{\mathrm{te}}\le\beta_{c_k}(\theta)$ then give the stated compatibility conditions. This argument permits singleton intervals: if $q_j=1/2$ and $\theta=0$, the acceptance interval is $\{1/2\}$ because $\psi(1/2,p)=0$.
%\end{proof}
%
%The upper rank constraints limit mass strictly below each observed residual, while the lower constraints guarantee mass at or below it. The two half-lines retain possible atoms at the observed values. The scalar endpoints can be found by one-dimensional bisection, after which the ambiguity set is described by linear constraints on the residual law.
%
%Define the CDF of the surrogate response by $H(t):=\mathbb P\left\{\hat g(\bm X)\le t\right\}$ for $t\in\mathbb R$. The observed residuals classify inputs by the location of the failure threshold $-\hat g(\bm x)$. Set $\xi_{(0)}^{\mathrm{te}}:=-\infty$ and $\xi_{(K+1)}^{\mathrm{te}}:=+\infty$, with $H(-\infty):=0$ and $H(+\infty):=1$. For $i=0,\ldots,K$, the input cells and their probabilities are
%\begin{equation}
%  \mathcal X_i:=\left\{\bm x\in\mathcal X:
%    \xi_{(i)}^{\mathrm{te}}\le-\hat g(\bm x)<\xi_{(i+1)}^{\mathrm{te}}\right\},
%  \qquad
%  w_i:=\mathbb P(\mathcal X_i)
%    =H\left(-\xi_{(i)}^{\mathrm{te}}\right)-H\left(-\xi_{(i+1)}^{\mathrm{te}}\right).
%  \label{eq-input-cells}
%\end{equation}
%These $K+1$ cells partition $\mathcal X$, and each threshold event is a union of consecutive cells:
%\[
%  \mathbf1_{\{\hat g(\bm x)\le-\xi_{(k)}^{\mathrm{te}}\}}
%    =\sum_{i=k}^K\mathbf1_{\mathcal X_i}(\bm x),
%  \qquad
%  H\left(-\xi_{(k)}^{\mathrm{te}}\right)=\sum_{i=k}^K w_i,
%  \qquad k=1,\ldots,K.
%\]
%The half-open convention assigns an input with $-\hat g(\bm x)=\xi_{(k)}^{\mathrm{te}}$ to $\mathcal X_k$, preserving threshold equalities even when $H$ has atoms.
%
%For every rank $k$, the failure indicator satisfies 
%\[
%  \mathbf1_{\{\xi\le\xi_{(k)}^{\mathrm{te}}\}}
%    +\mathbf1_{\{\hat g(\bm x)\le-\xi_{(k)}^{\mathrm{te}}\}}-1
%  \le\mathbf1_{\{\xi\le-\hat g(\bm x)\}}
%  \le\mathbf1_{\{\xi<\xi_{(k)}^{\mathrm{te}}\}}
%    +\mathbf1_{\{\hat g(\bm x)\le-\xi_{(k)}^{\mathrm{te}}\}}.
%\]
%The upper inequality covers failure by a union of two events. The lower inequality guarantees failure on their intersection. Integrating under any feasible joint law gives
%\begin{equation*}
%  l_k+H\left(-\xi_{(k)}^{\mathrm{te}}\right)-1\le p_{\hat f}(\Pi)\le u_k+H\left(-\xi_{(k)}^{\mathrm{te}}\right),
%  \qquad k=1,\ldots,K.
%\end{equation*}
%These candidate bounds hold for every admissible dependence structure. The duality argument below establishes that the best candidate in each direction is exact.
%
%Substituting \eqref{eq-finite-rank-constraints} into the upper failure probability problem gives the following linear program over $\mathcal M_+(\mathcal X\times\Xi)$:
%\begin{equation}
%  \begin{aligned}
%    \overline p_{\hat f}
%    =\sup_{\Pi\in\mathcal M_+(\mathcal X\times\Xi)}\quad&
%      \int_{\mathcal X\times\Xi}
%      \mathbf1_{\{\xi\le -\hat g(\bm x)\}}\,\Pi(\mathrm d\bm x,\mathrm d\xi)\\
%    \text{s.t.}\quad&
%      \int_{\mathcal X\times\Xi}
%      \mathbf1_{\mathcal U}(\bm x)\,\Pi(\mathrm d\bm x,\mathrm d\xi)
%      =\mathbb P(\bm X\in\mathcal U),
%      &&\forall \mathcal U\in\mathcal B(\mathcal X),\\
%    &\int_{\mathcal X\times\Xi}
%      \mathbf1_{\{\xi<\xi_{(k)}^{\mathrm{te}}\}}\,\Pi(\mathrm d\bm x,\mathrm d\xi)\le u_k,
%      &&k=1,\ldots,K,\\
%    &\int_{\mathcal X\times\Xi}
%      \mathbf1_{\{\xi\le \xi_{(k)}^{\mathrm{te}}\}}\,\Pi(\mathrm d\bm x,\mathrm d\xi)\ge l_k,
%      &&k=1,\ldots,K.
%  \end{aligned}
%  \label{eq-measure-upper}
%\end{equation}
%The lower problem minimizes the same integral subject to the same constraints:
%\begin{equation*}
%  \underline p_{\hat f}
%  =\inf_{\Pi\in\mathcal P(\mathcal Q^{\mathrm{DW}}(\theta))}
%    \int_{\mathcal X\times\Xi}
%    \mathbf1_{\{\xi\le -\hat g(\bm x)\}}\,\Pi(\mathrm d\bm x,\mathrm d\xi).
%\end{equation*}
%The marginal equality fixes the input law and total mass one. The rank inequalities restrict the residual marginal while allowing its dependence on the input to vary. For every $\theta\ge0$, the product of $\mathbb P$ and the empirical residual law is feasible when $\Xi=\mathbb R$.
%
%For $\bm\lambda,\bm\gamma\in\mathbb R_+^K$, collect the rank penalties in the residual step function
%\begin{equation}
%  \zeta_{\bm\lambda,\bm\gamma}(\xi)
%  :=\sum_{k=1}^K\lambda_k\mathbf1_{\{\xi<\xi_{(k)}^{\mathrm{te}}\}}
%    -\sum_{k=1}^K\gamma_k\mathbf1_{\{\xi\le\xi_{(k)}^{\mathrm{te}}\}}.
%  \label{eq-dual-step-function}
%\end{equation}
%
%\begin{lemma}[Lagrangian dual problems and weak duality]
%  \label{lem-weak-duality}
%  Consider the feasible upper measure program \eqref{eq-measure-upper} and its lower minimization counterpart. Write $\bm u=(u_1,\ldots,u_K)^\top$ and $\bm l=(l_1,\ldots,l_K)^\top$. The nonnegative multipliers $\bm\lambda,\bm\gamma\in\mathbb R_+^K$ correspond to the upper and lower rank constraints. The input marginal equality has an unrestricted multiplier $\eta\in\mathcal F_{\mathrm b}(\mathcal X)$, the space of bounded Borel measurable real-valued functions on $\mathcal X$. Eliminating $\eta$ pointwise gives the dual problems
%  \begin{align}
%    \overline d_\theta
%    &:=\inf_{\bm\lambda,\bm\gamma\in\mathbb R_+^K}
%    \left\{\bm\lambda^\top\bm u-\bm\gamma^\top\bm l
%    +\mathbb E_{\mathbb P}\left[\sup_{\xi\in\Xi}\left\{
%    \mathbf1_{\{\xi\le-\hat g(\bm X)\}}
%    -\zeta_{\bm\lambda,\bm\gamma}(\xi)
%    \right\}\right]\right\},
%    \label{eq-upper-measure-dual}\\
%    \underline d_\theta
%    &:=\sup_{\bm\lambda,\bm\gamma\in\mathbb R_+^K}
%    \left\{-\bm\lambda^\top\bm u+\bm\gamma^\top\bm l
%    +\mathbb E_{\mathbb P}\left[\inf_{\xi\in\Xi}\left\{
%    \mathbf1_{\{\xi\le-\hat g(\bm X)\}}
%    +\zeta_{\bm\lambda,\bm\gamma}(\xi)
%    \right\}\right]\right\}.
%    \label{eq-lower-measure-dual}
%  \end{align}
%  Their optimal values satisfy
%  \begin{equation}
%    \underline d_\theta\le\underline p_{\hat f}
%    \le\overline p_{\hat f}\le\overline d_\theta.
%    \label{eq-weak-duality}
%  \end{equation}
%\end{lemma}
%\begin{proof}
%  The input marginal equality is equivalent to
%  \begin{equation*}
%    \int_{\mathcal X\times\Xi}\eta(\bm x)\,\Pi(\mathrm d\bm x,\mathrm d\xi)
%    =\mathbb E_{\mathbb P}\left[\eta(\bm X)\right],
%    \qquad\forall\eta\in\mathcal F_{\mathrm b}(\mathcal X).
%  \end{equation*}
%  Relaxing the marginal equality and the rank inequalities gives the upper Lagrangian \citep[Section~5.1]{B.V04}
%  \begin{equation*}
%    \begin{aligned}
%      \mathcal L^+(\Pi,\eta,\bm\lambda,\bm\gamma) = & \int_{\mathcal X\times\Xi}\mathbf1_{\{\xi\le-\hat g(\bm x)\}}\,\Pi(\mathrm d\bm x,\mathrm d\xi) +\mathbb E_{\mathbb P}\left[\eta(\bm X)\right]-\int_{\mathcal X\times\Xi}\eta(\bm x)\,\Pi(\mathrm d\bm x,\mathrm d\xi)\\
%      &+\sum_{k=1}^K\lambda_k\left(u_k-\int_{\mathcal X\times\Xi}\mathbf1_{\{\xi<\xi_{(k)}^{\mathrm{te}}\}}\,\Pi(\mathrm d\bm x,\mathrm d\xi)\right) +\sum_{k=1}^K\gamma_k\left(\int_{\mathcal X\times\Xi}\mathbf1_{\{\xi\le\xi_{(k)}^{\mathrm{te}}\}}\,\Pi(\mathrm d\bm x,\mathrm d\xi)-l_k\right).
%    \end{aligned}
%  \end{equation*}
%  Collecting the terms involving $\Pi$ gives
%  \begin{equation*}
%    \mathcal L^+
%    =\bm\lambda^\top\bm u-\bm\gamma^\top\bm l
%      +\mathbb E_{\mathbb P}\left[\eta(\bm X)\right]
%      +\int_{\mathcal X\times\Xi}\left[
%        \mathbf1_{\{\xi\le-\hat g(\bm x)\}}-\eta(\bm x)-\zeta_{\bm\lambda,\bm\gamma}(\xi)
%      \right]\,\Pi(\mathrm d\bm x,\mathrm d\xi).
%  \end{equation*}
%  The Lagrangian dual function is defined as
%  \begin{equation*}
%    D^+(\eta,\bm\lambda,\bm\gamma)
%    :=\sup_{\Pi\in\mathcal M_+(\mathcal X\times\Xi)}\mathcal L^+(\Pi,\eta,\bm\lambda,\bm\gamma),
%  \end{equation*}
%  The relaxed measure has arbitrary finite nonnegative mass. If the integrand in the collected Lagrangian is positive at any point, scaling a point mass there by $M\to\infty$ gives $D^+=+\infty$. If the integrand is nonpositive everywhere, its supremum integral is zero, attained by the zero measure. The resulting functional dual is
%  \begin{equation*}
%    \begin{aligned}
%      \overline d_\theta
%      =\inf_{\substack{\eta\in\mathcal F_{\mathrm b}(\mathcal X)\\
%      \bm\lambda,\bm\gamma\in\mathbb R_+^K}}
%      \quad &\bm\lambda^\top\bm u-\bm\gamma^\top\bm l+\mathbb E_{\mathbb P}\left[\eta(\bm X)\right]\\
%      \text{s.t.}\quad & \eta(\bm x)\ge
%      \mathbf1_{\{\xi\le-\hat g(\bm x)\}}
%      -\zeta_{\bm\lambda,\bm\gamma}(\xi),
%      \qquad\forall(\bm x,\xi)\in\mathcal X\times\Xi.
%    \end{aligned}
%  \end{equation*}
%  For fixed rank multipliers, the pointwise smallest feasible marginal function is
%  \begin{equation*}
%    \eta^+_{\bm\lambda,\bm\gamma}(\bm x)
%    :=\sup_{\xi\in\Xi}\left\{
%      \mathbf1_{\{\xi\le-\hat g(\bm x)\}}
%      -\zeta_{\bm\lambda,\bm\gamma}(\xi)
%    \right\}.
%  \end{equation*}
%  The envelope is bounded in absolute value by $1+\|\bm\lambda\|_1+\|\bm\gamma\|_1$. As a function of the real threshold replacing $-\hat g(\bm x)$, it is nondecreasing and hence Borel measurable. Composing it with the measurable surrogate gives $\eta^+_{\bm\lambda,\bm\gamma}\in\mathcal F_{\mathrm b}(\mathcal X)$. This smallest feasible function minimizes the expectation and yields the upper dual \eqref{eq-upper-measure-dual}.
%
%  For the lower problem, reversing the inequality penalties gives
%  \begin{equation*}
%    \mathcal L^-
%    =-\bm\lambda^\top\bm u+\bm\gamma^\top\bm l
%      +\mathbb E_{\mathbb P}\left[\eta(\bm X)\right]
%      +\int_{\mathcal X\times\Xi}\left[
%        \mathbf1_{\{\xi\le-\hat g(\bm x)\}}-\eta(\bm x)+\zeta_{\bm\lambda,\bm\gamma}(\xi)
%      \right]\,\Pi(\mathrm d\bm x,\mathrm d\xi).
%  \end{equation*}
%  Its infimum over finite nonnegative measures is finite exactly when $\eta(\bm x)\le\mathbf1_{\{\xi\le-\hat g(\bm x)\}}+\zeta_{\bm\lambda,\bm\gamma}(\xi)$ for every $(\bm x,\xi)$. Under this condition, the infimum integral is zero. The pointwise largest feasible marginal function is therefore
%  \begin{equation*}
%    \eta^-_{\bm\lambda,\bm\gamma}(\bm x)
%    :=\inf_{\xi\in\Xi}\left\{
%      \mathbf1_{\{\xi\le-\hat g(\bm x)\}}
%      +\zeta_{\bm\lambda,\bm\gamma}(\xi)
%    \right\},
%  \end{equation*}
%  which is bounded and Borel measurable by the same argument. Substitution gives the lower dual \eqref{eq-lower-measure-dual}.
%
%  For a primal feasible $\Pi$, the upper Lagrangian is at least $p_{\hat f}(\Pi)$ and the lower Lagrangian is at most $p_{\hat f}(\Pi)$. Every feasible upper dual solution therefore bounds the primal objective from above, and every feasible lower dual solution bounds it from below. Taking the respective optima and using the common nonempty primal feasible set gives \eqref{eq-weak-duality}.
%\end{proof}
%
%The remaining supremum and infimum combine the residual penalty with the failure indicator. For $\Xi=\mathbb R$, the penalty in \eqref{eq-dual-step-function} changes only at observed residuals. Separate these values from the intervening intervals to form $J:=2K+1$ residual cells:
%\begin{equation}
%  \Xi_{2k}:=\left\{\xi_{(k)}^{\mathrm{te}}\right\}
%    \quad\left(k=1,\ldots,K\right),
%  \qquad
%  \Xi_{2k+1}:=\left(\xi_{(k)}^{\mathrm{te}},\xi_{(k+1)}^{\mathrm{te}}\right)
%    \quad\left(k=0,\ldots,K\right).
%  \label{eq-residual-cells}
%\end{equation}
%The input cells in \eqref{eq-input-cells} and the residual cells in \eqref{eq-residual-cells} satisfy
%\begin{equation*}
%  \mathcal X=\bigsqcup_{i=0}^K\mathcal X_i,
%  \qquad
%  \Xi=\bigsqcup_{j=1}^J\Xi_j.
%\end{equation*}
%The singleton cells retain residual atoms and distinguish the open and closed half-lines in the rank constraints. The penalty is constant on each $\Xi_j$, with value
%\begin{equation}
%  \zeta_j:=\sum_{k=1}^K\lambda_k\mathbf1_{\{j<2k\}}
%    -\sum_{k=1}^K\gamma_k\mathbf1_{\{j\le2k\}},
%    \qquad j=1,\ldots,J.
%  \label{eq-dual-cell-values}
%\end{equation}
%For $(\bm x,\xi)\in\mathcal X_i\times\Xi_j$, the failure indicator satisfies
%\begin{equation}
%  \underline h_{ij}:=\mathbf1_{\{j\le2i\}}
%  \le\mathbf1_{\{\xi\le-\hat g(\bm x)\}}
%  \le\overline h_{ij}:=\mathbf1_{\{j\le2i+1\}}.
%  \label{eq-cell-costs}
%\end{equation}
%For inputs in $\mathcal X_i$, residual cells with $j\le2i$ always cause failure, and those with $j\ge2i+2$ are always safe. The overlapping interval $\Xi_{2i+1}$ determines the difference between the two coefficients. When $-\hat g(\bm x)=\xi_{(i)}^{\mathrm{te}}$, that interval is safe; the upper construction assigns its mass to the left endpoint. The proof below verifies this assignment. The weights $w_i$, penalty values in \eqref{eq-dual-cell-values}, and failure coefficients now give finite linear programs for the duals.
%
%\begin{theorem}[Finite-dimensional duals and strong duality]
%  \label{thm-finite-dual-reformulations}
%  For $\Xi=\mathbb R$, $\mathcal Q=\mathcal Q^{\mathrm{DW}}(\theta)$, and every $\theta\ge0$, the dual values in Lemma~\ref{lem-weak-duality} equal the optimal values of
%  \begin{equation}
%    \begin{aligned}
%      \overline d_\theta
%      =\min_{\bm\lambda,\bm\gamma\ge\bm 0,\ \bm z\in\mathbb R^{K+1}}\quad&
%        \bm\lambda^\top\bm u-\bm\gamma^\top\bm l+\sum_{i=0}^K w_i z_i\\
%      \text{s.t.}\quad&
%        z_i+\zeta_j\ge\overline h_{ij},
%        \qquad i=0,\ldots,K,\quad j=1,\ldots,J,
%    \end{aligned}
%    \label{eq-finite-upper-dual}
%  \end{equation}
%  and
%  \begin{equation}
%    \begin{aligned}
%      \underline d_\theta
%      =\max_{\bm\lambda,\bm\gamma\ge\bm 0,\ \bm z\in\mathbb R^{K+1}}\quad&
%        -\bm\lambda^\top\bm u+\bm\gamma^\top\bm l+\sum_{i=0}^K w_i z_i\\
%      \text{s.t.}\quad&
%        z_i-\zeta_j\le\underline h_{ij},
%        \qquad i=0,\ldots,K,\quad j=1,\ldots,J.
%    \end{aligned}
%    \label{eq-finite-lower-dual}
%  \end{equation}
%  Moreover,
%  \begin{equation}
%    \overline p_{\hat f}=\overline d_\theta,\qquad
%    \underline p_{\hat f}=\underline d_\theta.
%    \label{eq-strong-duality}
%  \end{equation}
%  Each LP has $3K+1$ scalar variables and $(K+1)(2K+1)$ pairwise inequalities, together with multiplier nonnegativity. The free variables $z_i$ correspond to the input cell probability equalities. The finite LP optima and the corresponding measure primal and functional dual optima are attained.
%\end{theorem}
%\begin{proof}
%  For a feasible joint law, let $\bm\pi=(\pi_{ij})$ be its matrix of joint cell probabilities,
%  \begin{equation*}
%    \pi_{ij}:=\Pi(\mathcal X_i\times\Xi_j),
%    \qquad i=0,\ldots,K,\quad j=1,\ldots,J.
%  \end{equation*}
%  Its row sums are the fixed input cell probabilities, and its cumulative column sums satisfy the rank constraints. Thus $\bm\pi\in\mathcal C_\theta$, where
%  \begin{equation}
%    \mathcal C_\theta:=
%    \left\{\bm\pi\in\mathbb R_+^{(K+1)\times J}:
%    \begin{array}{ll}
%      \displaystyle\sum_{j=1}^J\pi_{ij}=w_i,
%        &i=0,\ldots,K,\\[1mm]
%      \displaystyle\sum_{i=0}^K\sum_{j<2k}\pi_{ij}\le u_k,
%        &k=1,\ldots,K,\\[1mm]
%      \displaystyle\sum_{i=0}^K\sum_{j\le2k}\pi_{ij}\ge l_k,
%        &k=1,\ldots,K
%    \end{array}\right\}.
%    \label{eq-cell-mass-polytope}
%  \end{equation}
%  By \eqref{eq-cell-costs}, maximizing $\sum_{i,j}\overline h_{ij}\pi_{ij}$ over this polytope bounds the upper measure problem from above, while minimizing $\sum_{i,j}\underline h_{ij}\pi_{ij}$ bounds the lower problem from below.
%
%  To establish the reverse inequalities, fix $\bm\pi\in\mathcal C_\theta$. Draw $\bm x$ from $\mathbb P$ and, on $\mathcal X_i$ with $w_i>0$, select a provisional residual cell $j$ with probability $\pi_{ij}/w_i$. These probabilities sum to one by the row equality. For $j\ne2i+1$, choose a fixed representative of $\Xi_j$; the two coefficients in \eqref{eq-cell-costs} then equal the actual failure indicator. The following assignments in the overlapping cell define the upper and lower joint laws $\Pi^+$ and $\Pi^-$, respectively.
%
%  For the upper problem, assign $\xi=-\hat g(\bm x)$ in the overlapping cell. This causes failure and lies in $\Xi_{2i+1}$ when $-\hat g(\bm x)>\xi_{(i)}^{\mathrm{te}}$. At equality, the assignment moves the provisional mass to the left endpoint. For $i\ge1$, any $\widetilde\xi\in\left(\xi_{(i)}^{\mathrm{te}},\xi_{(i+1)}^{\mathrm{te}}\right)$, and each observed threshold,
%  \begin{equation*}
%    \mathbf1_{\{\xi_{(i)}^{\mathrm{te}}<\xi_{(k)}^{\mathrm{te}}\}}=\mathbf1_{\{\widetilde\xi<\xi_{(k)}^{\mathrm{te}}\}},\qquad
%    \mathbf1_{\{\xi_{(i)}^{\mathrm{te}}\le \xi_{(k)}^{\mathrm{te}}\}}\ge\mathbf1_{\{\widetilde\xi\le \xi_{(k)}^{\mathrm{te}}\}}.
%  \end{equation*}
%  The assignment leaves every open upper probability unchanged and can only increase the closed lower probabilities. It therefore preserves all rank constraints and attains the upper coefficient. For $i=0$, the finite threshold $-\hat g(\bm x)$ is strictly above $\xi_{(0)}^{\mathrm{te}}=-\infty$.
%
%  For the lower problem, assign $\xi=\left(\xi_{(i+1)}^{\mathrm{te}}-\hat g(\bm x)\right)/2$ when $i<K$, and $\xi=-\hat g(\bm x)+1$ when $i=K$. Each value lies in $\Xi_{2i+1}$ and strictly exceeds $-\hat g(\bm x)$, including when the threshold equals $\xi_{(i)}^{\mathrm{te}}$. This attains the lower coefficient and preserves the provisional residual cell.
%
%  On an input cell with $w_i=0$, the row of $\bm\pi$ is zero; set $\xi=\xi_{(1)}^{\mathrm{te}}$ there. All residual assignments are Borel measurable functions of $\bm x$. Their finite mixtures define probability kernels, which, combined with $\mathbb P$, give joint probability laws with the prescribed input marginal. This construction also allows randomization at input atoms. Both joint laws satisfy the rank constraints and attain the assigned matrix objectives:
%  \begin{equation*}
%    p_{\hat f}(\Pi^+)=\sum_{i=0}^K\sum_{j=1}^J\overline h_{ij}\pi_{ij},
%    \qquad
%    p_{\hat f}(\Pi^-)=\sum_{i=0}^K\sum_{j=1}^J\underline h_{ij}\pi_{ij}.
%  \end{equation*}
%  Combining these identities with the initial bounds gives
%  \begin{equation}
%    \overline p_{\hat f}
%    =\max_{\bm\pi\in\mathcal C_\theta}\sum_{i=0}^K\sum_{j=1}^J
%      \overline h_{ij}\pi_{ij},\qquad
%    \underline p_{\hat f}
%    =\min_{\bm\pi\in\mathcal C_\theta}\sum_{i=0}^K\sum_{j=1}^J
%      \underline h_{ij}\pi_{ij}.
%    \label{eq-finite-primal-values}
%  \end{equation}
%  The lower assignment retains every joint cell mass $\pi_{ij}$. The upper assignment can transfer mass from an open interval to its left endpoint while preserving the rank inequalities and the assigned objective value; equality of the resulting cell matrix with $\bm\pi$ is not needed for \eqref{eq-finite-primal-values}.
%
%  For every $\theta\ge0$, a feasible matrix is induced by the empirical product law:
%  \begin{equation*}
%    \pi_{i,2k}^{\mathrm{emp}}:=w_i\frac{c_k-c_{k-1}}{N^{\mathrm{te}}},
%    \qquad i=0,\ldots,K,\quad k=1,\ldots,K,
%  \end{equation*}
%  with all odd columns zero. Its open and closed rank probabilities are $c_{k-1}/N^{\mathrm{te}}$ and $c_k/N^{\mathrm{te}}$, respectively, so \eqref{eq-dw-counts-and-levels} makes it feasible also at $\theta=0$. The set $\mathcal C_\theta$ is closed, and its nonnegative entries sum to one, so it is compact. Both matrix objectives lie in $[0,1]$ and attain their optima, which the preceding constructions realize by joint laws. Finite LP strong duality therefore applies \citep[Section~5.2.3]{B.V04}. The row equalities have free multipliers $z_i$, and the rank inequalities have nonnegative multipliers $\lambda_k,\gamma_k$. Their dual constraints are $z_i+\zeta_j\ge\overline h_{ij}$ and $z_i-\zeta_j\le\underline h_{ij}$, giving \eqref{eq-finite-upper-dual} and \eqref{eq-finite-lower-dual}.
%
%  An upper finite dual solution defines a bounded Borel function $\eta(\bm x)=z_i$ on $\mathcal X_i$. The indicator bounds \eqref{eq-cell-costs} make it feasible for the functional upper dual, so $\eta^+_{\bm\lambda,\bm\gamma}(\bm x)\le z_i$. A lower finite dual solution similarly gives $\eta^-_{\bm\lambda,\bm\gamma}(\bm x)\ge z_i$. The reduced upper objective is therefore no greater than the finite upper dual objective, and the reduced lower objective is no smaller than the finite lower dual objective.
%
%  Finite LP equality and Lemma~\ref{lem-weak-duality} give \eqref{eq-strong-duality}. Optimal finite dual solutions supply optimal rank multipliers and marginal functions for the functional duals. Feasibility and finite optimal values suffice for this LP argument, including at $\theta=0$ when some rank inequalities can be binding for every feasible law.
%\end{proof}
%
%The finite LPs establish exactness of the duals. Their ordered failure coefficients give a further simplification: the best single-threshold candidates derived above are already optimal.
%
%\begin{corollary}[Closed-form bounds and optimal certificates]
%  \label{thm-closed-form-bounds}
%  Under the conditions of Theorem~\ref{thm-finite-dual-reformulations},
%  \begin{align}
%    \overline p_{\hat f}
%    &=\min\left\{1,\min_{k=1,\ldots,K}[u_k+H\left(-\xi_{(k)}^{\mathrm{te}}\right)]\right\},
%    \label{eq-dw-closed-upper}\\
%    \underline p_{\hat f}
%    &=\max\left\{0,\max_{k=1,\ldots,K}[l_k+H\left(-\xi_{(k)}^{\mathrm{te}}\right)-1]\right\}.
%    \label{eq-dw-closed-lower}
%  \end{align}
%  Each dual optimum admits a certificate with at most one nonzero rank multiplier, of value one.
%\end{corollary}
%\begin{proof}
%  Fix nonnegative column masses $\bm y=(y_1,\ldots,y_J)^\top$ with $\sum_j y_j=1$, and write $Y_j:=\sum_{r=1}^j y_r$. Let $\mathcal T(\bm w,\bm y)$ be the set of transportation matrices with row sums $\bm w=(w_0,\ldots,w_K)^\top$ and column sums $\bm y$. It contains the product matrix $\pi_{ij}=w_i y_j$. The rank constraints depend only on $\bm y$, through $Y_{2k-1}\le u_k$ and $Y_{2k}\ge l_k$.
%
%  For fixed $\bm y$, the upper problem maximizes mass on failure pairs $j\le2i+1$; the lower problem maximizes mass on safe pairs $j>2i$ and subtracts it from one. Both are maximum-flow problems: source-to-row capacities are $w_i$, column-to-sink capacities are $y_j$, and a row is connected to each permitted column by an edge of capacity two. Restricting a transportation matrix to permitted pairs gives a feasible flow. Conversely, any flow can be completed to a transportation matrix because the remaining row and column masses are nonnegative and have equal totals. These two operations show that the maximum flow equals the maximum permitted mass.
%
%  A cut of capacity one is available, so a minimum cut crosses no edge of capacity two. Its source-side columns can be taken to be exactly the neighbors of its source-side rows. Because the neighborhoods are nested, adding any row whose neighbors are already included cannot increase the cut capacity. Thus, apart from empty or full row sets, upper-problem cuts reduce to the rows $i<k$ and columns $j<2k$, while safe-mass cuts reduce to the rows $i\ge k$ and columns $j>2k$. For $k=1,\ldots,K$, their capacities are, respectively,
%  \begin{align*}
%    \sum_{i=k}^K w_i+Y_{2k-1}
%      &=H\left(-\xi_{(k)}^{\mathrm{te}}\right)+Y_{2k-1},\\
%    \sum_{i=0}^{k-1}w_i+1-Y_{2k}
%      &=2-H\left(-\xi_{(k)}^{\mathrm{te}}\right)-Y_{2k}.
%  \end{align*}
%  Empty and full row sets give capacity one. Each remaining cut therefore tests a single separation between ordered cells. The max-flow/min-cut theorem \citep{F.F56} gives
%  \begin{align*}
%    \max_{\bm\pi\in\mathcal T(\bm w,\bm y)}\sum_{i,j}\overline h_{ij}\pi_{ij}
%      &=\min\left\{1,\min_k\left[H\left(-\xi_{(k)}^{\mathrm{te}}\right)+Y_{2k-1}\right]\right\},\\
%    \min_{\bm\pi\in\mathcal T(\bm w,\bm y)}\sum_{i,j}\underline h_{ij}\pi_{ij}
%      &=\max\left\{0,\max_k\left[H\left(-\xi_{(k)}^{\mathrm{te}}\right)+Y_{2k}-1\right]\right\}.
%  \end{align*}
%
%  The first value is nondecreasing in every open prefix $Y_{2k-1}$, and the second is nondecreasing in every closed prefix $Y_{2k}$. It remains to attain all upper prefix limits simultaneously for the maximization, and all lower prefix limits simultaneously for the minimization. With $u_{K+1}=1$ and $l_0=0$, choose
%  \begin{align*}
%    &y_1^+:=u_1,\qquad y_{2k}^+:=u_{k+1}-u_k,\qquad
%      y_{2k+1}^+:=0,\quad k=1,\ldots,K,\\
%    &y_{2k-1}^-:=0,\qquad y_{2k}^-:=l_k-l_{k-1},\quad k=1,\ldots,K,
%      \qquad y_J^-:=1-l_K.
%  \end{align*}
%  The monotonicity and compatibility conditions in Proposition~\ref{prop-finite-rank-constraints} make both vectors nonnegative with total mass one. Their prefixes satisfy
%  \[
%    Y_{2k-1}^+=u_k,\qquad Y_{2k}^+=u_{k+1}\ge l_k,
%    \qquad Y_{2k}^-=l_k,\qquad Y_{2k-1}^-=l_{k-1}\le u_k.
%  \]
%  Hence both allocations are feasible and attain the required prefix limits. Substituting them into the fixed-column values and using \eqref{eq-finite-primal-values} proves \eqref{eq-dw-closed-upper} and \eqref{eq-dw-closed-lower}. Zero masses and coincident prefix limits are allowed, so this argument includes $\theta=0$.
%
%  Finally, let $\bm e_k$ be the $k$th unit vector. The upper and lower dual certificates are, respectively,
%  \begin{align*}
%    &\bm\lambda=\bm e_k,\qquad \bm\gamma=\bm 0,\qquad
%      \eta^+_{\bm e_k,\bm 0}(\bm x)
%      =\mathbf1_{\{\hat g(\bm x)\le-\xi_{(k)}^{\mathrm{te}}\}},\\
%    &\bm\lambda=\bm 0,\qquad \bm\gamma=\bm e_k,\qquad
%      \eta^-_{\bm 0,\bm e_k}(\bm x)
%      =\mathbf1_{\{\hat g(\bm x)\le-\xi_{(k)}^{\mathrm{te}}\}}-1.
%  \end{align*}
%  The threshold indicator inequalities established before Lemma~\ref{lem-weak-duality} give feasibility. To verify the envelope values, take $\xi=\xi_{(k)}^{\mathrm{te}}$ in the upper problem when $-\hat g(\bm x)\ge\xi_{(k)}^{\mathrm{te}}$, and $\xi\in\left(-\hat g(\bm x),\xi_{(k)}^{\mathrm{te}}\right]$ in the lower problem when $-\hat g(\bm x)<\xi_{(k)}^{\mathrm{te}}$. These choices attain one and minus one, respectively; a residual above both thresholds attains zero in the complementary cases. At threshold equality, the upper and lower envelopes are therefore one and zero, and the input is included in $H\left(-\xi_{(k)}^{\mathrm{te}}\right)$. The certificate objectives are exactly the inner candidates in \eqref{eq-dw-closed-upper} and \eqref{eq-dw-closed-lower}. Zero rank multipliers give the constant marginal functions one and zero, attaining the outer candidates. Selecting the best candidate thus supplies an optimal full dual solution with at most one nonzero rank multiplier.
%\end{proof}
%
%The upper bound combines residual mass below an observed value with the probability that the surrogate threshold reaches that value. The lower bound measures overlap that every admissible dependence structure must retain. For a prescribed target $p_{\mathrm{tar}}\in(0,1)$, the upper bound is at most $p_{\mathrm{tar}}$ exactly when at least one rank satisfies $u_k+H\left(-\xi_{(k)}^{\mathrm{te}}\right)\le p_{\mathrm{tar}}$. Both the allowed lower-tail residual mass and the surrogate threshold probability must therefore be small. A small mean squared residual alone does not ensure this condition.
%
%For a known residual law, bounding the CDF of $\hat g(\bm X)+\xi$ over all couplings is the classical fixed-marginal sum problem studied by \citet{M82} and \citet{F.N.S87}. Corollary~\ref{thm-closed-form-bounds} gives its finite-band counterpart: the residual law varies within a statistically calibrated DW band, and the open and closed rank constraints retain atoms. The finite LP proof establishes sharpness for this marginal information model.
%
%Dependence uncertainty can remain even with a known residual marginal. For equally likely surrogate responses $0$ and $2$ and equally likely residuals $-1$ and $1$, failure occurs exactly when the response is $0$ and the residual is $-1$. Couplings with these marginals allow failure probabilities from $0$ to $1/2$, whereas independence gives $1/4$. The marginal model leaves the locations of errors unrestricted; paired observations can supply additional constraints on their joint behavior.
%
%\begin{remark}[Known interval support]
%  \label{rem-interval-support}
%  Suppose $\Xi=[a,b]$ with $-\infty<a\le b<\infty$ and the observed residuals lie in this interval. Then
%  \begin{align*}
%    \overline p_{\hat f}
%    &=\min\left\{H(-a),\min_{k=1,\ldots,K}[u_k+H\left(-\xi_{(k)}^{\mathrm{te}}\right)]\right\},\\
%    \underline p_{\hat f}
%    &=\max\left\{H(-b),\max_{k=1,\ldots,K}[l_k+H\left(-\xi_{(k)}^{\mathrm{te}}\right)-1]\right\}.
%  \end{align*}
%  To obtain these expressions, represent the support restriction on $\mathbb R$ by
%  $\mathbb Q\{\xi<a\}\le0$ and $\mathbb Q\{\xi\le b\}\ge1$.
%  Add $a,b$ to the ordered threshold list. At a newly added $a$, use upper and lower bounds zero; at a newly added $b$, use upper and lower bounds one. At coincident observed thresholds, tighten only the upper bound at $a$ and the lower bound at $b$. The enlarged sequences satisfy the compatibility conditions in Proposition~\ref{prop-finite-rank-constraints}, so the same finite-aggregation and ordered-cut arguments apply. The effective additional upper candidate is $H(-a)$, and the effective additional lower candidate is $H(-b)$; the other added candidates are redundant. If $a=b$, both bounds equal $H(-a)$. For the resulting partition on $\Xi=[a,b]$, intersect the singleton and open-interval residual cells from the augmented threshold list with $\Xi$ and omit empty cells. Including the support endpoints in this list retains their atoms and the corresponding changes in failure status.
%\end{remark}
%
%Finite calibration remains relevant in the rare-event regime. After zero failures in $N^{\mathrm{te}}$ independent tests, the one-sided exact binomial upper limit is $1-\delta^{1/N^{\mathrm{te}}}\approx\log(1/\delta)/N^{\mathrm{te}}$ \citep{C.P34}. The DW cutoff also depends on its simultaneous-band calibration: when $K=1$ and $H\left(-\xi_{(1)}^{\mathrm{te}}\right)=0$, the upper bound is $\beta_0(\theta)$. Known support can provide information beyond the finite test sample. In Remark~\ref{rem-interval-support}, $H(-a)=0$ certifies zero failure probability because the residual lower bound rules out failure.
%
%For fixed $\theta$, evaluation requires scalar DW inversion, sorting and counting the test residuals, and computing the $K$ values $H\left(-\xi_{(k)}^{\mathrm{te}}\right)$. The endpoint table depends only on $N^{\mathrm{te}}$ and $\theta$ and can be precomputed. Given this table and the surrogate CDF values, the final upper and lower scans cost $O(K)$.
%
%For $\theta\ge\theta_{N^{\mathrm{te}}}(\delta)$, these exact bounds inherit the finite-sample guarantee of Theorem~\ref{theorem-01}. If the surrogate CDF values are estimated by simulation, simultaneous upper CDF bounds can replace $H\left(-\xi_{(k)}^{\mathrm{te}}\right)$ in the upper expression, and simultaneous lower CDF bounds can replace $H\left(-\xi_{(k)}^{\mathrm{te}}\right)$ in the lower expression. An independent auxiliary simulation can provide these simultaneous bounds conditionally on the test data. Its error probability must be included in the overall confidence budget.
%
%
% END PRESERVED SECTIONS 2--4 BEFORE NARRATIVE REVISION

% BEGIN PRESERVED SECTION 4 BEFORE PARTITION-BASED DUALITY REWRITE (2026-11-25)
%\section{Tractable Reformulation}
%\label{sec-reformulation}
%
%Throughout this section, the training data, the test residuals, and $\theta\ge0$ are fixed. We first invert the DW constraint into linear rank restrictions, then derive the Lagrangian duals and prove their exactness through finite linear programs. The ordered constraints yield closed-form bounds. The main results use $\Xi=\mathbb R$; Remark~\ref{rem-interval-support} treats known interval support.
%
%For each empirical probability level $q_j:=j/N^{\mathrm{te}}$, define the scalar acceptance interval
%\begin{equation}
%  [\alpha_j(\theta),\beta_j(\theta)]
%  :=\left\{p\in[0,1]:
%  N^{\mathrm{te}}\phi(q_j,p)-\psi(q_j,p)\le\theta\right\},
%  \qquad j=0,\ldots,N^{\mathrm{te}}.
%  \label{eq-dw-rank-acceptance}
%\end{equation}
%Let $\xi_{(1)}^{\mathrm{te}}<\cdots<\xi_{(K)}^{\mathrm{te}}$ be the distinct ordered test residuals. Define
%\begin{equation} 
%    c_0 :=0, \quad
%    c_k :=\sum_{j=1}^{N^{\mathrm{te}}}\mathbf1_{\{\xi_j^{\mathrm{te}}\le\xi_{(k)}^{\mathrm{te}}\}}, \quad
%    l_k :=\alpha_{c_k}(\theta),  \quad
%    u_k :=\beta_{c_{k-1}}(\theta),\qquad k=1,\ldots,K.
%  \label{eq-dw-counts-and-levels}
%\end{equation}
%The cumulative counts retain all multiplicities, so $c_K=N^{\mathrm{te}}$ even when test residuals coincide.
%
%\begin{proposition}[CDF band and finite rank constraints]
%  \label{prop-finite-rank-constraints}
%  Each acceptance set in \eqref{eq-dw-rank-acceptance} is a nonempty closed interval containing $q_j$. Define
%  \[
%    L_\theta(t):=\alpha_{N^{\mathrm{te}}\widehat F^{\mathrm{te}}(t)}(\theta),
%    \qquad
%    U_\theta(t):=\beta_{N^{\mathrm{te}}\widehat F^{\mathrm{te}}(t)}(\theta).
%  \]
%  The DW ambiguity set has the equivalent representations
%  \begin{equation*}
%    \mathcal Q^{\mathrm{DW}}(\theta) =\left\{\mathbb Q\in\mathcal M_1(\Xi): L_\theta(t)\le F_{\mathbb Q}(t)\le U_\theta(t), \, \forall t\in\mathbb R\right\}
%  \end{equation*}
%  and
%  \begin{equation}
%    \mathcal Q^{\mathrm{DW}}(\theta)
%    =\left\{\mathbb Q\in\mathcal M_1(\Xi): \mathbb Q\{\xi<\xi_{(k)}^{\mathrm{te}}\}\le u_k, \, \mathbb Q\{\xi\le \xi_{(k)}^{\mathrm{te}}\}\ge l_k, \, k=1,\ldots,K \right\}.
%    \label{eq-finite-rank-constraints}
%  \end{equation}
%  Both endpoint sequences are nondecreasing in the empirical rank. Consequently, $l_k$ and $u_k$ are nondecreasing and satisfy
%  \begin{equation*}
%    0\le l_k,u_k\le1,\quad
%    l_{k-1}\le u_k,\quad
%    l_k\le u_{k+1},
%    \qquad k=1,\ldots,K,
%  \end{equation*}
%  where $l_0:=0$ and $u_{K+1}:=1$.
%\end{proposition}
%\begin{proof}
%  For fixed $q$, the Bernoulli divergence increases as $p$ moves away from $q$, while $\psi(q,p)$ cannot increase as the interval between $q$ and $p$ expands. The corrected discrepancy therefore decreases toward $q$ from below and increases above $q$. It is continuous for $p\in(0,1)$ and equals $-\psi_0(q)\le0$ at $p=q$. The boundary conventions in Section~\ref{sec-ambiguity-set} give a nonempty closed acceptance interval also for $q\in\{0,1\}$. Scalar inversion at each empirical CDF level yields the stated CDF band, following \citet[Supplement, Section~S.6]{L.J23}.
%
%  For $k=1,\ldots,K-1$, the empirical CDF on $\left[\xi_{(k)}^{\mathrm{te}},\xi_{(k+1)}^{\mathrm{te}}\right)$ equals $c_k/N^{\mathrm{te}}$. Monotonicity of a candidate CDF makes its lower band constraint equivalent to $F_{\mathbb Q}\left(\xi_{(k)}^{\mathrm{te}}\right)\ge\alpha_{c_k}(\theta)$ and its upper band constraint equivalent to $F_{\mathbb Q}\left(\xi_{(k+1)}^{\mathrm{te}}-\right)\le\beta_{c_k}(\theta)$. The left exterior interval contributes $F_{\mathbb Q}\left(\xi_{(1)}^{\mathrm{te}}-\right)\le\beta_0(\theta)$, and the right exterior interval contributes $F_{\mathbb Q}\left(\xi_{(K)}^{\mathrm{te}}\right)\ge\alpha_{N^{\mathrm{te}}}(\theta)$. The other exterior constraints are automatic because $\alpha_0(\theta)=0$ and $\beta_{N^{\mathrm{te}}}(\theta)=1$. Using $F_{\mathbb Q}\left(\xi_{(k)}^{\mathrm{te}}-\right)=\mathbb Q\{\xi<\xi_{(k)}^{\mathrm{te}}\}$ and $F_{\mathbb Q}\left(\xi_{(k)}^{\mathrm{te}}\right)=\mathbb Q\{\xi\le \xi_{(k)}^{\mathrm{te}}\}$ proves \eqref{eq-finite-rank-constraints}.
%
%  With $p$ fixed, the corrected discrepancy decreases as $q$ approaches $p$ from either side. For empirical levels $q<q'$, an accepted value $p<q$ at level $q'$ remains accepted at level $q$; an accepted value $p>q'$ at level $q$ remains accepted at level $q'$. Since each acceptance interval contains its empirical level, these implications make both endpoints nondecreasing. The inequalities $\alpha_{c_k}(\theta)\le c_k/N^{\mathrm{te}}\le\beta_{c_k}(\theta)$ then give the stated compatibility conditions. This argument permits singleton intervals: if $q_j=1/2$ and $\theta=0$, the acceptance interval is $\{1/2\}$ because $\psi(1/2,p)=0$.
%\end{proof}
%
%The open upper and closed lower half-lines distinguish mass at the observed residuals. Each scalar acceptance interval can be computed by one-dimensional bisection.
%
%Define the CDF of the surrogate response by $H(t):=\mathbb P\left\{\hat g(\bm X)\le t\right\}$ for $t\in\mathbb R$. The observed residuals classify inputs by the location of the failure threshold $-\hat g(\bm x)$. Set $\xi_{(0)}^{\mathrm{te}}:=-\infty$ and $\xi_{(K+1)}^{\mathrm{te}}:=+\infty$, with $H(-\infty):=0$ and $H(+\infty):=1$. For $i=0,\ldots,K$, the input cells and their probabilities are
%\begin{equation}
%  \mathcal X_i:=\left\{\bm x\in\mathcal X:
%    \xi_{(i)}^{\mathrm{te}}\le-\hat g(\bm x)<\xi_{(i+1)}^{\mathrm{te}}\right\},
%  \qquad
%  w_i:=\mathbb P(\mathcal X_i)
%    =H\left(-\xi_{(i)}^{\mathrm{te}}\right)-H\left(-\xi_{(i+1)}^{\mathrm{te}}\right).
%  \label{eq-input-cells}
%\end{equation}
%These $K+1$ cells partition $\mathcal X$, and each threshold event is a union of consecutive cells:
%\[
%  \mathbf1_{\{\hat g(\bm x)\le-\xi_{(k)}^{\mathrm{te}}\}}
%    =\sum_{i=k}^K\mathbf1_{\mathcal X_i}(\bm x),
%  \qquad
%  H\left(-\xi_{(k)}^{\mathrm{te}}\right)=\sum_{i=k}^K w_i,
%  \qquad k=1,\ldots,K.
%\]
%The half-open convention assigns an input with $-\hat g(\bm x)=\xi_{(k)}^{\mathrm{te}}$ to $\mathcal X_k$, preserving threshold equalities even when $H$ has atoms.
%
%For every rank $k$, the failure indicator satisfies 
%\[
%  \mathbf1_{\{\xi\le\xi_{(k)}^{\mathrm{te}}\}}
%    +\mathbf1_{\{\hat g(\bm x)\le-\xi_{(k)}^{\mathrm{te}}\}}-1
%  \le\mathbf1_{\{\xi\le-\hat g(\bm x)\}}
%  \le\mathbf1_{\{\xi<\xi_{(k)}^{\mathrm{te}}\}}
%    +\mathbf1_{\{\hat g(\bm x)\le-\xi_{(k)}^{\mathrm{te}}\}}.
%\]
%The upper inequality covers failure by a union of two events, while the lower inequality guarantees failure on their intersection. Integrating under any feasible joint law gives
%\begin{equation*}
%  l_k+H\left(-\xi_{(k)}^{\mathrm{te}}\right)-1\le p_{\hat f}(\Pi)\le u_k+H\left(-\xi_{(k)}^{\mathrm{te}}\right),
%  \qquad k=1,\ldots,K.
%\end{equation*}
%These threshold bounds will provide dual certificates; the finite reformulation will establish their sharpness.
%
%Substituting \eqref{eq-finite-rank-constraints} into the upper failure probability problem gives the following linear program over $\mathcal M_+(\mathcal X\times\Xi)$:
%\begin{equation}
%  \begin{aligned}
%    \overline p_{\hat f}
%    =\sup_{\Pi\in\mathcal M_+(\mathcal X\times\Xi)}\quad&
%      \int_{\mathcal X\times\Xi}
%      \mathbf1_{\{\xi\le -\hat g(\bm x)\}}\,\Pi(\mathrm d\bm x,\mathrm d\xi)\\
%    \text{s.t.}\quad&
%      \int_{\mathcal X\times\Xi}
%      \mathbf1_{\mathcal U}(\bm x)\,\Pi(\mathrm d\bm x,\mathrm d\xi)
%      =\mathbb P(\bm X\in\mathcal U),
%      &&\forall \mathcal U\in\mathcal B(\mathcal X),\\
%    &\int_{\mathcal X\times\Xi}
%      \mathbf1_{\{\xi<\xi_{(k)}^{\mathrm{te}}\}}\,\Pi(\mathrm d\bm x,\mathrm d\xi)\le u_k,
%      &&k=1,\ldots,K,\\
%    &\int_{\mathcal X\times\Xi}
%      \mathbf1_{\{\xi\le \xi_{(k)}^{\mathrm{te}}\}}\,\Pi(\mathrm d\bm x,\mathrm d\xi)\ge l_k,
%      &&k=1,\ldots,K.
%  \end{aligned}
%  \label{eq-measure-upper}
%\end{equation}
%The lower problem minimizes the same integral subject to the same constraints:
%\begin{equation*}
%  \underline p_{\hat f}
%  =\inf_{\Pi\in\mathcal P(\mathcal Q^{\mathrm{DW}}(\theta))}
%    \int_{\mathcal X\times\Xi}
%    \mathbf1_{\{\xi\le -\hat g(\bm x)\}}\,\Pi(\mathrm d\bm x,\mathrm d\xi).
%\end{equation*}
%The marginal equalities fix the input law and total mass one. For $\Xi=\mathbb R$, the product of $\mathbb P$ and the empirical residual law is feasible for every $\theta\ge0$.
%
%\begin{lemma}[Lagrangian dual problems and weak duality]
%  \label{lem-weak-duality}
%  Consider the feasible upper measure program \eqref{eq-measure-upper} and its lower minimization counterpart. Write $\bm u=(u_1,\ldots,u_K)^\top$ and $\bm l=(l_1,\ldots,l_K)^\top$. The nonnegative multipliers $\bm\lambda,\bm\gamma\in\mathbb R_+^K$ correspond to the upper and lower rank constraints. The input marginal equality has an unrestricted multiplier $\eta\in\mathcal F_{\mathrm b}(\mathcal X)$, the space of bounded Borel measurable real-valued functions on $\mathcal X$. For $\bm\lambda,\bm\gamma\in\mathbb R_+^K$, collect the rank penalties in the residual step function
%  \begin{equation}
%    \zeta_{\bm\lambda,\bm\gamma}(\xi)
%    :=\sum_{k=1}^K\lambda_k\mathbf1_{\{\xi<\xi_{(k)}^{\mathrm{te}}\}}
%      -\sum_{k=1}^K\gamma_k\mathbf1_{\{\xi\le\xi_{(k)}^{\mathrm{te}}\}}.
%    \label{eq-dual-step-function}
%  \end{equation}
%  Eliminating $\eta$ pointwise gives the dual problems
%  \begin{align}
%    \overline d_\theta
%    &:=\inf_{\bm\lambda,\bm\gamma\in\mathbb R_+^K}
%    \left\{\bm\lambda^\top\bm u-\bm\gamma^\top\bm l
%    +\mathbb E_{\mathbb P}\left[\sup_{\xi\in\Xi}\left\{
%    \mathbf1_{\{\xi\le-\hat g(\bm X)\}}
%    -\zeta_{\bm\lambda,\bm\gamma}(\xi)
%    \right\}\right]\right\},
%    \label{eq-upper-measure-dual}\\
%    \underline d_\theta
%    &:=\sup_{\bm\lambda,\bm\gamma\in\mathbb R_+^K}
%    \left\{-\bm\lambda^\top\bm u+\bm\gamma^\top\bm l
%    +\mathbb E_{\mathbb P}\left[\inf_{\xi\in\Xi}\left\{
%    \mathbf1_{\{\xi\le-\hat g(\bm X)\}}
%    +\zeta_{\bm\lambda,\bm\gamma}(\xi)
%    \right\}\right]\right\}.
%    \label{eq-lower-measure-dual}
%  \end{align}
%  Their optimal values satisfy
%  \begin{equation}
%    \underline d_\theta\le\underline p_{\hat f}
%    \le\overline p_{\hat f}\le\overline d_\theta.
%    \label{eq-weak-duality}
%  \end{equation}
%\end{lemma}
%\begin{proof}
%  The input marginal equality is equivalent to
%  \begin{equation*}
%    \int_{\mathcal X\times\Xi}\eta(\bm x)\,\Pi(\mathrm d\bm x,\mathrm d\xi)
%    =\mathbb E_{\mathbb P}\left[\eta(\bm X)\right],
%    \qquad\forall\eta\in\mathcal F_{\mathrm b}(\mathcal X).
%  \end{equation*}
%  Relaxing the marginal equality and the rank inequalities gives the upper Lagrangian \citep[Section~5.1]{B.V04}
%  \begin{equation*}
%    \begin{aligned}
%      \mathcal L^+(\Pi,\eta,\bm\lambda,\bm\gamma) = & \int_{\mathcal X\times\Xi}\mathbf1_{\{\xi\le-\hat g(\bm x)\}}\,\Pi(\mathrm d\bm x,\mathrm d\xi) +\mathbb E_{\mathbb P}\left[\eta(\bm X)\right]-\int_{\mathcal X\times\Xi}\eta(\bm x)\,\Pi(\mathrm d\bm x,\mathrm d\xi)\\
%      &+\sum_{k=1}^K\lambda_k\left(u_k-\int_{\mathcal X\times\Xi}\mathbf1_{\{\xi<\xi_{(k)}^{\mathrm{te}}\}}\,\Pi(\mathrm d\bm x,\mathrm d\xi)\right) +\sum_{k=1}^K\gamma_k\left(\int_{\mathcal X\times\Xi}\mathbf1_{\{\xi\le\xi_{(k)}^{\mathrm{te}}\}}\,\Pi(\mathrm d\bm x,\mathrm d\xi)-l_k\right).
%    \end{aligned}
%  \end{equation*}
%  Collecting the terms involving $\Pi$ gives
%  \begin{equation*}
%    \mathcal L^+
%    =\bm\lambda^\top\bm u-\bm\gamma^\top\bm l
%      +\mathbb E_{\mathbb P}\left[\eta(\bm X)\right]
%      +\int_{\mathcal X\times\Xi}\left[
%        \mathbf1_{\{\xi\le-\hat g(\bm x)\}}-\eta(\bm x)-\zeta_{\bm\lambda,\bm\gamma}(\xi)
%      \right]\,\Pi(\mathrm d\bm x,\mathrm d\xi).
%  \end{equation*}
%  The Lagrangian dual function is defined as
%  \begin{equation*}
%    D^+(\eta,\bm\lambda,\bm\gamma)
%    :=\sup_{\Pi\in\mathcal M_+(\mathcal X\times\Xi)}\mathcal L^+(\Pi,\eta,\bm\lambda,\bm\gamma),
%  \end{equation*}
%  The relaxed measure has arbitrary finite nonnegative mass. If the integrand in the collected Lagrangian is positive at any point, scaling a point mass there by $M\to\infty$ gives $D^+=+\infty$. If the integrand is nonpositive everywhere, its supremum integral is zero, attained by the zero measure. The resulting functional dual is
%  \begin{equation*}
%    \begin{aligned}
%      \overline d_\theta
%      =\inf_{\eta\in\mathcal F_{\mathrm b}(\mathcal X),\bm\lambda,\bm\gamma\in\mathbb R_+^K}
%      \quad &\bm\lambda^\top\bm u-\bm\gamma^\top\bm l+\mathbb E_{\mathbb P}\left[\eta(\bm X)\right]\\
%      \text{s.t.}\quad & \eta(\bm x)\ge
%      \mathbf1_{\{\xi\le-\hat g(\bm x)\}}
%      -\zeta_{\bm\lambda,\bm\gamma}(\xi),
%      && \forall(\bm x,\xi)\in\mathcal X\times\Xi.
%    \end{aligned}
%  \end{equation*}
%  For fixed rank multipliers, the pointwise smallest feasible marginal function is
%  \begin{equation*}
%    \eta^+_{\bm\lambda,\bm\gamma}(\bm x)
%    :=\sup_{\xi\in\Xi}\left\{
%      \mathbf1_{\{\xi\le-\hat g(\bm x)\}}
%      -\zeta_{\bm\lambda,\bm\gamma}(\xi)
%    \right\}.
%  \end{equation*}
%  The envelope is bounded in absolute value by $1+\|\bm\lambda\|_1+\|\bm\gamma\|_1$. As a function of the real threshold replacing $-\hat g(\bm x)$, it is nondecreasing and hence Borel measurable. Composing it with the measurable surrogate gives $\eta^+_{\bm\lambda,\bm\gamma}\in\mathcal F_{\mathrm b}(\mathcal X)$. This smallest feasible function minimizes the expectation and yields the upper dual \eqref{eq-upper-measure-dual}.
%
%  For the lower problem, reversing the inequality penalties gives
%  \begin{equation*}
%    \mathcal L^-
%    =-\bm\lambda^\top\bm u+\bm\gamma^\top\bm l
%      +\mathbb E_{\mathbb P}\left[\eta(\bm X)\right]
%      +\int_{\mathcal X\times\Xi}\left[
%        \mathbf1_{\{\xi\le-\hat g(\bm x)\}}-\eta(\bm x)+\zeta_{\bm\lambda,\bm\gamma}(\xi)
%      \right]\,\Pi(\mathrm d\bm x,\mathrm d\xi).
%  \end{equation*}
%  Its infimum over finite nonnegative measures is finite exactly when $\eta(\bm x)\le\mathbf1_{\{\xi\le-\hat g(\bm x)\}}+\zeta_{\bm\lambda,\bm\gamma}(\xi)$ for every $(\bm x,\xi)$. Under this condition, the infimum integral is zero. The pointwise largest feasible marginal function is therefore
%  \begin{equation*}
%    \eta^-_{\bm\lambda,\bm\gamma}(\bm x)
%    :=\inf_{\xi\in\Xi}\left\{
%      \mathbf1_{\{\xi\le-\hat g(\bm x)\}}
%      +\zeta_{\bm\lambda,\bm\gamma}(\xi)
%    \right\},
%  \end{equation*}
%  which is bounded and Borel measurable by the same argument. Substitution gives the lower dual \eqref{eq-lower-measure-dual}.
%
%  For a primal feasible $\Pi$, the upper Lagrangian is at least $p_{\hat f}(\Pi)$ and the lower Lagrangian is at most $p_{\hat f}(\Pi)$. Every feasible upper dual solution therefore bounds the primal objective from above, and every feasible lower dual solution bounds it from below. Taking the respective optima and using the common nonempty primal feasible set gives \eqref{eq-weak-duality}.
%\end{proof}
%
%To evaluate the envelopes in Lemma~\ref{lem-weak-duality}, use the fact that the residual penalty changes only at observed residuals. For $\Xi=\mathbb R$, separating these values from the intervening intervals gives $J:=2K+1$ residual cells:
%\begin{equation}
%  \Xi_{2k}:=\left\{\xi_{(k)}^{\mathrm{te}}\right\}
%    \quad\left(k=1,\ldots,K\right),
%  \qquad
%  \Xi_{2k+1}:=\left(\xi_{(k)}^{\mathrm{te}},\xi_{(k+1)}^{\mathrm{te}}\right)
%    \quad\left(k=0,\ldots,K\right).
%  \label{eq-residual-cells}
%\end{equation}
%The input cells in \eqref{eq-input-cells} and the residual cells in \eqref{eq-residual-cells} satisfy
%\begin{equation*}
%  \mathcal X=\bigsqcup_{i=0}^K\mathcal X_i,
%  \qquad
%  \Xi=\bigsqcup_{j=1}^J\Xi_j.
%\end{equation*}
%The singleton cells retain the distinction between the two rank events. The penalty is constant on each $\Xi_j$, with value
%\begin{equation}
%  \zeta_j:=\sum_{k=1}^K\lambda_k\mathbf1_{\{j<2k\}}
%    -\sum_{k=1}^K\gamma_k\mathbf1_{\{j\le2k\}},
%    \qquad j=1,\ldots,J.
%  \label{eq-dual-cell-values}
%\end{equation}
%For $(\bm x,\xi)\in\mathcal X_i\times\Xi_j$, the failure indicator satisfies
%\begin{equation}
%  \underline h_{ij}:=\mathbf1_{\{j\le2i\}}
%  \le\mathbf1_{\{\xi\le-\hat g(\bm x)\}}
%  \le\overline h_{ij}:=\mathbf1_{\{j\le2i+1\}}.
%  \label{eq-cell-costs}
%\end{equation}
%For inputs in $\mathcal X_i$, cells with $j\le2i$ always cause failure, and those with $j\ge2i+2$ are always safe. Only the overlapping interval $\Xi_{2i+1}$ gives different upper and lower coefficients. At threshold equality, the upper assignment may move mass to the left endpoint; the proof below establishes its feasibility. The cell weights and coefficients give the following finite duals.
%
%\begin{theorem}[Finite-dimensional duals and strong duality]
%  \label{thm-finite-dual-reformulations}
%  For $\Xi=\mathbb R$, $\mathcal Q=\mathcal Q^{\mathrm{DW}}(\theta)$, and every $\theta\ge0$, the dual values in Lemma~\ref{lem-weak-duality} equal the optimal values of
%  \begin{equation}
%    \begin{aligned}
%      \overline d_\theta
%      =\min_{\bm\lambda,\bm\gamma\ge\bm 0,\ \bm z\in\mathbb R^{K+1}}\quad&
%        \bm\lambda^\top\bm u-\bm\gamma^\top\bm l+\sum_{i=0}^K w_i z_i\\
%      \text{s.t.}\quad&
%        z_i+\zeta_j\ge\overline h_{ij},
%        \qquad i=0,\ldots,K,\quad j=1,\ldots,J,
%    \end{aligned}
%    \label{eq-finite-upper-dual}
%  \end{equation}
%  and
%  \begin{equation}
%    \begin{aligned}
%      \underline d_\theta
%      =\max_{\bm\lambda,\bm\gamma\ge\bm 0,\ \bm z\in\mathbb R^{K+1}}\quad&
%        -\bm\lambda^\top\bm u+\bm\gamma^\top\bm l+\sum_{i=0}^K w_i z_i\\
%      \text{s.t.}\quad&
%        z_i-\zeta_j\le\underline h_{ij},
%        \qquad i=0,\ldots,K,\quad j=1,\ldots,J.
%    \end{aligned}
%    \label{eq-finite-lower-dual}
%  \end{equation}
%  Moreover,
%  \begin{equation}
%    \overline p_{\hat f}=\overline d_\theta,\qquad
%    \underline p_{\hat f}=\underline d_\theta.
%    \label{eq-strong-duality}
%  \end{equation}
%  Each LP has $3K+1$ scalar variables and $(K+1)(2K+1)$ pairwise inequalities, together with multiplier nonnegativity. The free variables $z_i$ correspond to the input cell probability equalities. The finite LP optima and the corresponding measure primal and functional dual optima are attained.
%\end{theorem}
%\begin{proof}
%  For a feasible joint law, let $\bm\pi=(\pi_{ij})$ be its matrix of joint cell probabilities,
%  \begin{equation*}
%    \pi_{ij}:=\Pi(\mathcal X_i\times\Xi_j),
%    \qquad i=0,\ldots,K,\quad j=1,\ldots,J.
%  \end{equation*}
%  Its row sums are the fixed input cell probabilities, and its cumulative column sums satisfy the rank constraints. Thus $\bm\pi\in\mathcal C_\theta$, where
%  \begin{equation}
%    \mathcal C_\theta:=
%    \left\{\bm\pi\in\mathbb R_+^{(K+1)\times J}:
%    \begin{array}{ll}
%      \displaystyle\sum_{j=1}^J\pi_{ij}=w_i,
%        &i=0,\ldots,K,\\[1mm]
%      \displaystyle\sum_{i=0}^K\sum_{j<2k}\pi_{ij}\le u_k,
%        &k=1,\ldots,K,\\[1mm]
%      \displaystyle\sum_{i=0}^K\sum_{j\le2k}\pi_{ij}\ge l_k,
%        &k=1,\ldots,K
%    \end{array}\right\}.
%    \label{eq-cell-mass-polytope}
%  \end{equation}
%  By \eqref{eq-cell-costs}, maximizing $\sum_{i,j}\overline h_{ij}\pi_{ij}$ over this polytope bounds the upper measure problem from above, while minimizing $\sum_{i,j}\underline h_{ij}\pi_{ij}$ bounds the lower problem from below.
%
%  To establish the reverse inequalities, fix $\bm\pi\in\mathcal C_\theta$. Draw $\bm x$ from $\mathbb P$ and, on $\mathcal X_i$ with $w_i>0$, select a provisional residual cell $j$ with probability $\pi_{ij}/w_i$. These probabilities sum to one by the row equality. For $j\ne2i+1$, choose a fixed representative of $\Xi_j$; the two coefficients in \eqref{eq-cell-costs} then equal the actual failure indicator. The following assignments in the overlapping cell define the upper and lower joint laws $\Pi^+$ and $\Pi^-$, respectively.
%
%  For the upper problem, assign $\xi=-\hat g(\bm x)$ in the overlapping cell. This causes failure and lies in $\Xi_{2i+1}$ when $-\hat g(\bm x)>\xi_{(i)}^{\mathrm{te}}$. At equality, the assignment moves the provisional mass to the left endpoint. For $i\ge1$, any $\widetilde\xi\in\left(\xi_{(i)}^{\mathrm{te}},\xi_{(i+1)}^{\mathrm{te}}\right)$, and each observed threshold,
%  \begin{equation*}
%    \mathbf1_{\{\xi_{(i)}^{\mathrm{te}}<\xi_{(k)}^{\mathrm{te}}\}}=\mathbf1_{\{\widetilde\xi<\xi_{(k)}^{\mathrm{te}}\}},\qquad
%    \mathbf1_{\{\xi_{(i)}^{\mathrm{te}}\le \xi_{(k)}^{\mathrm{te}}\}}\ge\mathbf1_{\{\widetilde\xi\le \xi_{(k)}^{\mathrm{te}}\}}.
%  \end{equation*}
%  The assignment leaves every open upper probability unchanged and can only increase the closed lower probabilities. It therefore preserves all rank constraints and attains the upper coefficient. For $i=0$, the finite threshold $-\hat g(\bm x)$ is strictly above $\xi_{(0)}^{\mathrm{te}}=-\infty$.
%
%  For the lower problem, assign $\xi=\left(\xi_{(i+1)}^{\mathrm{te}}-\hat g(\bm x)\right)/2$ when $i<K$, and $\xi=-\hat g(\bm x)+1$ when $i=K$. Each value lies in $\Xi_{2i+1}$ and strictly exceeds $-\hat g(\bm x)$, including when the threshold equals $\xi_{(i)}^{\mathrm{te}}$. This attains the lower coefficient and preserves the provisional residual cell.
%
%  On an input cell with $w_i=0$, the row of $\bm\pi$ is zero; set $\xi=\xi_{(1)}^{\mathrm{te}}$ there. All residual assignments are Borel measurable functions of $\bm x$. Their finite mixtures define probability kernels, which, combined with $\mathbb P$, give joint probability laws with the prescribed input marginal. This construction also allows randomization at input atoms. Both joint laws satisfy the rank constraints and attain the assigned matrix objectives:
%  \begin{equation*}
%    p_{\hat f}(\Pi^+)=\sum_{i=0}^K\sum_{j=1}^J\overline h_{ij}\pi_{ij},
%    \qquad
%    p_{\hat f}(\Pi^-)=\sum_{i=0}^K\sum_{j=1}^J\underline h_{ij}\pi_{ij}.
%  \end{equation*}
%  Combining these identities with the initial bounds gives
%  \begin{equation}
%    \overline p_{\hat f}
%    =\max_{\bm\pi\in\mathcal C_\theta}\sum_{i=0}^K\sum_{j=1}^J
%      \overline h_{ij}\pi_{ij},\qquad
%    \underline p_{\hat f}
%    =\min_{\bm\pi\in\mathcal C_\theta}\sum_{i=0}^K\sum_{j=1}^J
%      \underline h_{ij}\pi_{ij}.
%    \label{eq-finite-primal-values}
%  \end{equation}
%  The lower assignment retains every joint cell mass $\pi_{ij}$. The upper assignment can transfer mass from an open interval to its left endpoint while preserving the rank inequalities and the assigned objective value; equality of the resulting cell matrix with $\bm\pi$ is not needed for \eqref{eq-finite-primal-values}.
%
%  For every $\theta\ge0$, a feasible matrix is induced by the empirical product law:
%  \begin{equation*}
%    \pi_{i,2k}^{\mathrm{emp}}:=w_i\frac{c_k-c_{k-1}}{N^{\mathrm{te}}},
%    \qquad i=0,\ldots,K,\quad k=1,\ldots,K,
%  \end{equation*}
%  with all odd columns zero. Its open and closed rank probabilities are $c_{k-1}/N^{\mathrm{te}}$ and $c_k/N^{\mathrm{te}}$, respectively, so \eqref{eq-dw-counts-and-levels} makes it feasible also at $\theta=0$. The set $\mathcal C_\theta$ is closed, and its nonnegative entries sum to one, so it is compact. Both matrix objectives lie in $[0,1]$ and attain their optima, which the preceding constructions realize by joint laws. Finite LP strong duality therefore applies \citep[Section~5.2.3]{B.V04}. The row equalities have free multipliers $z_i$, and the rank inequalities have nonnegative multipliers $\lambda_k,\gamma_k$. Their dual constraints are $z_i+\zeta_j\ge\overline h_{ij}$ and $z_i-\zeta_j\le\underline h_{ij}$, giving \eqref{eq-finite-upper-dual} and \eqref{eq-finite-lower-dual}.
%
%  An upper finite dual solution defines a bounded Borel function $\eta(\bm x)=z_i$ on $\mathcal X_i$. The indicator bounds \eqref{eq-cell-costs} make it feasible for the functional upper dual, so $\eta^+_{\bm\lambda,\bm\gamma}(\bm x)\le z_i$. A lower finite dual solution similarly gives $\eta^-_{\bm\lambda,\bm\gamma}(\bm x)\ge z_i$. The reduced upper objective is therefore no greater than the finite upper dual objective, and the reduced lower objective is no smaller than the finite lower dual objective.
%
%  Finite LP equality and Lemma~\ref{lem-weak-duality} give \eqref{eq-strong-duality}. Optimal finite dual solutions supply optimal rank multipliers and marginal functions for the functional duals. Feasibility and finite optimal values suffice for this LP argument, including at $\theta=0$ when some rank inequalities can be binding for every feasible law.
%\end{proof}
%
%The ordered coefficients in \eqref{eq-cell-costs} also identify optimal dual certificates, reducing both LPs to scans over the observed residual levels.
%
%\begin{corollary}[Closed-form bounds and optimal certificates]
%  \label{thm-closed-form-bounds}
%  Under the conditions of Theorem~\ref{thm-finite-dual-reformulations},
%  \begin{align}
%    \overline p_{\hat f}
%    &=\min\left\{1,\min_{k=1,\ldots,K}[u_k+H\left(-\xi_{(k)}^{\mathrm{te}}\right)]\right\},
%    \label{eq-dw-closed-upper}\\
%    \underline p_{\hat f}
%    &=\max\left\{0,\max_{k=1,\ldots,K}[l_k+H\left(-\xi_{(k)}^{\mathrm{te}}\right)-1]\right\}.
%    \label{eq-dw-closed-lower}
%  \end{align}
%  Each dual optimum admits a certificate with at most one nonzero rank multiplier, of value one.
%\end{corollary}
%\begin{proof}
%  Fix nonnegative column masses $\bm y=(y_1,\ldots,y_J)^\top$ with $\sum_j y_j=1$, and write $Y_j:=\sum_{r=1}^j y_r$. Let $\mathcal T(\bm w,\bm y)$ be the set of transportation matrices with row sums $\bm w=(w_0,\ldots,w_K)^\top$ and column sums $\bm y$. It contains the product matrix $\pi_{ij}=w_i y_j$. The rank constraints depend only on $\bm y$, through $Y_{2k-1}\le u_k$ and $Y_{2k}\ge l_k$.
%
%  For fixed $\bm y$, the upper problem maximizes mass on failure pairs $j\le2i+1$; the lower problem maximizes mass on safe pairs $j>2i$ and subtracts it from one. Both are maximum-flow problems: source-to-row capacities are $w_i$, column-to-sink capacities are $y_j$, and a row is connected to each permitted column by an edge of capacity two. Restricting a transportation matrix to permitted pairs gives a feasible flow. Conversely, any flow can be completed to a transportation matrix because the remaining row and column masses are nonnegative and have equal totals. These two operations show that the maximum flow equals the maximum permitted mass.
%
%  A cut of capacity one is available, so a minimum cut crosses no edge of capacity two. Its source-side columns can be taken to be exactly the neighbors of its source-side rows. Because the neighborhoods are nested, adding any row whose neighbors are already included cannot increase the cut capacity. Thus, apart from empty or full row sets, upper-problem cuts reduce to the rows $i<k$ and columns $j<2k$, while safe-mass cuts reduce to the rows $i\ge k$ and columns $j>2k$. For $k=1,\ldots,K$, their capacities are, respectively,
%  \begin{align*}
%    \sum_{i=k}^K w_i+Y_{2k-1}
%      &=H\left(-\xi_{(k)}^{\mathrm{te}}\right)+Y_{2k-1},\\
%    \sum_{i=0}^{k-1}w_i+1-Y_{2k}
%      &=2-H\left(-\xi_{(k)}^{\mathrm{te}}\right)-Y_{2k}.
%  \end{align*}
%  Empty and full row sets give capacity one. Each remaining cut therefore tests a single separation between ordered cells. The max-flow/min-cut theorem \citep{F.F56} gives
%  \begin{align*}
%    \max_{\bm\pi\in\mathcal T(\bm w,\bm y)}\sum_{i,j}\overline h_{ij}\pi_{ij}
%      &=\min\left\{1,\min_k\left[H\left(-\xi_{(k)}^{\mathrm{te}}\right)+Y_{2k-1}\right]\right\},\\
%    \min_{\bm\pi\in\mathcal T(\bm w,\bm y)}\sum_{i,j}\underline h_{ij}\pi_{ij}
%      &=\max\left\{0,\max_k\left[H\left(-\xi_{(k)}^{\mathrm{te}}\right)+Y_{2k}-1\right]\right\}.
%  \end{align*}
%
%  The first value is nondecreasing in every open prefix $Y_{2k-1}$, and the second is nondecreasing in every closed prefix $Y_{2k}$. It remains to attain all upper prefix limits simultaneously for the maximization, and all lower prefix limits simultaneously for the minimization. With $u_{K+1}=1$ and $l_0=0$, choose
%  \begin{align*}
%    &y_1^+:=u_1,\qquad y_{2k}^+:=u_{k+1}-u_k,\qquad
%      y_{2k+1}^+:=0,\quad k=1,\ldots,K,\\
%    &y_{2k-1}^-:=0,\qquad y_{2k}^-:=l_k-l_{k-1},\quad k=1,\ldots,K,
%      \qquad y_J^-:=1-l_K.
%  \end{align*}
%  The monotonicity and compatibility conditions in Proposition~\ref{prop-finite-rank-constraints} make both vectors nonnegative with total mass one. Their prefixes satisfy
%  \[
%    Y_{2k-1}^+=u_k,\qquad Y_{2k}^+=u_{k+1}\ge l_k,
%    \qquad Y_{2k}^-=l_k,\qquad Y_{2k-1}^-=l_{k-1}\le u_k.
%  \]
%  Hence both allocations are feasible and attain the required prefix limits. Substituting them into the fixed-column values and using \eqref{eq-finite-primal-values} proves \eqref{eq-dw-closed-upper} and \eqref{eq-dw-closed-lower}. Zero masses and coincident prefix limits are allowed, so this argument includes $\theta=0$.
%
%  Finally, let $\bm e_k$ be the $k$th unit vector. The upper and lower dual certificates are, respectively,
%  \begin{align*}
%    &\bm\lambda=\bm e_k,\qquad \bm\gamma=\bm 0,\qquad
%      \eta^+_{\bm e_k,\bm 0}(\bm x)
%      =\mathbf1_{\{\hat g(\bm x)\le-\xi_{(k)}^{\mathrm{te}}\}},\\
%    &\bm\lambda=\bm 0,\qquad \bm\gamma=\bm e_k,\qquad
%      \eta^-_{\bm 0,\bm e_k}(\bm x)
%      =\mathbf1_{\{\hat g(\bm x)\le-\xi_{(k)}^{\mathrm{te}}\}}-1.
%  \end{align*}
%  The threshold indicator inequalities established before Lemma~\ref{lem-weak-duality} give feasibility. To verify the envelope values, take $\xi=\xi_{(k)}^{\mathrm{te}}$ in the upper problem when $-\hat g(\bm x)\ge\xi_{(k)}^{\mathrm{te}}$, and $\xi\in\left(-\hat g(\bm x),\xi_{(k)}^{\mathrm{te}}\right]$ in the lower problem when $-\hat g(\bm x)<\xi_{(k)}^{\mathrm{te}}$. These choices attain one and minus one, respectively; a residual above both thresholds attains zero in the complementary cases. At threshold equality, the upper and lower envelopes are therefore one and zero, and the input is included in $H\left(-\xi_{(k)}^{\mathrm{te}}\right)$. The certificate objectives are exactly the inner candidates in \eqref{eq-dw-closed-upper} and \eqref{eq-dw-closed-lower}. Zero rank multipliers give the constant marginal functions one and zero, attaining the outer candidates. Selecting the best candidate thus supplies an optimal full dual solution with at most one nonzero rank multiplier.
%\end{proof}
%
%The same formulas give the response-threshold curves in Corollary~\ref{cor-simultaneous-threshold-bounds}. Replacing $g$ and $\hat g$ by $g-c$ and $\hat g-c$ leaves every residual, rank count, and endpoint $l_k,u_k$ unchanged. The surrogate CDF becomes $t\mapsto H(t+c)$, so Corollary~\ref{thm-closed-form-bounds} gives
%\begin{align}
%  \overline p_{\hat f}(c)
%    &=\min\left\{1,\min_{k=1,\ldots,K}\left[u_k+H\left(c-\xi_{(k)}^{\mathrm{te}}\right)\right]\right\},
%    \label{eq-threshold-closed-upper}\\
%  \underline p_{\hat f}(c)
%    &=\max\left\{0,\max_{k=1,\ldots,K}\left[l_k+H\left(c-\xi_{(k)}^{\mathrm{te}}\right)-1\right]\right\}.
%    \label{eq-threshold-closed-lower}
%\end{align}
%These identities hold for every $c\in\mathbb R$ and $\theta\ge0$, including $\theta=0$. When $\theta\ge\theta_{N^{\mathrm{te}}}(\delta)$, their coverage is simultaneous by \eqref{eq-threshold-curve-coverage}. An extremizing coupling may depend on $c$; the common feasible true joint law supplies the simultaneous guarantee.
%
%The upper bound combines residual mass below an observed value with the probability that the surrogate threshold reaches it. The lower bound measures the overlap required by every admissible coupling. For $p_{\mathrm{tar}}\in(0,1)$, the upper curve satisfies $\overline p_{\hat f}(c)\le p_{\mathrm{tar}}$ exactly when some rank satisfies $u_k+H\left(c-\xi_{(k)}^{\mathrm{te}}\right)\le p_{\mathrm{tar}}$. Both terms must therefore be small at the same residual level.
%
%With a known residual law, the coupling problem is the classical fixed-marginal sum problem studied by \citet{M82} and \citet{F.N.S87}. Here the residual law also varies within a calibrated DW band, and the open and closed rank constraints retain atoms.
%
%Dependence uncertainty can remain even with a known residual marginal. For equally likely surrogate responses $0$ and $2$ and equally likely residuals $-1$ and $1$, failure occurs exactly when the response is $0$ and the residual is $-1$. Its probability ranges from $0$ to $1/2$ over couplings, whereas independence gives $1/4$. Further marginal calibration alone does not resolve this dependence uncertainty.
%
%\begin{remark}[Known interval support]
%  \label{rem-interval-support}
%  Suppose $\Xi=[a,b]$ with $-\infty<a\le b<\infty$ and the observed residuals lie in this interval. Then
%  \begin{align*}
%    \overline p_{\hat f}
%    &=\min\left\{H(-a),\min_{k=1,\ldots,K}[u_k+H\left(-\xi_{(k)}^{\mathrm{te}}\right)]\right\},\\
%    \underline p_{\hat f}
%    &=\max\left\{H(-b),\max_{k=1,\ldots,K}[l_k+H\left(-\xi_{(k)}^{\mathrm{te}}\right)-1]\right\}.
%  \end{align*}
%  To obtain these expressions, represent the support restriction on $\mathbb R$ by
%  $\mathbb Q\{\xi<a\}\le0$ and $\mathbb Q\{\xi\le b\}\ge1$.
%  Add $a,b$ to the ordered threshold list. At a newly added $a$, use upper and lower bounds zero; at a newly added $b$, use upper and lower bounds one. At coincident observed thresholds, tighten only the upper bound at $a$ and the lower bound at $b$. The enlarged sequences satisfy the compatibility conditions in Proposition~\ref{prop-finite-rank-constraints}, so the same finite-aggregation and ordered-cut arguments apply. The effective additional upper candidate is $H(-a)$, and the effective additional lower candidate is $H(-b)$; the other added candidates are redundant. If $a=b$, both bounds equal $H(-a)$. For the resulting partition on $\Xi=[a,b]$, intersect the singleton and open-interval residual cells from the augmented threshold list with $\Xi$ and omit empty cells. Including the support endpoints in this list retains their atoms and the corresponding changes in failure status. For a response threshold $c$, replace $H(-a)$, $H(-b)$, and $H\left(-\xi_{(k)}^{\mathrm{te}}\right)$ by $H(c-a)$, $H(c-b)$, and $H\left(c-\xi_{(k)}^{\mathrm{te}}\right)$, respectively.
%\end{remark}
%
%After zero failures in $N^{\mathrm{te}}$ independent tests, the one-sided exact binomial upper limit is $1-\delta^{1/N^{\mathrm{te}}}\approx\log(1/\delta)/N^{\mathrm{te}}$ \citep{C.P34}. At $c=0$, the simple DW case $K=1$ and $H\left(-\xi_{(1)}^{\mathrm{te}}\right)=0$ gives an upper bound $\beta_0(\theta)$ determined by the simultaneous-band calibration. Verified support can exclude otherwise unobserved failures: in Remark~\ref{rem-interval-support}, $H(-a)=0$ rules out failure.
%
%Evaluation requires scalar DW inversion, sorting and counting the test residuals, and computing the $K$ values $H\left(c-\xi_{(k)}^{\mathrm{te}}\right)$ at each response threshold $c$. The endpoint table depends only on $N^{\mathrm{te}}$ and $\theta$ and can be precomputed. Given this table and the surrogate CDF values, the two scans cost $O(K)$ at each response threshold.
%
%For $\theta\ge\theta_{N^{\mathrm{te}}}(\delta)$, the exact expressions retain the coverage in Theorem~\ref{theorem-01} and Corollary~\ref{cor-simultaneous-threshold-bounds}. If $H$ is evaluated by independent auxiliary simulation, upper CDF bounds can replace its values in the upper expression, and lower CDF bounds can replace them in the lower expression. For one fixed response threshold, simultaneous bounds at its $K$ evaluation points suffice. To retain coverage of the entire response-threshold curve, the auxiliary bounds must hold over all arguments of $H$. In either case, they must hold conditionally on the test data, with their error probability included in the overall confidence budget.
%
%
% END PRESERVED SECTION 4 BEFORE PARTITION-BASED DUALITY REWRITE

% BEGIN PRESERVED RANK-PROBABILITY EXPOSITION (2026-11-25)
%The singleton cells distinguish probability strictly below an observed residual from probability at or below it.
%
%Define the surrogate-response CDF by $H(t):=\mathbb P\left\{\hat g(\bm X)\le t\right\}$, with $H(-\infty):=0$ and $H(+\infty):=1$. For $i=0,\ldots,K$, the same residual levels classify inputs by the location of the failure threshold $-\hat g(\bm x)$:
%\begin{equation}
%  \mathcal X_i:=\left\{\bm x\in\mathcal X:
%    \xi_{(i)}^{\mathrm{te}}\le-\hat g(\bm x)<\xi_{(i+1)}^{\mathrm{te}}\right\},
%  \qquad
%  w_i:=\mathbb P(\mathcal X_i)
%    =H\left(-\xi_{(i)}^{\mathrm{te}}\right)-H\left(-\xi_{(i+1)}^{\mathrm{te}}\right).
%  \label{eq-input-cells}
%\end{equation}
%Thus the two spaces are partitioned as
%\[
%  \mathcal X=\bigsqcup_{i=0}^{K}\mathcal X_i,
%  \qquad
%  \Xi=\bigsqcup_{j=1}^{J}\Xi_j,
%  \qquad
%  H\left(-\xi_{(k)}^{\mathrm{te}}\right)=\sum_{i=k}^{K}w_i,
%  \quad k=1,\ldots,K.
%\]
%The half-open input convention assigns $-\hat g(\bm x)=\xi_{(k)}^{\mathrm{te}}$ to $\mathcal X_k$. The tail sum therefore includes threshold equalities even when $H$ has atoms.
%
%\begin{proposition}[CDF band and residual cell probabilities]
%  \label{prop-finite-rank-constraints}
%  Each acceptance set in \eqref{eq-dw-rank-acceptance} is a nonempty closed interval containing $q_j$. Define
%  \[
%    L_\theta(t):=\alpha_{N^{\mathrm{te}}\widehat F^{\mathrm{te}}(t)}(\theta),
%    \qquad
%    U_\theta(t):=\beta_{N^{\mathrm{te}}\widehat F^{\mathrm{te}}(t)}(\theta).
%  \]
%  The DW ambiguity set has the equivalent representations
%  \[
%    \mathcal Q^{\mathrm{DW}}(\theta)
%    =\left\{\mathbb Q\in\mathcal M_1(\Xi):
%      L_\theta(t)\le F_{\mathbb Q}(t)\le U_\theta(t),\ \forall t\in\mathbb R\right\}
%  \]
%  and
%  \begin{equation}
%    \mathcal Q^{\mathrm{DW}}(\theta)
%    =\left\{\mathbb Q\in\mathcal M_1(\Xi):
%      \sum_{j<2k}\mathbb Q(\Xi_j)\le u_k,\quad
%      \sum_{j\le2k}\mathbb Q(\Xi_j)\ge l_k,\quad k=1,\ldots,K\right\}.
%    \label{eq-finite-rank-constraints}
%  \end{equation}
%  Equivalently, the rank probabilities satisfy
%  \[
%    \mathbb Q\left\{\xi<\xi_{(k)}^{\mathrm{te}}\right\}\le u_k,
%    \qquad
%    \mathbb Q\left\{\xi\le\xi_{(k)}^{\mathrm{te}}\right\}\ge l_k,
%    \qquad k=1,\ldots,K.
%  \]
%  The sequences $l_k$ and $u_k$ are nondecreasing and satisfy
%  \[
%    0\le l_k,u_k\le1,\qquad l_{k-1}\le u_k,\qquad l_k\le u_{k+1},
%    \qquad k=1,\ldots,K,
%  \]
%  where $l_0:=0$ and $u_{K+1}:=1$.
%\end{proposition}
%\begin{proof}
%  For fixed $q$, the Bernoulli divergence increases as $p$ moves away from $q$, while $\psi(q,p)$ cannot increase as the interval between $q$ and $p$ expands. The corrected discrepancy therefore decreases toward $q$ from below and increases above $q$. It is continuous for $p\in(0,1)$ and equals $-\psi_0(q)\le0$ at $p=q$. The boundary conventions in Section~\ref{sec-ambiguity-set} give a nonempty closed acceptance interval also for $q\in\{0,1\}$. Scalar inversion at each empirical CDF level yields the stated CDF band, following \citet[Supplement, Section~S.6]{L.J23}.
%
%  For $k=1,\ldots,K-1$, the empirical CDF on $\left[\xi_{(k)}^{\mathrm{te}},\xi_{(k+1)}^{\mathrm{te}}\right)$ equals $c_k/N^{\mathrm{te}}$. Monotonicity of a candidate CDF makes its lower band constraint equivalent to $F_{\mathbb Q}\left(\xi_{(k)}^{\mathrm{te}}\right)\ge\alpha_{c_k}(\theta)$ and its upper band constraint equivalent to $F_{\mathbb Q}\left(\xi_{(k+1)}^{\mathrm{te}}-\right)\le\beta_{c_k}(\theta)$. The exterior intervals give $F_{\mathbb Q}\left(\xi_{(1)}^{\mathrm{te}}-\right)\le\beta_0(\theta)$ and $F_{\mathbb Q}\left(\xi_{(K)}^{\mathrm{te}}\right)\ge\alpha_{N^{\mathrm{te}}}(\theta)$; their other band constraints are automatic because $\alpha_0(\theta)=0$ and $\beta_{N^{\mathrm{te}}}(\theta)=1$. The residual partition gives
%  \[
%    F_{\mathbb Q}\left(\xi_{(k)}^{\mathrm{te}}-\right)=\sum_{j<2k}\mathbb Q(\Xi_j),
%    \qquad
%    F_{\mathbb Q}\left(\xi_{(k)}^{\mathrm{te}}\right)=\sum_{j\le2k}\mathbb Q(\Xi_j),
%  \]
%  proving \eqref{eq-finite-rank-constraints} and its half-line representation.
%
%  With $p$ fixed, the corrected discrepancy decreases as $q$ approaches $p$ from either side. For empirical levels $q<q'$, an accepted value $p<q$ at level $q'$ remains accepted at level $q$; an accepted value $p>q'$ at level $q$ remains accepted at level $q'$. Since each acceptance interval contains its empirical level, both endpoints are nondecreasing. The inequalities $\alpha_{c_k}(\theta)\le c_k/N^{\mathrm{te}}\le\beta_{c_k}(\theta)$ give the stated compatibility conditions. Singleton intervals are allowed: if $q_j=1/2$ and $\theta=0$, the interval is $\{1/2\}$ because $\psi(1/2,p)=0$.
%\end{proof}
% END PRESERVED RANK-PROBABILITY EXPOSITION

% BEGIN PRESERVED SECTION 4 BEFORE DIRECT CELL-PROBABILITY NOTATION (2026-11-25)
%\section{Tractable Reformulation}
%\label{sec-reformulation}
%
%Throughout this section, the training data, the test residuals, and $\theta\ge0$ are fixed. We use the observed residuals to partition both the residual and input spaces. These partitions express the DW restrictions as cumulative cell probabilities and reduce the Lagrangian duals to finite linear programs. A joint-distribution construction establishes exact finite primal representations, after which LP duality and the ordering of the cells yield strong duality and closed-form bounds. The main development takes $\Xi=\mathbb R$; Remark~\ref{rem-interval-support} treats known interval support.
%
%For each empirical probability level $q_j:=j/N^{\mathrm{te}}$, define the scalar acceptance interval
%\begin{equation}
%  [\alpha_j(\theta),\beta_j(\theta)]
%  :=\left\{p\in[0,1]:
%  N^{\mathrm{te}}\phi(q_j,p)-\psi(q_j,p)\le\theta\right\},
%  \qquad j=0,\ldots,N^{\mathrm{te}}.
%  \label{eq-dw-rank-acceptance}
%\end{equation}
%Let $\xi_{(1)}^{\mathrm{te}}<\cdots<\xi_{(K)}^{\mathrm{te}}$ be the distinct ordered test residuals, and set
%\begin{equation}
%  c_0:=0,\quad
%  c_k:=\sum_{j=1}^{N^{\mathrm{te}}}\mathbf1_{\{\xi_j^{\mathrm{te}}\le\xi_{(k)}^{\mathrm{te}}\}},\quad
%  l_k:=\alpha_{c_k}(\theta),\quad
%  u_k:=\beta_{c_{k-1}}(\theta),\qquad k=1,\ldots,K.
%  \label{eq-dw-counts-and-levels}
%\end{equation}
%The counts retain all multiplicities, so $c_K=N^{\mathrm{te}}$ even when test residuals coincide. Set $\xi_{(0)}^{\mathrm{te}}:=-\infty$ and $\xi_{(K+1)}^{\mathrm{te}}:=+\infty$. Separating the observed values from the intervening intervals gives $J:=2K+1$ residual cells:
%\begin{equation}
%  \Xi_{2k}:=\left\{\xi_{(k)}^{\mathrm{te}}\right\}
%    \quad\left(k=1,\ldots,K\right),
%  \qquad
%  \Xi_{2k+1}:=\left(\xi_{(k)}^{\mathrm{te}},\xi_{(k+1)}^{\mathrm{te}}\right)
%    \quad\left(k=0,\ldots,K\right).
%  \label{eq-residual-cells}
%\end{equation}
%For any candidate residual law $\mathbb Q\in\mathcal M_1(\Xi)$, the open rank probability measures mass strictly below an observed residual, while the closed rank probability also includes mass at that value. For $k=1,\ldots,K$, the partition gives
%\[
%  \begin{aligned}
%    \mathbb Q\left\{\xi<\xi_{(k)}^{\mathrm{te}}\right\}
%      &=F_{\mathbb Q}\left(\xi_{(k)}^{\mathrm{te}}-\right)
%        =\sum_{j<2k}\mathbb Q(\Xi_j),\\
%    \mathbb Q\left\{\xi\le\xi_{(k)}^{\mathrm{te}}\right\}
%      &=F_{\mathbb Q}\left(\xi_{(k)}^{\mathrm{te}}\right)
%        =\sum_{j\le2k}\mathbb Q(\Xi_j).
%  \end{aligned}
%\]
%Here $F_{\mathbb Q}(t-):=\lim_{s\uparrow t}F_{\mathbb Q}(s)$ is the CDF left limit. The two probabilities differ by the atom mass $\mathbb Q(\Xi_{2k})$.
%
%
%Define the surrogate-response CDF by $H(t):=\mathbb P\left\{\hat g(\bm X)\le t\right\}$, with $H(-\infty):=0$ and $H(+\infty):=1$. For $i=0,\ldots,K$, the same residual levels classify inputs by the location of the failure threshold $-\hat g(\bm x)$:
%\begin{equation}
%  \mathcal X_i:=\left\{\bm x\in\mathcal X:
%    \xi_{(i)}^{\mathrm{te}}\le-\hat g(\bm x)<\xi_{(i+1)}^{\mathrm{te}}\right\},
%  \qquad
%  w_i:=\mathbb P(\mathcal X_i)
%    =H\left(-\xi_{(i)}^{\mathrm{te}}\right)-H\left(-\xi_{(i+1)}^{\mathrm{te}}\right).
%  \label{eq-input-cells}
%\end{equation}
%Thus the two spaces are partitioned as
%\[
%  \mathcal X=\bigsqcup_{i=0}^{K}\mathcal X_i,
%  \qquad
%  \Xi=\bigsqcup_{j=1}^{J}\Xi_j,
%  \qquad
%  H\left(-\xi_{(k)}^{\mathrm{te}}\right)=\sum_{i=k}^{K}w_i,
%  \quad k=1,\ldots,K.
%\]
%The half-open input convention assigns $-\hat g(\bm x)=\xi_{(k)}^{\mathrm{te}}$ to $\mathcal X_k$. The tail sum therefore includes threshold equalities even when $H$ has atoms.
%
%\begin{proposition}[Finite representation of the DW ambiguity set]
%  \label{prop-finite-rank-constraints}
%  Each acceptance set in \eqref{eq-dw-rank-acceptance} is a nonempty closed interval containing $q_j$. The DW ambiguity set has the finite representation
%  \begin{equation}
%    \mathcal Q^{\mathrm{DW}}(\theta)
%    =\left\{\mathbb Q\in\mathcal M_1(\Xi):
%      \sum_{j<2k}\mathbb Q(\Xi_j)\le u_k,\quad
%      \sum_{j\le2k}\mathbb Q(\Xi_j)\ge l_k,\quad k=1,\ldots,K\right\}.
%    \label{eq-finite-rank-constraints}
%  \end{equation}
%  The sequences $l_k$ and $u_k$ are nondecreasing and satisfy
%  \[
%    0\le l_k,u_k\le1,\qquad l_{k-1}\le u_k,\qquad l_k\le u_{k+1},
%    \qquad k=1,\ldots,K,
%  \]
%  where $l_0:=0$ and $u_{K+1}:=1$.
%\end{proposition}
%\begin{proof}
%  For fixed $q$, the Bernoulli divergence increases as $p$ moves away from $q$, while $\psi(q,p)$ cannot increase as the interval between $q$ and $p$ expands. The corrected discrepancy therefore decreases toward $q$ from below and increases above $q$. It is continuous for $p\in(0,1)$ and equals $-\psi_0(q)\le0$ at $p=q$. The boundary conventions in Section~\ref{sec-ambiguity-set} give a nonempty closed acceptance interval also for $q\in\{0,1\}$. Scalar inversion of the DW constraint therefore gives the CDF band \citep[Supplement, Section~S.6]{L.J23}
%  \[
%    \alpha_{N^{\mathrm{te}}\widehat F^{\mathrm{te}}(t)}(\theta)
%      \le F_{\mathbb Q}(t)
%      \le\beta_{N^{\mathrm{te}}\widehat F^{\mathrm{te}}(t)}(\theta),
%    \qquad\forall t\in\mathbb R.
%  \]
%
%  For $k=1,\ldots,K-1$, the empirical CDF on $\left[\xi_{(k)}^{\mathrm{te}},\xi_{(k+1)}^{\mathrm{te}}\right)$ equals $c_k/N^{\mathrm{te}}$. Monotonicity of a candidate CDF makes its lower band constraint equivalent to $F_{\mathbb Q}\left(\xi_{(k)}^{\mathrm{te}}\right)\ge\alpha_{c_k}(\theta)$ and its upper band constraint equivalent to $F_{\mathbb Q}\left(\xi_{(k+1)}^{\mathrm{te}}-\right)\le\beta_{c_k}(\theta)$. The exterior intervals give $F_{\mathbb Q}\left(\xi_{(1)}^{\mathrm{te}}-\right)\le\beta_0(\theta)$ and $F_{\mathbb Q}\left(\xi_{(K)}^{\mathrm{te}}\right)\ge\alpha_{N^{\mathrm{te}}}(\theta)$; their other band constraints are automatic because $\alpha_0(\theta)=0$ and $\beta_{N^{\mathrm{te}}}(\theta)=1$. The rank-probability identities above then yield \eqref{eq-finite-rank-constraints}.
%
%  With $p$ fixed, the corrected discrepancy decreases as $q$ approaches $p$ from either side. For empirical levels $q<q'$, an accepted value $p<q$ at level $q'$ remains accepted at level $q$; an accepted value $p>q'$ at level $q$ remains accepted at level $q'$. Since each acceptance interval contains its empirical level, both endpoints are nondecreasing. The inequalities $\alpha_{c_k}(\theta)\le c_k/N^{\mathrm{te}}\le\beta_{c_k}(\theta)$ give the stated compatibility conditions. Singleton intervals are allowed: if $q_j=1/2$ and $\theta=0$, the interval is $\{1/2\}$ because $\psi(1/2,p)=0$.
%\end{proof}
%
%Each scalar interval can be computed by one-dimensional bisection. Proposition~\ref{prop-finite-rank-constraints} then restricts only the probabilities of the residual cells; the distribution within each cell remains free.
%
%For a feasible joint law, define its matrix of joint cell probabilities by $\pi_{ij}:=\Pi(\mathcal X_i\times\Xi_j)$. The fixed input marginal and the residual restrictions imply $\bm\pi=(\pi_{ij})\in\mathcal C_\theta$, where
%\begin{equation}
%  \mathcal C_\theta:=\left\{\bm\pi\in\mathbb R_+^{(K+1)\times J}:
%  \begin{array}{ll}
%    \displaystyle\sum_{j=1}^{J}\pi_{ij}=w_i,&i=0,\ldots,K,\\[1mm]
%    \displaystyle\sum_{i=0}^{K}\sum_{j<2k}\pi_{ij}\le u_k,&k=1,\ldots,K,\\[1mm]
%    \displaystyle\sum_{i=0}^{K}\sum_{j\le2k}\pi_{ij}\ge l_k,&k=1,\ldots,K
%  \end{array}\right\}.
%  \label{eq-cell-mass-polytope}
%\end{equation}
%On each product cell, the failure indicator satisfies
%\begin{equation}
%  \underline h_{ij}:=\mathbf1_{\{j\le2i\}}
%  \le\mathbf1_{\{\xi\le-\hat g(\bm x)\}}
%  \le\overline h_{ij}:=\mathbf1_{\{j\le2i+1\}},
%  \qquad (\bm x,\xi)\in\mathcal X_i\times\Xi_j.
%  \label{eq-cell-costs}
%\end{equation}
%Residual cells with $j\le2i$ always cause failure for inputs in $\mathcal X_i$, and those with $j\ge2i+2$ are always safe. Only the overlapping interval $\Xi_{2i+1}$ has different upper and lower coefficients. Thus the cell matrix records the rank constraints exactly, while the failure objective requires the constructions below.
%
%The upper problem remains a linear program over joint measures, with the full input marginal fixed:
%\begin{equation}
%  \begin{aligned}
%    \overline p_{\hat f}
%    =\sup_{\Pi\in\mathcal M_+(\mathcal X\times\Xi)}\quad&
%      \int_{\mathcal X\times\Xi}\mathbf1_{\{\xi\le-\hat g(\bm x)\}}\,\Pi(\mathrm d\bm x,\mathrm d\xi)\\
%    \text{s.t.}\quad&
%      \int_{\mathcal X\times\Xi}\mathbf1_{\mathcal U}(\bm x)\,\Pi(\mathrm d\bm x,\mathrm d\xi)
%        =\mathbb P(\bm X\in\mathcal U),&&\forall\mathcal U\in\mathcal B(\mathcal X),\\
%    &\int_{\mathcal X\times\Xi}\mathbf1_{\{\xi<\xi_{(k)}^{\mathrm{te}}\}}\,\Pi(\mathrm d\bm x,\mathrm d\xi)\le u_k,
%        &&k=1,\ldots,K,\\
%    &\int_{\mathcal X\times\Xi}\mathbf1_{\{\xi\le\xi_{(k)}^{\mathrm{te}}\}}\,\Pi(\mathrm d\bm x,\mathrm d\xi)\ge l_k,
%        &&k=1,\ldots,K.
%  \end{aligned}
%  \label{eq-measure-upper}
%\end{equation}
%The lower problem minimizes the same integral under the same constraints. The marginal equalities fix both the input law and total mass one. The product of $\mathbb P$ and the empirical residual law is feasible for every $\theta\ge0$.
%
%\begin{lemma}[Finite-dimensional Lagrangian duals and weak duality]
%  \label{lem-weak-duality}
%  Write $\bm u=(u_1,\ldots,u_K)^\top$ and $\bm l=(l_1,\ldots,l_K)^\top$. The nonnegative multipliers $\bm\lambda,\bm\gamma\in\mathbb R_+^K$ correspond to the upper and lower rank constraints. The input marginal equality has an unrestricted multiplier $\eta\in\mathcal F_{\mathrm b}(\mathcal X)$, the space of bounded Borel measurable real-valued functions. Define the residual penalty
%  \begin{equation}
%    \zeta_{\bm\lambda,\bm\gamma}(\xi)
%    :=\sum_{k=1}^{K}\lambda_k\mathbf1_{\{\xi<\xi_{(k)}^{\mathrm{te}}\}}
%      -\sum_{k=1}^{K}\gamma_k\mathbf1_{\{\xi\le\xi_{(k)}^{\mathrm{te}}\}},
%    \label{eq-dual-step-function}
%  \end{equation}
%  whose value on $\Xi_j$ is
%  \begin{equation}
%    \zeta_j:=\sum_{k=1}^{K}\lambda_k\mathbf1_{\{j<2k\}}
%      -\sum_{k=1}^{K}\gamma_k\mathbf1_{\{j\le2k\}},
%    \qquad j=1,\ldots,J.
%    \label{eq-dual-cell-values}
%  \end{equation}
%  The upper and lower Lagrangian dual values are given by the finite LPs
%  \begin{equation}
%    \begin{aligned}
%      \overline d_\theta
%      :=\inf_{\bm\lambda,\bm\gamma\ge\bm0,\ \bm z\in\mathbb R^{K+1}}\quad&
%        \bm\lambda^\top\bm u-\bm\gamma^\top\bm l+\sum_{i=0}^{K}w_i z_i\\
%      \text{s.t.}\quad&z_i+\zeta_j\ge\overline h_{ij},
%        \qquad i=0,\ldots,K,\quad j=1,\ldots,J,
%    \end{aligned}
%    \label{eq-finite-upper-dual}
%  \end{equation}
%  and
%  \begin{equation}
%    \begin{aligned}
%      \underline d_\theta
%      :=\sup_{\bm\lambda,\bm\gamma\ge\bm0,\ \bm z\in\mathbb R^{K+1}}\quad&
%        -\bm\lambda^\top\bm u+\bm\gamma^\top\bm l+\sum_{i=0}^{K}w_i z_i\\
%      \text{s.t.}\quad&z_i-\zeta_j\le\underline h_{ij},
%        \qquad i=0,\ldots,K,\quad j=1,\ldots,J.
%    \end{aligned}
%    \label{eq-finite-lower-dual}
%  \end{equation}
%  Their optimal values satisfy
%  \begin{equation}
%    \underline d_\theta\le\underline p_{\hat f}
%      \le\overline p_{\hat f}\le\overline d_\theta.
%    \label{eq-weak-duality}
%  \end{equation}
%\end{lemma}
%\begin{proof}
%  The input marginal equality is equivalent to
%  \[
%    \int_{\mathcal X\times\Xi}\eta(\bm x)\,\Pi(\mathrm d\bm x,\mathrm d\xi)
%      =\mathbb E_{\mathbb P}\left[\eta(\bm X)\right],
%    \qquad\forall\eta\in\mathcal F_{\mathrm b}(\mathcal X).
%  \]
%  Relaxing this equality and the rank inequalities gives the upper Lagrangian \citep[Section~5.1]{B.V04}
%  \[
%    \begin{aligned}
%      \mathcal L^+(\Pi,\eta,\bm\lambda,\bm\gamma)
%      ={}&\int_{\mathcal X\times\Xi}\mathbf1_{\{\xi\le-\hat g(\bm x)\}}\,\Pi(\mathrm d\bm x,\mathrm d\xi)
%        +\mathbb E_{\mathbb P}\left[\eta(\bm X)\right]
%        -\int_{\mathcal X\times\Xi}\eta(\bm x)\,\Pi(\mathrm d\bm x,\mathrm d\xi)\\
%      &+\sum_{k=1}^{K}\lambda_k\left(u_k-\int_{\mathcal X\times\Xi}
%          \mathbf1_{\{\xi<\xi_{(k)}^{\mathrm{te}}\}}\,\Pi(\mathrm d\bm x,\mathrm d\xi)\right)\\
%      &+\sum_{k=1}^{K}\gamma_k\left(\int_{\mathcal X\times\Xi}
%          \mathbf1_{\{\xi\le\xi_{(k)}^{\mathrm{te}}\}}\,\Pi(\mathrm d\bm x,\mathrm d\xi)-l_k\right)\\
%      ={}&\mathbb E_{\mathbb P}\left[\eta(\bm X)\right]
%        +\bm\lambda^\top\bm u-\bm\gamma^\top\bm l\\
%      &+\int_{\mathcal X\times\Xi}\left[
%          \mathbf1_{\{\xi\le-\hat g(\bm x)\}}-\eta(\bm x)-\zeta_{\bm\lambda,\bm\gamma}(\xi)
%        \right]\,\Pi(\mathrm d\bm x,\mathrm d\xi).
%    \end{aligned}
%  \]
%  Its supremum over finite nonnegative measures is finite exactly when the integrand is nonpositive at every $(\bm x,\xi)$. A positive value at any point permits an arbitrarily large point mass there; otherwise the supremum integral is zero, attained by the zero measure. The functional upper dual therefore minimizes $\mathbb E_{\mathbb P}[\eta(\bm X)]+\bm\lambda^\top\bm u-\bm\gamma^\top\bm l$ subject to $\eta(\bm x)\ge\mathbf1_{\{\xi\le-\hat g(\bm x)\}}-\zeta_{\bm\lambda,\bm\gamma}(\xi)$ at every pair. For the lower problem, reversing the rank penalties gives
%  \[
%    \mathcal L^-
%      =\mathbb E_{\mathbb P}\left[\eta(\bm X)\right]
%        -\bm\lambda^\top\bm u+\bm\gamma^\top\bm l
%        +\int_{\mathcal X\times\Xi}\left[
%          \mathbf1_{\{\xi\le-\hat g(\bm x)\}}-\eta(\bm x)+\zeta_{\bm\lambda,\bm\gamma}(\xi)
%        \right]\,\Pi(\mathrm d\bm x,\mathrm d\xi).
%  \]
%  Its infimum over finite nonnegative measures is finite exactly when the integrand is nonnegative everywhere, in which case the infimum integral is zero. The lower functional dual maximizes $\mathbb E_{\mathbb P}[\eta(\bm X)]-\bm\lambda^\top\bm u+\bm\gamma^\top\bm l$ subject to $\eta(\bm x)\le\mathbf1_{\{\xi\le-\hat g(\bm x)\}}+\zeta_{\bm\lambda,\bm\gamma}(\xi)$ at every pair.
%
%  For fixed rank multipliers, the best marginal functions are the pointwise envelopes. For $\bm x\in\mathcal X_i$, they satisfy the identities
%  \begin{equation}
%    \begin{aligned}
%      \eta^+_{\bm\lambda,\bm\gamma}(\bm x)
%        &:=\sup_{\xi\in\Xi}\left\{\mathbf1_{\{\xi\le-\hat g(\bm x)\}}
%          -\zeta_{\bm\lambda,\bm\gamma}(\xi)\right\}
%          =\max_{1\le j\le J}\left\{\overline h_{ij}-\zeta_j\right\},\\
%      \eta^-_{\bm\lambda,\bm\gamma}(\bm x)
%        &:=\inf_{\xi\in\Xi}\left\{\mathbf1_{\{\xi\le-\hat g(\bm x)\}}
%          +\zeta_{\bm\lambda,\bm\gamma}(\xi)\right\}
%          =\min_{1\le j\le J}\left\{\underline h_{ij}+\zeta_j\right\}.
%    \end{aligned}
%    \label{eq-dual-envelope-identities}
%  \end{equation}
%  To verify them, note that outside $\Xi_{2i+1}$ the two cell coefficients equal the actual failure indicator. In the overlapping interval, the lower coefficient zero is attained by $\xi=\left(\xi_{(i+1)}^{\mathrm{te}}-\hat g(\bm x)\right)/2$ if $i<K$, and by $\xi=-\hat g(\bm x)+1$ if $i=K$. The upper coefficient one is attained by $\xi=-\hat g(\bm x)$ whenever $-\hat g(\bm x)>\xi_{(i)}^{\mathrm{te}}$. The only remaining case is $-\hat g(\bm x)=\xi_{(i)}^{\mathrm{te}}$ with $i\ge1$. Then the overlapping interval is safe, but
%  \[
%    \zeta_{2i}=\zeta_{2i+1}-\gamma_i\le\zeta_{2i+1},
%    \qquad
%    1-\zeta_{2i}\ge1-\zeta_{2i+1}.
%  \]
%  Choosing the left endpoint therefore attains an upper-envelope value at least as large as the candidate from the overlapping interval. Together with \eqref{eq-cell-costs}, this proves both identities, including threshold equalities.
%
%  The envelopes are bounded and constant on the Borel cells $\mathcal X_i$, so they belong to $\mathcal F_{\mathrm b}(\mathcal X)$. They respectively minimize and maximize the expectation over feasible marginal functions. Substitution yields
%  \begin{align}
%    \overline d_\theta
%      &=\inf_{\bm\lambda,\bm\gamma\ge\bm0}\left\{
%        \bm\lambda^\top\bm u-\bm\gamma^\top\bm l
%        +\sum_{i=0}^{K}w_i\max_{1\le j\le J}
%          \left\{\overline h_{ij}-\zeta_j\right\}\right\},
%      \label{eq-upper-measure-dual}\\
%    \underline d_\theta
%      &=\sup_{\bm\lambda,\bm\gamma\ge\bm0}\left\{
%        -\bm\lambda^\top\bm u+\bm\gamma^\top\bm l
%        +\sum_{i=0}^{K}w_i\min_{1\le j\le J}
%          \left\{\underline h_{ij}+\zeta_j\right\}\right\}.
%      \label{eq-lower-measure-dual}
%  \end{align}
%  Replacing each maximum by an epigraph variable $z_i$, and each minimum by a hypograph variable $z_i$, gives \eqref{eq-finite-upper-dual} and \eqref{eq-finite-lower-dual}. This remains exact when $w_i=0$, since $z_i$ can still be chosen as the corresponding finite envelope.
%
%  For a feasible joint law, the signs of the rank penalties give $\mathcal L^-\le p_{\hat f}(\Pi)\le\mathcal L^+$. Feasible lower dual objectives therefore bound the primal objectives from below, and feasible upper dual objectives bound them from above. Taking the respective optima over the common nonempty primal feasible set proves \eqref{eq-weak-duality}.
%\end{proof}
%
%Both duals have $3K+1$ scalar variables and $(K+1)(2K+1)$ pairwise inequalities, together with multiplier nonnegativity. The variables $z_i$ represent the optimal marginal function on each input cell. The envelope identities show that these finite LPs are the full Lagrangian duals; their derivation does not require strong duality.
%
%\begin{theorem}[Exact finite reformulations and strong duality]
%  \label{thm-finite-dual-reformulations}
%  For $\mathcal Q=\mathcal Q^{\mathrm{DW}}(\theta)$ and every $\theta\ge0$, the failure probability bounds admit the exact representations
%  \begin{equation}
%    \overline p_{\hat f}
%      =\max_{\bm\pi\in\mathcal C_\theta}\sum_{i=0}^{K}\sum_{j=1}^{J}\overline h_{ij}\pi_{ij},
%    \qquad
%    \underline p_{\hat f}
%      =\min_{\bm\pi\in\mathcal C_\theta}\sum_{i=0}^{K}\sum_{j=1}^{J}\underline h_{ij}\pi_{ij}.
%    \label{eq-finite-primal-values}
%  \end{equation}
%  Their Lagrangian dual values in Lemma~\ref{lem-weak-duality} satisfy
%  \begin{equation}
%    \overline p_{\hat f}=\overline d_\theta,
%    \qquad
%    \underline p_{\hat f}=\underline d_\theta.
%    \label{eq-strong-duality}
%  \end{equation}
%  The finite primal and dual optima, the measure primal optima, and the functional dual optima are attained.
%\end{theorem}
%\begin{proof}
%  Any feasible joint law induces a matrix in $\mathcal C_\theta$. Integrating \eqref{eq-cell-costs} therefore bounds the upper measure problem by the finite maximization from above and the lower measure problem by the finite minimization from below.
%
%  For the reverse inequalities, fix $\bm\pi\in\mathcal C_\theta$. Draw $\bm x$ from $\mathbb P$. On $\mathcal X_i$ with $w_i>0$, select a provisional residual cell $j$ with probability $\pi_{ij}/w_i$; the row equality makes these probabilities sum to one. For $j\ne2i+1$, choose a fixed representative of $\Xi_j$, where the actual failure indicator equals both coefficients. In the overlapping cell, use the following assignments to define joint laws $\Pi^+$ and $\Pi^-$.
%
%  For the upper problem, assign $\xi=-\hat g(\bm x)$. This causes failure and remains in $\Xi_{2i+1}$ unless $-\hat g(\bm x)=\xi_{(i)}^{\mathrm{te}}$. At equality, which requires $i\ge1$, the assignment transfers the provisional mass to the left endpoint. For any $\widetilde\xi\in\left(\xi_{(i)}^{\mathrm{te}},\xi_{(i+1)}^{\mathrm{te}}\right)$ and every observed threshold,
%  \[
%    \mathbf1_{\{\xi_{(i)}^{\mathrm{te}}<\xi_{(k)}^{\mathrm{te}}\}}
%      =\mathbf1_{\{\widetilde\xi<\xi_{(k)}^{\mathrm{te}}\}},
%    \qquad
%    \mathbf1_{\{\xi_{(i)}^{\mathrm{te}}\le\xi_{(k)}^{\mathrm{te}}\}}
%      \ge\mathbf1_{\{\widetilde\xi\le\xi_{(k)}^{\mathrm{te}}\}}.
%  \]
%  Thus every open upper probability is unchanged and every closed lower probability can only increase. The upper assignment preserves all rank inequalities and attains the coefficient $\overline h_{i,2i+1}=1$.
%
%  For the lower problem, assign $\xi=\left(\xi_{(i+1)}^{\mathrm{te}}-\hat g(\bm x)\right)/2$ if $i<K$, and $\xi=-\hat g(\bm x)+1$ if $i=K$. These values remain in $\Xi_{2i+1}$ and strictly exceed the failure threshold, including at threshold equality. Hence the lower assignment retains the cell matrix and attains $\underline h_{i,2i+1}=0$.
%
%  If $w_i=0$, the corresponding row is zero; set $\xi=\xi_{(1)}^{\mathrm{te}}$ on that input cell. All assignments are Borel measurable, and their finite mixtures define probability kernels. Combining them with $\mathbb P$ gives joint laws with the full prescribed input marginal, including randomization at input atoms. Both laws are feasible and satisfy
%  \[
%    p_{\hat f}(\Pi^+)=\sum_{i=0}^{K}\sum_{j=1}^{J}\overline h_{ij}\pi_{ij},
%    \qquad
%    p_{\hat f}(\Pi^-)=\sum_{i=0}^{K}\sum_{j=1}^{J}\underline h_{ij}\pi_{ij}.
%  \]
%  The lower law retains every provisional cell mass. The upper law may change the cell matrix at threshold equalities, but it preserves feasibility and the assigned upper objective. These identities establish the reverse inequalities required for \eqref{eq-finite-primal-values}.
%
%  For every $\theta\ge0$, the empirical product law induces the feasible matrix
%  \[
%    \pi_{i,2k}^{\mathrm{emp}}:=w_i\frac{c_k-c_{k-1}}{N^{\mathrm{te}}},
%    \qquad i=0,\ldots,K,\quad k=1,\ldots,K,
%  \]
%  with all odd columns zero. Its open and closed rank probabilities are $c_{k-1}/N^{\mathrm{te}}$ and $c_k/N^{\mathrm{te}}$, respectively, so feasibility includes $\theta=0$. The nonnegative entries of a matrix in $\mathcal C_\theta$ sum to one, and the constraints are closed. Thus $\mathcal C_\theta$ is nonempty and compact. The finite objectives lie in $[0,1]$ and attain their optima, which the joint-law constructions realize in the measure problems.
%
%  The finite primal LPs in \eqref{eq-finite-primal-values} have row multipliers $z_i$ and rank multipliers $\lambda_k,\gamma_k$. Their dual constraints are precisely $z_i+\zeta_j\ge\overline h_{ij}$ and $z_i-\zeta_j\le\underline h_{ij}$. Finite LP strong duality therefore gives \eqref{eq-strong-duality} and dual attainment \citep[Section~5.2.3]{B.V04}. By \eqref{eq-dual-envelope-identities}, optimal rank multipliers also yield optimal bounded marginal functions for the original functional duals. Feasibility and finite optimal values suffice; strict rank inequalities are unnecessary, including at $\theta=0$.
%\end{proof}
%
%The finite formulations are exact, and their ordered failure coefficients give a further simplification. The closed-form bounds below follow from cuts between consecutive cells and admit certificates using at most one rank multiplier.
%
%\begin{corollary}[Closed-form bounds and optimal certificates]
%  \label{thm-closed-form-bounds}
%  Under the conditions of Theorem~\ref{thm-finite-dual-reformulations},
%  \begin{align}
%    \overline p_{\hat f}
%      &=\min\left\{1,\min_{k=1,\ldots,K}\left[u_k+H\left(-\xi_{(k)}^{\mathrm{te}}\right)\right]\right\},
%      \label{eq-dw-closed-upper}\\
%    \underline p_{\hat f}
%      &=\max\left\{0,\max_{k=1,\ldots,K}\left[l_k+H\left(-\xi_{(k)}^{\mathrm{te}}\right)-1\right]\right\}.
%      \label{eq-dw-closed-lower}
%  \end{align}
%  Each dual optimum admits a certificate with at most one nonzero rank multiplier, of value one.
%\end{corollary}
%\begin{proof}
%  Fix nonnegative column masses $\bm y=(y_1,\ldots,y_J)^\top$ with $\sum_j y_j=1$, and write $Y_j:=\sum_{r=1}^{j}y_r$. Let $\mathcal T(\bm w,\bm y)$ be the transportation matrices with row sums $\bm w=(w_0,\ldots,w_K)^\top$ and column sums $\bm y$. This set contains the product matrix $\pi_{ij}=w_i y_j$. The residual constraints depend only on the column masses, through $Y_{2k-1}\le u_k$ and $Y_{2k}\ge l_k$.
%
%  For fixed $\bm y$, the upper problem maximizes mass on pairs $j\le2i+1$; the lower problem maximizes safe mass on pairs $j>2i$ and subtracts it from one. Both are maximum-flow problems with source-to-row capacities $w_i$, column-to-sink capacities $y_j$, and capacity two on each permitted row-to-column edge. A transportation matrix restricts to a feasible flow on permitted pairs. Conversely, any flow can be completed to a transportation matrix because the remaining row and column masses are nonnegative with equal totals. Thus the maximum flow equals the maximum permitted mass.
%
%  A cut of capacity one is available, so a minimum cut crosses no edge of capacity two. Its source-side columns can be taken to be exactly the neighbors of its source-side rows. Since the neighborhoods are nested, adding rows whose neighbors are already included cannot increase the capacity. Apart from empty and full row sets, upper-problem cuts therefore reduce to rows $i<k$ and columns $j<2k$, while safe-mass cuts reduce to rows $i\ge k$ and columns $j>2k$. Their capacities, for $k=1,\ldots,K$, are respectively
%  \begin{align*}
%    \sum_{i=k}^{K}w_i+Y_{2k-1}
%      &=H\left(-\xi_{(k)}^{\mathrm{te}}\right)+Y_{2k-1},\\
%    \sum_{i=0}^{k-1}w_i+1-Y_{2k}
%      &=2-H\left(-\xi_{(k)}^{\mathrm{te}}\right)-Y_{2k}.
%  \end{align*}
%  Empty and full row sets give capacity one. The max-flow/min-cut theorem \citep{F.F56} consequently yields
%  \begin{align*}
%    \max_{\bm\pi\in\mathcal T(\bm w,\bm y)}\sum_{i,j}\overline h_{ij}\pi_{ij}
%      &=\min\left\{1,\min_k\left[H\left(-\xi_{(k)}^{\mathrm{te}}\right)+Y_{2k-1}\right]\right\},\\
%    \min_{\bm\pi\in\mathcal T(\bm w,\bm y)}\sum_{i,j}\underline h_{ij}\pi_{ij}
%      &=\max\left\{0,\max_k\left[H\left(-\xi_{(k)}^{\mathrm{te}}\right)+Y_{2k}-1\right]\right\}.
%  \end{align*}
%
%  These values are nondecreasing in their respective open and closed prefixes. To attain all upper prefix limits for the maximization and all lower prefix limits for the minimization, choose
%  \begin{align*}
%    &y_1^+:=u_1,\qquad y_{2k}^+:=u_{k+1}-u_k,\qquad
%      y_{2k+1}^+:=0,\quad k=1,\ldots,K,\\
%    &y_{2k-1}^-:=0,\qquad y_{2k}^-:=l_k-l_{k-1},\quad k=1,\ldots,K,
%      \qquad y_J^-:=1-l_K.
%  \end{align*}
%  The monotonicity and compatibility conditions in Proposition~\ref{prop-finite-rank-constraints} make both vectors nonnegative with total mass one, and
%  \[
%    Y_{2k-1}^+=u_k,\qquad Y_{2k}^+=u_{k+1}\ge l_k,
%    \qquad Y_{2k}^-=l_k,\qquad Y_{2k-1}^-=l_{k-1}\le u_k.
%  \]
%  Both allocations are feasible and attain the required prefix limits. Substitution into the fixed-column values, followed by \eqref{eq-finite-primal-values}, proves \eqref{eq-dw-closed-upper} and \eqref{eq-dw-closed-lower}. The construction allows zero masses and coincident prefix limits, including at $\theta=0$.
%
%  Finally, let $\bm e_k$ be the $k$th unit vector. Upper and lower dual certificates are respectively
%  \begin{align*}
%    &\bm\lambda=\bm e_k,\qquad\bm\gamma=\bm0,\qquad
%      z_i=\mathbf1_{\{i\ge k\}},\\
%    &\bm\lambda=\bm0,\qquad\bm\gamma=\bm e_k,\qquad
%      z_i=\mathbf1_{\{i\ge k\}}-1.
%  \end{align*}
%  Their feasibility follows from the index inequalities
%  \[
%    \mathbf1_{\{i\ge k\}}+\mathbf1_{\{j<2k\}}\ge\overline h_{ij},
%    \qquad
%    \mathbf1_{\{i\ge k\}}-1+\mathbf1_{\{j\le2k\}}\le\underline h_{ij}.
%  \]
%  Their objectives are $u_k+\sum_{i=k}^{K}w_i$ and $l_k+\sum_{i=k}^{K}w_i-1$, the inner candidates in the closed-form expressions. Zero rank multipliers with $z_i=1$ or $z_i=0$ supply the outer candidates. Selecting the best candidate gives an optimal full dual solution. The envelope identities recover its functional marginal multiplier on each input cell; the tail sum includes $i=k$ and therefore uses the closed CDF value $H\left(-\xi_{(k)}^{\mathrm{te}}\right)$ also at threshold equality.
%\end{proof}
%
%The same formulas give the response-threshold curves in Corollary~\ref{cor-simultaneous-threshold-bounds}. Replacing $g$ and $\hat g$ by $g-c$ and $\hat g-c$ leaves every residual, rank count, and endpoint $l_k,u_k$ unchanged. The surrogate CDF becomes $t\mapsto H(t+c)$, so Corollary~\ref{thm-closed-form-bounds} gives
%\begin{align}
%  \overline p_{\hat f}(c)
%    &=\min\left\{1,\min_{k=1,\ldots,K}\left[u_k+H\left(c-\xi_{(k)}^{\mathrm{te}}\right)\right]\right\},
%    \label{eq-threshold-closed-upper}\\
%  \underline p_{\hat f}(c)
%    &=\max\left\{0,\max_{k=1,\ldots,K}\left[l_k+H\left(c-\xi_{(k)}^{\mathrm{te}}\right)-1\right]\right\}.
%    \label{eq-threshold-closed-lower}
%\end{align}
%These identities hold for every $c\in\mathbb R$ and $\theta\ge0$, including $\theta=0$. When $\theta\ge\theta_{N^{\mathrm{te}}}(\delta)$, their coverage is simultaneous by \eqref{eq-threshold-curve-coverage}. An extremizing coupling may depend on $c$; the common feasible true joint law supplies the simultaneous guarantee.
%
%The upper bound combines residual mass below an observed value with the probability that the surrogate threshold reaches it. The lower bound measures the overlap required by every admissible coupling. For $p_{\mathrm{tar}}\in(0,1)$, the upper curve satisfies $\overline p_{\hat f}(c)\le p_{\mathrm{tar}}$ exactly when some rank satisfies $u_k+H\left(c-\xi_{(k)}^{\mathrm{te}}\right)\le p_{\mathrm{tar}}$. Both terms must therefore be small at the same residual level.
%
%With a known residual law, the coupling problem is the classical fixed-marginal sum problem studied by \citet{M82} and \citet{F.N.S87}. Here the residual law also varies within a calibrated DW band, and the open and closed rank constraints retain atoms.
%
%Dependence uncertainty can remain even with a known residual marginal. For equally likely surrogate responses $0$ and $2$ and equally likely residuals $-1$ and $1$, failure occurs exactly when the response is $0$ and the residual is $-1$. Its probability ranges from $0$ to $1/2$ over couplings, whereas independence gives $1/4$. Further marginal calibration alone does not resolve this dependence uncertainty.
%
%\begin{remark}[Known interval support]
%  \label{rem-interval-support}
%  Suppose $\Xi=[a,b]$ with $-\infty<a\le b<\infty$ and the observed residuals lie in this interval. Then
%  \begin{align*}
%    \overline p_{\hat f}
%    &=\min\left\{H(-a),\min_{k=1,\ldots,K}[u_k+H\left(-\xi_{(k)}^{\mathrm{te}}\right)]\right\},\\
%    \underline p_{\hat f}
%    &=\max\left\{H(-b),\max_{k=1,\ldots,K}[l_k+H\left(-\xi_{(k)}^{\mathrm{te}}\right)-1]\right\}.
%  \end{align*}
%  To obtain these expressions, represent the support restriction on $\mathbb R$ by
%  $\mathbb Q\{\xi<a\}\le0$ and $\mathbb Q\{\xi\le b\}\ge1$.
%  Add $a,b$ to the ordered threshold list. At a newly added $a$, use upper and lower bounds zero; at a newly added $b$, use upper and lower bounds one. At coincident observed thresholds, tighten only the upper bound at $a$ and the lower bound at $b$. The enlarged sequences satisfy the compatibility conditions in Proposition~\ref{prop-finite-rank-constraints}, so the same finite-aggregation and ordered-cut arguments apply. The effective additional upper candidate is $H(-a)$, and the effective additional lower candidate is $H(-b)$; the other added candidates are redundant. If $a=b$, both bounds equal $H(-a)$. For the resulting partition on $\Xi=[a,b]$, intersect the singleton and open-interval residual cells from the augmented threshold list with $\Xi$ and omit empty cells. Including the support endpoints in this list retains their atoms and the corresponding changes in failure status. For a response threshold $c$, replace $H(-a)$, $H(-b)$, and $H\left(-\xi_{(k)}^{\mathrm{te}}\right)$ by $H(c-a)$, $H(c-b)$, and $H\left(c-\xi_{(k)}^{\mathrm{te}}\right)$, respectively.
%\end{remark}
%
%After zero failures in $N^{\mathrm{te}}$ independent tests, the one-sided exact binomial upper limit is $1-\delta^{1/N^{\mathrm{te}}}\approx\log(1/\delta)/N^{\mathrm{te}}$ \citep{C.P34}. At $c=0$, the simple DW case $K=1$ and $H\left(-\xi_{(1)}^{\mathrm{te}}\right)=0$ gives an upper bound $\beta_0(\theta)$ determined by the simultaneous-band calibration. Verified support can exclude otherwise unobserved failures: in Remark~\ref{rem-interval-support}, $H(-a)=0$ rules out failure.
%
%Evaluation requires scalar DW inversion, sorting and counting the test residuals, and computing the $K$ values $H\left(c-\xi_{(k)}^{\mathrm{te}}\right)$ at each response threshold $c$. The endpoint table depends only on $N^{\mathrm{te}}$ and $\theta$ and can be precomputed. Given this table and the surrogate CDF values, the two scans cost $O(K)$ at each response threshold.
%
%For $\theta\ge\theta_{N^{\mathrm{te}}}(\delta)$, the exact expressions retain the coverage in Theorem~\ref{theorem-01} and Corollary~\ref{cor-simultaneous-threshold-bounds}. If $H$ is evaluated by independent auxiliary simulation, upper CDF bounds can replace its values in the upper expression, and lower CDF bounds can replace them in the lower expression. For one fixed response threshold, simultaneous bounds at its $K$ evaluation points suffice. To retain coverage of the entire response-threshold curve, the auxiliary bounds must hold over all arguments of $H$. In either case, they must hold conditionally on the test data, with their error probability included in the overall confidence budget.
%
%
% END PRESERVED SECTION 4 BEFORE DIRECT CELL-PROBABILITY NOTATION

% BEGIN PRESERVED SECTION 4 BEFORE PARTITION AND EVENT PROBABILITIES (2026-11-25)
%\section{Tractable Reformulation}
%\label{sec-reformulation}
%
%Throughout this section, the training data, the test residuals, and $\theta\ge0$ are fixed. We use the observed residuals to partition both the residual and input spaces. These partitions express the DW restrictions as cumulative cell probabilities and reduce the Lagrangian duals to finite linear programs. A joint-distribution construction establishes exact finite primal representations, after which LP duality and the ordering of the cells yield strong duality and closed-form bounds. The main development takes $\Xi=\mathbb R$; Remark~\ref{rem-interval-support} treats known interval support.
%
%For each empirical probability level $q_j:=j/N^{\mathrm{te}}$, define the scalar acceptance interval
%\begin{equation}
%  [\alpha_j(\theta),\beta_j(\theta)]
%  :=\left\{p\in[0,1]:
%  N^{\mathrm{te}}\phi(q_j,p)-\psi(q_j,p)\le\theta\right\},
%  \qquad j=0,\ldots,N^{\mathrm{te}}.
%  \label{eq-dw-rank-acceptance}
%\end{equation}
%Let $\xi_{(1)}^{\mathrm{te}}<\cdots<\xi_{(K)}^{\mathrm{te}}$ be the distinct ordered test residuals, and set
%\begin{equation}
%  c_0:=0,\quad
%  c_k:=\sum_{j=1}^{N^{\mathrm{te}}}\mathbf1_{\{\xi_j^{\mathrm{te}}\le\xi_{(k)}^{\mathrm{te}}\}},\quad
%  l_k:=\alpha_{c_k}(\theta),\quad
%  u_k:=\beta_{c_{k-1}}(\theta),\qquad k=1,\ldots,K.
%  \label{eq-dw-counts-and-levels}
%\end{equation}
%The counts retain all multiplicities, so $c_K=N^{\mathrm{te}}$ even when test residuals coincide. Set $\xi_{(0)}^{\mathrm{te}}:=-\infty$ and $\xi_{(K+1)}^{\mathrm{te}}:=+\infty$. Separating the observed values from the intervening intervals gives $J:=2K+1$ residual cells:
%\begin{equation}
%  \Xi_{2k}:=\left\{\xi_{(k)}^{\mathrm{te}}\right\}
%    \quad\left(k=1,\ldots,K\right),
%  \qquad
%  \Xi_{2k+1}:=\left(\xi_{(k)}^{\mathrm{te}},\xi_{(k+1)}^{\mathrm{te}}\right)
%    \quad\left(k=0,\ldots,K\right).
%  \label{eq-residual-cells}
%\end{equation}
%For any candidate residual law $\mathbb Q\in\mathcal M_1(\Xi)$, the open rank probability measures mass strictly below an observed residual, while the closed rank probability also includes mass at that value. For $k=1,\ldots,K$, the partition gives
%\[
%  \mathbb Q\left\{\xi<\xi_{(k)}^{\mathrm{te}}\right\}
%    =\sum_{j<2k}\mathbb Q(\Xi_j),
%  \qquad
%  \mathbb Q\left\{\xi\le\xi_{(k)}^{\mathrm{te}}\right\}
%    =\sum_{j\le2k}\mathbb Q(\Xi_j).
%\]
%The two probabilities differ by the atom mass $\mathbb Q(\Xi_{2k})$.
%
%Define the surrogate-response CDF by $H(t):=\mathbb P\left\{\hat g(\bm X)\le t\right\}$, with $H(-\infty):=0$ and $H(+\infty):=1$. For $i=0,\ldots,K$, the same residual levels classify inputs by the location of the failure threshold $-\hat g(\bm x)$:
%\begin{equation}
%  \mathcal X_i:=\left\{\bm x\in\mathcal X:
%    \xi_{(i)}^{\mathrm{te}}\le-\hat g(\bm x)<\xi_{(i+1)}^{\mathrm{te}}\right\},
%  \qquad
%  \mathbb P(\mathcal X_i)
%    =H\left(-\xi_{(i)}^{\mathrm{te}}\right)-H\left(-\xi_{(i+1)}^{\mathrm{te}}\right).
%  \label{eq-input-cells}
%\end{equation}
%Thus the two spaces are partitioned as
%\[
%  \mathcal X=\bigsqcup_{i=0}^{K}\mathcal X_i,
%  \qquad
%  \Xi=\bigsqcup_{j=1}^{J}\Xi_j,
%  \qquad
%  H\left(-\xi_{(k)}^{\mathrm{te}}\right)=\sum_{i=k}^{K}\mathbb P(\mathcal X_i),
%  \quad k=1,\ldots,K.
%\]
%The half-open input convention assigns $-\hat g(\bm x)=\xi_{(k)}^{\mathrm{te}}$ to $\mathcal X_k$. The tail sum therefore includes threshold equalities even when $H$ has atoms.
%
%\begin{proposition}[Finite representation of the DW ambiguity set]
%  \label{prop-finite-rank-constraints}
%  Each acceptance set in \eqref{eq-dw-rank-acceptance} is a nonempty closed interval containing $q_j$. The DW ambiguity set has the finite representation
%  \begin{equation}
%    \mathcal Q^{\mathrm{DW}}(\theta)
%    =\left\{\mathbb Q\in\mathcal M_1(\Xi):
%      \sum_{j<2k}\mathbb Q(\Xi_j)\le u_k,\quad
%      \sum_{j\le2k}\mathbb Q(\Xi_j)\ge l_k,\quad k=1,\ldots,K\right\}.
%    \label{eq-finite-rank-constraints}
%  \end{equation}
%  The sequences $l_k$ and $u_k$ are nondecreasing and satisfy
%  \[
%    0\le l_k,u_k\le1,\qquad l_{k-1}\le u_k,\qquad l_k\le u_{k+1},
%    \qquad k=1,\ldots,K,
%  \]
%  where $l_0:=0$ and $u_{K+1}:=1$.
%\end{proposition}
%\begin{proof}
%  For fixed $q$, the Bernoulli divergence increases as $p$ moves away from $q$, while $\psi(q,p)$ cannot increase as the interval between $q$ and $p$ expands. The corrected discrepancy therefore decreases toward $q$ from below and increases above $q$. It is continuous for $p\in(0,1)$ and equals $-\psi_0(q)\le0$ at $p=q$. The boundary conventions in Section~\ref{sec-ambiguity-set} give a nonempty closed acceptance interval also for $q\in\{0,1\}$. Scalar inversion of the DW constraint therefore gives the CDF band \citep[Supplement, Section~S.6]{L.J23}
%  \[
%    \alpha_{N^{\mathrm{te}}\widehat F^{\mathrm{te}}(t)}(\theta)
%      \le F_{\mathbb Q}(t)
%      \le\beta_{N^{\mathrm{te}}\widehat F^{\mathrm{te}}(t)}(\theta),
%    \qquad\forall t\in\mathbb R.
%  \]
%
%  For $k=1,\ldots,K-1$, the empirical CDF on $\left[\xi_{(k)}^{\mathrm{te}},\xi_{(k+1)}^{\mathrm{te}}\right)$ equals $c_k/N^{\mathrm{te}}$. Monotonicity makes the lower band constraint on this interval equivalent to $\mathbb Q\left\{\xi\le\xi_{(k)}^{\mathrm{te}}\right\}\ge\alpha_{c_k}(\theta)$. By continuity of probability from below, the upper band constraint is equivalent to $\mathbb Q\left\{\xi<\xi_{(k+1)}^{\mathrm{te}}\right\}\le\beta_{c_k}(\theta)$. The exterior intervals give $\mathbb Q\left\{\xi<\xi_{(1)}^{\mathrm{te}}\right\}\le\beta_0(\theta)$ and $\mathbb Q\left\{\xi\le\xi_{(K)}^{\mathrm{te}}\right\}\ge\alpha_{N^{\mathrm{te}}}(\theta)$; their other band constraints are automatic because $\alpha_0(\theta)=0$ and $\beta_{N^{\mathrm{te}}}(\theta)=1$. The partition identities above then yield \eqref{eq-finite-rank-constraints}.
%
%  With $p$ fixed, the corrected discrepancy decreases as $q$ approaches $p$ from either side. For empirical levels $q<q'$, an accepted value $p<q$ at level $q'$ remains accepted at level $q$; an accepted value $p>q'$ at level $q$ remains accepted at level $q'$. Since each acceptance interval contains its empirical level, both endpoints are nondecreasing. The inequalities $\alpha_{c_k}(\theta)\le c_k/N^{\mathrm{te}}\le\beta_{c_k}(\theta)$ give the stated compatibility conditions. Singleton intervals are allowed: if $q_j=1/2$ and $\theta=0$, the interval is $\{1/2\}$ because $\psi(1/2,p)=0$.
%\end{proof}
%
%Each scalar interval can be computed by one-dimensional bisection. Proposition~\ref{prop-finite-rank-constraints} then restricts only the probabilities of the residual cells; the distribution within each cell remains free.
%
%For a feasible joint law, define its matrix of joint cell probabilities by $\pi_{ij}:=\Pi(\mathcal X_i\times\Xi_j)$. The fixed input marginal and the residual restrictions imply $\bm\pi=(\pi_{ij})\in\mathcal C_\theta$, where
%\begin{equation}
%  \mathcal C_\theta:=\left\{\bm\pi\in\mathbb R_+^{(K+1)\times J}:
%  \begin{array}{ll}
%    \displaystyle\sum_{j=1}^{J}\pi_{ij}=\mathbb P(\mathcal X_i),&i=0,\ldots,K,\\[1mm]
%    \displaystyle\sum_{i=0}^{K}\sum_{j<2k}\pi_{ij}\le u_k,&k=1,\ldots,K,\\[1mm]
%    \displaystyle\sum_{i=0}^{K}\sum_{j\le2k}\pi_{ij}\ge l_k,&k=1,\ldots,K
%  \end{array}\right\}.
%  \label{eq-cell-mass-polytope}
%\end{equation}
%On each product cell, the failure indicator satisfies
%\begin{equation}
%  \underline h_{ij}:=\mathbf1_{\{j\le2i\}}
%  \le\mathbf1_{\{\xi\le-\hat g(\bm x)\}}
%  \le\overline h_{ij}:=\mathbf1_{\{j\le2i+1\}},
%  \qquad (\bm x,\xi)\in\mathcal X_i\times\Xi_j.
%  \label{eq-cell-costs}
%\end{equation}
%Residual cells with $j\le2i$ always cause failure for inputs in $\mathcal X_i$, and those with $j\ge2i+2$ are always safe. Only the overlapping interval $\Xi_{2i+1}$ has different upper and lower coefficients. Thus the cell matrix records the rank constraints exactly, while the failure objective requires the constructions below.
%
%The upper problem remains a linear program over joint measures, with the full input marginal fixed:
%\begin{equation}
%  \begin{aligned}
%    \overline p_{\hat f}
%    =\sup_{\Pi\in\mathcal M_+(\mathcal X\times\Xi)}\quad&
%      \int_{\mathcal X\times\Xi}\mathbf1_{\{\xi\le-\hat g(\bm x)\}}\,\Pi(\mathrm d\bm x,\mathrm d\xi)\\
%    \text{s.t.}\quad&
%      \int_{\mathcal X\times\Xi}\mathbf1_{\mathcal U}(\bm x)\,\Pi(\mathrm d\bm x,\mathrm d\xi)
%        =\mathbb P(\bm X\in\mathcal U),&&\forall\mathcal U\in\mathcal B(\mathcal X),\\
%    &\int_{\mathcal X\times\Xi}\mathbf1_{\{\xi<\xi_{(k)}^{\mathrm{te}}\}}\,\Pi(\mathrm d\bm x,\mathrm d\xi)\le u_k,
%        &&k=1,\ldots,K,\\
%    &\int_{\mathcal X\times\Xi}\mathbf1_{\{\xi\le\xi_{(k)}^{\mathrm{te}}\}}\,\Pi(\mathrm d\bm x,\mathrm d\xi)\ge l_k,
%        &&k=1,\ldots,K.
%  \end{aligned}
%  \label{eq-measure-upper}
%\end{equation}
%The lower problem minimizes the same integral under the same constraints. The marginal equalities fix both the input law and total mass one. The product of $\mathbb P$ and the empirical residual law is feasible for every $\theta\ge0$.
%
%\begin{lemma}[Finite-dimensional Lagrangian duals and weak duality]
%  \label{lem-weak-duality}
%  Write $\bm u=(u_1,\ldots,u_K)^\top$ and $\bm l=(l_1,\ldots,l_K)^\top$. The nonnegative multipliers $\bm\lambda,\bm\gamma\in\mathbb R_+^K$ correspond to the upper and lower rank constraints. The input marginal equality has an unrestricted multiplier $\eta\in\mathcal F_{\mathrm b}(\mathcal X)$, the space of bounded Borel measurable real-valued functions. Define the residual penalty
%  \begin{equation}
%    \zeta_{\bm\lambda,\bm\gamma}(\xi)
%    :=\sum_{k=1}^{K}\lambda_k\mathbf1_{\{\xi<\xi_{(k)}^{\mathrm{te}}\}}
%      -\sum_{k=1}^{K}\gamma_k\mathbf1_{\{\xi\le\xi_{(k)}^{\mathrm{te}}\}},
%    \label{eq-dual-step-function}
%  \end{equation}
%  whose value on $\Xi_j$ is
%  \begin{equation}
%    \zeta_j:=\sum_{k=1}^{K}\lambda_k\mathbf1_{\{j<2k\}}
%      -\sum_{k=1}^{K}\gamma_k\mathbf1_{\{j\le2k\}},
%    \qquad j=1,\ldots,J.
%    \label{eq-dual-cell-values}
%  \end{equation}
%  The upper and lower Lagrangian dual values are given by the finite LPs
%  \begin{equation}
%    \begin{aligned}
%      \overline d_\theta
%      :=\inf_{\bm\lambda,\bm\gamma\ge\bm0,\ \bm z\in\mathbb R^{K+1}}\quad&
%        \bm\lambda^\top\bm u-\bm\gamma^\top\bm l+\sum_{i=0}^{K}\mathbb P(\mathcal X_i) z_i\\
%      \text{s.t.}\quad&z_i+\zeta_j\ge\overline h_{ij},
%        \qquad i=0,\ldots,K,\quad j=1,\ldots,J,
%    \end{aligned}
%    \label{eq-finite-upper-dual}
%  \end{equation}
%  and
%  \begin{equation}
%    \begin{aligned}
%      \underline d_\theta
%      :=\sup_{\bm\lambda,\bm\gamma\ge\bm0,\ \bm z\in\mathbb R^{K+1}}\quad&
%        -\bm\lambda^\top\bm u+\bm\gamma^\top\bm l+\sum_{i=0}^{K}\mathbb P(\mathcal X_i) z_i\\
%      \text{s.t.}\quad&z_i-\zeta_j\le\underline h_{ij},
%        \qquad i=0,\ldots,K,\quad j=1,\ldots,J.
%    \end{aligned}
%    \label{eq-finite-lower-dual}
%  \end{equation}
%  Their optimal values satisfy
%  \begin{equation}
%    \underline d_\theta\le\underline p_{\hat f}
%      \le\overline p_{\hat f}\le\overline d_\theta.
%    \label{eq-weak-duality}
%  \end{equation}
%\end{lemma}
%\begin{proof}
%  The input marginal equality is equivalent to
%  \[
%    \int_{\mathcal X\times\Xi}\eta(\bm x)\,\Pi(\mathrm d\bm x,\mathrm d\xi)
%      =\mathbb E_{\mathbb P}\left[\eta(\bm X)\right],
%    \qquad\forall\eta\in\mathcal F_{\mathrm b}(\mathcal X).
%  \]
%  Relaxing this equality and the rank inequalities gives the upper Lagrangian \citep[Section~5.1]{B.V04}
%  \[
%    \begin{aligned}
%      \mathcal L^+(\Pi,\eta,\bm\lambda,\bm\gamma)
%      ={}&\int_{\mathcal X\times\Xi}\mathbf1_{\{\xi\le-\hat g(\bm x)\}}\,\Pi(\mathrm d\bm x,\mathrm d\xi)
%        +\mathbb E_{\mathbb P}\left[\eta(\bm X)\right]
%        -\int_{\mathcal X\times\Xi}\eta(\bm x)\,\Pi(\mathrm d\bm x,\mathrm d\xi)\\
%      &+\sum_{k=1}^{K}\lambda_k\left(u_k-\int_{\mathcal X\times\Xi}
%          \mathbf1_{\{\xi<\xi_{(k)}^{\mathrm{te}}\}}\,\Pi(\mathrm d\bm x,\mathrm d\xi)\right)\\
%      &+\sum_{k=1}^{K}\gamma_k\left(\int_{\mathcal X\times\Xi}
%          \mathbf1_{\{\xi\le\xi_{(k)}^{\mathrm{te}}\}}\,\Pi(\mathrm d\bm x,\mathrm d\xi)-l_k\right)\\
%      ={}&\mathbb E_{\mathbb P}\left[\eta(\bm X)\right]
%        +\bm\lambda^\top\bm u-\bm\gamma^\top\bm l\\
%      &+\int_{\mathcal X\times\Xi}\left[
%          \mathbf1_{\{\xi\le-\hat g(\bm x)\}}-\eta(\bm x)-\zeta_{\bm\lambda,\bm\gamma}(\xi)
%        \right]\,\Pi(\mathrm d\bm x,\mathrm d\xi).
%    \end{aligned}
%  \]
%  Its supremum over finite nonnegative measures is finite exactly when the integrand is nonpositive at every $(\bm x,\xi)$. A positive value at any point permits an arbitrarily large point mass there; otherwise the supremum integral is zero, attained by the zero measure. The functional upper dual therefore minimizes $\mathbb E_{\mathbb P}[\eta(\bm X)]+\bm\lambda^\top\bm u-\bm\gamma^\top\bm l$ subject to $\eta(\bm x)\ge\mathbf1_{\{\xi\le-\hat g(\bm x)\}}-\zeta_{\bm\lambda,\bm\gamma}(\xi)$ at every pair. For the lower problem, reversing the rank penalties gives
%  \[
%    \mathcal L^-
%      =\mathbb E_{\mathbb P}\left[\eta(\bm X)\right]
%        -\bm\lambda^\top\bm u+\bm\gamma^\top\bm l
%        +\int_{\mathcal X\times\Xi}\left[
%          \mathbf1_{\{\xi\le-\hat g(\bm x)\}}-\eta(\bm x)+\zeta_{\bm\lambda,\bm\gamma}(\xi)
%        \right]\,\Pi(\mathrm d\bm x,\mathrm d\xi).
%  \]
%  Its infimum over finite nonnegative measures is finite exactly when the integrand is nonnegative everywhere, in which case the infimum integral is zero. The lower functional dual maximizes $\mathbb E_{\mathbb P}[\eta(\bm X)]-\bm\lambda^\top\bm u+\bm\gamma^\top\bm l$ subject to $\eta(\bm x)\le\mathbf1_{\{\xi\le-\hat g(\bm x)\}}+\zeta_{\bm\lambda,\bm\gamma}(\xi)$ at every pair.
%
%  For fixed rank multipliers, the best marginal functions are the pointwise envelopes. For $\bm x\in\mathcal X_i$, they satisfy the identities
%  \begin{equation}
%    \begin{aligned}
%      \eta^+_{\bm\lambda,\bm\gamma}(\bm x)
%        &:=\sup_{\xi\in\Xi}\left\{\mathbf1_{\{\xi\le-\hat g(\bm x)\}}
%          -\zeta_{\bm\lambda,\bm\gamma}(\xi)\right\}
%          =\max_{1\le j\le J}\left\{\overline h_{ij}-\zeta_j\right\},\\
%      \eta^-_{\bm\lambda,\bm\gamma}(\bm x)
%        &:=\inf_{\xi\in\Xi}\left\{\mathbf1_{\{\xi\le-\hat g(\bm x)\}}
%          +\zeta_{\bm\lambda,\bm\gamma}(\xi)\right\}
%          =\min_{1\le j\le J}\left\{\underline h_{ij}+\zeta_j\right\}.
%    \end{aligned}
%    \label{eq-dual-envelope-identities}
%  \end{equation}
%  To verify them, note that outside $\Xi_{2i+1}$ the two cell coefficients equal the actual failure indicator. In the overlapping interval, the lower coefficient zero is attained by $\xi=\left(\xi_{(i+1)}^{\mathrm{te}}-\hat g(\bm x)\right)/2$ if $i<K$, and by $\xi=-\hat g(\bm x)+1$ if $i=K$. The upper coefficient one is attained by $\xi=-\hat g(\bm x)$ whenever $-\hat g(\bm x)>\xi_{(i)}^{\mathrm{te}}$. The only remaining case is $-\hat g(\bm x)=\xi_{(i)}^{\mathrm{te}}$ with $i\ge1$. Then the overlapping interval is safe, but
%  \[
%    \zeta_{2i}=\zeta_{2i+1}-\gamma_i\le\zeta_{2i+1},
%    \qquad
%    1-\zeta_{2i}\ge1-\zeta_{2i+1}.
%  \]
%  Choosing the left endpoint therefore attains an upper-envelope value at least as large as the candidate from the overlapping interval. Together with \eqref{eq-cell-costs}, this proves both identities, including threshold equalities.
%
%  The envelopes are bounded and constant on the Borel cells $\mathcal X_i$, so they belong to $\mathcal F_{\mathrm b}(\mathcal X)$. They respectively minimize and maximize the expectation over feasible marginal functions. Substitution yields
%  \begin{align}
%    \overline d_\theta
%      &=\inf_{\bm\lambda,\bm\gamma\ge\bm0}\left\{
%        \bm\lambda^\top\bm u-\bm\gamma^\top\bm l
%        +\sum_{i=0}^{K}\mathbb P(\mathcal X_i)\max_{1\le j\le J}
%          \left\{\overline h_{ij}-\zeta_j\right\}\right\},
%      \label{eq-upper-measure-dual}\\
%    \underline d_\theta
%      &=\sup_{\bm\lambda,\bm\gamma\ge\bm0}\left\{
%        -\bm\lambda^\top\bm u+\bm\gamma^\top\bm l
%        +\sum_{i=0}^{K}\mathbb P(\mathcal X_i)\min_{1\le j\le J}
%          \left\{\underline h_{ij}+\zeta_j\right\}\right\}.
%      \label{eq-lower-measure-dual}
%  \end{align}
%  Replacing each maximum by an epigraph variable $z_i$, and each minimum by a hypograph variable $z_i$, gives \eqref{eq-finite-upper-dual} and \eqref{eq-finite-lower-dual}. This remains exact when $\mathbb P(\mathcal X_i)=0$, since $z_i$ can still be chosen as the corresponding finite envelope.
%
%  For a feasible joint law, the signs of the rank penalties give $\mathcal L^-\le p_{\hat f}(\Pi)\le\mathcal L^+$. Feasible lower dual objectives therefore bound the primal objectives from below, and feasible upper dual objectives bound them from above. Taking the respective optima over the common nonempty primal feasible set proves \eqref{eq-weak-duality}.
%\end{proof}
%
%Both duals have $3K+1$ scalar variables and $(K+1)(2K+1)$ pairwise inequalities, together with multiplier nonnegativity. The variables $z_i$ represent the optimal marginal function on each input cell. The envelope identities show that these finite LPs are the full Lagrangian duals; their derivation does not require strong duality.
%
%\begin{theorem}[Exact finite reformulations and strong duality]
%  \label{thm-finite-dual-reformulations}
%  For $\mathcal Q=\mathcal Q^{\mathrm{DW}}(\theta)$ and every $\theta\ge0$, the failure probability bounds admit the exact representations
%  \begin{equation}
%    \overline p_{\hat f}
%      =\max_{\bm\pi\in\mathcal C_\theta}\sum_{i=0}^{K}\sum_{j=1}^{J}\overline h_{ij}\pi_{ij},
%    \qquad
%    \underline p_{\hat f}
%      =\min_{\bm\pi\in\mathcal C_\theta}\sum_{i=0}^{K}\sum_{j=1}^{J}\underline h_{ij}\pi_{ij}.
%    \label{eq-finite-primal-values}
%  \end{equation}
%  Their Lagrangian dual values in Lemma~\ref{lem-weak-duality} satisfy
%  \begin{equation}
%    \overline p_{\hat f}=\overline d_\theta,
%    \qquad
%    \underline p_{\hat f}=\underline d_\theta.
%    \label{eq-strong-duality}
%  \end{equation}
%  The finite primal and dual optima, the measure primal optima, and the functional dual optima are attained.
%\end{theorem}
%\begin{proof}
%  Any feasible joint law induces a matrix in $\mathcal C_\theta$. Integrating \eqref{eq-cell-costs} therefore bounds the upper measure problem by the finite maximization from above and the lower measure problem by the finite minimization from below.
%
%  For the reverse inequalities, fix $\bm\pi\in\mathcal C_\theta$. Draw $\bm x$ from $\mathbb P$. On $\mathcal X_i$ with $\mathbb P(\mathcal X_i)>0$, select a provisional residual cell $j$ with probability $\pi_{ij}/\mathbb P(\mathcal X_i)$; the row equality makes these probabilities sum to one. For $j\ne2i+1$, choose a fixed representative of $\Xi_j$, where the actual failure indicator equals both coefficients. In the overlapping cell, use the following assignments to define joint laws $\Pi^+$ and $\Pi^-$.
%
%  For the upper problem, assign $\xi=-\hat g(\bm x)$. This causes failure and remains in $\Xi_{2i+1}$ unless $-\hat g(\bm x)=\xi_{(i)}^{\mathrm{te}}$. At equality, which requires $i\ge1$, the assignment transfers the provisional mass to the left endpoint. For any $\widetilde\xi\in\left(\xi_{(i)}^{\mathrm{te}},\xi_{(i+1)}^{\mathrm{te}}\right)$ and every observed threshold,
%  \[
%    \mathbf1_{\{\xi_{(i)}^{\mathrm{te}}<\xi_{(k)}^{\mathrm{te}}\}}
%      =\mathbf1_{\{\widetilde\xi<\xi_{(k)}^{\mathrm{te}}\}},
%    \qquad
%    \mathbf1_{\{\xi_{(i)}^{\mathrm{te}}\le\xi_{(k)}^{\mathrm{te}}\}}
%      \ge\mathbf1_{\{\widetilde\xi\le\xi_{(k)}^{\mathrm{te}}\}}.
%  \]
%  Thus every open upper probability is unchanged and every closed lower probability can only increase. The upper assignment preserves all rank inequalities and attains the coefficient $\overline h_{i,2i+1}=1$.
%
%  For the lower problem, assign $\xi=\left(\xi_{(i+1)}^{\mathrm{te}}-\hat g(\bm x)\right)/2$ if $i<K$, and $\xi=-\hat g(\bm x)+1$ if $i=K$. These values remain in $\Xi_{2i+1}$ and strictly exceed the failure threshold, including at threshold equality. Hence the lower assignment retains the cell matrix and attains $\underline h_{i,2i+1}=0$.
%
%  If $\mathbb P(\mathcal X_i)=0$, the corresponding row is zero; set $\xi=\xi_{(1)}^{\mathrm{te}}$ on that input cell. All assignments are Borel measurable, and their finite mixtures define probability kernels. Combining them with $\mathbb P$ gives joint laws with the full prescribed input marginal, including randomization at input atoms. Both laws are feasible and satisfy
%  \[
%    p_{\hat f}(\Pi^+)=\sum_{i=0}^{K}\sum_{j=1}^{J}\overline h_{ij}\pi_{ij},
%    \qquad
%    p_{\hat f}(\Pi^-)=\sum_{i=0}^{K}\sum_{j=1}^{J}\underline h_{ij}\pi_{ij}.
%  \]
%  The lower law retains every provisional cell mass. The upper law may change the cell matrix at threshold equalities, but it preserves feasibility and the assigned upper objective. These identities establish the reverse inequalities required for \eqref{eq-finite-primal-values}.
%
%  For every $\theta\ge0$, the empirical product law induces the feasible matrix
%  \[
%    \pi_{i,2k}^{\mathrm{emp}}:=\mathbb P(\mathcal X_i)\frac{c_k-c_{k-1}}{N^{\mathrm{te}}},
%    \qquad i=0,\ldots,K,\quad k=1,\ldots,K,
%  \]
%  with all odd columns zero. Its open and closed rank probabilities are $c_{k-1}/N^{\mathrm{te}}$ and $c_k/N^{\mathrm{te}}$, respectively, so feasibility includes $\theta=0$. The nonnegative entries of a matrix in $\mathcal C_\theta$ sum to one, and the constraints are closed. Thus $\mathcal C_\theta$ is nonempty and compact. The finite objectives lie in $[0,1]$ and attain their optima, which the joint-law constructions realize in the measure problems.
%
%  The finite primal LPs in \eqref{eq-finite-primal-values} have row multipliers $z_i$ and rank multipliers $\lambda_k,\gamma_k$. Their dual constraints are precisely $z_i+\zeta_j\ge\overline h_{ij}$ and $z_i-\zeta_j\le\underline h_{ij}$. Finite LP strong duality therefore gives \eqref{eq-strong-duality} and dual attainment \citep[Section~5.2.3]{B.V04}. By \eqref{eq-dual-envelope-identities}, optimal rank multipliers also yield optimal bounded marginal functions for the original functional duals. Feasibility and finite optimal values suffice; strict rank inequalities are unnecessary, including at $\theta=0$.
%\end{proof}
%
%The finite formulations are exact, and their ordered failure coefficients give a further simplification. The closed-form bounds below follow from cuts between consecutive cells and admit certificates using at most one rank multiplier.
%
%\begin{corollary}[Closed-form bounds and optimal certificates]
%  \label{thm-closed-form-bounds}
%  Under the conditions of Theorem~\ref{thm-finite-dual-reformulations},
%  \begin{align}
%    \overline p_{\hat f}
%      &=\min\left\{1,\min_{k=1,\ldots,K}\left[u_k+H\left(-\xi_{(k)}^{\mathrm{te}}\right)\right]\right\},
%      \label{eq-dw-closed-upper}\\
%    \underline p_{\hat f}
%      &=\max\left\{0,\max_{k=1,\ldots,K}\left[l_k+H\left(-\xi_{(k)}^{\mathrm{te}}\right)-1\right]\right\}.
%      \label{eq-dw-closed-lower}
%  \end{align}
%  Each dual optimum admits a certificate with at most one nonzero rank multiplier, of value one.
%\end{corollary}
%\begin{proof}
%  Fix nonnegative column masses $\bm y=(y_1,\ldots,y_J)^\top$ with $\sum_j y_j=1$, and write $Y_j:=\sum_{r=1}^{j}y_r$. Let $\mathcal T(\bm y)$ be the set of nonnegative transportation matrices with row sums $\mathbb P(\mathcal X_i)$ and column sums $y_j$. This set contains the product matrix $\pi_{ij}=\mathbb P(\mathcal X_i) y_j$. The residual constraints depend only on the column masses, through $Y_{2k-1}\le u_k$ and $Y_{2k}\ge l_k$.
%
%  For fixed $\bm y$, the upper problem maximizes mass on pairs $j\le2i+1$; the lower problem maximizes safe mass on pairs $j>2i$ and subtracts it from one. Both are maximum-flow problems with source-to-row capacities $\mathbb P(\mathcal X_i)$, column-to-sink capacities $y_j$, and capacity two on each permitted row-to-column edge. A transportation matrix restricts to a feasible flow on permitted pairs. Conversely, any flow can be completed to a transportation matrix because the remaining row and column masses are nonnegative with equal totals. Thus the maximum flow equals the maximum permitted mass.
%
%  A cut of capacity one is available, so a minimum cut crosses no edge of capacity two. Its source-side columns can be taken to be exactly the neighbors of its source-side rows. Since the neighborhoods are nested, adding rows whose neighbors are already included cannot increase the capacity. Apart from empty and full row sets, upper-problem cuts therefore reduce to rows $i<k$ and columns $j<2k$, while safe-mass cuts reduce to rows $i\ge k$ and columns $j>2k$. Their capacities, for $k=1,\ldots,K$, are respectively
%  \begin{align*}
%    \sum_{i=k}^{K}\mathbb P(\mathcal X_i)+Y_{2k-1}
%      &=H\left(-\xi_{(k)}^{\mathrm{te}}\right)+Y_{2k-1},\\
%    \sum_{i=0}^{k-1}\mathbb P(\mathcal X_i)+1-Y_{2k}
%      &=2-H\left(-\xi_{(k)}^{\mathrm{te}}\right)-Y_{2k}.
%  \end{align*}
%  Empty and full row sets give capacity one. The max-flow/min-cut theorem \citep{F.F56} consequently yields
%  \begin{align*}
%    \max_{\bm\pi\in\mathcal T(\bm y)}\sum_{i,j}\overline h_{ij}\pi_{ij}
%      &=\min\left\{1,\min_k\left[H\left(-\xi_{(k)}^{\mathrm{te}}\right)+Y_{2k-1}\right]\right\},\\
%    \min_{\bm\pi\in\mathcal T(\bm y)}\sum_{i,j}\underline h_{ij}\pi_{ij}
%      &=\max\left\{0,\max_k\left[H\left(-\xi_{(k)}^{\mathrm{te}}\right)+Y_{2k}-1\right]\right\}.
%  \end{align*}
%
%  These values are nondecreasing in their respective open and closed prefixes. To attain all upper prefix limits for the maximization and all lower prefix limits for the minimization, choose
%  \begin{align*}
%    &y_1^+:=u_1,\qquad y_{2k}^+:=u_{k+1}-u_k,\qquad
%      y_{2k+1}^+:=0,\quad k=1,\ldots,K,\\
%    &y_{2k-1}^-:=0,\qquad y_{2k}^-:=l_k-l_{k-1},\quad k=1,\ldots,K,
%      \qquad y_J^-:=1-l_K.
%  \end{align*}
%  The monotonicity and compatibility conditions in Proposition~\ref{prop-finite-rank-constraints} make both vectors nonnegative with total mass one, and
%  \[
%    Y_{2k-1}^+=u_k,\qquad Y_{2k}^+=u_{k+1}\ge l_k,
%    \qquad Y_{2k}^-=l_k,\qquad Y_{2k-1}^-=l_{k-1}\le u_k.
%  \]
%  Both allocations are feasible and attain the required prefix limits. Substitution into the fixed-column values, followed by \eqref{eq-finite-primal-values}, proves \eqref{eq-dw-closed-upper} and \eqref{eq-dw-closed-lower}. The construction allows zero masses and coincident prefix limits, including at $\theta=0$.
%
%  Finally, let $\bm e_k$ be the $k$th unit vector. Upper and lower dual certificates are respectively
%  \begin{align*}
%    &\bm\lambda=\bm e_k,\qquad\bm\gamma=\bm0,\qquad
%      z_i=\mathbf1_{\{i\ge k\}},\\
%    &\bm\lambda=\bm0,\qquad\bm\gamma=\bm e_k,\qquad
%      z_i=\mathbf1_{\{i\ge k\}}-1.
%  \end{align*}
%  Their feasibility follows from the index inequalities
%  \[
%    \mathbf1_{\{i\ge k\}}+\mathbf1_{\{j<2k\}}\ge\overline h_{ij},
%    \qquad
%    \mathbf1_{\{i\ge k\}}-1+\mathbf1_{\{j\le2k\}}\le\underline h_{ij}.
%  \]
%  Their objectives are $u_k+\sum_{i=k}^{K}\mathbb P(\mathcal X_i)$ and $l_k+\sum_{i=k}^{K}\mathbb P(\mathcal X_i)-1$, the inner candidates in the closed-form expressions. Zero rank multipliers with $z_i=1$ or $z_i=0$ supply the outer candidates. Selecting the best candidate gives an optimal full dual solution. The envelope identities recover its functional marginal multiplier on each input cell; the tail sum includes $i=k$ and therefore uses the closed CDF value $H\left(-\xi_{(k)}^{\mathrm{te}}\right)$ also at threshold equality.
%\end{proof}
%
%The same formulas give the response-threshold curves in Corollary~\ref{cor-simultaneous-threshold-bounds}. Replacing $g$ and $\hat g$ by $g-c$ and $\hat g-c$ leaves every residual, rank count, and endpoint $l_k,u_k$ unchanged. The surrogate CDF becomes $t\mapsto H(t+c)$, so Corollary~\ref{thm-closed-form-bounds} gives
%\begin{align}
%  \overline p_{\hat f}(c)
%    &=\min\left\{1,\min_{k=1,\ldots,K}\left[u_k+H\left(c-\xi_{(k)}^{\mathrm{te}}\right)\right]\right\},
%    \label{eq-threshold-closed-upper}\\
%  \underline p_{\hat f}(c)
%    &=\max\left\{0,\max_{k=1,\ldots,K}\left[l_k+H\left(c-\xi_{(k)}^{\mathrm{te}}\right)-1\right]\right\}.
%    \label{eq-threshold-closed-lower}
%\end{align}
%These identities hold for every $c\in\mathbb R$ and $\theta\ge0$, including $\theta=0$. When $\theta\ge\theta_{N^{\mathrm{te}}}(\delta)$, their coverage is simultaneous by \eqref{eq-threshold-curve-coverage}. An extremizing coupling may depend on $c$; the common feasible true joint law supplies the simultaneous guarantee.
%
%The upper bound combines residual mass below an observed value with the probability that the surrogate threshold reaches it. The lower bound measures the overlap required by every admissible coupling. For $p_{\mathrm{tar}}\in(0,1)$, the upper curve satisfies $\overline p_{\hat f}(c)\le p_{\mathrm{tar}}$ exactly when some rank satisfies $u_k+H\left(c-\xi_{(k)}^{\mathrm{te}}\right)\le p_{\mathrm{tar}}$. Both terms must therefore be small at the same residual level.
%
%With a known residual law, the coupling problem is the classical fixed-marginal sum problem studied by \citet{M82} and \citet{F.N.S87}. Here the residual law also varies within a calibrated DW band, and the open and closed rank constraints retain atoms.
%
%Dependence uncertainty can remain even with a known residual marginal. For equally likely surrogate responses $0$ and $2$ and equally likely residuals $-1$ and $1$, failure occurs exactly when the response is $0$ and the residual is $-1$. Its probability ranges from $0$ to $1/2$ over couplings, whereas independence gives $1/4$. Further marginal calibration alone does not resolve this dependence uncertainty.
%
%\begin{remark}[Known interval support]
%  \label{rem-interval-support}
%  Suppose $\Xi=[a,b]$ with $-\infty<a\le b<\infty$ and the observed residuals lie in this interval. Then
%  \begin{align*}
%    \overline p_{\hat f}
%    &=\min\left\{H(-a),\min_{k=1,\ldots,K}[u_k+H\left(-\xi_{(k)}^{\mathrm{te}}\right)]\right\},\\
%    \underline p_{\hat f}
%    &=\max\left\{H(-b),\max_{k=1,\ldots,K}[l_k+H\left(-\xi_{(k)}^{\mathrm{te}}\right)-1]\right\}.
%  \end{align*}
%  To obtain these expressions, represent the support restriction on $\mathbb R$ by
%  $\mathbb Q\{\xi<a\}\le0$ and $\mathbb Q\{\xi\le b\}\ge1$.
%  Add $a,b$ to the ordered threshold list. At a newly added $a$, use upper and lower bounds zero; at a newly added $b$, use upper and lower bounds one. At coincident observed thresholds, tighten only the upper bound at $a$ and the lower bound at $b$. The enlarged sequences satisfy the compatibility conditions in Proposition~\ref{prop-finite-rank-constraints}, so the same finite-aggregation and ordered-cut arguments apply. The effective additional upper candidate is $H(-a)$, and the effective additional lower candidate is $H(-b)$; the other added candidates are redundant. If $a=b$, both bounds equal $H(-a)$. For the resulting partition on $\Xi=[a,b]$, intersect the singleton and open-interval residual cells from the augmented threshold list with $\Xi$ and omit empty cells. Including the support endpoints in this list retains their atoms and the corresponding changes in failure status. For a response threshold $c$, replace $H(-a)$, $H(-b)$, and $H\left(-\xi_{(k)}^{\mathrm{te}}\right)$ by $H(c-a)$, $H(c-b)$, and $H\left(c-\xi_{(k)}^{\mathrm{te}}\right)$, respectively.
%\end{remark}
%
%After zero failures in $N^{\mathrm{te}}$ independent tests, the one-sided exact binomial upper limit is $1-\delta^{1/N^{\mathrm{te}}}\approx\log(1/\delta)/N^{\mathrm{te}}$ \citep{C.P34}. At $c=0$, the simple DW case $K=1$ and $H\left(-\xi_{(1)}^{\mathrm{te}}\right)=0$ gives an upper bound $\beta_0(\theta)$ determined by the simultaneous-band calibration. Verified support can exclude otherwise unobserved failures: in Remark~\ref{rem-interval-support}, $H(-a)=0$ rules out failure.
%
%Evaluation requires scalar DW inversion, sorting and counting the test residuals, and computing the $K$ values $H\left(c-\xi_{(k)}^{\mathrm{te}}\right)$ at each response threshold $c$. The endpoint table depends only on $N^{\mathrm{te}}$ and $\theta$ and can be precomputed. Given this table and the surrogate CDF values, the two scans cost $O(K)$ at each response threshold.
%
%For $\theta\ge\theta_{N^{\mathrm{te}}}(\delta)$, the exact expressions retain the coverage in Theorem~\ref{theorem-01} and Corollary~\ref{cor-simultaneous-threshold-bounds}. If $H$ is evaluated by independent auxiliary simulation, upper CDF bounds can replace its values in the upper expression, and lower CDF bounds can replace them in the lower expression. For one fixed response threshold, simultaneous bounds at its $K$ evaluation points suffice. To retain coverage of the entire response-threshold curve, the auxiliary bounds must hold over all arguments of $H$. In either case, they must hold conditionally on the test data, with their error probability included in the overall confidence budget.
%
%
% END PRESERVED SECTION 4 BEFORE PARTITION AND EVENT PROBABILITIES

% BEGIN PRESERVED LEMMA BEFORE ELIMINATED-DUAL STATEMENT (2026-11-25)
%\begin{lemma}[Finite-dimensional Lagrangian duals and weak duality]
%  \label{lem-weak-duality}
%  Write $\bm u=(u_1,\ldots,u_K)^\top$ and $\bm l=(l_1,\ldots,l_K)^\top$. The nonnegative multipliers $\bm\lambda,\bm\gamma\in\mathbb R_+^K$ correspond to the upper and lower rank constraints. The input marginal equality has an unrestricted multiplier $\eta\in\mathcal F_{\mathrm b}(\mathcal X)$, the space of bounded Borel measurable real-valued functions. Define the residual penalty
%  \begin{equation}
%    \zeta_{\bm\lambda,\bm\gamma}(\xi)
%    :=\sum_{k=1}^{K}\lambda_k\mathbf1_{\{\xi<\xi_{(k)}^{\mathrm{te}}\}}
%      -\sum_{k=1}^{K}\gamma_k\mathbf1_{\{\xi\le\xi_{(k)}^{\mathrm{te}}\}},
%    \label{eq-dual-step-function}
%  \end{equation}
%  whose value on $\Xi_j$ is
%  \begin{equation}
%    \zeta_j:=\sum_{k=1}^{K}\lambda_k\mathbf1_{\{j<2k\}}
%      -\sum_{k=1}^{K}\gamma_k\mathbf1_{\{j\le2k\}},
%    \qquad j=1,\ldots,J.
%    \label{eq-dual-cell-values}
%  \end{equation}
%  The upper and lower Lagrangian dual values are given by the finite LPs
%  \begin{equation}
%    \begin{aligned}
%      \overline d_\theta
%      :=\inf_{\bm\lambda,\bm\gamma\ge\bm0,\ \bm z\in\mathbb R^{K+1}}\quad&
%        \bm\lambda^\top\bm u-\bm\gamma^\top\bm l+\sum_{i=0}^{K}\mathbb P(\mathcal X_i) z_i\\
%      \text{s.t.}\quad&z_i+\zeta_j\ge\overline h_{ij},
%        \qquad i=0,\ldots,K,\quad j=1,\ldots,J,
%    \end{aligned}
%    \label{eq-finite-upper-dual}
%  \end{equation}
%  and
%  \begin{equation}
%    \begin{aligned}
%      \underline d_\theta
%      :=\sup_{\bm\lambda,\bm\gamma\ge\bm0,\ \bm z\in\mathbb R^{K+1}}\quad&
%        -\bm\lambda^\top\bm u+\bm\gamma^\top\bm l+\sum_{i=0}^{K}\mathbb P(\mathcal X_i) z_i\\
%      \text{s.t.}\quad&z_i-\zeta_j\le\underline h_{ij},
%        \qquad i=0,\ldots,K,\quad j=1,\ldots,J.
%    \end{aligned}
%    \label{eq-finite-lower-dual}
%  \end{equation}
%  Their optimal values satisfy
%  \begin{equation}
%    \underline d_\theta\le\underline p_{\hat f}
%      \le\overline p_{\hat f}\le\overline d_\theta.
%    \label{eq-weak-duality}
%  \end{equation}
%\end{lemma}
%\begin{proof}
%  The input marginal equality is equivalent to
%  \[
%    \int_{\mathcal X\times\Xi}\eta(\bm x)\,\Pi(\mathrm d\bm x,\mathrm d\xi)
%      =\mathbb E_{\mathbb P}\left[\eta(\bm X)\right],
%    \qquad\forall\eta\in\mathcal F_{\mathrm b}(\mathcal X).
%  \]
%  Relaxing this equality and the rank inequalities gives the upper Lagrangian \citep[Section~5.1]{B.V04}
%  \[
%    \begin{aligned}
%      \mathcal L^+(\Pi,\eta,\bm\lambda,\bm\gamma)
%      ={}&\int_{\mathcal X\times\Xi}\mathbf1_{\{\xi\le-\hat g(\bm x)\}}\,\Pi(\mathrm d\bm x,\mathrm d\xi)
%        +\mathbb E_{\mathbb P}\left[\eta(\bm X)\right]
%        -\int_{\mathcal X\times\Xi}\eta(\bm x)\,\Pi(\mathrm d\bm x,\mathrm d\xi)\\
%      &+\sum_{k=1}^{K}\lambda_k\left(u_k-\int_{\mathcal X\times\Xi}
%          \mathbf1_{\{\xi<\xi_{(k)}^{\mathrm{te}}\}}\,\Pi(\mathrm d\bm x,\mathrm d\xi)\right)
%        +\sum_{k=1}^{K}\gamma_k\left(\int_{\mathcal X\times\Xi}
%          \mathbf1_{\{\xi\le\xi_{(k)}^{\mathrm{te}}\}}\,\Pi(\mathrm d\bm x,\mathrm d\xi)-l_k\right)\\
%      ={}&\mathbb E_{\mathbb P}\left[\eta(\bm X)\right]
%        +\bm\lambda^\top\bm u-\bm\gamma^\top\bm l
%        +\int_{\mathcal X\times\Xi}\left[
%          \mathbf1_{\{\xi\le-\hat g(\bm x)\}}-\eta(\bm x)-\zeta_{\bm\lambda,\bm\gamma}(\xi)
%        \right]\,\Pi(\mathrm d\bm x,\mathrm d\xi).
%    \end{aligned}
%  \]
%  Its supremum over finite nonnegative measures is finite exactly when the integrand is nonpositive at every $(\bm x,\xi)$. A positive value at any point permits an arbitrarily large point mass there; otherwise the supremum integral is zero, attained by the zero measure. The functional upper dual therefore minimizes $\mathbb E_{\mathbb P}[\eta(\bm X)]+\bm\lambda^\top\bm u-\bm\gamma^\top\bm l$ subject to $\eta(\bm x)\ge\mathbf1_{\{\xi\le-\hat g(\bm x)\}}-\zeta_{\bm\lambda,\bm\gamma}(\xi)$ at every pair. For the lower problem, reversing the rank penalties gives
%  \[
%    \mathcal L^-
%      =\mathbb E_{\mathbb P}\left[\eta(\bm X)\right]
%        -\bm\lambda^\top\bm u+\bm\gamma^\top\bm l
%        +\int_{\mathcal X\times\Xi}\left[
%          \mathbf1_{\{\xi\le-\hat g(\bm x)\}}-\eta(\bm x)+\zeta_{\bm\lambda,\bm\gamma}(\xi)
%        \right]\,\Pi(\mathrm d\bm x,\mathrm d\xi).
%  \]
%  Its infimum over finite nonnegative measures is finite exactly when the integrand is nonnegative everywhere, in which case the infimum integral is zero. The lower functional dual maximizes $\mathbb E_{\mathbb P}[\eta(\bm X)]-\bm\lambda^\top\bm u+\bm\gamma^\top\bm l$ subject to $\eta(\bm x)\le\mathbf1_{\{\xi\le-\hat g(\bm x)\}}+\zeta_{\bm\lambda,\bm\gamma}(\xi)$ at every pair.
%
%  For fixed rank multipliers, the best marginal functions are the pointwise envelopes. For $\bm x\in\mathcal X_i$, they satisfy the identities
%  \begin{equation}
%    \begin{aligned}
%      \eta^+_{\bm\lambda,\bm\gamma}(\bm x)
%        &:=\sup_{\xi\in\Xi}\left\{\mathbf1_{\{\xi\le-\hat g(\bm x)\}}
%          -\zeta_{\bm\lambda,\bm\gamma}(\xi)\right\}
%          =\max_{1\le j\le J}\left\{\overline h_{ij}-\zeta_j\right\},\\
%      \eta^-_{\bm\lambda,\bm\gamma}(\bm x)
%        &:=\inf_{\xi\in\Xi}\left\{\mathbf1_{\{\xi\le-\hat g(\bm x)\}}
%          +\zeta_{\bm\lambda,\bm\gamma}(\xi)\right\}
%          =\min_{1\le j\le J}\left\{\underline h_{ij}+\zeta_j\right\}.
%    \end{aligned}
%    \label{eq-dual-envelope-identities}
%  \end{equation}
%  To verify them, note that outside $\Xi_{2i+1}$ the two cell coefficients equal the actual failure indicator. In the overlapping interval, the lower coefficient zero is attained by $\xi=\left(\xi_{(i+1)}^{\mathrm{te}}-\hat g(\bm x)\right)/2$ if $i<K$, and by $\xi=-\hat g(\bm x)+1$ if $i=K$. The upper coefficient one is attained by $\xi=-\hat g(\bm x)$ whenever $-\hat g(\bm x)>\xi_{(i)}^{\mathrm{te}}$. The only remaining case is $-\hat g(\bm x)=\xi_{(i)}^{\mathrm{te}}$ with $i\ge1$. Then the overlapping interval is safe, but
%  \[
%    \zeta_{2i}=\zeta_{2i+1}-\gamma_i\le\zeta_{2i+1},
%    \qquad
%    1-\zeta_{2i}\ge1-\zeta_{2i+1}.
%  \]
%  Choosing the left endpoint therefore attains an upper-envelope value at least as large as the candidate from the overlapping interval. Together with \eqref{eq-cell-costs}, this proves both identities, including threshold equalities.
%
%  The envelopes are bounded and constant on the Borel cells $\mathcal X_i$, so they belong to $\mathcal F_{\mathrm b}(\mathcal X)$. They respectively minimize and maximize the expectation over feasible marginal functions. Substitution yields
%  \begin{align}
%    \overline d_\theta
%      &=\inf_{\bm\lambda,\bm\gamma\ge\bm0}\left\{
%        \bm\lambda^\top\bm u-\bm\gamma^\top\bm l
%        +\sum_{i=0}^{K}\mathbb P(\mathcal X_i)\max_{1\le j\le J}
%          \left\{\overline h_{ij}-\zeta_j\right\}\right\},
%      \label{eq-upper-measure-dual}\\
%    \underline d_\theta
%      &=\sup_{\bm\lambda,\bm\gamma\ge\bm0}\left\{
%        -\bm\lambda^\top\bm u+\bm\gamma^\top\bm l
%        +\sum_{i=0}^{K}\mathbb P(\mathcal X_i)\min_{1\le j\le J}
%          \left\{\underline h_{ij}+\zeta_j\right\}\right\}.
%      \label{eq-lower-measure-dual}
%  \end{align}
%  Replacing each maximum by an epigraph variable $z_i$, and each minimum by a hypograph variable $z_i$, gives \eqref{eq-finite-upper-dual} and \eqref{eq-finite-lower-dual}. This remains exact when $\mathbb P(\mathcal X_i)=0$, since $z_i$ can still be chosen as the corresponding finite envelope.
%
%  For a feasible joint law, the signs of the rank penalties give $\mathcal L^-\le p_{\hat f}(\Pi)\le\mathcal L^+$. Feasible lower dual objectives therefore bound the primal objectives from below, and feasible upper dual objectives bound them from above. Taking the respective optima over the common nonempty primal feasible set proves \eqref{eq-weak-duality}.
%\end{proof}
%
%Both duals have $3K+1$ scalar variables and $(K+1)(2K+1)$ pairwise inequalities, together with multiplier nonnegativity. The variables $z_i$ represent the optimal marginal function on each input cell. The envelope identities show that these finite LPs are the full Lagrangian duals; their derivation does not require strong duality.
%
% END PRESERVED LEMMA BEFORE ELIMINATED-DUAL STATEMENT

% BEGIN PRESERVED BODY BEFORE RESULT-EXPOSITION STYLE REVISION (2026-11-25)
%\section{Problem Description}
%\label{sec-problem}
%
%Let $\bm X\sim\mathbb P$ be a random input taking values in $\mathcal X\subseteq\mathbb R^n$, and let $g:\mathcal X\to\mathbb R$ be a measurable limit-state function. The failure probability is defined as
%\[
%  p_f:=\Pr\{g(\bm X)\le0\}=\mathbb E_{\bm X\sim\mathbb P}\left[\mathbf 1_{\{g(\bm X)\le0\}}\right]=\int_{\mathcal X}\mathbf 1_{\{g(\bm x)\le0\}}\,\mathbb P(\mathrm d\bm x).
%\]
%
%Evaluating $g$ can require an expensive high-fidelity simulation. We use the training dataset
%\[
%  \mathcal D^{\mathrm{tr}}:=\left\{\left(\bm x_i^{\mathrm{tr}},g(\bm x_i^{\mathrm{tr}})\right)\right\}_{i=1}^{N^{\mathrm{tr}}}
%\]
%to fit a computationally efficient, measurable surrogate $\hat g:=\hat g_{\mathcal D^{\mathrm{tr}}}$. Model fitting and selection use only the training data and are completed before the test data are used. We condition on $\mathcal D^{\mathrm{tr}}$, treat $\hat g$ as fixed, and suppress this dependence in the notation.
%
%Define the signed residual function
%\[
%  \xi(\bm x):=g(\bm x)-\hat g(\bm x),\qquad \forall \bm x\in\mathcal X.
%\]
%For a new input $\bm X\sim\mathbb P$ independent of the training data, let $\mathbb Q^\star$ be the distribution of $\xi^\star:=\xi(\bm X)$. Let $\Xi\subseteq\mathbb R$ be a known Borel set with $\mathbb Q^\star(\Xi)=1$; the default choice is $\Xi=\mathbb R$. For any Borel set $\mathcal S$ in a Euclidean space, write $\mathcal B(\mathcal S)$ for its Borel sigma-algebra. Let $\mathcal M_+(\mathcal S)$ denote the finite nonnegative Borel measures concentrated on $\mathcal S$, and let $\mathcal M_1(\mathcal S)$ denote the probability measures.
%
%For a candidate residual law $\mathbb Q\in\mathcal M_1(\Xi)$, let $\Pi$ be a joint distribution on $\mathcal X\times\Xi$ with marginals $\Pi_{\bm X}=\mathbb P$ and $\Pi_\xi=\mathbb Q$. Its surrogate-based failure probability is
%\begin{equation}
%  p_{\hat f}(\Pi):=\Pr\{\hat g(\bm X)+\xi\le0\}=\int_{\mathcal X\times\Xi}\mathbf 1_{\{\hat g(\bm x)+\xi\le0\}}\,\Pi(\mathrm d\bm x,\mathrm d\xi).
%  \label{eq-surrogate-failure-probability}
%\end{equation}
%For fixed $\hat g$, this functional is linear in $\Pi$. In particular, the true joint distribution $\Pi^\star$ satisfies $p_{\hat f}(\Pi^\star)=p_f$.
%
%The residual marginal and its dependence on the input jointly determine the failure probability. We represent residual-marginal uncertainty by an ambiguity set $\mathcal Q$ calibrated to contain $\mathbb Q^\star$ with high confidence. For a given $\mathcal Q$, define the joint ambiguity set
%\[
%  \mathcal P(\mathcal Q) :=\left\{\Pi\in\mathcal M_1(\mathcal X\times\Xi):\Pi_{\bm X}=\mathbb P,\ \Pi_\xi\in\mathcal Q\right\}.
%\]
%The true residual is determined by the input, but its unknown functional relationship is not encoded in $\mathcal Q$. The set $\mathcal P(\mathcal Q)$ therefore allows all couplings with the specified marginals, including randomized conditional residual laws. The upper and lower failure probability bounds are
%\begin{equation}
%  \overline p_{\hat f}:=\sup_{\Pi\in\mathcal P(\mathcal Q)}p_{\hat f}(\Pi),
%  \label{eq-upper-bound}
%\end{equation}
%and
%\begin{equation}
%  \underline p_{\hat f}:=\inf_{\Pi\in\mathcal P(\mathcal Q)}p_{\hat f}(\Pi),
%  \label{eq-lower-bound}
%\end{equation}
%respectively. If $\mathbb Q^\star\in\mathcal Q$, then $\Pi^\star\in\mathcal P(\mathcal Q)$ and hence $\underline p_{\hat f}\le p_f\le\overline p_{\hat f}$.
%
%
%Section~\ref{sec-ambiguity-set} constructs the residual ambiguity set and establishes finite-sample coverage of the failure probability bounds, simultaneously across response thresholds. Section~\ref{sec-reformulation} then derives their Lagrangian duals, proves strong duality through a finite reformulation, and obtains closed-form expressions.
%
%
%\section{Ambiguity Set Design}
%\label{sec-ambiguity-set}
%
%Failure occurs when $\xi\le-\hat g(\bm X)$, so the residual cumulative distribution function (CDF) gives the marginal mass at or below each failure threshold. We use the D\"umbgen--Wellner (DW) construction \citep{L.J23} to control empirical CDF deviations simultaneously over all thresholds, with tolerances that vary across probability levels.
%
%Let $q$ and $p$ denote an empirical probability and a candidate probability. Following \citet{B.J79} and the CDF confidence bands of \citet{O95}, we measure their local discrepancy by the Bernoulli Kullback--Leibler divergence
%\begin{equation}
%  \phi(q,p):=q\log\frac{q}{p}+(1-q)\log\frac{1-q}{1-p},\qquad \forall q\in[0,1],\, p\in(0,1),
%  \label{eq-bernoulli-kl}
%\end{equation}
%where logarithms are natural and a term with zero numerator is interpreted as zero. For $q$ close to an interior $p$, the approximation $\phi(q,p)\approx(q-p)^2/[2p(1-p)]$ scales the deviation by the Bernoulli variance. Thus, the same absolute deviation receives a larger weight near the probability boundaries. The full divergence also covers zero and full empirical counts through $\phi(0,p)=-\log(1-p)$ and $\phi(1,p)=-\log p$. At the boundaries of the candidate probability, we set $\phi(0,0)=\phi(1,1)=0$, $\phi(q,0)=+\infty$ for $q>0$, and $\phi(q,1)=+\infty$ for $q<1$.
%
%To compare deviations simultaneously across probability levels, we use the additive correction of \citet{L.J23}. For $v\in(0,1)$, define
%\begin{equation*}
%  \rho(v):=\log\log\frac{e}{4v(1-v)},\qquad \omega(v):=\log\!\left(1+\rho(v)^2\right).
%\end{equation*}
%Before sampling, we fix the divergence parameter $s=1$, which gives \eqref{eq-bernoulli-kl}, and the correction parameter $\nu=1$. Set $\psi_0(v):=\rho(v)+\omega(v)$, with $\psi_0(0)=\psi_0(1)=+\infty$, and define the bivariate correction
%\begin{equation*}
%  \psi(q,p):=\min_{v\in[\min\{q,p\},\,\max\{q,p\}]}\psi_0(v).
%\end{equation*}
%The functions $\rho$ and $\omega$ vanish at $v=1/2$ and increase toward the probability boundaries. The minimum over the interval between $q$ and $p$ is therefore large only when both probabilities lie in the same tail. This correction balances fluctuations across probability scales in the simultaneous CDF constraint.
%
%We construct the empirical CDF from a test dataset
%\[
%  \mathcal D^{\mathrm{te}}:=\left\{\left(\bm x_j^{\mathrm{te}},g(\bm x_j^{\mathrm{te}})\right)\right\}_{j=1}^{N^{\mathrm{te}}}.
%\]
%The test inputs are independent draws from $\mathbb P$ and are independent of $\mathcal D^{\mathrm{tr}}$. We fix $N^{\mathrm{te}}\ge1$ before sampling and use these data to calibrate the residual ambiguity set.
%
%Define the signed test residuals
%\[
%  \xi_j^{\mathrm{te}}:=\xi(\bm x_j^{\mathrm{te}})=g(\bm x_j^{\mathrm{te}})-\hat g(\bm x_j^{\mathrm{te}}),\qquad j=1,\ldots,N^{\mathrm{te}}.
%\]
%Conditional on $\mathcal D^{\mathrm{tr}}$, these residuals are independent and identically distributed according to $\mathbb Q^\star$. Their empirical CDF is
%\begin{equation}
%  \widehat F^{\mathrm{te}}(t):=\frac{1}{N^{\mathrm{te}}}\sum_{j=1}^{N^{\mathrm{te}}}\mathbf 1_{\{\xi_j^{\mathrm{te}}\le t\}},\qquad \forall t\in\mathbb R.
%  \label{eq-empirical-cdf}
%\end{equation}
%For $\mathbb Q\in\mathcal M_1(\Xi)$, write $F_{\mathbb Q}(t):=\mathbb Q\left(({-\infty},t]\right)$. At a threshold $t$, $q=\widehat F^{\mathrm{te}}(t)$ is the mean of the Bernoulli indicators in \eqref{eq-empirical-cdf}, and $p=F_{\mathbb Q}(t)$ is the candidate probability. Thus $N^{\mathrm{te}}\phi(q,p)$ is the Bernoulli log-likelihood difference between the empirical and candidate probabilities. Combining it with the multiscale correction gives $N^{\mathrm{te}}\phi(q,p)-\psi(q,p)$.
%
%At the probability boundaries, the corrected discrepancy is $-\infty$ for $q=p\in\{0,1\}$ and $+\infty$ for $p=0<q$ or $p=1>q$. These conventions extend the constraint to every residual threshold.
%
%\begin{definition}[DW-based CDF ambiguity set]
%  \label{definition-dw-ambiguity-set}
%  Given the empirical CDF $\widehat F^{\mathrm{te}}$ and a threshold $\theta\ge0$, define
%  \begin{equation} 
%    \label{eq-cdf-ambiguity-set}
%      \mathcal Q^{\mathrm{DW}}(\theta):=\left\{\mathbb Q\in\mathcal M_1(\Xi):\ 
%      \sup_{t\in\mathbb R}\left[ N^{\mathrm{te}}\phi\!\left(\widehat F^{\mathrm{te}}(t),F_{\mathbb Q}(t)\right) 
%       {}-\psi\!\left(\widehat F^{\mathrm{te}}(t),F_{\mathbb Q}(t)\right)\right]\le\theta\right\}. 
%  \end{equation} 
%\end{definition}
%
%The threshold $\theta$ controls the corrected likelihood discrepancy. The induced CDF tolerances vary across probability levels and can be asymmetric about the empirical CDF. Statistical calibration assigns a confidence level to this set.
%
%We calibrate $\theta$ against a distribution-free reference law. Following \citet{L.J23}, let $V_1,\ldots,V_{N^{\mathrm{te}}}$ be independent standard uniform variables, independent of the training and test data. Their empirical CDF is
%\begin{equation*}
%  G_{N^{\mathrm{te}}}(u):=\frac{1}{N^{\mathrm{te}}}\sum_{j=1}^{N^{\mathrm{te}}}\mathbf 1_{\{V_j\le u\}},\qquad \forall u\in[0,1].
%\end{equation*}
%Define the uniform reference statistic
%\begin{equation}
%  T_{N^{\mathrm{te}}}^{\mathrm U}:=\sup_{u\in(0,1)}\left[N^{\mathrm{te}}\phi\!\left(G_{N^{\mathrm{te}}}(u),u\right)-\psi\!\left(G_{N^{\mathrm{te}}}(u),u\right)\right].
%  \label{eq-dw-reference-statistic}
%\end{equation}
%
%
%\begin{proposition}[Finite-sample threshold calibration]
%  \label{proposition-01}
%  Under the sampling conditions stated above, let $\delta\in(0,1)$ and define
%  \begin{equation}
%    \theta_{N^{\mathrm{te}}}(\delta):=\inf\left\{\tau\ge0:\Pr\left\{T_{N^{\mathrm{te}}}^{\mathrm U}\le\tau\right\}\ge1-\delta\right\}.
%    \label{eq-dw-calibrated-threshold}
%  \end{equation}
%
%  Then, for every $\theta\ge\theta_{N^{\mathrm{te}}}(\delta)$, the DW-based CDF ambiguity set satisfies
%  \begin{equation}
%    \Pr\left\{\mathbb Q^\star\in\mathcal Q^{\mathrm{DW}}(\theta)\,\middle|\,\mathcal D^{\mathrm{tr}}\right\}\ge1-\delta.
%    \label{eq-cdf-coverage}
%  \end{equation}
%\end{proposition}
%\begin{proof}
%  The reference statistic is finite almost surely, and its distribution function is right-continuous. Hence \eqref{eq-dw-calibrated-threshold} gives
%  \[
%    \Pr\left\{T_{N^{\mathrm{te}}}^{\mathrm U}>\theta_{N^{\mathrm{te}}}(\delta)\right\}\le\delta.
%  \]
%  Conditioning on $\mathcal D^{\mathrm{tr}}$ fixes $\hat g$ and $\mathbb Q^\star$. Define the generalized quantile function $F_{\mathbb Q^\star}^{-1}(u):=\inf\{t\in\mathbb R:F_{\mathbb Q^\star}(t)\ge u\}$ for all $u\in(0,1)$. The sample $F_{\mathbb Q^\star}^{-1}(V_j)$ for $j=1,\ldots,N^{\mathrm{te}}$, has the same conditional joint law as the test residuals. Thus, given $\mathcal D^{\mathrm{tr}}$,
%  \[
%    \left(\widehat F^{\mathrm{te}}(t)\right)_{t\in\mathbb R}\overset{\mathrm d}{=}\left(G_{N^{\mathrm{te}}}\!\left(F_{\mathbb Q^\star}(t)\right)\right)_{t\in\mathbb R}.
%  \]
%  For $0<F_{\mathbb Q^\star}(t)<1$, the discrepancy at $\mathbb Q^\star$ is evaluated at a probability level included in \eqref{eq-dw-reference-statistic}. When $F_{\mathbb Q^\star}(t)$ is zero or one, the empirical CDF equals that boundary almost surely, giving discrepancy $-\infty$. The supremum in \eqref{eq-cdf-ambiguity-set}, evaluated at $\mathbb Q^\star$, is therefore bounded by $T_{N^{\mathrm{te}}}^{\mathrm U}$ under this coupling. Since $\mathbb Q^\star(\Xi)=1$, we obtain
%  \[
%    \Pr\left\{\mathbb Q^\star\notin\mathcal Q^{\mathrm{DW}}\!\left(\theta_{N^{\mathrm{te}}}(\delta)\right)\,\middle|\,\mathcal D^{\mathrm{tr}}\right\}
%    \le\Pr\left\{T_{N^{\mathrm{te}}}^{\mathrm U}>\theta_{N^{\mathrm{te}}}(\delta)\right\}\le\delta.
%  \]
%  For every test sample, $\theta\ge\theta_{N^{\mathrm{te}}}(\delta)$ implies $\mathcal Q^{\mathrm{DW}}\!\left(\theta_{N^{\mathrm{te}}}(\delta)\right)\subseteq\mathcal Q^{\mathrm{DW}}(\theta)$. This inclusion proves \eqref{eq-cdf-coverage} for every such $\theta$.
%\end{proof}
%
%Proposition~\ref{proposition-01} applies to arbitrary residual distributions, including laws with atoms or infinite moments. For fixed $\phi$ and $\psi$, the reference law depends only on $N^{\mathrm{te}}$, so $\theta_{N^{\mathrm{te}}}(\delta)$ can be calibrated offline. \citet{L.J23} describe finite-sample evaluation using uniform order-statistic boundary probabilities. A conservative upper bound on the quantile in \eqref{eq-dw-calibrated-threshold} preserves \eqref{eq-cdf-coverage}.
%% If calibration is performed by simulation, quantile estimation error must be accounted for in the confidence guarantee. A threshold chosen from the test data also preserves \eqref{eq-cdf-coverage}, provided $\theta\ge\theta_{N^{\mathrm{te}}}(\delta)$ almost surely, by the same nesting argument.
%
%The residual-law coverage in Proposition~\ref{proposition-01} transfers to the failure probability bounds defined in Section~\ref{sec-problem}.
%
%\begin{theorem}[Finite-sample failure probability bounds]
%  \label{theorem-01}
%  Under the sampling conditions of Proposition~\ref{proposition-01}, let $\delta\in(0,1)$ and $\theta\ge\theta_{N^{\mathrm{te}}}(\delta)$, with the calibrated threshold defined in \eqref{eq-dw-calibrated-threshold}. Then the bounds defined in \eqref{eq-upper-bound} and \eqref{eq-lower-bound}, evaluated with $\mathcal Q=\mathcal Q^{\mathrm{DW}}(\theta)$, satisfy
%  \begin{equation}
%    \Pr\left\{\underline p_{\hat f}\le p_f\le\overline p_{\hat f}\,\middle|\,\mathcal D^{\mathrm{tr}}\right\}\ge1-\delta.
%    \label{eq-failure-probability-coverage}
%  \end{equation}
%\end{theorem}
%
%\begin{proof}
%  On the event $\mathbb Q^\star\in\mathcal Q^{\mathrm{DW}}(\theta)$, the true joint distribution $\Pi^\star$ belongs to $\mathcal P\left(\mathcal Q^{\mathrm{DW}}(\theta)\right)$. Hence
%  \[
%    \underline p_{\hat f}\le p_{\hat f}(\Pi^\star)=p_f\le\overline p_{\hat f}.
%  \]
%  By Proposition~\ref{proposition-01}, this event has conditional probability at least $1-\delta$ given $\mathcal D^{\mathrm{tr}}$, proving \eqref{eq-failure-probability-coverage}.
%\end{proof}
%
%
%The same event gives both inequalities, so $1-\delta$ is the confidence level for the entire interval. For fixed test data, $0\le\theta_1\le\theta_2$ implies $\mathcal Q^{\mathrm{DW}}(\theta_1)\subseteq\mathcal Q^{\mathrm{DW}}(\theta_2)$, making the lower bound nonincreasing and the upper bound nondecreasing in $\theta$. A higher prescribed confidence level requires a nondecreasing calibrated threshold and therefore gives a wider or unchanged interval.
%
%The same calibrated set can be reused across response thresholds.
%
%\begin{corollary}[Simultaneous bounds across response thresholds]
%  \label{cor-simultaneous-threshold-bounds}
%  Under the conditions of Theorem~\ref{theorem-01}, let $\mathcal Q=\mathcal Q^{\mathrm{DW}}(\theta)$. For $c\in\mathbb R$, write $p_f(c):=\Pr \left\{g(\bm X)\le c\right\}$ and define
%  \[
%    p_{\hat f}(\Pi;c):=\int_{\mathcal X\times\Xi}
%      \mathbf1_{\{\hat g(\bm x)+\xi\le c\}}\,\Pi(\mathrm d\bm x,\mathrm d\xi),
%  \]
%  with lower and upper bounds
%  \[
%    \underline p_{\hat f}(c):=\inf_{\Pi\in\mathcal P(\mathcal Q)}p_{\hat f}(\Pi;c),
%    \qquad
%    \overline p_{\hat f}(c):=\sup_{\Pi\in\mathcal P(\mathcal Q)}p_{\hat f}(\Pi;c).
%  \]
%  Then
%  \begin{equation}
%    \Pr\left\{\underline p_{\hat f}(c)\le p_f(c)\le\overline p_{\hat f}(c),
%      \ \forall c\in\mathbb R\,\middle|\,\mathcal D^{\mathrm{tr}}\right\}\ge1-\delta.
%    \label{eq-threshold-curve-coverage}
%  \end{equation}
%  At $c=0$, these definitions recover the original failure probability and its bounds.
%\end{corollary}
%\begin{proof}
%  On $\mathbb Q^\star\in\mathcal Q^{\mathrm{DW}}(\theta)$, the true joint law $\Pi^\star$ belongs to the same set $\mathcal P(\mathcal Q)$ for every $c$. Consequently,
%  \[ 
%    \underline p_{\hat f}(c) =\inf_{\Pi\in\mathcal P(\mathcal Q)} p_{\hat f}(\Pi;c) \le p_{\hat f}(\Pi^\star;c) =p_f(c) \le\sup_{\Pi\in\mathcal P(\mathcal Q)} p_{\hat f}(\Pi;c) =\overline p_{\hat f}(c), \qquad\forall c\in\mathbb R.
%  \]
%  Both bound curves are nondecreasing. Monotonicity and the CDF limits of $p_f$ imply that any violation is detected at a rational threshold or an atom of $g(\bm X)$, a countable set. The simultaneous event is therefore measurable. Proposition~\ref{proposition-01} gives the common residual-coverage event conditional probability at least $1-\delta$, proving \eqref{eq-threshold-curve-coverage}.
%\end{proof}
%
%One calibration therefore covers the entire response-threshold curve, including thresholds selected after observing the test data, without an additional confidence allocation across $c$. Complementing the endpoints also bounds the strict exceedance probabilities $\Pr \left\{g(\bm X)>c\right\}$ simultaneously. Here $g$, $\hat g$, and the input law $\mathbb P$ remain fixed. Fragility curves formed by conditioning on different hazard intensities require coverage for the corresponding conditional residual laws.
%
%The coverage argument also applies to any residual ambiguity set containing $\mathbb Q^\star$ with conditional probability at least $1-\delta$. It remains to evaluate the DW extrema in \eqref{eq-upper-bound} and \eqref{eq-lower-bound} and their response-threshold extensions.
%
%
%\section{Tractable Reformulation}
%\label{sec-reformulation}
%
%Throughout this section, the training data, the test residuals, and $\theta\ge0$ are fixed. We use the observed residuals to partition both the residual and input spaces. These partitions express the DW restrictions as cumulative cell probabilities and reduce the Lagrangian duals to finite linear programs. A joint-distribution construction establishes exact finite primal representations, after which LP duality and the ordering of the cells yield strong duality and closed-form bounds. The main development takes $\Xi=\mathbb R$; Remark~\ref{rem-interval-support} treats known interval support.
%
%For each empirical probability level $q_j:=j/N^{\mathrm{te}}$, define the scalar acceptance interval
%\begin{equation}
%  [\alpha_j(\theta),\beta_j(\theta)]
%  :=\left\{p\in[0,1]:
%  N^{\mathrm{te}}\phi(q_j,p)-\psi(q_j,p)\le\theta\right\},
%  \qquad j=0,\ldots,N^{\mathrm{te}}.
%  \label{eq-dw-rank-acceptance}
%\end{equation}
%Let $\xi_{(1)}^{\mathrm{te}}<\cdots<\xi_{(K)}^{\mathrm{te}}$ be the distinct ordered test residuals, and set
%\begin{equation}
%  c_0:=0,\quad
%  c_k:=\sum_{j=1}^{N^{\mathrm{te}}}\mathbf1_{\{\xi_j^{\mathrm{te}}\le\xi_{(k)}^{\mathrm{te}}\}},\quad
%  l_k:=\alpha_{c_k}(\theta),\quad
%  u_k:=\beta_{c_{k-1}}(\theta),\qquad k=1,\ldots,K.
%  \label{eq-dw-counts-and-levels}
%\end{equation}
%The counts retain all multiplicities, so $c_K=N^{\mathrm{te}}$ even when test residuals coincide. Set $\xi_{(0)}^{\mathrm{te}}:=-\infty$ and $\xi_{(K+1)}^{\mathrm{te}}:=+\infty$. Separating the observed values from the intervening intervals gives $J:=2K+1$ residual cells:
%\begin{equation}
%  \Xi_{2k}:=\left\{\xi_{(k)}^{\mathrm{te}}\right\}
%    \quad\left(k=1,\ldots,K\right),
%  \qquad
%  \Xi_{2k+1}:=\left(\xi_{(k)}^{\mathrm{te}},\xi_{(k+1)}^{\mathrm{te}}\right)
%    \quad\left(k=0,\ldots,K\right).
%  \label{eq-residual-cells}
%\end{equation}
%For any candidate residual law $\mathbb Q\in\mathcal M_1(\Xi)$, the open rank probability measures mass strictly below an observed residual, while the closed rank probability also includes mass at that value. For $k=1,\ldots,K$, the partition gives
%\[
%  \mathbb Q\left\{\xi<\xi_{(k)}^{\mathrm{te}}\right\}
%    =\sum_{j<2k}\mathbb Q(\Xi_j),
%  \qquad
%  \mathbb Q\left\{\xi\le\xi_{(k)}^{\mathrm{te}}\right\}
%    =\sum_{j\le2k}\mathbb Q(\Xi_j).
%\]
%The two probabilities differ by the atom mass $\mathbb Q(\Xi_{2k})$.
%
%For $i=0,\ldots,K$, the observed residual levels classify inputs by the location of the failure threshold $-\hat g(\bm x)$:
%\begin{equation}
%  \mathcal X_i:=\left\{\bm x\in\mathcal X:
%    \xi_{(i)}^{\mathrm{te}}\le-\hat g(\bm x)<\xi_{(i+1)}^{\mathrm{te}}\right\}.
%  \label{eq-input-cells}
%\end{equation}
%The input cell probabilities give the surrogate-response threshold probabilities through
%\[
%  \mathbb P\left\{\hat g(\bm X)\le-\xi_{(k)}^{\mathrm{te}}\right\}
%    =\sum_{i=k}^{K}\mathbb P(\mathcal X_i),
%  \qquad k=1,\ldots,K.
%\]
%The half-open convention assigns $-\hat g(\bm x)=\xi_{(k)}^{\mathrm{te}}$ to $\mathcal X_k$. The tail sum therefore includes threshold equalities even when the surrogate response has atoms. 
%The two partitions satisfy
%\[
%  \mathcal X=\bigsqcup_{i=0}^{K}\mathcal X_i,
%  \qquad
%  \Xi=\bigsqcup_{j=1}^{J}\Xi_j.
%\]
%Figure~\ref{fig-residual-input-cells} illustrates the two partitions for $K=3$.
%\begin{figure}[!htbp]
%  \centering
%  \includegraphics[width=0.5\columnwidth]{fig_residual_input_cells.png}
%  \caption{Illustration of the residual and input-threshold partitions for $K=3$. The residual partition consists of $K$ singleton cells and $K+1$ open intervals, whereas the input partition is induced by half-open intervals of the failure threshold $-\hat g(\bm x)$. Here, $w_i=\mathbb P(\mathcal X_i)$ denotes the probability mass of input cell $\mathcal X_i$.}
%  \label{fig-residual-input-cells}
%\end{figure}
%
%\begin{proposition}[Finite representation of the DW ambiguity set]
%  \label{prop-finite-rank-constraints}
%  Each acceptance set in \eqref{eq-dw-rank-acceptance} is a nonempty closed interval containing $q_j$. The DW ambiguity set has the finite representation
%  \begin{equation}
%    \mathcal Q^{\mathrm{DW}}(\theta)
%    =\left\{\mathbb Q\in\mathcal M_1(\Xi):
%      \sum_{j<2k}\mathbb Q(\Xi_j)\le u_k,\quad
%      \sum_{j\le2k}\mathbb Q(\Xi_j)\ge l_k,\quad k=1,\ldots,K\right\}.
%    \label{eq-finite-rank-constraints}
%  \end{equation}
%  The sequences $l_k$ and $u_k$ are nondecreasing and satisfy
%  \[
%    0\le l_k,u_k\le1,\qquad l_{k-1}\le u_k,\qquad l_k\le u_{k+1},
%    \qquad k=1,\ldots,K,
%  \]
%  where $l_0:=0$ and $u_{K+1}:=1$.
%\end{proposition}
%\begin{proof}
%  For fixed $q$, the Bernoulli divergence increases as $p$ moves away from $q$, while $\psi(q,p)$ cannot increase as the interval between $q$ and $p$ expands. The corrected discrepancy therefore decreases toward $q$ from below and increases above $q$. It is continuous for $p\in(0,1)$ and equals $-\psi_0(q)\le0$ at $p=q$. The boundary conventions in Section~\ref{sec-ambiguity-set} give a nonempty closed acceptance interval also for $q\in\{0,1\}$. Scalar inversion of the DW constraint therefore gives the CDF band \citep[Supplement, Section~S.6]{L.J23}
%  \[
%    \alpha_{N^{\mathrm{te}}\widehat F^{\mathrm{te}}(t)}(\theta)
%      \le F_{\mathbb Q}(t)
%      \le\beta_{N^{\mathrm{te}}\widehat F^{\mathrm{te}}(t)}(\theta),
%    \qquad\forall t\in\mathbb R.
%  \]
%
%  For $k=1,\ldots,K-1$, the empirical CDF on $\left[\xi_{(k)}^{\mathrm{te}},\xi_{(k+1)}^{\mathrm{te}}\right)$ equals $c_k/N^{\mathrm{te}}$. Monotonicity makes the lower band constraint on this interval equivalent to $\mathbb Q\left\{\xi\le\xi_{(k)}^{\mathrm{te}}\right\}\ge\alpha_{c_k}(\theta)$. By continuity of probability from below, the upper band constraint is equivalent to $\mathbb Q\left\{\xi<\xi_{(k+1)}^{\mathrm{te}}\right\}\le\beta_{c_k}(\theta)$. The exterior intervals give $\mathbb Q\left\{\xi<\xi_{(1)}^{\mathrm{te}}\right\}\le\beta_0(\theta)$ and $\mathbb Q\left\{\xi\le\xi_{(K)}^{\mathrm{te}}\right\}\ge\alpha_{N^{\mathrm{te}}}(\theta)$; their other band constraints are automatic because $\alpha_0(\theta)=0$ and $\beta_{N^{\mathrm{te}}}(\theta)=1$. The partition identities above then yield \eqref{eq-finite-rank-constraints}.
%
%  With $p$ fixed, the corrected discrepancy decreases as $q$ approaches $p$ from either side. For empirical levels $q<q'$, an accepted value $p<q$ at level $q'$ remains accepted at level $q$; an accepted value $p>q'$ at level $q$ remains accepted at level $q'$. Since each acceptance interval contains its empirical level, both endpoints are nondecreasing. The inequalities $\alpha_{c_k}(\theta)\le c_k/N^{\mathrm{te}}\le\beta_{c_k}(\theta)$ give the stated compatibility conditions. Singleton intervals are allowed: if $q_j=1/2$ and $\theta=0$, the interval is $\{1/2\}$ because $\psi(1/2,p)=0$.
%\end{proof}
%
%Each scalar interval can be computed by one-dimensional bisection. Proposition~\ref{prop-finite-rank-constraints} then restricts only the probabilities of the residual cells; the distribution within each cell remains free.
%
%For a feasible joint law, define its matrix of joint cell probabilities by $\pi_{ij}:=\Pi(\mathcal X_i\times\Xi_j)$. The fixed input marginal and the residual restrictions imply $\bm\pi=(\pi_{ij})\in\mathcal C_\theta$, where
%\begin{equation}
%  \mathcal C_\theta:=\left\{\bm\pi\in\mathbb R_+^{(K+1)\times J}:
%  \begin{array}{ll}
%    \displaystyle\sum_{j=1}^{J}\pi_{ij}=\mathbb P(\mathcal X_i),&i=0,\ldots,K,\\[1mm]
%    \displaystyle\sum_{i=0}^{K}\sum_{j<2k}\pi_{ij}\le u_k,&k=1,\ldots,K,\\[1mm]
%    \displaystyle\sum_{i=0}^{K}\sum_{j\le2k}\pi_{ij}\ge l_k,&k=1,\ldots,K
%  \end{array}\right\}.
%  \label{eq-cell-mass-polytope}
%\end{equation}
%On each product cell, the failure indicator satisfies
%\begin{equation}
%  \underline h_{ij}:=\mathbf1_{\{j\le2i\}}
%  \le\mathbf1_{\{\xi\le-\hat g(\bm x)\}}
%  \le\overline h_{ij}:=\mathbf1_{\{j\le2i+1\}},
%  \qquad \forall (\bm x,\xi)\in\mathcal X_i\times\Xi_j.
%  \label{eq-cell-costs}
%\end{equation}
%Residual cells with $j\le2i$ always cause failure for inputs in $\mathcal X_i$, and those with $j\ge2i+2$ are always safe. Only the overlapping interval $\Xi_{2i+1}$ has different upper and lower coefficients. Figure~\ref{fig-joint-partition} illustrates the resulting joint cell structure and the classification of product cells for $K=3$. Thus the cell matrix records the rank constraints exactly, while the failure objective requires the constructions below.
%\begin{figure}[!htbp]
%  \centering
%  \includegraphics[width=0.5\columnwidth]{fig_joint_partition.png}
%  \caption{Joint partition of the input and residual spaces for $K=3$. Each product cell has probability mass $\pi_{ij}=\Pi(\mathcal X_i\times\Xi_j)$. The blue, yellow, and red cells correspond to $(\underline h_{ij},\overline h_{ij})=(0,0)$, $(0,1)$, and $(1,1)$, respectively.}
%  \label{fig-joint-partition}
%\end{figure}
%
%The upper problem remains a linear program over joint measures, with the full input marginal fixed:
%\begin{equation}
%  \begin{aligned}
%    \overline p_{\hat f}
%    =\sup_{\Pi\in\mathcal M_+(\mathcal X\times\Xi)}\quad&
%      \int_{\mathcal X\times\Xi}\mathbf1_{\{\xi\le-\hat g(\bm x)\}}\,\Pi(\mathrm d\bm x,\mathrm d\xi)\\
%    \text{s.t.}\quad&
%      \int_{\mathcal X\times\Xi}\mathbf1_{\mathcal U}(\bm x)\,\Pi(\mathrm d\bm x,\mathrm d\xi)
%        =\mathbb P(\bm X\in\mathcal U),&&\forall\mathcal U\in\mathcal B(\mathcal X),\\
%    &\int_{\mathcal X\times\Xi}\mathbf1_{\{\xi<\xi_{(k)}^{\mathrm{te}}\}}\,\Pi(\mathrm d\bm x,\mathrm d\xi)\le u_k,
%        &&k=1,\ldots,K,\\
%    &\int_{\mathcal X\times\Xi}\mathbf1_{\{\xi\le\xi_{(k)}^{\mathrm{te}}\}}\,\Pi(\mathrm d\bm x,\mathrm d\xi)\ge l_k,
%        &&k=1,\ldots,K.
%  \end{aligned}
%  \label{eq-measure-upper}
%\end{equation}
%The lower problem minimizes the same integral under the same constraints. The marginal equalities fix both the input law and total mass one. The product of $\mathbb P$ and the empirical residual law is feasible for every $\theta\ge0$.
%
%\begin{lemma}[Finite-dimensional Lagrangian duals and weak duality]
%  \label{lem-weak-duality}
%  Write $\bm u=(u_1,\ldots,u_K)^\top$ and $\bm l=(l_1,\ldots,l_K)^\top$. The nonnegative multipliers $\bm\lambda,\bm\gamma\in\mathbb R_+^K$ correspond to the upper and lower rank constraints. The input marginal equality has an unrestricted multiplier $\eta\in\mathcal F_{\mathrm b}(\mathcal X)$, the space of bounded Borel measurable real-valued functions. Define the residual penalty
%  \begin{equation}
%    \zeta_{\bm\lambda,\bm\gamma}(\xi)
%    :=\sum_{k=1}^{K}\lambda_k\mathbf1_{\{\xi<\xi_{(k)}^{\mathrm{te}}\}}
%      -\sum_{k=1}^{K}\gamma_k\mathbf1_{\{\xi\le\xi_{(k)}^{\mathrm{te}}\}},
%    \label{eq-dual-step-function}
%  \end{equation}
%  whose value on $\Xi_j$ is
%  \begin{equation}
%    \zeta_j:=\sum_{k=1}^{K}\lambda_k\mathbf1_{\{j<2k\}}
%      -\sum_{k=1}^{K}\gamma_k\mathbf1_{\{j\le2k\}},
%    \qquad j=1,\ldots,J.
%    \label{eq-dual-cell-values}
%  \end{equation}
%  Eliminating the input marginal multiplier gives the upper and lower Lagrangian dual problems
%  \begin{align}
%    \overline d_\theta
%      &:=\inf_{\bm\lambda,\bm\gamma\ge\bm0}\left\{
%        \bm\lambda^\top\bm u-\bm\gamma^\top\bm l
%        +\sum_{i=0}^{K}\mathbb P(\mathcal X_i)\max_{1\le j\le J}
%          \left\{\overline h_{ij}-\zeta_j\right\}\right\},
%      \label{eq-upper-measure-dual}\\
%    \underline d_\theta
%      &:=\sup_{\bm\lambda,\bm\gamma\ge\bm0}\left\{
%        -\bm\lambda^\top\bm u+\bm\gamma^\top\bm l
%        +\sum_{i=0}^{K}\mathbb P(\mathcal X_i)\min_{1\le j\le J}
%          \left\{\underline h_{ij}+\zeta_j\right\}\right\}.
%      \label{eq-lower-measure-dual}
%  \end{align}
%  Their optimal values satisfy
%  \begin{equation}
%    \underline d_\theta\le\underline p_{\hat f}
%      \le\overline p_{\hat f}\le\overline d_\theta.
%    \label{eq-weak-duality}
%  \end{equation}
%\end{lemma}
%\begin{proof}
%  The input marginal equality is equivalent to
%  \[
%    \int_{\mathcal X\times\Xi}\eta(\bm x)\,\Pi(\mathrm d\bm x,\mathrm d\xi)
%      =\mathbb E_{\mathbb P}\left[\eta(\bm X)\right],
%    \qquad\forall\eta\in\mathcal F_{\mathrm b}(\mathcal X).
%  \]
%  Relaxing this equality and the rank inequalities gives the upper Lagrangian \citep[Section~5.1]{B.V04}
%  \[
%    \begin{aligned}
%      \mathcal L^+(\Pi,\eta,\bm\lambda,\bm\gamma)
%      ={}&\int_{\mathcal X\times\Xi}\mathbf1_{\{\xi\le-\hat g(\bm x)\}}\,\Pi(\mathrm d\bm x,\mathrm d\xi)
%        +\mathbb E_{\mathbb P}\left[\eta(\bm X)\right]
%        -\int_{\mathcal X\times\Xi}\eta(\bm x)\,\Pi(\mathrm d\bm x,\mathrm d\xi)\\
%      &+\sum_{k=1}^{K}\lambda_k\left(u_k-\int_{\mathcal X\times\Xi}
%          \mathbf1_{\{\xi<\xi_{(k)}^{\mathrm{te}}\}}\,\Pi(\mathrm d\bm x,\mathrm d\xi)\right)
%        +\sum_{k=1}^{K}\gamma_k\left(\int_{\mathcal X\times\Xi}
%          \mathbf1_{\{\xi\le\xi_{(k)}^{\mathrm{te}}\}}\,\Pi(\mathrm d\bm x,\mathrm d\xi)-l_k\right)\\
%      ={}&\mathbb E_{\mathbb P}\left[\eta(\bm X)\right]
%        +\bm\lambda^\top\bm u-\bm\gamma^\top\bm l
%        +\int_{\mathcal X\times\Xi}\left[
%          \mathbf1_{\{\xi\le-\hat g(\bm x)\}}-\eta(\bm x)-\zeta_{\bm\lambda,\bm\gamma}(\xi)
%        \right]\,\Pi(\mathrm d\bm x,\mathrm d\xi).
%    \end{aligned}
%  \]
%  Its supremum over finite nonnegative measures is finite exactly when the integrand is nonpositive at every $(\bm x,\xi)$. A positive value at any point permits an arbitrarily large point mass there; otherwise the supremum integral is zero, attained by the zero measure. The functional upper dual therefore minimizes $\mathbb E_{\mathbb P}[\eta(\bm X)]+\bm\lambda^\top\bm u-\bm\gamma^\top\bm l$ subject to $\eta(\bm x)\ge\mathbf1_{\{\xi\le-\hat g(\bm x)\}}-\zeta_{\bm\lambda,\bm\gamma}(\xi)$ at every pair. For the lower problem, reversing the rank penalties gives
%  \[
%    \mathcal L^-
%      =\mathbb E_{\mathbb P}\left[\eta(\bm X)\right]
%        -\bm\lambda^\top\bm u+\bm\gamma^\top\bm l
%        +\int_{\mathcal X\times\Xi}\left[
%          \mathbf1_{\{\xi\le-\hat g(\bm x)\}}-\eta(\bm x)+\zeta_{\bm\lambda,\bm\gamma}(\xi)
%        \right]\,\Pi(\mathrm d\bm x,\mathrm d\xi).
%  \]
%  Its infimum over finite nonnegative measures is finite exactly when the integrand is nonnegative everywhere, in which case the infimum integral is zero. The lower functional dual maximizes $\mathbb E_{\mathbb P}[\eta(\bm X)]-\bm\lambda^\top\bm u+\bm\gamma^\top\bm l$ subject to $\eta(\bm x)\le\mathbf1_{\{\xi\le-\hat g(\bm x)\}}+\zeta_{\bm\lambda,\bm\gamma}(\xi)$ at every pair.
%
%  For fixed rank multipliers, the best marginal functions are the pointwise envelopes. For $\bm x\in\mathcal X_i$, they satisfy the identities
%  \begin{equation}
%    \begin{aligned}
%      \eta^+_{\bm\lambda,\bm\gamma}(\bm x)
%        &:=\sup_{\xi\in\Xi}\left\{\mathbf1_{\{\xi\le-\hat g(\bm x)\}}
%          -\zeta_{\bm\lambda,\bm\gamma}(\xi)\right\}
%          =\max_{1\le j\le J}\left\{\overline h_{ij}-\zeta_j\right\},\\
%      \eta^-_{\bm\lambda,\bm\gamma}(\bm x)
%        &:=\inf_{\xi\in\Xi}\left\{\mathbf1_{\{\xi\le-\hat g(\bm x)\}}
%          +\zeta_{\bm\lambda,\bm\gamma}(\xi)\right\}
%          =\min_{1\le j\le J}\left\{\underline h_{ij}+\zeta_j\right\}.
%    \end{aligned}
%    \label{eq-dual-envelope-identities}
%  \end{equation}
%  To verify them, note that outside $\Xi_{2i+1}$ the two cell coefficients equal the actual failure indicator. In the overlapping interval, the lower coefficient zero is attained by $\xi=\left(\xi_{(i+1)}^{\mathrm{te}}-\hat g(\bm x)\right)/2$ if $i<K$, and by $\xi=-\hat g(\bm x)+1$ if $i=K$. The upper coefficient one is attained by $\xi=-\hat g(\bm x)$ whenever $-\hat g(\bm x)>\xi_{(i)}^{\mathrm{te}}$. The only remaining case is $-\hat g(\bm x)=\xi_{(i)}^{\mathrm{te}}$ with $i\ge1$. Then the overlapping interval is safe, but
%  \[
%    \zeta_{2i}=\zeta_{2i+1}-\gamma_i\le\zeta_{2i+1},
%    \qquad
%    1-\zeta_{2i}\ge1-\zeta_{2i+1}.
%  \]
%  Choosing the left endpoint therefore attains an upper-envelope value at least as large as the candidate from the overlapping interval. Together with \eqref{eq-cell-costs}, this proves both identities, including threshold equalities.
%
%  The envelopes are bounded and constant on the Borel cells $\mathcal X_i$, so they belong to $\mathcal F_{\mathrm b}(\mathcal X)$. They respectively minimize and maximize the expectation over feasible marginal functions. Substituting them into the functional dual objectives gives \eqref{eq-upper-measure-dual} and \eqref{eq-lower-measure-dual}. Introducing epigraph variables $z_i$ for the upper maxima and hypograph variables $z_i$ for the lower minima gives the equivalent LPs
%  \begin{equation}
%    \begin{aligned}
%      \overline d_\theta
%      =\inf_{\bm\lambda,\bm\gamma\ge\bm0,\ \bm z\in\mathbb R^{K+1}}\quad&
%        \bm\lambda^\top\bm u-\bm\gamma^\top\bm l+\sum_{i=0}^{K}\mathbb P(\mathcal X_i) z_i\\
%      \text{s.t.}\quad&z_i+\zeta_j\ge\overline h_{ij},
%        \qquad i=0,\ldots,K,\quad j=1,\ldots,J,
%    \end{aligned}
%    \label{eq-finite-upper-dual}
%  \end{equation}
%  and
%  \begin{equation}
%    \begin{aligned}
%      \underline d_\theta
%      =\sup_{\bm\lambda,\bm\gamma\ge\bm0,\ \bm z\in\mathbb R^{K+1}}\quad&
%        -\bm\lambda^\top\bm u+\bm\gamma^\top\bm l+\sum_{i=0}^{K}\mathbb P(\mathcal X_i) z_i\\
%      \text{s.t.}\quad&z_i-\zeta_j\le\underline h_{ij},
%        \qquad i=0,\ldots,K,\quad j=1,\ldots,J.
%    \end{aligned}
%    \label{eq-finite-lower-dual}
%  \end{equation}
%  Each $z_i$ can be chosen as its corresponding finite envelope, including on input cells with $\mathbb P(\mathcal X_i)=0$. Thus introducing these variables preserves both optimal values.
%
%  For a feasible joint law, the signs of the rank penalties give $\mathcal L^-\le p_{\hat f}(\Pi)\le\mathcal L^+$. Feasible lower dual objectives therefore bound the primal objectives from below, and feasible upper dual objectives bound them from above. Taking the respective optima over the common nonempty primal feasible set proves \eqref{eq-weak-duality}.
%\end{proof}
%
%The duals in \eqref{eq-upper-measure-dual} and \eqref{eq-lower-measure-dual} involve $2K$ nonnegative rank multipliers. Their equivalent LP forms \eqref{eq-finite-upper-dual} and \eqref{eq-finite-lower-dual} add $K+1$ free variables $z_i$ and $(K+1)(2K+1)$ pairwise inequalities. The envelope identities establish these representations before strong duality is invoked.
%
%\begin{theorem}[Exact finite reformulations and strong duality]
%  \label{thm-finite-dual-reformulations}
%  For $\mathcal Q=\mathcal Q^{\mathrm{DW}}(\theta)$ and every $\theta\ge0$, the failure probability bounds admit the exact representations
%  \begin{equation}
%    \overline p_{\hat f}
%      =\max_{\bm\pi\in\mathcal C_\theta}\sum_{i=0}^{K}\sum_{j=1}^{J}\overline h_{ij}\pi_{ij},
%    \qquad
%    \underline p_{\hat f}
%      =\min_{\bm\pi\in\mathcal C_\theta}\sum_{i=0}^{K}\sum_{j=1}^{J}\underline h_{ij}\pi_{ij}.
%    \label{eq-finite-primal-values}
%  \end{equation}
%  Their Lagrangian dual values in Lemma~\ref{lem-weak-duality} satisfy
%  \begin{equation}
%    \overline p_{\hat f}=\overline d_\theta,
%    \qquad
%    \underline p_{\hat f}=\underline d_\theta.
%    \label{eq-strong-duality}
%  \end{equation}
%  The finite primal and dual optima, the measure primal optima, and the functional dual optima are attained.
%\end{theorem}
%\begin{proof}
%  Any feasible joint law induces a matrix in $\mathcal C_\theta$. Integrating \eqref{eq-cell-costs} therefore bounds the upper measure problem by the finite maximization from above and the lower measure problem by the finite minimization from below.
%
%  For the reverse inequalities, fix $\bm\pi\in\mathcal C_\theta$. Draw $\bm x$ from $\mathbb P$. On $\mathcal X_i$ with $\mathbb P(\mathcal X_i)>0$, select a provisional residual cell $j$ with probability $\pi_{ij}/\mathbb P(\mathcal X_i)$; the row equality makes these probabilities sum to one. For $j\ne2i+1$, choose a fixed representative of $\Xi_j$, where the actual failure indicator equals both coefficients. In the overlapping cell, use the following assignments to define joint laws $\Pi^+$ and $\Pi^-$.
%
%  For the upper problem, assign $\xi=-\hat g(\bm x)$. This causes failure and remains in $\Xi_{2i+1}$ unless $-\hat g(\bm x)=\xi_{(i)}^{\mathrm{te}}$. At equality, which requires $i\ge1$, the assignment transfers the provisional mass to the left endpoint. For any $\widetilde\xi\in\left(\xi_{(i)}^{\mathrm{te}},\xi_{(i+1)}^{\mathrm{te}}\right)$ and every observed threshold,
%  \[
%    \mathbf1_{\{\xi_{(i)}^{\mathrm{te}}<\xi_{(k)}^{\mathrm{te}}\}}
%      =\mathbf1_{\{\widetilde\xi<\xi_{(k)}^{\mathrm{te}}\}},
%    \qquad
%    \mathbf1_{\{\xi_{(i)}^{\mathrm{te}}\le\xi_{(k)}^{\mathrm{te}}\}}
%      \ge\mathbf1_{\{\widetilde\xi\le\xi_{(k)}^{\mathrm{te}}\}}.
%  \]
%  Thus every open upper probability is unchanged and every closed lower probability can only increase. The upper assignment preserves all rank inequalities and attains the coefficient $\overline h_{i,2i+1}=1$.
%
%  For the lower problem, assign $\xi=\left(\xi_{(i+1)}^{\mathrm{te}}-\hat g(\bm x)\right)/2$ if $i<K$, and $\xi=-\hat g(\bm x)+1$ if $i=K$. These values remain in $\Xi_{2i+1}$ and strictly exceed the failure threshold, including at threshold equality. Hence the lower assignment retains the cell matrix and attains $\underline h_{i,2i+1}=0$.
%
%  If $\mathbb P(\mathcal X_i)=0$, the corresponding row is zero; set $\xi=\xi_{(1)}^{\mathrm{te}}$ on that input cell. All assignments are Borel measurable, and their finite mixtures define probability kernels. Combining them with $\mathbb P$ gives joint laws with the full prescribed input marginal, including randomization at input atoms. Both laws are feasible and satisfy
%  \[
%    p_{\hat f}(\Pi^+)=\sum_{i=0}^{K}\sum_{j=1}^{J}\overline h_{ij}\pi_{ij},
%    \qquad
%    p_{\hat f}(\Pi^-)=\sum_{i=0}^{K}\sum_{j=1}^{J}\underline h_{ij}\pi_{ij}.
%  \]
%  The lower law retains every provisional cell mass. The upper law may change the cell matrix at threshold equalities, but it preserves feasibility and the assigned upper objective. These identities establish the reverse inequalities required for \eqref{eq-finite-primal-values}.
%
%  For every $\theta\ge0$, the empirical product law induces the feasible matrix
%  \[
%    \pi_{i,2k}^{\mathrm{emp}}:=\mathbb P(\mathcal X_i)\frac{c_k-c_{k-1}}{N^{\mathrm{te}}},
%    \qquad i=0,\ldots,K,\quad k=1,\ldots,K,
%  \]
%  with all odd columns zero. Its open and closed rank probabilities are $c_{k-1}/N^{\mathrm{te}}$ and $c_k/N^{\mathrm{te}}$, respectively, so feasibility includes $\theta=0$. The nonnegative entries of a matrix in $\mathcal C_\theta$ sum to one, and the constraints are closed. Thus $\mathcal C_\theta$ is nonempty and compact. The finite objectives lie in $[0,1]$ and attain their optima, which the joint-law constructions realize in the measure problems.
%
%  The finite primal LPs in \eqref{eq-finite-primal-values} have row multipliers $z_i$ and rank multipliers $\lambda_k,\gamma_k$. Their dual constraints are precisely $z_i+\zeta_j\ge\overline h_{ij}$ and $z_i-\zeta_j\le\underline h_{ij}$. Finite LP strong duality therefore gives \eqref{eq-strong-duality} and dual attainment \citep[Section~5.2.3]{B.V04}. By \eqref{eq-dual-envelope-identities}, optimal rank multipliers also yield optimal bounded marginal functions for the original functional duals. Feasibility and finite optimal values suffice; strict rank inequalities are unnecessary, including at $\theta=0$.
%\end{proof}
%
%The finite formulations are exact, and their ordered failure coefficients give a further simplification. The closed-form bounds below follow from cuts between consecutive cells and admit certificates using at most one rank multiplier.
%
%\begin{corollary}[Closed-form bounds and optimal certificates]
%  \label{thm-closed-form-bounds}
%  Under the conditions of Theorem~\ref{thm-finite-dual-reformulations},
%  \begin{align}
%    \overline p_{\hat f}
%      &=\min\left\{1,\min_{k=1,\ldots,K}\left[u_k+\sum_{i=k}^{K}\mathbb P(\mathcal X_i)\right]\right\},
%      \label{eq-dw-closed-upper}\\
%    \underline p_{\hat f}
%      &=\max\left\{0,\max_{k=1,\ldots,K}\left[l_k+\sum_{i=k}^{K}\mathbb P(\mathcal X_i)-1\right]\right\}.
%      \label{eq-dw-closed-lower}
%  \end{align}
%  Each dual optimum admits a certificate with at most one nonzero rank multiplier, of value one.
%\end{corollary}
%\begin{proof}
%  For a candidate residual law $\mathbb Q\in\mathcal M_1(\Xi)$, define the finite transportation set
%  \[
%    \mathcal T(\mathbb Q):=\left\{\bm\pi\in\mathbb R_+^{(K+1)\times J}:
%    \begin{array}{ll}
%      \displaystyle\sum_{j=1}^{J}\pi_{ij}=\mathbb P(\mathcal X_i),&i=0,\ldots,K,\\[1mm]
%      \displaystyle\sum_{i=0}^{K}\pi_{ij}=\mathbb Q(\Xi_j),&j=1,\ldots,J
%    \end{array}\right\}.
%  \]
%  It contains the product matrix $\pi_{ij}=\mathbb P(\mathcal X_i)\mathbb Q(\Xi_j)$. Every residual cell is nonempty, so any nonnegative column allocation with total mass one is realized by a residual law. Proposition~\ref{prop-finite-rank-constraints} consequently gives
%  \[
%    \mathcal C_\theta
%      =\bigcup_{\mathbb Q\in\mathcal Q^{\mathrm{DW}}(\theta)}\mathcal T(\mathbb Q).
%  \]
%  The parameter $\mathbb Q$ fixes only residual partition masses in this finite matrix problem. The joint-law construction in Theorem~\ref{thm-finite-dual-reformulations} preserves feasibility and the assigned objective while allowing the residual marginal to change, including endpoint transfers in the upper problem.
%
%  For fixed cell masses $\mathbb Q(\Xi_j)$, the upper problem maximizes mass on pairs $j\le2i+1$; the lower problem maximizes safe mass on pairs $j>2i$ and subtracts it from one. Both are maximum-flow problems with source-to-row capacities $\mathbb P(\mathcal X_i)$, column-to-sink capacities $\mathbb Q(\Xi_j)$, and capacity two on each permitted row-to-column edge. A transportation matrix restricts to a feasible flow. Conversely, any flow can be completed to a transportation matrix because the remaining row and column masses are nonnegative with equal totals. Thus the maximum flow equals the maximum permitted mass.
%
%  A cut of capacity one is available, so a minimum cut crosses no edge of capacity two. Its source-side columns can be taken to be exactly the neighbors of its source-side rows. Since the neighborhoods are nested, adding rows whose neighbors are already included cannot increase the capacity. Apart from empty and full row sets, upper-problem cuts reduce to rows $i<k$ and columns $j<2k$, while safe-mass cuts reduce to rows $i\ge k$ and columns $j>2k$. For $k=1,\ldots,K$, their capacities are respectively
%  \[
%    \sum_{i=k}^{K}\mathbb P(\mathcal X_i)+\sum_{j<2k}\mathbb Q(\Xi_j),
%    \qquad
%    \sum_{i=0}^{k-1}\mathbb P(\mathcal X_i)+\sum_{j>2k}\mathbb Q(\Xi_j).
%  \]
%  Empty and full row sets give capacity one. The max-flow/min-cut theorem \citep{F.F56} therefore yields
%  \begin{align*}
%    \max_{\bm\pi\in\mathcal T(\mathbb Q)}\sum_{i,j}\overline h_{ij}\pi_{ij}
%      &=\min\left\{1,\min_k\left[\sum_{i=k}^{K}\mathbb P(\mathcal X_i)
%        +\sum_{j<2k}\mathbb Q(\Xi_j)\right]\right\},\\
%    \min_{\bm\pi\in\mathcal T(\mathbb Q)}\sum_{i,j}\underline h_{ij}\pi_{ij}
%      &=\max\left\{0,\max_k\left[\sum_{i=k}^{K}\mathbb P(\mathcal X_i)
%        +\sum_{j\le2k}\mathbb Q(\Xi_j)-1\right]\right\}.
%  \end{align*}
%
%  These values are nondecreasing in the respective open and closed cumulative residual probabilities. The rank constraints bound the former from above by $u_k$ and the latter from below by $l_k$. To attain all relevant limits simultaneously, choose laws $\mathbb Q^+$ and $\mathbb Q^-$ with cell probabilities
%  \begin{align*}
%    &\mathbb Q^+(\Xi_1):=u_1,\quad
%      \mathbb Q^+(\Xi_{2k}):=u_{k+1}-u_k,\quad
%      \mathbb Q^+(\Xi_{2k+1}):=0,\\
%    &\mathbb Q^-(\Xi_{2k-1}):=0,\quad
%      \mathbb Q^-(\Xi_{2k}):=l_k-l_{k-1},\quad
%      \mathbb Q^-(\Xi_J):=1-l_K,
%      \qquad k=1,\ldots,K.
%  \end{align*}
%  By Proposition~\ref{prop-finite-rank-constraints}, the prescribed masses are nonnegative and sum to one. Choosing a representative in each residual cell realizes both laws. Their cumulative probabilities satisfy
%  \begin{align*}
%    &\sum_{j<2k}\mathbb Q^+(\Xi_j)=u_k,
%      \qquad \sum_{j\le2k}\mathbb Q^+(\Xi_j)=u_{k+1}\ge l_k,\\
%    &\sum_{j\le2k}\mathbb Q^-(\Xi_j)=l_k,
%      \qquad \sum_{j<2k}\mathbb Q^-(\Xi_j)=l_{k-1}\le u_k.
%  \end{align*}
%  Both laws therefore belong to $\mathcal Q^{\mathrm{DW}}(\theta)$. The union representation of $\mathcal C_\theta$ and \eqref{eq-finite-primal-values} show that these allocations attain the upper and lower bounds in \eqref{eq-dw-closed-upper} and \eqref{eq-dw-closed-lower}. Zero masses and coincident cumulative limits are allowed, including at $\theta=0$.
%
%  Finally, let $\bm e_k$ be the $k$th unit vector. Certificates for the upper and lower LPs \eqref{eq-finite-upper-dual} and \eqref{eq-finite-lower-dual} are respectively
%  \begin{align*}
%    &\bm\lambda=\bm e_k,\qquad\bm\gamma=\bm0,\qquad
%      z_i=\mathbf1_{\{i\ge k\}},\\
%    &\bm\lambda=\bm0,\qquad\bm\gamma=\bm e_k,\qquad
%      z_i=\mathbf1_{\{i\ge k\}}-1.
%  \end{align*}
%  Their feasibility follows from
%  \[
%    \mathbf1_{\{i\ge k\}}+\mathbf1_{\{j<2k\}}\ge\overline h_{ij},
%    \qquad
%    \mathbf1_{\{i\ge k\}}-1+\mathbf1_{\{j\le2k\}}\le\underline h_{ij}.
%  \]
%  Their objectives are $u_k+\sum_{i=k}^{K}\mathbb P(\mathcal X_i)$ and $l_k+\sum_{i=k}^{K}\mathbb P(\mathcal X_i)-1$, the inner candidates in the closed-form expressions. Zero rank multipliers with $z_i=1$ or $z_i=0$ supply the outer candidates. Selecting the best candidate gives an optimal full dual solution, whose functional marginal multiplier follows from \eqref{eq-dual-envelope-identities}. The inclusion of $i=k$ in the tail sum retains inputs whose failure threshold equals the observed residual.
%\end{proof}
%
%For the response-threshold curves in Corollary~\ref{cor-simultaneous-threshold-bounds}, apply the same construction to $g-c$ and $\hat g-c$. Every residual, rank count, and endpoint $l_k,u_k$ is unchanged. The input partition now depends on $c$, and its cumulative probability at rank $k$ is $\mathbb P\left\{\hat g(\bm X)\le c-\xi_{(k)}^{\mathrm{te}}\right\}$. Thus Corollary~\ref{thm-closed-form-bounds} gives
%\begin{align}
%  \overline p_{\hat f}(c)
%    &=\min\left\{1,\min_{k=1,\ldots,K}\left[u_k
%      +\mathbb P\left\{\hat g(\bm X)\le c-\xi_{(k)}^{\mathrm{te}}\right\}\right]\right\},
%    \label{eq-threshold-closed-upper}\\
%  \underline p_{\hat f}(c)
%    &=\max\left\{0,\max_{k=1,\ldots,K}\left[l_k
%      +\mathbb P\left\{\hat g(\bm X)\le c-\xi_{(k)}^{\mathrm{te}}\right\}-1\right]\right\}.
%    \label{eq-threshold-closed-lower}
%\end{align}
%These identities hold for every $c\in\mathbb R$ and $\theta\ge0$, including $\theta=0$. When $\theta\ge\theta_{N^{\mathrm{te}}}(\delta)$, their coverage is simultaneous by \eqref{eq-threshold-curve-coverage}. An extremizing coupling may depend on $c$; the common feasible true joint law supplies the simultaneous guarantee.
%
%The upper bound combines residual mass below an observed value with the probability that the surrogate threshold reaches it. The lower bound measures the overlap required by every admissible coupling. For $p_{\mathrm{tar}}\in(0,1)$, the upper curve satisfies $\overline p_{\hat f}(c)\le p_{\mathrm{tar}}$ exactly when some rank satisfies
%\[
%  u_k+\mathbb P\left\{\hat g(\bm X)\le c-\xi_{(k)}^{\mathrm{te}}\right\}\le p_{\mathrm{tar}}.
%\]
%Both terms must therefore be small at the same residual level.
%
%With a known residual law, the coupling problem is the classical fixed-marginal sum problem studied by \citet{M82} and \citet{F.N.S87}. Here the residual law also varies within a calibrated DW band, and the open and closed rank constraints retain atoms.
%
%Dependence uncertainty can remain even with a known residual marginal. For equally likely surrogate responses $0$ and $2$ and equally likely residuals $-1$ and $1$, failure occurs exactly when the response is $0$ and the residual is $-1$. Its probability ranges from $0$ to $1/2$ over couplings, whereas independence gives $1/4$. Further marginal calibration alone does not resolve this dependence uncertainty.
%
%\begin{remark}[Known interval support]
%  \label{rem-interval-support}
%  Suppose $\Xi=[a,b]$ with $-\infty<a\le b<\infty$ and the observed residuals lie in this interval. Then
%  \begin{align*}
%    \overline p_{\hat f}
%      &=\min\left\{\mathbb P\left\{\hat g(\bm X)\le-a\right\},
%        \min_{k=1,\ldots,K}\left[u_k+\mathbb P\left\{\hat g(\bm X)\le-\xi_{(k)}^{\mathrm{te}}\right\}\right]\right\},\\
%    \underline p_{\hat f}
%      &=\max\left\{\mathbb P\left\{\hat g(\bm X)\le-b\right\},
%        \max_{k=1,\ldots,K}\left[l_k+\mathbb P\left\{\hat g(\bm X)\le-\xi_{(k)}^{\mathrm{te}}\right\}-1\right]\right\}.
%  \end{align*}
%  Represent the support restriction on $\mathbb R$ by $\mathbb Q\{\xi<a\}\le0$ and $\mathbb Q\{\xi\le b\}\ge1$, and add $a,b$ to the ordered threshold list. At a newly added $a$, use upper and lower bounds zero; at a newly added $b$, use upper and lower bounds one. At coincident observed thresholds, tighten only the upper bound at $a$ and the lower bound at $b$. The enlarged sequences satisfy the compatibility conditions in Proposition~\ref{prop-finite-rank-constraints}, so the same finite-aggregation and ordered-cut arguments apply. The effective additional upper candidate is $\mathbb P\left\{\hat g(\bm X)\le-a\right\}$, and the effective additional lower candidate is $\mathbb P\left\{\hat g(\bm X)\le-b\right\}$; the other added candidates are redundant. If $a=b$, both bounds equal $\mathbb P\left\{\hat g(\bm X)\le-a\right\}$.
%
%  For the partition on $\Xi=[a,b]$, intersect the residual cells from the augmented threshold list with $\Xi$ and omit empty cells. Including the support endpoints retains their atoms and the corresponding changes in failure status. For a response threshold $c$, shift each surrogate-response threshold in the displayed event probabilities by $c$.
%\end{remark}
%
%After zero failures in $N^{\mathrm{te}}$ independent tests, the one-sided exact binomial upper limit is $1-\delta^{1/N^{\mathrm{te}}}\approx\log(1/\delta)/N^{\mathrm{te}}$ \citep{C.P34}. At $c=0$, the simple DW case $K=1$ and $\mathbb P(\mathcal X_1)=0$ gives an upper bound $\beta_0(\theta)$ determined by the simultaneous-band calibration. Verified support can exclude otherwise unobserved failures: in Remark~\ref{rem-interval-support}, $\mathbb P\left\{\hat g(\bm X)\le-a\right\}=0$ rules out failure.
%
%Evaluation requires scalar DW inversion, sorting and counting the test residuals, and computing the $K$ input event probabilities $\mathbb P\left\{\hat g(\bm X)\le c-\xi_{(k)}^{\mathrm{te}}\right\}$ at each response threshold $c$. The endpoint table depends only on $N^{\mathrm{te}}$ and $\theta$ and can be precomputed. Given this table and the input event probabilities, the two scans cost $O(K)$ at each response threshold.
%
%For $\theta\ge\theta_{N^{\mathrm{te}}}(\delta)$, the exact expressions retain the coverage in Theorem~\ref{theorem-01} and Corollary~\ref{cor-simultaneous-threshold-bounds}. Independent auxiliary simulation of $\bm X\sim\mathbb P$ can supply simultaneous upper and lower confidence bounds on the input event probabilities. Substituting these bounds in the respective upper and lower expressions preserves conservative failure probability bounds. For one fixed $c$, simultaneous bounds at its $K$ evaluation thresholds suffice. To retain coverage of the entire response-threshold curve, the auxiliary bounds must hold for the surrogate-response CDF at every real threshold. They must hold conditionally on the test data, with their error probability included in the overall confidence budget.
%
%
% END PRESERVED BODY BEFORE RESULT-EXPOSITION STYLE REVISION

% BEGIN PRESERVED BODY BEFORE CONCISE RESULT TRANSITIONS (2026-11-26)
%\section{Problem Description}
%\label{sec-problem}
%
%Let $\bm X\sim\mathbb P$ be a random input taking values in $\mathcal X\subseteq\mathbb R^n$, and let $g:\mathcal X\to\mathbb R$ be a measurable limit-state function. The failure probability is defined as
%\[
%  p_f:=\Pr\{g(\bm X)\le0\}=\mathbb E_{\bm X\sim\mathbb P}\left[\mathbf 1_{\{g(\bm X)\le0\}}\right]=\int_{\mathcal X}\mathbf 1_{\{g(\bm x)\le0\}}\,\mathbb P(\mathrm d\bm x).
%\]
%
%Evaluating $g$ can require an expensive high-fidelity simulation. We use the training dataset
%\[
%  \mathcal D^{\mathrm{tr}}:=\left\{\left(\bm x_i^{\mathrm{tr}},g(\bm x_i^{\mathrm{tr}})\right)\right\}_{i=1}^{N^{\mathrm{tr}}}
%\]
%to fit a computationally efficient, measurable surrogate $\hat g:=\hat g_{\mathcal D^{\mathrm{tr}}}$. Model fitting and selection use only the training data and are completed before the test data are used. We condition on $\mathcal D^{\mathrm{tr}}$, treat $\hat g$ as fixed, and suppress this dependence in the notation.
%
%Define the signed residual function
%\[
%  \xi(\bm x):=g(\bm x)-\hat g(\bm x),\qquad \forall \bm x\in\mathcal X.
%\]
%For a new input $\bm X\sim\mathbb P$ independent of the training data, let $\mathbb Q^\star$ be the distribution of $\xi^\star:=\xi(\bm X)$. Let $\Xi\subseteq\mathbb R$ be a known Borel set with $\mathbb Q^\star(\Xi)=1$; the default choice is $\Xi=\mathbb R$. For any Borel set $\mathcal S$ in a Euclidean space, write $\mathcal B(\mathcal S)$ for its Borel sigma-algebra. Let $\mathcal M_+(\mathcal S)$ denote the finite nonnegative Borel measures concentrated on $\mathcal S$, and let $\mathcal M_1(\mathcal S)$ denote the probability measures.
%
%For a candidate residual law $\mathbb Q\in\mathcal M_1(\Xi)$, let $\Pi$ be a joint distribution on $\mathcal X\times\Xi$ with marginals $\Pi_{\bm X}=\mathbb P$ and $\Pi_\xi=\mathbb Q$. Its surrogate-based failure probability is
%\begin{equation}
%  p_{\hat f}(\Pi):=\Pr\{\hat g(\bm X)+\xi\le0\}=\int_{\mathcal X\times\Xi}\mathbf 1_{\{\hat g(\bm x)+\xi\le0\}}\,\Pi(\mathrm d\bm x,\mathrm d\xi).
%  \label{eq-surrogate-failure-probability}
%\end{equation}
%For fixed $\hat g$, this functional is linear in $\Pi$. In particular, the true joint distribution $\Pi^\star$ satisfies $p_{\hat f}(\Pi^\star)=p_f$.
%
%The residual marginal and its dependence on the input jointly determine the failure probability. We represent residual-marginal uncertainty by an ambiguity set $\mathcal Q$ calibrated to contain $\mathbb Q^\star$ with high confidence. For a given $\mathcal Q$, define the joint ambiguity set
%\[
%  \mathcal P(\mathcal Q) :=\left\{\Pi\in\mathcal M_1(\mathcal X\times\Xi):\Pi_{\bm X}=\mathbb P,\ \Pi_\xi\in\mathcal Q\right\}.
%\]
%The true residual is determined by the input, but its unknown functional relationship is not encoded in $\mathcal Q$. The set $\mathcal P(\mathcal Q)$ therefore allows all couplings with the specified marginals, including randomized conditional residual laws. The upper and lower failure probability bounds are
%\begin{equation}
%  \overline p_{\hat f}:=\sup_{\Pi\in\mathcal P(\mathcal Q)}p_{\hat f}(\Pi),
%  \label{eq-upper-bound}
%\end{equation}
%and
%\begin{equation}
%  \underline p_{\hat f}:=\inf_{\Pi\in\mathcal P(\mathcal Q)}p_{\hat f}(\Pi),
%  \label{eq-lower-bound}
%\end{equation}
%respectively. If $\mathbb Q^\star\in\mathcal Q$, then $\Pi^\star\in\mathcal P(\mathcal Q)$ and hence $\underline p_{\hat f}\le p_f\le\overline p_{\hat f}$.
%
%
%Section~\ref{sec-ambiguity-set} constructs the residual ambiguity set and establishes finite-sample coverage of the failure probability bounds, simultaneously across response thresholds. Section~\ref{sec-reformulation} then derives their Lagrangian duals, proves strong duality through a finite reformulation, and obtains closed-form expressions.
%
%
%\section{Ambiguity Set Design}
%\label{sec-ambiguity-set}
%
%Failure occurs when $\xi\le-\hat g(\bm X)$, so the residual cumulative distribution function (CDF) gives the marginal mass at or below each failure threshold. We use the D\"umbgen--Wellner (DW) construction \citep{L.J23} to control empirical CDF deviations simultaneously over all thresholds, with tolerances that vary across probability levels.
%
%Let $q$ and $p$ denote an empirical probability and a candidate probability. Following \citet{B.J79} and the CDF confidence bands of \citet{O95}, we measure their local discrepancy by the Bernoulli Kullback--Leibler divergence
%\begin{equation}
%  \phi(q,p):=q\log\frac{q}{p}+(1-q)\log\frac{1-q}{1-p},\qquad \forall q\in[0,1],\, p\in(0,1),
%  \label{eq-bernoulli-kl}
%\end{equation}
%where logarithms are natural and a term with zero numerator is interpreted as zero. For $q$ close to an interior $p$, the approximation $\phi(q,p)\approx(q-p)^2/[2p(1-p)]$ scales the deviation by the Bernoulli variance. Thus, the same absolute deviation receives a larger weight near the probability boundaries. The full divergence also covers zero and full empirical counts through $\phi(0,p)=-\log(1-p)$ and $\phi(1,p)=-\log p$. At the boundaries of the candidate probability, we set $\phi(0,0)=\phi(1,1)=0$, $\phi(q,0)=+\infty$ for $q>0$, and $\phi(q,1)=+\infty$ for $q<1$.
%
%To compare deviations simultaneously across probability levels, we use the additive correction of \citet{L.J23}. For $v\in(0,1)$, define
%\begin{equation*}
%  \rho(v):=\log\log\frac{e}{4v(1-v)},\qquad \omega(v):=\log\!\left(1+\rho(v)^2\right).
%\end{equation*}
%Before sampling, we fix the divergence parameter $s=1$, which gives \eqref{eq-bernoulli-kl}, and the correction parameter $\nu=1$. Set $\psi_0(v):=\rho(v)+\omega(v)$, with $\psi_0(0)=\psi_0(1)=+\infty$, and define the bivariate correction
%\begin{equation*}
%  \psi(q,p):=\min_{v\in[\min\{q,p\},\,\max\{q,p\}]}\psi_0(v).
%\end{equation*}
%The functions $\rho$ and $\omega$ vanish at $v=1/2$ and increase toward the probability boundaries. The minimum over the interval between $q$ and $p$ is therefore large only when both probabilities lie in the same tail. This correction balances fluctuations across probability scales in the simultaneous CDF constraint.
%
%We construct the empirical CDF from a test dataset
%\[
%  \mathcal D^{\mathrm{te}}:=\left\{\left(\bm x_j^{\mathrm{te}},g(\bm x_j^{\mathrm{te}})\right)\right\}_{j=1}^{N^{\mathrm{te}}}.
%\]
%The test inputs are independent draws from $\mathbb P$ and are independent of $\mathcal D^{\mathrm{tr}}$. We fix $N^{\mathrm{te}}\ge1$ before sampling and use these data to calibrate the residual ambiguity set.
%
%Define the signed test residuals
%\[
%  \xi_j^{\mathrm{te}}:=\xi(\bm x_j^{\mathrm{te}})=g(\bm x_j^{\mathrm{te}})-\hat g(\bm x_j^{\mathrm{te}}),\qquad j=1,\ldots,N^{\mathrm{te}}.
%\]
%Conditional on $\mathcal D^{\mathrm{tr}}$, these residuals are independent and identically distributed according to $\mathbb Q^\star$. Their empirical CDF is
%\begin{equation}
%  \widehat F^{\mathrm{te}}(t):=\frac{1}{N^{\mathrm{te}}}\sum_{j=1}^{N^{\mathrm{te}}}\mathbf 1_{\{\xi_j^{\mathrm{te}}\le t\}},\qquad \forall t\in\mathbb R.
%  \label{eq-empirical-cdf}
%\end{equation}
%For $\mathbb Q\in\mathcal M_1(\Xi)$, write $F_{\mathbb Q}(t):=\mathbb Q\left(({-\infty},t]\right)$. At a threshold $t$, $q=\widehat F^{\mathrm{te}}(t)$ is the mean of the Bernoulli indicators in \eqref{eq-empirical-cdf}, and $p=F_{\mathbb Q}(t)$ is the candidate probability. Thus $N^{\mathrm{te}}\phi(q,p)$ is the Bernoulli log-likelihood difference between the empirical and candidate probabilities. Combining it with the multiscale correction gives $N^{\mathrm{te}}\phi(q,p)-\psi(q,p)$.
%
%At the probability boundaries, the corrected discrepancy is $-\infty$ for $q=p\in\{0,1\}$ and $+\infty$ for $p=0<q$ or $p=1>q$. These conventions extend the constraint to every residual threshold.
%
%\Needspace{5\baselineskip}
%The following definition uses the corrected discrepancy to constrain the residual CDF simultaneously over all thresholds.
%
%\begin{definition}[DW-based CDF ambiguity set]
%  \label{definition-dw-ambiguity-set}
%  Given the empirical CDF $\widehat F^{\mathrm{te}}$ and a threshold $\theta\ge0$, define
%  \begin{equation} 
%    \label{eq-cdf-ambiguity-set}
%      \mathcal Q^{\mathrm{DW}}(\theta):=\left\{\mathbb Q\in\mathcal M_1(\Xi):\ 
%      \sup_{t\in\mathbb R}\left[ N^{\mathrm{te}}\phi\!\left(\widehat F^{\mathrm{te}}(t),F_{\mathbb Q}(t)\right) 
%       {}-\psi\!\left(\widehat F^{\mathrm{te}}(t),F_{\mathbb Q}(t)\right)\right]\le\theta\right\}. 
%  \end{equation} 
%\end{definition}
%
%In Definition~\ref{definition-dw-ambiguity-set}, $\theta$ controls the allowed corrected likelihood discrepancy. The induced CDF tolerances vary across probability levels and can be asymmetric about the empirical CDF. Calibrating $\theta$ assigns a confidence level to this simultaneous band.
%
%We calibrate $\theta$ against a distribution-free reference law. Following \citet{L.J23}, let $V_1,\ldots,V_{N^{\mathrm{te}}}$ be independent standard uniform variables, independent of the training and test data. Their empirical CDF is
%\begin{equation*}
%  G_{N^{\mathrm{te}}}(u):=\frac{1}{N^{\mathrm{te}}}\sum_{j=1}^{N^{\mathrm{te}}}\mathbf 1_{\{V_j\le u\}},\qquad \forall u\in[0,1].
%\end{equation*}
%Define the uniform reference statistic
%\begin{equation}
%  T_{N^{\mathrm{te}}}^{\mathrm U}:=\sup_{u\in(0,1)}\left[N^{\mathrm{te}}\phi\!\left(G_{N^{\mathrm{te}}}(u),u\right)-\psi\!\left(G_{N^{\mathrm{te}}}(u),u\right)\right].
%  \label{eq-dw-reference-statistic}
%\end{equation}
%
%
%\Needspace{5\baselineskip}
%The following proposition calibrates $\theta$ using the uniform reference statistic and establishes conditional coverage of the true residual law.
%
%\begin{proposition}[Finite-sample threshold calibration]
%  \label{proposition-01}
%  Under the sampling conditions stated above, let $\delta\in(0,1)$ and define
%  \begin{equation}
%    \theta_{N^{\mathrm{te}}}(\delta):=\inf\left\{\tau\ge0:\Pr\left\{T_{N^{\mathrm{te}}}^{\mathrm U}\le\tau\right\}\ge1-\delta\right\}.
%    \label{eq-dw-calibrated-threshold}
%  \end{equation}
%
%  Then, for every $\theta\ge\theta_{N^{\mathrm{te}}}(\delta)$, the DW-based CDF ambiguity set satisfies
%  \begin{equation}
%    \Pr\left\{\mathbb Q^\star\in\mathcal Q^{\mathrm{DW}}(\theta)\,\middle|\,\mathcal D^{\mathrm{tr}}\right\}\ge1-\delta.
%    \label{eq-cdf-coverage}
%  \end{equation}
%\end{proposition}
%\begin{proof}
%  The reference statistic is finite almost surely, and its distribution function is right-continuous. Hence \eqref{eq-dw-calibrated-threshold} gives
%  \[
%    \Pr\left\{T_{N^{\mathrm{te}}}^{\mathrm U}>\theta_{N^{\mathrm{te}}}(\delta)\right\}\le\delta.
%  \]
%  Conditioning on $\mathcal D^{\mathrm{tr}}$ fixes $\hat g$ and $\mathbb Q^\star$. Define the generalized quantile function $F_{\mathbb Q^\star}^{-1}(u):=\inf\{t\in\mathbb R:F_{\mathbb Q^\star}(t)\ge u\}$ for all $u\in(0,1)$. The sample $F_{\mathbb Q^\star}^{-1}(V_j)$ for $j=1,\ldots,N^{\mathrm{te}}$, has the same conditional joint law as the test residuals. Thus, given $\mathcal D^{\mathrm{tr}}$,
%  \[
%    \left(\widehat F^{\mathrm{te}}(t)\right)_{t\in\mathbb R}\overset{\mathrm d}{=}\left(G_{N^{\mathrm{te}}}\!\left(F_{\mathbb Q^\star}(t)\right)\right)_{t\in\mathbb R}.
%  \]
%  For $0<F_{\mathbb Q^\star}(t)<1$, the discrepancy at $\mathbb Q^\star$ is evaluated at a probability level included in \eqref{eq-dw-reference-statistic}. When $F_{\mathbb Q^\star}(t)$ is zero or one, the empirical CDF equals that boundary almost surely, giving discrepancy $-\infty$. The supremum in \eqref{eq-cdf-ambiguity-set}, evaluated at $\mathbb Q^\star$, is therefore bounded by $T_{N^{\mathrm{te}}}^{\mathrm U}$ under this coupling. Since $\mathbb Q^\star(\Xi)=1$, we obtain
%  \[
%    \Pr\left\{\mathbb Q^\star\notin\mathcal Q^{\mathrm{DW}}\!\left(\theta_{N^{\mathrm{te}}}(\delta)\right)\,\middle|\,\mathcal D^{\mathrm{tr}}\right\}
%    \le\Pr\left\{T_{N^{\mathrm{te}}}^{\mathrm U}>\theta_{N^{\mathrm{te}}}(\delta)\right\}\le\delta.
%  \]
%  For every test sample, $\theta\ge\theta_{N^{\mathrm{te}}}(\delta)$ implies $\mathcal Q^{\mathrm{DW}}\!\left(\theta_{N^{\mathrm{te}}}(\delta)\right)\subseteq\mathcal Q^{\mathrm{DW}}(\theta)$. This inclusion proves \eqref{eq-cdf-coverage} for every such $\theta$.
%\end{proof}
%
%Proposition~\ref{proposition-01} applies to arbitrary residual distributions, including laws with atoms or infinite moments. For fixed $\phi$ and $\psi$, the reference law depends only on $N^{\mathrm{te}}$, so $\theta_{N^{\mathrm{te}}}(\delta)$ can be calibrated offline. \citet{L.J23} describe finite-sample evaluation using uniform order-statistic boundary probabilities. A conservative upper bound on the quantile in \eqref{eq-dw-calibrated-threshold} preserves \eqref{eq-cdf-coverage}.
%% If calibration is performed by simulation, quantile estimation error must be accounted for in the confidence guarantee. A threshold chosen from the test data also preserves \eqref{eq-cdf-coverage}, provided $\theta\ge\theta_{N^{\mathrm{te}}}(\delta)$ almost surely, by the same nesting argument.
%
%\Needspace{5\baselineskip}
%The following theorem transfers the residual-law coverage in Proposition~\ref{proposition-01} to the upper and lower failure probability bounds.
%
%\begin{theorem}[Finite-sample failure probability bounds]
%  \label{theorem-01}
%  Under the sampling conditions of Proposition~\ref{proposition-01}, let $\delta\in(0,1)$ and $\theta\ge\theta_{N^{\mathrm{te}}}(\delta)$, with the calibrated threshold defined in \eqref{eq-dw-calibrated-threshold}. Then the bounds defined in \eqref{eq-upper-bound} and \eqref{eq-lower-bound}, evaluated with $\mathcal Q=\mathcal Q^{\mathrm{DW}}(\theta)$, satisfy
%  \begin{equation}
%    \Pr\left\{\underline p_{\hat f}\le p_f\le\overline p_{\hat f}\,\middle|\,\mathcal D^{\mathrm{tr}}\right\}\ge1-\delta.
%    \label{eq-failure-probability-coverage}
%  \end{equation}
%\end{theorem}
%
%\begin{proof}
%  On the event $\mathbb Q^\star\in\mathcal Q^{\mathrm{DW}}(\theta)$, the true joint distribution $\Pi^\star$ belongs to $\mathcal P\left(\mathcal Q^{\mathrm{DW}}(\theta)\right)$. Hence
%  \[
%    \underline p_{\hat f}\le p_{\hat f}(\Pi^\star)=p_f\le\overline p_{\hat f}.
%  \]
%  By Proposition~\ref{proposition-01}, this event has conditional probability at least $1-\delta$ given $\mathcal D^{\mathrm{tr}}$, proving \eqref{eq-failure-probability-coverage}.
%\end{proof}
%
%
%Theorem~\ref{theorem-01} uses the same residual-coverage event for both endpoints, so $1-\delta$ is the confidence level for the entire interval. For fixed test data, $0\le\theta_1\le\theta_2$ implies $\mathcal Q^{\mathrm{DW}}(\theta_1)\subseteq\mathcal Q^{\mathrm{DW}}(\theta_2)$, making the lower bound nonincreasing and the upper bound nondecreasing in $\theta$. A higher prescribed confidence level requires a nondecreasing calibrated threshold and therefore gives a wider or unchanged interval.
%
%\Needspace{5\baselineskip}
%The following corollary extends this common coverage event to all response thresholds using the same calibrated ambiguity set.
%
%\begin{corollary}[Simultaneous bounds across response thresholds]
%  \label{cor-simultaneous-threshold-bounds}
%  Under the conditions of Theorem~\ref{theorem-01}, let $\mathcal Q=\mathcal Q^{\mathrm{DW}}(\theta)$. For $c\in\mathbb R$, write $p_f(c):=\Pr \left\{g(\bm X)\le c\right\}$ and define
%  \[
%    p_{\hat f}(\Pi;c):=\int_{\mathcal X\times\Xi}
%      \mathbf1_{\{\hat g(\bm x)+\xi\le c\}}\,\Pi(\mathrm d\bm x,\mathrm d\xi),
%  \]
%  with lower and upper bounds
%  \[
%    \underline p_{\hat f}(c):=\inf_{\Pi\in\mathcal P(\mathcal Q)}p_{\hat f}(\Pi;c),
%    \qquad
%    \overline p_{\hat f}(c):=\sup_{\Pi\in\mathcal P(\mathcal Q)}p_{\hat f}(\Pi;c).
%  \]
%  Then
%  \begin{equation}
%    \Pr\left\{\underline p_{\hat f}(c)\le p_f(c)\le\overline p_{\hat f}(c),
%      \ \forall c\in\mathbb R\,\middle|\,\mathcal D^{\mathrm{tr}}\right\}\ge1-\delta.
%    \label{eq-threshold-curve-coverage}
%  \end{equation}
%  At $c=0$, these definitions recover the original failure probability and its bounds.
%\end{corollary}
%\begin{proof}
%  On $\mathbb Q^\star\in\mathcal Q^{\mathrm{DW}}(\theta)$, the true joint law $\Pi^\star$ belongs to the same set $\mathcal P(\mathcal Q)$ for every $c$. Consequently,
%  \[ 
%    \underline p_{\hat f}(c) =\inf_{\Pi\in\mathcal P(\mathcal Q)} p_{\hat f}(\Pi;c) \le p_{\hat f}(\Pi^\star;c) =p_f(c) \le\sup_{\Pi\in\mathcal P(\mathcal Q)} p_{\hat f}(\Pi;c) =\overline p_{\hat f}(c), \qquad\forall c\in\mathbb R.
%  \]
%  Both bound curves are nondecreasing. Monotonicity and the CDF limits of $p_f$ imply that any violation is detected at a rational threshold or an atom of $g(\bm X)$, a countable set. The simultaneous event is therefore measurable. Proposition~\ref{proposition-01} gives the common residual-coverage event conditional probability at least $1-\delta$, proving \eqref{eq-threshold-curve-coverage}.
%\end{proof}
%
%Corollary~\ref{cor-simultaneous-threshold-bounds} gives the entire response-threshold curve from one residual calibration, including thresholds selected after observing the test data, with a common confidence level across $c$. Complementing the endpoints also bounds the strict exceedance probabilities $\Pr \left\{g(\bm X)>c\right\}$ simultaneously. Here $g$, $\hat g$, and the input law $\mathbb P$ remain fixed. Fragility curves formed by conditioning on different hazard intensities require coverage for the corresponding conditional residual laws.
%
%The coverage argument also applies to any residual ambiguity set containing $\mathbb Q^\star$ with conditional probability at least $1-\delta$. It remains to evaluate the DW extrema in \eqref{eq-upper-bound} and \eqref{eq-lower-bound} and their response-threshold extensions.
%
%
%\section{Tractable Reformulation}
%\label{sec-reformulation}
%
%Throughout this section, the training data, the test residuals, and $\theta\ge0$ are fixed. We use the observed residuals to partition both the residual and input spaces. These partitions express the DW restrictions as cumulative cell probabilities and reduce the Lagrangian duals to finite linear programs. A joint-distribution construction establishes exact finite primal representations, after which LP duality and the ordering of the cells yield strong duality and closed-form bounds. The main development takes $\Xi=\mathbb R$; Remark~\ref{rem-interval-support} treats known interval support.
%
%For each empirical probability level $q_j:=j/N^{\mathrm{te}}$, define the scalar acceptance interval
%\begin{equation}
%  [\alpha_j(\theta),\beta_j(\theta)]
%  :=\left\{p\in[0,1]:
%  N^{\mathrm{te}}\phi(q_j,p)-\psi(q_j,p)\le\theta\right\},
%  \qquad j=0,\ldots,N^{\mathrm{te}}.
%  \label{eq-dw-rank-acceptance}
%\end{equation}
%Let $\xi_{(1)}^{\mathrm{te}}<\cdots<\xi_{(K)}^{\mathrm{te}}$ be the distinct ordered test residuals, and set
%\begin{equation}
%  c_0:=0,\quad
%  c_k:=\sum_{j=1}^{N^{\mathrm{te}}}\mathbf1_{\{\xi_j^{\mathrm{te}}\le\xi_{(k)}^{\mathrm{te}}\}},\quad
%  l_k:=\alpha_{c_k}(\theta),\quad
%  u_k:=\beta_{c_{k-1}}(\theta),\qquad k=1,\ldots,K.
%  \label{eq-dw-counts-and-levels}
%\end{equation}
%The counts retain all multiplicities, so $c_K=N^{\mathrm{te}}$ even when test residuals coincide. Set $\xi_{(0)}^{\mathrm{te}}:=-\infty$ and $\xi_{(K+1)}^{\mathrm{te}}:=+\infty$. Separating the observed values from the intervening intervals gives $J:=2K+1$ residual cells:
%\begin{equation}
%  \Xi_{2k}:=\left\{\xi_{(k)}^{\mathrm{te}}\right\}
%    \quad\left(k=1,\ldots,K\right),
%  \qquad
%  \Xi_{2k+1}:=\left(\xi_{(k)}^{\mathrm{te}},\xi_{(k+1)}^{\mathrm{te}}\right)
%    \quad\left(k=0,\ldots,K\right).
%  \label{eq-residual-cells}
%\end{equation}
%For any candidate residual law $\mathbb Q\in\mathcal M_1(\Xi)$, the open rank probability measures mass strictly below an observed residual, while the closed rank probability also includes mass at that value. For $k=1,\ldots,K$, the partition gives
%\[
%  \mathbb Q\left\{\xi<\xi_{(k)}^{\mathrm{te}}\right\}
%    =\sum_{j<2k}\mathbb Q(\Xi_j),
%  \qquad
%  \mathbb Q\left\{\xi\le\xi_{(k)}^{\mathrm{te}}\right\}
%    =\sum_{j\le2k}\mathbb Q(\Xi_j).
%\]
%The two probabilities differ by the atom mass $\mathbb Q(\Xi_{2k})$.
%
%For $i=0,\ldots,K$, the observed residual levels classify inputs by the location of the failure threshold $-\hat g(\bm x)$:
%\begin{equation}
%  \mathcal X_i:=\left\{\bm x\in\mathcal X:
%    \xi_{(i)}^{\mathrm{te}}\le-\hat g(\bm x)<\xi_{(i+1)}^{\mathrm{te}}\right\}.
%  \label{eq-input-cells}
%\end{equation}
%The input cell probabilities give the surrogate-response threshold probabilities through
%\[
%  \mathbb P\left\{\hat g(\bm X)\le-\xi_{(k)}^{\mathrm{te}}\right\}
%    =\sum_{i=k}^{K}\mathbb P(\mathcal X_i),
%  \qquad k=1,\ldots,K.
%\]
%The half-open convention assigns $-\hat g(\bm x)=\xi_{(k)}^{\mathrm{te}}$ to $\mathcal X_k$. The tail sum therefore includes threshold equalities even when the surrogate response has atoms. 
%The two partitions satisfy
%\[
%  \mathcal X=\bigsqcup_{i=0}^{K}\mathcal X_i,
%  \qquad
%  \Xi=\bigsqcup_{j=1}^{J}\Xi_j.
%\]
%Figure~\ref{fig-residual-input-cells} illustrates the two partitions for $K=3$.
%\begin{figure}[!htbp]
%  \centering
%  \includegraphics[width=0.5\columnwidth]{fig_residual_input_cells.png}
%  \caption{Illustration of the residual and input-threshold partitions for $K=3$. The residual partition consists of $K$ singleton cells and $K+1$ open intervals, whereas the input partition is induced by half-open intervals of the failure threshold $-\hat g(\bm x)$. Here, $w_i=\mathbb P(\mathcal X_i)$ denotes the probability mass of input cell $\mathcal X_i$.}
%  \label{fig-residual-input-cells}
%\end{figure}
%
%\Needspace{5\baselineskip}
%The following proposition expresses the DW ambiguity set through linear constraints on cumulative residual cell probabilities.
%
%\begin{proposition}[Finite representation of the DW ambiguity set]
%  \label{prop-finite-rank-constraints}
%  Each acceptance set in \eqref{eq-dw-rank-acceptance} is a nonempty closed interval containing $q_j$. The DW ambiguity set has the finite representation
%  \begin{equation}
%    \mathcal Q^{\mathrm{DW}}(\theta)
%    =\left\{\mathbb Q\in\mathcal M_1(\Xi):
%      \sum_{j<2k}\mathbb Q(\Xi_j)\le u_k,\quad
%      \sum_{j\le2k}\mathbb Q(\Xi_j)\ge l_k,\quad k=1,\ldots,K\right\}.
%    \label{eq-finite-rank-constraints}
%  \end{equation}
%  The sequences $l_k$ and $u_k$ are nondecreasing and satisfy
%  \[
%    0\le l_k,u_k\le1,\qquad l_{k-1}\le u_k,\qquad l_k\le u_{k+1},
%    \qquad k=1,\ldots,K,
%  \]
%  where $l_0:=0$ and $u_{K+1}:=1$.
%\end{proposition}
%\begin{proof}
%  For fixed $q$, the Bernoulli divergence increases as $p$ moves away from $q$, while $\psi(q,p)$ cannot increase as the interval between $q$ and $p$ expands. The corrected discrepancy therefore decreases toward $q$ from below and increases above $q$. It is continuous for $p\in(0,1)$ and equals $-\psi_0(q)\le0$ at $p=q$. The boundary conventions in Section~\ref{sec-ambiguity-set} give a nonempty closed acceptance interval also for $q\in\{0,1\}$. Scalar inversion of the DW constraint therefore gives the CDF band \citep[Supplement, Section~S.6]{L.J23}
%  \[
%    \alpha_{N^{\mathrm{te}}\widehat F^{\mathrm{te}}(t)}(\theta)
%      \le F_{\mathbb Q}(t)
%      \le\beta_{N^{\mathrm{te}}\widehat F^{\mathrm{te}}(t)}(\theta),
%    \qquad\forall t\in\mathbb R.
%  \]
%
%  For $k=1,\ldots,K-1$, the empirical CDF on $\left[\xi_{(k)}^{\mathrm{te}},\xi_{(k+1)}^{\mathrm{te}}\right)$ equals $c_k/N^{\mathrm{te}}$. Monotonicity makes the lower band constraint on this interval equivalent to $\mathbb Q\left\{\xi\le\xi_{(k)}^{\mathrm{te}}\right\}\ge\alpha_{c_k}(\theta)$. By continuity of probability from below, the upper band constraint is equivalent to $\mathbb Q\left\{\xi<\xi_{(k+1)}^{\mathrm{te}}\right\}\le\beta_{c_k}(\theta)$. The exterior intervals give $\mathbb Q\left\{\xi<\xi_{(1)}^{\mathrm{te}}\right\}\le\beta_0(\theta)$ and $\mathbb Q\left\{\xi\le\xi_{(K)}^{\mathrm{te}}\right\}\ge\alpha_{N^{\mathrm{te}}}(\theta)$; their other band constraints are automatic because $\alpha_0(\theta)=0$ and $\beta_{N^{\mathrm{te}}}(\theta)=1$. The partition identities above then yield \eqref{eq-finite-rank-constraints}.
%
%  With $p$ fixed, the corrected discrepancy decreases as $q$ approaches $p$ from either side. For empirical levels $q<q'$, an accepted value $p<q$ at level $q'$ remains accepted at level $q$; an accepted value $p>q'$ at level $q$ remains accepted at level $q'$. Since each acceptance interval contains its empirical level, both endpoints are nondecreasing. The inequalities $\alpha_{c_k}(\theta)\le c_k/N^{\mathrm{te}}\le\beta_{c_k}(\theta)$ give the stated compatibility conditions. Singleton intervals are allowed: if $q_j=1/2$ and $\theta=0$, the interval is $\{1/2\}$ because $\psi(1/2,p)=0$.
%\end{proof}
%
%Proposition~\ref{prop-finite-rank-constraints} replaces the CDF constraint over all thresholds by $2K$ linear inequalities on residual cell probabilities. The distribution within each cell remains free, and the scalar interval endpoints can be computed by one-dimensional bisection.
%
%For a feasible joint law, define its matrix of joint cell probabilities by $\pi_{ij}:=\Pi(\mathcal X_i\times\Xi_j)$. The fixed input marginal and the residual restrictions imply $\bm\pi=(\pi_{ij})\in\mathcal C_\theta$, where
%\begin{equation}
%  \mathcal C_\theta:=\left\{\bm\pi\in\mathbb R_+^{(K+1)\times J}:
%  \begin{array}{ll}
%    \displaystyle\sum_{j=1}^{J}\pi_{ij}=\mathbb P(\mathcal X_i),&i=0,\ldots,K,\\[1mm]
%    \displaystyle\sum_{i=0}^{K}\sum_{j<2k}\pi_{ij}\le u_k,&k=1,\ldots,K,\\[1mm]
%    \displaystyle\sum_{i=0}^{K}\sum_{j\le2k}\pi_{ij}\ge l_k,&k=1,\ldots,K
%  \end{array}\right\}.
%  \label{eq-cell-mass-polytope}
%\end{equation}
%The row equalities in $\mathcal C_\theta$ retain the input cell probabilities, and the cumulative column inequalities impose the DW band. On each product cell, the failure indicator satisfies
%\begin{equation}
%  \underline h_{ij}:=\mathbf1_{\{j\le2i\}}
%  \le\mathbf1_{\{\xi\le-\hat g(\bm x)\}}
%  \le\overline h_{ij}:=\mathbf1_{\{j\le2i+1\}},
%  \qquad \forall (\bm x,\xi)\in\mathcal X_i\times\Xi_j.
%  \label{eq-cell-costs}
%\end{equation}
%Residual cells with $j\le2i$ always cause failure for inputs in $\mathcal X_i$, and those with $j\ge2i+2$ are always safe. Only the overlapping interval $\Xi_{2i+1}$ has different upper and lower coefficients. Figure~\ref{fig-joint-partition} illustrates the resulting joint cell structure and the classification of product cells for $K=3$. Thus the cell matrix records the rank constraints exactly, while the failure objective requires the constructions below.
%\begin{figure}[!htbp]
%  \centering
%  \includegraphics[width=0.5\columnwidth]{fig_joint_partition.png}
%  \caption{Joint partition of the input and residual spaces for $K=3$. Each product cell has probability mass $\pi_{ij}=\Pi(\mathcal X_i\times\Xi_j)$. The blue, yellow, and red cells correspond to $(\underline h_{ij},\overline h_{ij})=(0,0)$, $(0,1)$, and $(1,1)$, respectively.}
%  \label{fig-joint-partition}
%\end{figure}
%
%The upper problem remains a linear program over joint measures, with the full input marginal fixed:
%\begin{equation}
%  \begin{aligned}
%    \overline p_{\hat f}
%    =\sup_{\Pi\in\mathcal M_+(\mathcal X\times\Xi)}\quad&
%      \int_{\mathcal X\times\Xi}\mathbf1_{\{\xi\le-\hat g(\bm x)\}}\,\Pi(\mathrm d\bm x,\mathrm d\xi)\\
%    \text{s.t.}\quad&
%      \int_{\mathcal X\times\Xi}\mathbf1_{\mathcal U}(\bm x)\,\Pi(\mathrm d\bm x,\mathrm d\xi)
%        =\mathbb P(\bm X\in\mathcal U),&&\forall\mathcal U\in\mathcal B(\mathcal X),\\
%    &\int_{\mathcal X\times\Xi}\mathbf1_{\{\xi<\xi_{(k)}^{\mathrm{te}}\}}\,\Pi(\mathrm d\bm x,\mathrm d\xi)\le u_k,
%        &&k=1,\ldots,K,\\
%    &\int_{\mathcal X\times\Xi}\mathbf1_{\{\xi\le\xi_{(k)}^{\mathrm{te}}\}}\,\Pi(\mathrm d\bm x,\mathrm d\xi)\ge l_k,
%        &&k=1,\ldots,K.
%  \end{aligned}
%  \label{eq-measure-upper}
%\end{equation}
%The lower problem minimizes the same integral under the same constraints. The marginal equalities fix both the input law and total mass one. The product of $\mathbb P$ and the empirical residual law is feasible for every $\theta\ge0$.
%
%\Needspace{5\baselineskip}
%The following lemma reduces the Lagrangian duals to $2K$ nonnegative rank multipliers and establishes weak duality.
%
%\begin{lemma}[Finite-dimensional Lagrangian duals and weak duality]
%  \label{lem-weak-duality}
%  Write $\bm u=(u_1,\ldots,u_K)^\top$ and $\bm l=(l_1,\ldots,l_K)^\top$. The nonnegative multipliers $\bm\lambda,\bm\gamma\in\mathbb R_+^K$ correspond to the upper and lower rank constraints. The input marginal equality has an unrestricted multiplier $\eta\in\mathcal F_{\mathrm b}(\mathcal X)$, the space of bounded Borel measurable real-valued functions. Define the residual penalty
%  \begin{equation}
%    \zeta_{\bm\lambda,\bm\gamma}(\xi)
%    :=\sum_{k=1}^{K}\lambda_k\mathbf1_{\{\xi<\xi_{(k)}^{\mathrm{te}}\}}
%      -\sum_{k=1}^{K}\gamma_k\mathbf1_{\{\xi\le\xi_{(k)}^{\mathrm{te}}\}},
%    \label{eq-dual-step-function}
%  \end{equation}
%  whose value on $\Xi_j$ is
%  \begin{equation}
%    \zeta_j:=\sum_{k=1}^{K}\lambda_k\mathbf1_{\{j<2k\}}
%      -\sum_{k=1}^{K}\gamma_k\mathbf1_{\{j\le2k\}},
%    \qquad j=1,\ldots,J.
%    \label{eq-dual-cell-values}
%  \end{equation}
%  Eliminating the input marginal multiplier gives the upper and lower Lagrangian dual problems
%  \begin{align}
%    \overline d_\theta
%      &:=\inf_{\bm\lambda,\bm\gamma\ge\bm0}\left\{
%        \bm\lambda^\top\bm u-\bm\gamma^\top\bm l
%        +\sum_{i=0}^{K}\mathbb P(\mathcal X_i)\max_{1\le j\le J}
%          \left\{\overline h_{ij}-\zeta_j\right\}\right\},
%      \label{eq-upper-measure-dual}\\
%    \underline d_\theta
%      &:=\sup_{\bm\lambda,\bm\gamma\ge\bm0}\left\{
%        -\bm\lambda^\top\bm u+\bm\gamma^\top\bm l
%        +\sum_{i=0}^{K}\mathbb P(\mathcal X_i)\min_{1\le j\le J}
%          \left\{\underline h_{ij}+\zeta_j\right\}\right\}.
%      \label{eq-lower-measure-dual}
%  \end{align}
%  Their optimal values satisfy
%  \begin{equation}
%    \underline d_\theta\le\underline p_{\hat f}
%      \le\overline p_{\hat f}\le\overline d_\theta.
%    \label{eq-weak-duality}
%  \end{equation}
%\end{lemma}
%\begin{proof}
%  The input marginal equality is equivalent to
%  \[
%    \int_{\mathcal X\times\Xi}\eta(\bm x)\,\Pi(\mathrm d\bm x,\mathrm d\xi)
%      =\mathbb E_{\mathbb P}\left[\eta(\bm X)\right],
%    \qquad\forall\eta\in\mathcal F_{\mathrm b}(\mathcal X).
%  \]
%  Relaxing this equality and the rank inequalities gives the upper Lagrangian \citep[Section~5.1]{B.V04}
%  \[
%    \begin{aligned}
%      \mathcal L^+(\Pi,\eta,\bm\lambda,\bm\gamma)
%      ={}&\int_{\mathcal X\times\Xi}\mathbf1_{\{\xi\le-\hat g(\bm x)\}}\,\Pi(\mathrm d\bm x,\mathrm d\xi)
%        +\mathbb E_{\mathbb P}\left[\eta(\bm X)\right]
%        -\int_{\mathcal X\times\Xi}\eta(\bm x)\,\Pi(\mathrm d\bm x,\mathrm d\xi)\\
%      &+\sum_{k=1}^{K}\lambda_k\left(u_k-\int_{\mathcal X\times\Xi}
%          \mathbf1_{\{\xi<\xi_{(k)}^{\mathrm{te}}\}}\,\Pi(\mathrm d\bm x,\mathrm d\xi)\right)
%        +\sum_{k=1}^{K}\gamma_k\left(\int_{\mathcal X\times\Xi}
%          \mathbf1_{\{\xi\le\xi_{(k)}^{\mathrm{te}}\}}\,\Pi(\mathrm d\bm x,\mathrm d\xi)-l_k\right)\\
%      ={}&\mathbb E_{\mathbb P}\left[\eta(\bm X)\right]
%        +\bm\lambda^\top\bm u-\bm\gamma^\top\bm l
%        +\int_{\mathcal X\times\Xi}\left[
%          \mathbf1_{\{\xi\le-\hat g(\bm x)\}}-\eta(\bm x)-\zeta_{\bm\lambda,\bm\gamma}(\xi)
%        \right]\,\Pi(\mathrm d\bm x,\mathrm d\xi).
%    \end{aligned}
%  \]
%  Its supremum over finite nonnegative measures is finite exactly when the integrand is nonpositive at every $(\bm x,\xi)$. A positive value at any point permits an arbitrarily large point mass there; otherwise the supremum integral is zero, attained by the zero measure. The functional upper dual therefore minimizes $\mathbb E_{\mathbb P}[\eta(\bm X)]+\bm\lambda^\top\bm u-\bm\gamma^\top\bm l$ subject to $\eta(\bm x)\ge\mathbf1_{\{\xi\le-\hat g(\bm x)\}}-\zeta_{\bm\lambda,\bm\gamma}(\xi)$ at every pair. For the lower problem, reversing the rank penalties gives
%  \[
%    \mathcal L^-
%      =\mathbb E_{\mathbb P}\left[\eta(\bm X)\right]
%        -\bm\lambda^\top\bm u+\bm\gamma^\top\bm l
%        +\int_{\mathcal X\times\Xi}\left[
%          \mathbf1_{\{\xi\le-\hat g(\bm x)\}}-\eta(\bm x)+\zeta_{\bm\lambda,\bm\gamma}(\xi)
%        \right]\,\Pi(\mathrm d\bm x,\mathrm d\xi).
%  \]
%  Its infimum over finite nonnegative measures is finite exactly when the integrand is nonnegative everywhere, in which case the infimum integral is zero. The lower functional dual maximizes $\mathbb E_{\mathbb P}[\eta(\bm X)]-\bm\lambda^\top\bm u+\bm\gamma^\top\bm l$ subject to $\eta(\bm x)\le\mathbf1_{\{\xi\le-\hat g(\bm x)\}}+\zeta_{\bm\lambda,\bm\gamma}(\xi)$ at every pair.
%
%  For fixed rank multipliers, the best marginal functions are the pointwise envelopes. For $\bm x\in\mathcal X_i$, they satisfy the identities
%  \begin{equation}
%    \begin{aligned}
%      \eta^+_{\bm\lambda,\bm\gamma}(\bm x)
%        &:=\sup_{\xi\in\Xi}\left\{\mathbf1_{\{\xi\le-\hat g(\bm x)\}}
%          -\zeta_{\bm\lambda,\bm\gamma}(\xi)\right\}
%          =\max_{1\le j\le J}\left\{\overline h_{ij}-\zeta_j\right\},\\
%      \eta^-_{\bm\lambda,\bm\gamma}(\bm x)
%        &:=\inf_{\xi\in\Xi}\left\{\mathbf1_{\{\xi\le-\hat g(\bm x)\}}
%          +\zeta_{\bm\lambda,\bm\gamma}(\xi)\right\}
%          =\min_{1\le j\le J}\left\{\underline h_{ij}+\zeta_j\right\}.
%    \end{aligned}
%    \label{eq-dual-envelope-identities}
%  \end{equation}
%  To verify them, note that outside $\Xi_{2i+1}$ the two cell coefficients equal the actual failure indicator. In the overlapping interval, the lower coefficient zero is attained by $\xi=\left(\xi_{(i+1)}^{\mathrm{te}}-\hat g(\bm x)\right)/2$ if $i<K$, and by $\xi=-\hat g(\bm x)+1$ if $i=K$. The upper coefficient one is attained by $\xi=-\hat g(\bm x)$ whenever $-\hat g(\bm x)>\xi_{(i)}^{\mathrm{te}}$. The only remaining case is $-\hat g(\bm x)=\xi_{(i)}^{\mathrm{te}}$ with $i\ge1$. Then the overlapping interval is safe, but
%  \[
%    \zeta_{2i}=\zeta_{2i+1}-\gamma_i\le\zeta_{2i+1},
%    \qquad
%    1-\zeta_{2i}\ge1-\zeta_{2i+1}.
%  \]
%  Choosing the left endpoint therefore attains an upper-envelope value at least as large as the candidate from the overlapping interval. Together with \eqref{eq-cell-costs}, this proves both identities, including threshold equalities.
%
%  The envelopes are bounded and constant on the Borel cells $\mathcal X_i$, so they belong to $\mathcal F_{\mathrm b}(\mathcal X)$. They respectively minimize and maximize the expectation over feasible marginal functions. Substituting them into the functional dual objectives gives \eqref{eq-upper-measure-dual} and \eqref{eq-lower-measure-dual}. Introducing epigraph variables $z_i$ for the upper maxima and hypograph variables $z_i$ for the lower minima gives the equivalent LPs
%  \begin{equation}
%    \begin{aligned}
%      \overline d_\theta
%      =\inf_{\bm\lambda,\bm\gamma\ge\bm0,\ \bm z\in\mathbb R^{K+1}}\quad&
%        \bm\lambda^\top\bm u-\bm\gamma^\top\bm l+\sum_{i=0}^{K}\mathbb P(\mathcal X_i) z_i\\
%      \text{s.t.}\quad&z_i+\zeta_j\ge\overline h_{ij},
%        \qquad i=0,\ldots,K,\quad j=1,\ldots,J,
%    \end{aligned}
%    \label{eq-finite-upper-dual}
%  \end{equation}
%  and
%  \begin{equation}
%    \begin{aligned}
%      \underline d_\theta
%      =\sup_{\bm\lambda,\bm\gamma\ge\bm0,\ \bm z\in\mathbb R^{K+1}}\quad&
%        -\bm\lambda^\top\bm u+\bm\gamma^\top\bm l+\sum_{i=0}^{K}\mathbb P(\mathcal X_i) z_i\\
%      \text{s.t.}\quad&z_i-\zeta_j\le\underline h_{ij},
%        \qquad i=0,\ldots,K,\quad j=1,\ldots,J.
%    \end{aligned}
%    \label{eq-finite-lower-dual}
%  \end{equation}
%  Each $z_i$ can be chosen as its corresponding finite envelope, including on input cells with $\mathbb P(\mathcal X_i)=0$. Thus introducing these variables preserves both optimal values.
%
%  For a feasible joint law, the signs of the rank penalties give $\mathcal L^-\le p_{\hat f}(\Pi)\le\mathcal L^+$. Feasible lower dual objectives therefore bound the primal objectives from below, and feasible upper dual objectives bound them from above. Taking the respective optima over the common nonempty primal feasible set proves \eqref{eq-weak-duality}.
%\end{proof}
%
%Lemma~\ref{lem-weak-duality} gives certificates for the primal extrema through optimization over $2K$ nonnegative rank multipliers. Its equivalent LP forms \eqref{eq-finite-upper-dual} and \eqref{eq-finite-lower-dual} add $K+1$ free variables $z_i$ and $(K+1)(2K+1)$ pairwise inequalities. Each feasible dual choice supplies a bound in the corresponding direction through \eqref{eq-weak-duality}.
%
%\Needspace{5\baselineskip}
%The following theorem shows that these dual bounds are exact by constructing feasible joint laws that attain the finite primal objectives.
%
%\begin{theorem}[Exact finite reformulations and strong duality]
%  \label{thm-finite-dual-reformulations}
%  For $\mathcal Q=\mathcal Q^{\mathrm{DW}}(\theta)$ and every $\theta\ge0$, the failure probability bounds admit the exact representations
%  \begin{equation}
%    \overline p_{\hat f}
%      =\max_{\bm\pi\in\mathcal C_\theta}\sum_{i=0}^{K}\sum_{j=1}^{J}\overline h_{ij}\pi_{ij},
%    \qquad
%    \underline p_{\hat f}
%      =\min_{\bm\pi\in\mathcal C_\theta}\sum_{i=0}^{K}\sum_{j=1}^{J}\underline h_{ij}\pi_{ij}.
%    \label{eq-finite-primal-values}
%  \end{equation}
%  Their Lagrangian dual values in Lemma~\ref{lem-weak-duality} satisfy
%  \begin{equation}
%    \overline p_{\hat f}=\overline d_\theta,
%    \qquad
%    \underline p_{\hat f}=\underline d_\theta.
%    \label{eq-strong-duality}
%  \end{equation}
%  The finite primal and dual optima, the measure primal optima, and the functional dual optima are attained.
%\end{theorem}
%\begin{proof}
%  Any feasible joint law induces a matrix in $\mathcal C_\theta$. Integrating \eqref{eq-cell-costs} therefore bounds the upper measure problem by the finite maximization from above and the lower measure problem by the finite minimization from below.
%
%  For the reverse inequalities, fix $\bm\pi\in\mathcal C_\theta$. Draw $\bm x$ from $\mathbb P$. On $\mathcal X_i$ with $\mathbb P(\mathcal X_i)>0$, select a provisional residual cell $j$ with probability $\pi_{ij}/\mathbb P(\mathcal X_i)$; the row equality makes these probabilities sum to one. For $j\ne2i+1$, choose a fixed representative of $\Xi_j$, where the actual failure indicator equals both coefficients. In the overlapping cell, use the following assignments to define joint laws $\Pi^+$ and $\Pi^-$.
%
%  For the upper problem, assign $\xi=-\hat g(\bm x)$. This causes failure and remains in $\Xi_{2i+1}$ unless $-\hat g(\bm x)=\xi_{(i)}^{\mathrm{te}}$. At equality, which requires $i\ge1$, the assignment transfers the provisional mass to the left endpoint. For any $\widetilde\xi\in\left(\xi_{(i)}^{\mathrm{te}},\xi_{(i+1)}^{\mathrm{te}}\right)$ and every observed threshold,
%  \[
%    \mathbf1_{\{\xi_{(i)}^{\mathrm{te}}<\xi_{(k)}^{\mathrm{te}}\}}
%      =\mathbf1_{\{\widetilde\xi<\xi_{(k)}^{\mathrm{te}}\}},
%    \qquad
%    \mathbf1_{\{\xi_{(i)}^{\mathrm{te}}\le\xi_{(k)}^{\mathrm{te}}\}}
%      \ge\mathbf1_{\{\widetilde\xi\le\xi_{(k)}^{\mathrm{te}}\}}.
%  \]
%  Thus every open upper probability is unchanged and every closed lower probability can only increase. The upper assignment preserves all rank inequalities and attains the coefficient $\overline h_{i,2i+1}=1$.
%
%  For the lower problem, assign $\xi=\left(\xi_{(i+1)}^{\mathrm{te}}-\hat g(\bm x)\right)/2$ if $i<K$, and $\xi=-\hat g(\bm x)+1$ if $i=K$. These values remain in $\Xi_{2i+1}$ and strictly exceed the failure threshold, including at threshold equality. Hence the lower assignment retains the cell matrix and attains $\underline h_{i,2i+1}=0$.
%
%  If $\mathbb P(\mathcal X_i)=0$, the corresponding row is zero; set $\xi=\xi_{(1)}^{\mathrm{te}}$ on that input cell. All assignments are Borel measurable, and their finite mixtures define probability kernels. Combining them with $\mathbb P$ gives joint laws with the full prescribed input marginal, including randomization at input atoms. Both laws are feasible and satisfy
%  \[
%    p_{\hat f}(\Pi^+)=\sum_{i=0}^{K}\sum_{j=1}^{J}\overline h_{ij}\pi_{ij},
%    \qquad
%    p_{\hat f}(\Pi^-)=\sum_{i=0}^{K}\sum_{j=1}^{J}\underline h_{ij}\pi_{ij}.
%  \]
%  The lower law retains every provisional cell mass. The upper law may change the cell matrix at threshold equalities, but it preserves feasibility and the assigned upper objective. These identities establish the reverse inequalities required for \eqref{eq-finite-primal-values}.
%
%  For every $\theta\ge0$, the empirical product law induces the feasible matrix
%  \[
%    \pi_{i,2k}^{\mathrm{emp}}:=\mathbb P(\mathcal X_i)\frac{c_k-c_{k-1}}{N^{\mathrm{te}}},
%    \qquad i=0,\ldots,K,\quad k=1,\ldots,K,
%  \]
%  with all odd columns zero. Its open and closed rank probabilities are $c_{k-1}/N^{\mathrm{te}}$ and $c_k/N^{\mathrm{te}}$, respectively, so feasibility includes $\theta=0$. The nonnegative entries of a matrix in $\mathcal C_\theta$ sum to one, and the constraints are closed. Thus $\mathcal C_\theta$ is nonempty and compact. The finite objectives lie in $[0,1]$ and attain their optima, which the joint-law constructions realize in the measure problems.
%
%  The finite primal LPs in \eqref{eq-finite-primal-values} have row multipliers $z_i$ and rank multipliers $\lambda_k,\gamma_k$. Their dual constraints are precisely $z_i+\zeta_j\ge\overline h_{ij}$ and $z_i-\zeta_j\le\underline h_{ij}$. Finite LP strong duality therefore gives \eqref{eq-strong-duality} and dual attainment \citep[Section~5.2.3]{B.V04}. By \eqref{eq-dual-envelope-identities}, optimal rank multipliers also yield optimal bounded marginal functions for the original functional duals. Feasibility and finite optimal values suffice; strict rank inequalities are unnecessary, including at $\theta=0$.
%\end{proof}
%
%Theorem~\ref{thm-finite-dual-reformulations} replaces the search over joint distributions by exact allocations of joint cell probability. The conditional residual assignments retain the full input law, and the equalities also hold with residual atoms, zero-probability input cells, and $\theta=0$.
%
%\Needspace{5\baselineskip}
%The following corollary uses the ordering of the cells to evaluate both extrema explicitly and identify optimal certificates with at most one nonzero rank multiplier.
%
%\begin{corollary}[Closed-form bounds and optimal certificates]
%  \label{thm-closed-form-bounds}
%  Under the conditions of Theorem~\ref{thm-finite-dual-reformulations},
%  \begin{align}
%    \overline p_{\hat f}
%      &=\min\left\{1,\min_{k=1,\ldots,K}\left[u_k+\sum_{i=k}^{K}\mathbb P(\mathcal X_i)\right]\right\},
%      \label{eq-dw-closed-upper}\\
%    \underline p_{\hat f}
%      &=\max\left\{0,\max_{k=1,\ldots,K}\left[l_k+\sum_{i=k}^{K}\mathbb P(\mathcal X_i)-1\right]\right\}.
%      \label{eq-dw-closed-lower}
%  \end{align}
%  Each dual optimum admits a certificate with at most one nonzero rank multiplier, of value one.
%\end{corollary}
%\begin{proof}
%  For a candidate residual law $\mathbb Q\in\mathcal M_1(\Xi)$, define the finite transportation set
%  \[
%    \mathcal T(\mathbb Q):=\left\{\bm\pi\in\mathbb R_+^{(K+1)\times J}:
%    \begin{array}{ll}
%      \displaystyle\sum_{j=1}^{J}\pi_{ij}=\mathbb P(\mathcal X_i),&i=0,\ldots,K,\\[1mm]
%      \displaystyle\sum_{i=0}^{K}\pi_{ij}=\mathbb Q(\Xi_j),&j=1,\ldots,J
%    \end{array}\right\}.
%  \]
%  It contains the product matrix $\pi_{ij}=\mathbb P(\mathcal X_i)\mathbb Q(\Xi_j)$. Every residual cell is nonempty, so any nonnegative column allocation with total mass one is realized by a residual law. Proposition~\ref{prop-finite-rank-constraints} consequently gives
%  \[
%    \mathcal C_\theta
%      =\bigcup_{\mathbb Q\in\mathcal Q^{\mathrm{DW}}(\theta)}\mathcal T(\mathbb Q).
%  \]
%  The parameter $\mathbb Q$ fixes only residual partition masses in this finite matrix problem. The joint-law construction in Theorem~\ref{thm-finite-dual-reformulations} preserves feasibility and the assigned objective while allowing the residual marginal to change, including endpoint transfers in the upper problem.
%
%  For fixed cell masses $\mathbb Q(\Xi_j)$, the upper problem maximizes mass on pairs $j\le2i+1$; the lower problem maximizes safe mass on pairs $j>2i$ and subtracts it from one. Both are maximum-flow problems with source-to-row capacities $\mathbb P(\mathcal X_i)$, column-to-sink capacities $\mathbb Q(\Xi_j)$, and capacity two on each permitted row-to-column edge. A transportation matrix restricts to a feasible flow. Conversely, any flow can be completed to a transportation matrix because the remaining row and column masses are nonnegative with equal totals. Thus the maximum flow equals the maximum permitted mass.
%
%  A cut of capacity one is available, so a minimum cut crosses no edge of capacity two. Its source-side columns can be taken to be exactly the neighbors of its source-side rows. Since the neighborhoods are nested, adding rows whose neighbors are already included cannot increase the capacity. Apart from empty and full row sets, upper-problem cuts reduce to rows $i<k$ and columns $j<2k$, while safe-mass cuts reduce to rows $i\ge k$ and columns $j>2k$. For $k=1,\ldots,K$, their capacities are respectively
%  \[
%    \sum_{i=k}^{K}\mathbb P(\mathcal X_i)+\sum_{j<2k}\mathbb Q(\Xi_j),
%    \qquad
%    \sum_{i=0}^{k-1}\mathbb P(\mathcal X_i)+\sum_{j>2k}\mathbb Q(\Xi_j).
%  \]
%  Empty and full row sets give capacity one. The max-flow/min-cut theorem \citep{F.F56} therefore yields
%  \begin{align*}
%    \max_{\bm\pi\in\mathcal T(\mathbb Q)}\sum_{i,j}\overline h_{ij}\pi_{ij}
%      &=\min\left\{1,\min_k\left[\sum_{i=k}^{K}\mathbb P(\mathcal X_i)
%        +\sum_{j<2k}\mathbb Q(\Xi_j)\right]\right\},\\
%    \min_{\bm\pi\in\mathcal T(\mathbb Q)}\sum_{i,j}\underline h_{ij}\pi_{ij}
%      &=\max\left\{0,\max_k\left[\sum_{i=k}^{K}\mathbb P(\mathcal X_i)
%        +\sum_{j\le2k}\mathbb Q(\Xi_j)-1\right]\right\}.
%  \end{align*}
%
%  These values are nondecreasing in the respective open and closed cumulative residual probabilities. The rank constraints bound the former from above by $u_k$ and the latter from below by $l_k$. To attain all relevant limits simultaneously, choose laws $\mathbb Q^+$ and $\mathbb Q^-$ with cell probabilities
%  \begin{align*}
%    &\mathbb Q^+(\Xi_1):=u_1,\quad
%      \mathbb Q^+(\Xi_{2k}):=u_{k+1}-u_k,\quad
%      \mathbb Q^+(\Xi_{2k+1}):=0,\\
%    &\mathbb Q^-(\Xi_{2k-1}):=0,\quad
%      \mathbb Q^-(\Xi_{2k}):=l_k-l_{k-1},\quad
%      \mathbb Q^-(\Xi_J):=1-l_K,
%      \qquad k=1,\ldots,K.
%  \end{align*}
%  By Proposition~\ref{prop-finite-rank-constraints}, the prescribed masses are nonnegative and sum to one. Choosing a representative in each residual cell realizes both laws. Their cumulative probabilities satisfy
%  \begin{align*}
%    &\sum_{j<2k}\mathbb Q^+(\Xi_j)=u_k,
%      \qquad \sum_{j\le2k}\mathbb Q^+(\Xi_j)=u_{k+1}\ge l_k,\\
%    &\sum_{j\le2k}\mathbb Q^-(\Xi_j)=l_k,
%      \qquad \sum_{j<2k}\mathbb Q^-(\Xi_j)=l_{k-1}\le u_k.
%  \end{align*}
%  Both laws therefore belong to $\mathcal Q^{\mathrm{DW}}(\theta)$. The union representation of $\mathcal C_\theta$ and \eqref{eq-finite-primal-values} show that these allocations attain the upper and lower bounds in \eqref{eq-dw-closed-upper} and \eqref{eq-dw-closed-lower}. Zero masses and coincident cumulative limits are allowed, including at $\theta=0$.
%
%  Finally, let $\bm e_k$ be the $k$th unit vector. Certificates for the upper and lower LPs \eqref{eq-finite-upper-dual} and \eqref{eq-finite-lower-dual} are respectively
%  \begin{align*}
%    &\bm\lambda=\bm e_k,\qquad\bm\gamma=\bm0,\qquad
%      z_i=\mathbf1_{\{i\ge k\}},\\
%    &\bm\lambda=\bm0,\qquad\bm\gamma=\bm e_k,\qquad
%      z_i=\mathbf1_{\{i\ge k\}}-1.
%  \end{align*}
%  Their feasibility follows from
%  \[
%    \mathbf1_{\{i\ge k\}}+\mathbf1_{\{j<2k\}}\ge\overline h_{ij},
%    \qquad
%    \mathbf1_{\{i\ge k\}}-1+\mathbf1_{\{j\le2k\}}\le\underline h_{ij}.
%  \]
%  Their objectives are $u_k+\sum_{i=k}^{K}\mathbb P(\mathcal X_i)$ and $l_k+\sum_{i=k}^{K}\mathbb P(\mathcal X_i)-1$, the inner candidates in the closed-form expressions. Zero rank multipliers with $z_i=1$ or $z_i=0$ supply the outer candidates. Selecting the best candidate gives an optimal full dual solution, whose functional marginal multiplier follows from \eqref{eq-dual-envelope-identities}. The inclusion of $i=k$ in the tail sum retains inputs whose failure threshold equals the observed residual.
%\end{proof}
%
%In Corollary~\ref{thm-closed-form-bounds}, the upper bound combines residual mass below an observed value with the probability that the surrogate threshold reaches it. The lower bound measures the overlap required by every admissible coupling.
%
%For the response-threshold curves in Corollary~\ref{cor-simultaneous-threshold-bounds}, apply the same construction to $g-c$ and $\hat g-c$. Every residual, rank count, and endpoint $l_k,u_k$ is unchanged. The input partition now depends on $c$, and its cumulative probability at rank $k$ is $\mathbb P\left\{\hat g(\bm X)\le c-\xi_{(k)}^{\mathrm{te}}\right\}$. Thus Corollary~\ref{thm-closed-form-bounds} gives
%\begin{align}
%  \overline p_{\hat f}(c)
%    &=\min\left\{1,\min_{k=1,\ldots,K}\left[u_k
%      +\mathbb P\left\{\hat g(\bm X)\le c-\xi_{(k)}^{\mathrm{te}}\right\}\right]\right\},
%    \label{eq-threshold-closed-upper}\\
%  \underline p_{\hat f}(c)
%    &=\max\left\{0,\max_{k=1,\ldots,K}\left[l_k
%      +\mathbb P\left\{\hat g(\bm X)\le c-\xi_{(k)}^{\mathrm{te}}\right\}-1\right]\right\}.
%    \label{eq-threshold-closed-lower}
%\end{align}
%These identities hold for every $c\in\mathbb R$ and $\theta\ge0$, including $\theta=0$. When $\theta\ge\theta_{N^{\mathrm{te}}}(\delta)$, their coverage is simultaneous by \eqref{eq-threshold-curve-coverage}. An extremizing coupling may depend on $c$; the common feasible true joint law supplies the simultaneous guarantee.
%
%For $p_{\mathrm{tar}}\in(0,1)$, the upper curve satisfies $\overline p_{\hat f}(c)\le p_{\mathrm{tar}}$ exactly when some rank satisfies
%\[
%  u_k+\mathbb P\left\{\hat g(\bm X)\le c-\xi_{(k)}^{\mathrm{te}}\right\}\le p_{\mathrm{tar}}.
%\]
%Both terms must therefore be small at the same residual level.
%
%With a known residual law, the coupling problem is the classical fixed-marginal sum problem studied by \citet{M82} and \citet{F.N.S87}. Here the residual law also varies within a calibrated DW band, and the open and closed rank constraints retain atoms.
%
%Dependence uncertainty can remain even with a known residual marginal. For equally likely surrogate responses $0$ and $2$ and equally likely residuals $-1$ and $1$, failure occurs exactly when the response is $0$ and the residual is $-1$. Its probability ranges from $0$ to $1/2$ over couplings, whereas independence gives $1/4$. Further marginal calibration alone does not resolve this dependence uncertainty.
%
%\Needspace{7\baselineskip}
%The following remark incorporates known interval support through additional constraints at its endpoints.
%
%\begin{remark}[Known interval support]
%  \label{rem-interval-support}
%  Suppose $\Xi=[a,b]$ with $-\infty<a\le b<\infty$ and the observed residuals lie in this interval. Then
%  \begin{align*}
%    \overline p_{\hat f}
%      &=\min\left\{\mathbb P\left\{\hat g(\bm X)\le-a\right\},
%        \min_{k=1,\ldots,K}\left[u_k+\mathbb P\left\{\hat g(\bm X)\le-\xi_{(k)}^{\mathrm{te}}\right\}\right]\right\},\\
%    \underline p_{\hat f}
%      &=\max\left\{\mathbb P\left\{\hat g(\bm X)\le-b\right\},
%        \max_{k=1,\ldots,K}\left[l_k+\mathbb P\left\{\hat g(\bm X)\le-\xi_{(k)}^{\mathrm{te}}\right\}-1\right]\right\}.
%  \end{align*}
%  Represent the support restriction on $\mathbb R$ by $\mathbb Q\{\xi<a\}\le0$ and $\mathbb Q\{\xi\le b\}\ge1$, and add $a,b$ to the ordered threshold list. At a newly added $a$, use upper and lower bounds zero; at a newly added $b$, use upper and lower bounds one. At coincident observed thresholds, tighten only the upper bound at $a$ and the lower bound at $b$. The enlarged sequences satisfy the compatibility conditions in Proposition~\ref{prop-finite-rank-constraints}, so the same finite-aggregation and ordered-cut arguments apply. The effective additional upper candidate is $\mathbb P\left\{\hat g(\bm X)\le-a\right\}$, and the effective additional lower candidate is $\mathbb P\left\{\hat g(\bm X)\le-b\right\}$; the other added candidates are redundant. If $a=b$, both bounds equal $\mathbb P\left\{\hat g(\bm X)\le-a\right\}$.
%
%  For the partition on $\Xi=[a,b]$, intersect the residual cells from the augmented threshold list with $\Xi$ and omit empty cells. Including the support endpoints retains their atoms and the corresponding changes in failure status. For a response threshold $c$, shift each surrogate-response threshold in the displayed event probabilities by $c$.
%\end{remark}
%
%Remark~\ref{rem-interval-support} separates inputs that fail for every admissible residual, $\hat g(\bm X)\le-b$, from inputs that are always safe, $\hat g(\bm X)>-a$. In particular, $\mathbb P\left\{\hat g(\bm X)\le-a\right\}=0$ rules out failure.
%
%After zero failures in $N^{\mathrm{te}}$ independent tests, the one-sided exact binomial upper limit is $1-\delta^{1/N^{\mathrm{te}}}\approx\log(1/\delta)/N^{\mathrm{te}}$ \citep{C.P34}. At $c=0$, the simple DW case $K=1$ and $\mathbb P(\mathcal X_1)=0$ gives an upper bound $\beta_0(\theta)$ determined by the simultaneous-band calibration.
%
%Evaluation requires scalar DW inversion, sorting and counting the test residuals, and computing the $K$ input event probabilities $\mathbb P\left\{\hat g(\bm X)\le c-\xi_{(k)}^{\mathrm{te}}\right\}$ at each response threshold $c$. The endpoint table depends only on $N^{\mathrm{te}}$ and $\theta$ and can be precomputed. Given this table and the input event probabilities, the two scans cost $O(K)$ at each response threshold.
%
%For $\theta\ge\theta_{N^{\mathrm{te}}}(\delta)$, the exact expressions retain the coverage in Theorem~\ref{theorem-01} and Corollary~\ref{cor-simultaneous-threshold-bounds}. Independent auxiliary simulation of $\bm X\sim\mathbb P$ can supply simultaneous upper and lower confidence bounds on the input event probabilities. Substituting these bounds in the respective upper and lower expressions preserves conservative failure probability bounds. For one fixed $c$, simultaneous bounds at its $K$ evaluation thresholds suffice. To retain coverage of the entire response-threshold curve, the auxiliary bounds must hold for the surrogate-response CDF at every real threshold. They must hold conditionally on the test data, with their error probability included in the overall confidence budget.
%
%
% END PRESERVED BODY BEFORE CONCISE RESULT TRANSITIONS

% BEGIN PRESERVED BODY BEFORE NAMED RESULT TRANSITIONS (2026-11-26)
%\section{Problem Description}
%\label{sec-problem}
%
%Let $\bm X\sim\mathbb P$ be a random input taking values in $\mathcal X\subseteq\mathbb R^n$, and let $g:\mathcal X\to\mathbb R$ be a measurable limit-state function. The failure probability is defined as
%\[
%  p_f:=\Pr\{g(\bm X)\le0\}=\mathbb E_{\bm X\sim\mathbb P}\left[\mathbf 1_{\{g(\bm X)\le0\}}\right]=\int_{\mathcal X}\mathbf 1_{\{g(\bm x)\le0\}}\,\mathbb P(\mathrm d\bm x).
%\]
%
%Evaluating $g$ can require an expensive high-fidelity simulation. We use the training dataset
%\[
%  \mathcal D^{\mathrm{tr}}:=\left\{\left(\bm x_i^{\mathrm{tr}},g(\bm x_i^{\mathrm{tr}})\right)\right\}_{i=1}^{N^{\mathrm{tr}}}
%\]
%to fit a computationally efficient, measurable surrogate $\hat g:=\hat g_{\mathcal D^{\mathrm{tr}}}$. Model fitting and selection use only the training data and are completed before the test data are used. We condition on $\mathcal D^{\mathrm{tr}}$, treat $\hat g$ as fixed, and suppress this dependence in the notation.
%
%Define the signed residual function
%\[
%  \xi(\bm x):=g(\bm x)-\hat g(\bm x),\qquad \forall \bm x\in\mathcal X.
%\]
%For a new input $\bm X\sim\mathbb P$ independent of the training data, let $\mathbb Q^\star$ be the distribution of $\xi^\star:=\xi(\bm X)$. Let $\Xi\subseteq\mathbb R$ be a known Borel set with $\mathbb Q^\star(\Xi)=1$; the default choice is $\Xi=\mathbb R$. For any Borel set $\mathcal S$ in a Euclidean space, write $\mathcal B(\mathcal S)$ for its Borel sigma-algebra. Let $\mathcal M_+(\mathcal S)$ denote the finite nonnegative Borel measures concentrated on $\mathcal S$, and let $\mathcal M_1(\mathcal S)$ denote the probability measures.
%
%For a candidate residual law $\mathbb Q\in\mathcal M_1(\Xi)$, let $\Pi$ be a joint distribution on $\mathcal X\times\Xi$ with marginals $\Pi_{\bm X}=\mathbb P$ and $\Pi_\xi=\mathbb Q$. Its surrogate-based failure probability is
%\begin{equation}
%  p_{\hat f}(\Pi):=\Pr\{\hat g(\bm X)+\xi\le0\}=\int_{\mathcal X\times\Xi}\mathbf 1_{\{\hat g(\bm x)+\xi\le0\}}\,\Pi(\mathrm d\bm x,\mathrm d\xi).
%  \label{eq-surrogate-failure-probability}
%\end{equation}
%For fixed $\hat g$, this functional is linear in $\Pi$. In particular, the true joint distribution $\Pi^\star$ satisfies $p_{\hat f}(\Pi^\star)=p_f$.
%
%The residual marginal and its dependence on the input jointly determine the failure probability. We represent residual-marginal uncertainty by an ambiguity set $\mathcal Q$ calibrated to contain $\mathbb Q^\star$ with high confidence. For a given $\mathcal Q$, define the joint ambiguity set
%\[
%  \mathcal P(\mathcal Q) :=\left\{\Pi\in\mathcal M_1(\mathcal X\times\Xi):\Pi_{\bm X}=\mathbb P,\ \Pi_\xi\in\mathcal Q\right\}.
%\]
%The true residual is determined by the input, but its unknown functional relationship is not encoded in $\mathcal Q$. The set $\mathcal P(\mathcal Q)$ therefore allows all couplings with the specified marginals, including randomized conditional residual laws. The upper and lower failure probability bounds are
%\begin{equation}
%  \overline p_{\hat f}:=\sup_{\Pi\in\mathcal P(\mathcal Q)}p_{\hat f}(\Pi),
%  \label{eq-upper-bound}
%\end{equation}
%and
%\begin{equation}
%  \underline p_{\hat f}:=\inf_{\Pi\in\mathcal P(\mathcal Q)}p_{\hat f}(\Pi),
%  \label{eq-lower-bound}
%\end{equation}
%respectively. If $\mathbb Q^\star\in\mathcal Q$, then $\Pi^\star\in\mathcal P(\mathcal Q)$ and hence $\underline p_{\hat f}\le p_f\le\overline p_{\hat f}$.
%
%
%Section~\ref{sec-ambiguity-set} constructs the residual ambiguity set and establishes finite-sample coverage of the failure probability bounds, simultaneously across response thresholds. Section~\ref{sec-reformulation} then derives their Lagrangian duals, proves strong duality through a finite reformulation, and obtains closed-form expressions.
%
%
%\section{Ambiguity Set Design}
%\label{sec-ambiguity-set}
%
%Failure occurs when $\xi\le-\hat g(\bm X)$, so the residual cumulative distribution function (CDF) gives the marginal mass at or below each failure threshold. We use the D\"umbgen--Wellner (DW) construction \citep{L.J23} to control empirical CDF deviations simultaneously over all thresholds, with tolerances that vary across probability levels.
%
%Let $q$ and $p$ denote an empirical probability and a candidate probability. Following \citet{B.J79} and the CDF confidence bands of \citet{O95}, we measure their local discrepancy by the Bernoulli Kullback--Leibler divergence
%\begin{equation}
%  \phi(q,p):=q\log\frac{q}{p}+(1-q)\log\frac{1-q}{1-p},\qquad \forall q\in[0,1],\, p\in(0,1),
%  \label{eq-bernoulli-kl}
%\end{equation}
%where logarithms are natural and a term with zero numerator is interpreted as zero. For $q$ close to an interior $p$, the approximation $\phi(q,p)\approx(q-p)^2/[2p(1-p)]$ scales the deviation by the Bernoulli variance. Thus, the same absolute deviation receives a larger weight near the probability boundaries. The full divergence also covers zero and full empirical counts through $\phi(0,p)=-\log(1-p)$ and $\phi(1,p)=-\log p$. At the boundaries of the candidate probability, we set $\phi(0,0)=\phi(1,1)=0$, $\phi(q,0)=+\infty$ for $q>0$, and $\phi(q,1)=+\infty$ for $q<1$.
%
%To compare deviations simultaneously across probability levels, we use the additive correction of \citet{L.J23}. For $v\in(0,1)$, define
%\begin{equation*}
%  \rho(v):=\log\log\frac{e}{4v(1-v)},\qquad \omega(v):=\log\!\left(1+\rho(v)^2\right).
%\end{equation*}
%Before sampling, we fix the divergence parameter $s=1$, which gives \eqref{eq-bernoulli-kl}, and the correction parameter $\nu=1$. Set $\psi_0(v):=\rho(v)+\omega(v)$, with $\psi_0(0)=\psi_0(1)=+\infty$, and define the bivariate correction
%\begin{equation*}
%  \psi(q,p):=\min_{v\in[\min\{q,p\},\,\max\{q,p\}]}\psi_0(v).
%\end{equation*}
%The functions $\rho$ and $\omega$ vanish at $v=1/2$ and increase toward the probability boundaries. The minimum over the interval between $q$ and $p$ is therefore large only when both probabilities lie in the same tail. This correction balances fluctuations across probability scales in the simultaneous CDF constraint.
%
%We construct the empirical CDF from a test dataset
%\[
%  \mathcal D^{\mathrm{te}}:=\left\{\left(\bm x_j^{\mathrm{te}},g(\bm x_j^{\mathrm{te}})\right)\right\}_{j=1}^{N^{\mathrm{te}}}.
%\]
%The test inputs are independent draws from $\mathbb P$ and are independent of $\mathcal D^{\mathrm{tr}}$. We fix $N^{\mathrm{te}}\ge1$ before sampling and use these data to calibrate the residual ambiguity set.
%
%Define the signed test residuals
%\[
%  \xi_j^{\mathrm{te}}:=\xi(\bm x_j^{\mathrm{te}})=g(\bm x_j^{\mathrm{te}})-\hat g(\bm x_j^{\mathrm{te}}),\qquad j=1,\ldots,N^{\mathrm{te}}.
%\]
%Conditional on $\mathcal D^{\mathrm{tr}}$, these residuals are independent and identically distributed according to $\mathbb Q^\star$. Their empirical CDF is
%\begin{equation}
%  \widehat F^{\mathrm{te}}(t):=\frac{1}{N^{\mathrm{te}}}\sum_{j=1}^{N^{\mathrm{te}}}\mathbf 1_{\{\xi_j^{\mathrm{te}}\le t\}},\qquad \forall t\in\mathbb R.
%  \label{eq-empirical-cdf}
%\end{equation}
%For $\mathbb Q\in\mathcal M_1(\Xi)$, write $F_{\mathbb Q}(t):=\mathbb Q\left(({-\infty},t]\right)$. At a threshold $t$, $q=\widehat F^{\mathrm{te}}(t)$ is the mean of the Bernoulli indicators in \eqref{eq-empirical-cdf}, and $p=F_{\mathbb Q}(t)$ is the candidate probability. Thus $N^{\mathrm{te}}\phi(q,p)$ is the Bernoulli log-likelihood difference between the empirical and candidate probabilities. Combining it with the multiscale correction gives $N^{\mathrm{te}}\phi(q,p)-\psi(q,p)$.
%
%At the probability boundaries, the corrected discrepancy is $-\infty$ for $q=p\in\{0,1\}$ and $+\infty$ for $p=0<q$ or $p=1>q$. These conventions extend the constraint to every residual threshold.
%
%\Needspace{5\baselineskip}
%We use this discrepancy to define the residual ambiguity set.
%
%\begin{definition}[DW-based CDF ambiguity set]
%  \label{definition-dw-ambiguity-set}
%  Given the empirical CDF $\widehat F^{\mathrm{te}}$ and a threshold $\theta\ge0$, define
%  \begin{equation} 
%    \label{eq-cdf-ambiguity-set}
%      \mathcal Q^{\mathrm{DW}}(\theta):=\left\{\mathbb Q\in\mathcal M_1(\Xi):\ 
%      \sup_{t\in\mathbb R}\left[ N^{\mathrm{te}}\phi\!\left(\widehat F^{\mathrm{te}}(t),F_{\mathbb Q}(t)\right) 
%       {}-\psi\!\left(\widehat F^{\mathrm{te}}(t),F_{\mathbb Q}(t)\right)\right]\le\theta\right\}. 
%  \end{equation} 
%\end{definition}
%
%We calibrate $\theta$ against a distribution-free reference law. Following \citet{L.J23}, let $V_1,\ldots,V_{N^{\mathrm{te}}}$ be independent standard uniform variables, independent of the training and test data. Their empirical CDF is
%\begin{equation*}
%  G_{N^{\mathrm{te}}}(u):=\frac{1}{N^{\mathrm{te}}}\sum_{j=1}^{N^{\mathrm{te}}}\mathbf 1_{\{V_j\le u\}},\qquad \forall u\in[0,1].
%\end{equation*}
%Define the uniform reference statistic
%\begin{equation}
%  T_{N^{\mathrm{te}}}^{\mathrm U}:=\sup_{u\in(0,1)}\left[N^{\mathrm{te}}\phi\!\left(G_{N^{\mathrm{te}}}(u),u\right)-\psi\!\left(G_{N^{\mathrm{te}}}(u),u\right)\right].
%  \label{eq-dw-reference-statistic}
%\end{equation}
%
%
%\Needspace{5\baselineskip}
%We choose $\theta$ from a quantile of this statistic.
%
%\begin{proposition}[Finite-sample threshold calibration]
%  \label{proposition-01}
%  Under the sampling conditions stated above, let $\delta\in(0,1)$ and define
%  \begin{equation}
%    \theta_{N^{\mathrm{te}}}(\delta):=\inf\left\{\tau\ge0:\Pr\left\{T_{N^{\mathrm{te}}}^{\mathrm U}\le\tau\right\}\ge1-\delta\right\}.
%    \label{eq-dw-calibrated-threshold}
%  \end{equation}
%
%  Then, for every $\theta\ge\theta_{N^{\mathrm{te}}}(\delta)$, the DW-based CDF ambiguity set satisfies
%  \begin{equation}
%    \Pr\left\{\mathbb Q^\star\in\mathcal Q^{\mathrm{DW}}(\theta)\,\middle|\,\mathcal D^{\mathrm{tr}}\right\}\ge1-\delta.
%    \label{eq-cdf-coverage}
%  \end{equation}
%\end{proposition}
%\begin{proof}
%  The reference statistic is finite almost surely, and its distribution function is right-continuous. Hence \eqref{eq-dw-calibrated-threshold} gives
%  \[
%    \Pr\left\{T_{N^{\mathrm{te}}}^{\mathrm U}>\theta_{N^{\mathrm{te}}}(\delta)\right\}\le\delta.
%  \]
%  Conditioning on $\mathcal D^{\mathrm{tr}}$ fixes $\hat g$ and $\mathbb Q^\star$. Define the generalized quantile function $F_{\mathbb Q^\star}^{-1}(u):=\inf\{t\in\mathbb R:F_{\mathbb Q^\star}(t)\ge u\}$ for all $u\in(0,1)$. The sample $F_{\mathbb Q^\star}^{-1}(V_j)$ for $j=1,\ldots,N^{\mathrm{te}}$, has the same conditional joint law as the test residuals. Thus, given $\mathcal D^{\mathrm{tr}}$,
%  \[
%    \left(\widehat F^{\mathrm{te}}(t)\right)_{t\in\mathbb R}\overset{\mathrm d}{=}\left(G_{N^{\mathrm{te}}}\!\left(F_{\mathbb Q^\star}(t)\right)\right)_{t\in\mathbb R}.
%  \]
%  For $0<F_{\mathbb Q^\star}(t)<1$, the discrepancy at $\mathbb Q^\star$ is evaluated at a probability level included in \eqref{eq-dw-reference-statistic}. When $F_{\mathbb Q^\star}(t)$ is zero or one, the empirical CDF equals that boundary almost surely, giving discrepancy $-\infty$. The supremum in \eqref{eq-cdf-ambiguity-set}, evaluated at $\mathbb Q^\star$, is therefore bounded by $T_{N^{\mathrm{te}}}^{\mathrm U}$ under this coupling. Since $\mathbb Q^\star(\Xi)=1$, we obtain
%  \[
%    \Pr\left\{\mathbb Q^\star\notin\mathcal Q^{\mathrm{DW}}\!\left(\theta_{N^{\mathrm{te}}}(\delta)\right)\,\middle|\,\mathcal D^{\mathrm{tr}}\right\}
%    \le\Pr\left\{T_{N^{\mathrm{te}}}^{\mathrm U}>\theta_{N^{\mathrm{te}}}(\delta)\right\}\le\delta.
%  \]
%  For every test sample, $\theta\ge\theta_{N^{\mathrm{te}}}(\delta)$ implies $\mathcal Q^{\mathrm{DW}}\!\left(\theta_{N^{\mathrm{te}}}(\delta)\right)\subseteq\mathcal Q^{\mathrm{DW}}(\theta)$. This inclusion proves \eqref{eq-cdf-coverage} for every such $\theta$.
%\end{proof}
%
%The guarantee allows atoms and infinite moments. For fixed $\phi$ and $\psi$, the reference law depends only on $N^{\mathrm{te}}$, so thresholds can be precomputed. \citet{L.J23} give finite-sample evaluation through uniform order-statistic boundary probabilities. A conservative quantile upper bound preserves \eqref{eq-cdf-coverage}.
%% If calibration is performed by simulation, quantile estimation error must be accounted for in the confidence guarantee. A threshold chosen from the test data also preserves \eqref{eq-cdf-coverage}, provided $\theta\ge\theta_{N^{\mathrm{te}}}(\delta)$ almost surely, by the same nesting argument.
%
%\Needspace{5\baselineskip}
%Residual-law coverage also controls the failure probability interval.
%
%\begin{theorem}[Finite-sample failure probability bounds]
%  \label{theorem-01}
%  Under the sampling conditions of Proposition~\ref{proposition-01}, let $\delta\in(0,1)$ and $\theta\ge\theta_{N^{\mathrm{te}}}(\delta)$, with the calibrated threshold defined in \eqref{eq-dw-calibrated-threshold}. Then the bounds defined in \eqref{eq-upper-bound} and \eqref{eq-lower-bound}, evaluated with $\mathcal Q=\mathcal Q^{\mathrm{DW}}(\theta)$, satisfy
%  \begin{equation}
%    \Pr\left\{\underline p_{\hat f}\le p_f\le\overline p_{\hat f}\,\middle|\,\mathcal D^{\mathrm{tr}}\right\}\ge1-\delta.
%    \label{eq-failure-probability-coverage}
%  \end{equation}
%\end{theorem}
%
%\begin{proof}
%  On the event $\mathbb Q^\star\in\mathcal Q^{\mathrm{DW}}(\theta)$, the true joint distribution $\Pi^\star$ belongs to $\mathcal P\left(\mathcal Q^{\mathrm{DW}}(\theta)\right)$. Hence
%  \[
%    \underline p_{\hat f}\le p_{\hat f}(\Pi^\star)=p_f\le\overline p_{\hat f}.
%  \]
%  By Proposition~\ref{proposition-01}, this event has conditional probability at least $1-\delta$ given $\mathcal D^{\mathrm{tr}}$, proving \eqref{eq-failure-probability-coverage}.
%\end{proof}
%
%
%For fixed test data, $0\le\theta_1\le\theta_2$ implies $\mathcal Q^{\mathrm{DW}}(\theta_1)\subseteq\mathcal Q^{\mathrm{DW}}(\theta_2)$. The lower bound is nonincreasing and the upper bound nondecreasing in $\theta$, so higher confidence gives a wider or unchanged interval.
%
%\Needspace{5\baselineskip}
%The same coverage event can be used across response thresholds.
%
%\begin{corollary}[Simultaneous bounds across response thresholds]
%  \label{cor-simultaneous-threshold-bounds}
%  Under the conditions of Theorem~\ref{theorem-01}, let $\mathcal Q=\mathcal Q^{\mathrm{DW}}(\theta)$. For $c\in\mathbb R$, write $p_f(c):=\Pr \left\{g(\bm X)\le c\right\}$ and define
%  \[
%    p_{\hat f}(\Pi;c):=\int_{\mathcal X\times\Xi}
%      \mathbf1_{\{\hat g(\bm x)+\xi\le c\}}\,\Pi(\mathrm d\bm x,\mathrm d\xi),
%  \]
%  with lower and upper bounds
%  \[
%    \underline p_{\hat f}(c):=\inf_{\Pi\in\mathcal P(\mathcal Q)}p_{\hat f}(\Pi;c),
%    \qquad
%    \overline p_{\hat f}(c):=\sup_{\Pi\in\mathcal P(\mathcal Q)}p_{\hat f}(\Pi;c).
%  \]
%  Then
%  \begin{equation}
%    \Pr\left\{\underline p_{\hat f}(c)\le p_f(c)\le\overline p_{\hat f}(c),
%      \ \forall c\in\mathbb R\,\middle|\,\mathcal D^{\mathrm{tr}}\right\}\ge1-\delta.
%    \label{eq-threshold-curve-coverage}
%  \end{equation}
%  At $c=0$, these definitions recover the original failure probability and its bounds.
%\end{corollary}
%\begin{proof}
%  On $\mathbb Q^\star\in\mathcal Q^{\mathrm{DW}}(\theta)$, the true joint law $\Pi^\star$ belongs to the same set $\mathcal P(\mathcal Q)$ for every $c$. Consequently,
%  \[ 
%    \underline p_{\hat f}(c) =\inf_{\Pi\in\mathcal P(\mathcal Q)} p_{\hat f}(\Pi;c) \le p_{\hat f}(\Pi^\star;c) =p_f(c) \le\sup_{\Pi\in\mathcal P(\mathcal Q)} p_{\hat f}(\Pi;c) =\overline p_{\hat f}(c), \qquad\forall c\in\mathbb R.
%  \]
%  Both bound curves are nondecreasing. Monotonicity and the CDF limits of $p_f$ imply that any violation is detected at a rational threshold or an atom of $g(\bm X)$, a countable set. The simultaneous event is therefore measurable. Proposition~\ref{proposition-01} gives the common residual-coverage event conditional probability at least $1-\delta$, proving \eqref{eq-threshold-curve-coverage}.
%\end{proof}
%
%Thresholds may therefore be selected after observing the test data. Complementing the endpoints gives simultaneous bounds on $\Pr \left\{g(\bm X)>c\right\}$. The functions $g$, $\hat g$, and the law $\mathbb P$ remain fixed; fragility curves conditioned on hazard intensity require coverage for the corresponding conditional residual laws.
%
%The coverage argument also applies to any residual ambiguity set containing $\mathbb Q^\star$ with conditional probability at least $1-\delta$. It remains to evaluate the DW extrema in \eqref{eq-upper-bound} and \eqref{eq-lower-bound} and their response-threshold extensions.
%
%
%\section{Tractable Reformulation}
%\label{sec-reformulation}
%
%Throughout this section, the training data, test residuals, and $\theta\ge0$ are fixed. Input and residual cell probabilities give finite primal and dual problems and closed-form bounds. The main results use $\Xi=\mathbb R$; Remark~\ref{rem-interval-support} treats known interval support.
%
%For each empirical probability level $q_j:=j/N^{\mathrm{te}}$, define the scalar acceptance interval
%\begin{equation}
%  [\alpha_j(\theta),\beta_j(\theta)]
%  :=\left\{p\in[0,1]:
%  N^{\mathrm{te}}\phi(q_j,p)-\psi(q_j,p)\le\theta\right\},
%  \qquad j=0,\ldots,N^{\mathrm{te}}.
%  \label{eq-dw-rank-acceptance}
%\end{equation}
%Let $\xi_{(1)}^{\mathrm{te}}<\cdots<\xi_{(K)}^{\mathrm{te}}$ be the distinct ordered test residuals, and set
%\begin{equation}
%  c_0:=0,\quad
%  c_k:=\sum_{j=1}^{N^{\mathrm{te}}}\mathbf1_{\{\xi_j^{\mathrm{te}}\le\xi_{(k)}^{\mathrm{te}}\}},\quad
%  l_k:=\alpha_{c_k}(\theta),\quad
%  u_k:=\beta_{c_{k-1}}(\theta),\qquad k=1,\ldots,K.
%  \label{eq-dw-counts-and-levels}
%\end{equation}
%The counts retain all multiplicities, so $c_K=N^{\mathrm{te}}$ even when test residuals coincide. Set $\xi_{(0)}^{\mathrm{te}}:=-\infty$ and $\xi_{(K+1)}^{\mathrm{te}}:=+\infty$. Separating the observed values from the intervening intervals gives $J:=2K+1$ residual cells:
%\begin{equation}
%  \Xi_{2k}:=\left\{\xi_{(k)}^{\mathrm{te}}\right\}
%    \quad\left(k=1,\ldots,K\right),
%  \qquad
%  \Xi_{2k+1}:=\left(\xi_{(k)}^{\mathrm{te}},\xi_{(k+1)}^{\mathrm{te}}\right)
%    \quad\left(k=0,\ldots,K\right).
%  \label{eq-residual-cells}
%\end{equation}
%For any candidate residual law $\mathbb Q\in\mathcal M_1(\Xi)$, the open rank probability measures mass strictly below an observed residual, while the closed rank probability also includes mass at that value. For $k=1,\ldots,K$, the partition gives
%\[
%  \mathbb Q\left\{\xi<\xi_{(k)}^{\mathrm{te}}\right\}
%    =\sum_{j<2k}\mathbb Q(\Xi_j),
%  \qquad
%  \mathbb Q\left\{\xi\le\xi_{(k)}^{\mathrm{te}}\right\}
%    =\sum_{j\le2k}\mathbb Q(\Xi_j).
%\]
%The two probabilities differ by the atom mass $\mathbb Q(\Xi_{2k})$.
%
%For $i=0,\ldots,K$, the observed residual levels classify inputs by the location of the failure threshold $-\hat g(\bm x)$:
%\begin{equation}
%  \mathcal X_i:=\left\{\bm x\in\mathcal X:
%    \xi_{(i)}^{\mathrm{te}}\le-\hat g(\bm x)<\xi_{(i+1)}^{\mathrm{te}}\right\}.
%  \label{eq-input-cells}
%\end{equation}
%The input cell probabilities give the surrogate-response threshold probabilities through
%\[
%  \mathbb P\left\{\hat g(\bm X)\le-\xi_{(k)}^{\mathrm{te}}\right\}
%    =\sum_{i=k}^{K}\mathbb P(\mathcal X_i),
%  \qquad k=1,\ldots,K.
%\]
%The half-open convention assigns $-\hat g(\bm x)=\xi_{(k)}^{\mathrm{te}}$ to $\mathcal X_k$. The tail sum therefore includes threshold equalities even when the surrogate response has atoms. 
%The two partitions satisfy
%\[
%  \mathcal X=\bigsqcup_{i=0}^{K}\mathcal X_i,
%  \qquad
%  \Xi=\bigsqcup_{j=1}^{J}\Xi_j.
%\]
%Figure~\ref{fig-residual-input-cells} illustrates the two partitions for $K=3$.
%\begin{figure}[!htbp]
%  \centering
%  \includegraphics[width=0.5\columnwidth]{fig_residual_input_cells.png}
%  \caption{Illustration of the residual and input-threshold partitions for $K=3$. The residual partition consists of $K$ singleton cells and $K+1$ open intervals, whereas the input partition is induced by half-open intervals of the failure threshold $-\hat g(\bm x)$. Here, $w_i=\mathbb P(\mathcal X_i)$ denotes the probability mass of input cell $\mathcal X_i$.}
%  \label{fig-residual-input-cells}
%\end{figure}
%
%\Needspace{5\baselineskip}
%The DW band can now be expressed through cell masses.
%
%\begin{proposition}[Finite representation of the DW ambiguity set]
%  \label{prop-finite-rank-constraints}
%  Each acceptance set in \eqref{eq-dw-rank-acceptance} is a nonempty closed interval containing $q_j$. The DW ambiguity set has the finite representation
%  \begin{equation}
%    \mathcal Q^{\mathrm{DW}}(\theta)
%    =\left\{\mathbb Q\in\mathcal M_1(\Xi):
%      \sum_{j<2k}\mathbb Q(\Xi_j)\le u_k,\quad
%      \sum_{j\le2k}\mathbb Q(\Xi_j)\ge l_k,\quad k=1,\ldots,K\right\}.
%    \label{eq-finite-rank-constraints}
%  \end{equation}
%  The sequences $l_k$ and $u_k$ are nondecreasing and satisfy
%  \[
%    0\le l_k,u_k\le1,\qquad l_{k-1}\le u_k,\qquad l_k\le u_{k+1},
%    \qquad k=1,\ldots,K,
%  \]
%  where $l_0:=0$ and $u_{K+1}:=1$.
%\end{proposition}
%\begin{proof}
%  For fixed $q$, the Bernoulli divergence increases as $p$ moves away from $q$, while $\psi(q,p)$ cannot increase as the interval between $q$ and $p$ expands. The corrected discrepancy therefore decreases toward $q$ from below and increases above $q$. It is continuous for $p\in(0,1)$ and equals $-\psi_0(q)\le0$ at $p=q$. The boundary conventions in Section~\ref{sec-ambiguity-set} give a nonempty closed acceptance interval also for $q\in\{0,1\}$. Scalar inversion of the DW constraint therefore gives the CDF band \citep[Supplement, Section~S.6]{L.J23}
%  \[
%    \alpha_{N^{\mathrm{te}}\widehat F^{\mathrm{te}}(t)}(\theta)
%      \le F_{\mathbb Q}(t)
%      \le\beta_{N^{\mathrm{te}}\widehat F^{\mathrm{te}}(t)}(\theta),
%    \qquad\forall t\in\mathbb R.
%  \]
%
%  For $k=1,\ldots,K-1$, the empirical CDF on $\left[\xi_{(k)}^{\mathrm{te}},\xi_{(k+1)}^{\mathrm{te}}\right)$ equals $c_k/N^{\mathrm{te}}$. Monotonicity makes the lower band constraint on this interval equivalent to $\mathbb Q\left\{\xi\le\xi_{(k)}^{\mathrm{te}}\right\}\ge\alpha_{c_k}(\theta)$. By continuity of probability from below, the upper band constraint is equivalent to $\mathbb Q\left\{\xi<\xi_{(k+1)}^{\mathrm{te}}\right\}\le\beta_{c_k}(\theta)$. The exterior intervals give $\mathbb Q\left\{\xi<\xi_{(1)}^{\mathrm{te}}\right\}\le\beta_0(\theta)$ and $\mathbb Q\left\{\xi\le\xi_{(K)}^{\mathrm{te}}\right\}\ge\alpha_{N^{\mathrm{te}}}(\theta)$; their other band constraints are automatic because $\alpha_0(\theta)=0$ and $\beta_{N^{\mathrm{te}}}(\theta)=1$. The partition identities above then yield \eqref{eq-finite-rank-constraints}.
%
%  With $p$ fixed, the corrected discrepancy decreases as $q$ approaches $p$ from either side. For empirical levels $q<q'$, an accepted value $p<q$ at level $q'$ remains accepted at level $q$; an accepted value $p>q'$ at level $q$ remains accepted at level $q'$. Since each acceptance interval contains its empirical level, both endpoints are nondecreasing. The inequalities $\alpha_{c_k}(\theta)\le c_k/N^{\mathrm{te}}\le\beta_{c_k}(\theta)$ give the stated compatibility conditions. Singleton intervals are allowed: if $q_j=1/2$ and $\theta=0$, the interval is $\{1/2\}$ because $\psi(1/2,p)=0$.
%\end{proof}
%
%Within each cell, the residual distribution remains free. The scalar endpoints can be computed by one-dimensional bisection.
%
%For a feasible joint law, define its matrix of joint cell probabilities by $\pi_{ij}:=\Pi(\mathcal X_i\times\Xi_j)$. The fixed input marginal and the residual restrictions imply $\bm\pi=(\pi_{ij})\in\mathcal C_\theta$, where
%\begin{equation}
%  \mathcal C_\theta:=\left\{\bm\pi\in\mathbb R_+^{(K+1)\times J}:
%  \begin{array}{ll}
%    \displaystyle\sum_{j=1}^{J}\pi_{ij}=\mathbb P(\mathcal X_i),&i=0,\ldots,K,\\[1mm]
%    \displaystyle\sum_{i=0}^{K}\sum_{j<2k}\pi_{ij}\le u_k,&k=1,\ldots,K,\\[1mm]
%    \displaystyle\sum_{i=0}^{K}\sum_{j\le2k}\pi_{ij}\ge l_k,&k=1,\ldots,K
%  \end{array}\right\}.
%  \label{eq-cell-mass-polytope}
%\end{equation}
%To bound the failure objective, compare the indicator on each product cell:
%\begin{equation}
%  \underline h_{ij}:=\mathbf1_{\{j\le2i\}}
%  \le\mathbf1_{\{\xi\le-\hat g(\bm x)\}}
%  \le\overline h_{ij}:=\mathbf1_{\{j\le2i+1\}},
%  \qquad \forall (\bm x,\xi)\in\mathcal X_i\times\Xi_j.
%  \label{eq-cell-costs}
%\end{equation}
%For inputs in $\mathcal X_i$, the failure status can vary only within $\Xi_{2i+1}$. The other cells are uniformly failing or safe, as illustrated in Figure~\ref{fig-joint-partition} for $K=3$.
%\begin{figure}[!htbp]
%  \centering
%  \includegraphics[width=0.5\columnwidth]{fig_joint_partition.png}
%  \caption{Joint partition of the input and residual spaces for $K=3$. Each product cell has probability mass $\pi_{ij}=\Pi(\mathcal X_i\times\Xi_j)$. The blue, yellow, and red cells correspond to $(\underline h_{ij},\overline h_{ij})=(0,0)$, $(0,1)$, and $(1,1)$, respectively.}
%  \label{fig-joint-partition}
%\end{figure}
%
%The upper problem remains a linear program over joint measures, with the full input marginal fixed:
%\begin{equation}
%  \begin{aligned}
%    \overline p_{\hat f}
%    =\sup_{\Pi\in\mathcal M_+(\mathcal X\times\Xi)}\quad&
%      \int_{\mathcal X\times\Xi}\mathbf1_{\{\xi\le-\hat g(\bm x)\}}\,\Pi(\mathrm d\bm x,\mathrm d\xi)\\
%    \text{s.t.}\quad&
%      \int_{\mathcal X\times\Xi}\mathbf1_{\mathcal U}(\bm x)\,\Pi(\mathrm d\bm x,\mathrm d\xi)
%        =\mathbb P(\bm X\in\mathcal U),&&\forall\mathcal U\in\mathcal B(\mathcal X),\\
%    &\int_{\mathcal X\times\Xi}\mathbf1_{\{\xi<\xi_{(k)}^{\mathrm{te}}\}}\,\Pi(\mathrm d\bm x,\mathrm d\xi)\le u_k,
%        &&k=1,\ldots,K,\\
%    &\int_{\mathcal X\times\Xi}\mathbf1_{\{\xi\le\xi_{(k)}^{\mathrm{te}}\}}\,\Pi(\mathrm d\bm x,\mathrm d\xi)\ge l_k,
%        &&k=1,\ldots,K.
%  \end{aligned}
%  \label{eq-measure-upper}
%\end{equation}
%The lower problem minimizes the same integral under the same constraints. The marginal equalities fix both the input law and total mass one. The product of $\mathbb P$ and the empirical residual law is feasible for every $\theta\ge0$.
%
%\Needspace{5\baselineskip}
%Lagrangian duality provides finite-dimensional certificates for these bounds.
%
%\begin{lemma}[Finite-dimensional Lagrangian duals and weak duality]
%  \label{lem-weak-duality}
%  Write $\bm u=(u_1,\ldots,u_K)^\top$ and $\bm l=(l_1,\ldots,l_K)^\top$. The nonnegative multipliers $\bm\lambda,\bm\gamma\in\mathbb R_+^K$ correspond to the upper and lower rank constraints. The input marginal equality has an unrestricted multiplier $\eta\in\mathcal F_{\mathrm b}(\mathcal X)$, the space of bounded Borel measurable real-valued functions. Define the residual penalty
%  \begin{equation}
%    \zeta_{\bm\lambda,\bm\gamma}(\xi)
%    :=\sum_{k=1}^{K}\lambda_k\mathbf1_{\{\xi<\xi_{(k)}^{\mathrm{te}}\}}
%      -\sum_{k=1}^{K}\gamma_k\mathbf1_{\{\xi\le\xi_{(k)}^{\mathrm{te}}\}},
%    \label{eq-dual-step-function}
%  \end{equation}
%  whose value on $\Xi_j$ is
%  \begin{equation}
%    \zeta_j:=\sum_{k=1}^{K}\lambda_k\mathbf1_{\{j<2k\}}
%      -\sum_{k=1}^{K}\gamma_k\mathbf1_{\{j\le2k\}},
%    \qquad j=1,\ldots,J.
%    \label{eq-dual-cell-values}
%  \end{equation}
%  Eliminating the input marginal multiplier gives the upper and lower Lagrangian dual problems
%  \begin{align}
%    \overline d_\theta
%      &:=\inf_{\bm\lambda,\bm\gamma\ge\bm0}\left\{
%        \bm\lambda^\top\bm u-\bm\gamma^\top\bm l
%        +\sum_{i=0}^{K}\mathbb P(\mathcal X_i)\max_{1\le j\le J}
%          \left\{\overline h_{ij}-\zeta_j\right\}\right\},
%      \label{eq-upper-measure-dual}\\
%    \underline d_\theta
%      &:=\sup_{\bm\lambda,\bm\gamma\ge\bm0}\left\{
%        -\bm\lambda^\top\bm u+\bm\gamma^\top\bm l
%        +\sum_{i=0}^{K}\mathbb P(\mathcal X_i)\min_{1\le j\le J}
%          \left\{\underline h_{ij}+\zeta_j\right\}\right\}.
%      \label{eq-lower-measure-dual}
%  \end{align}
%  Their optimal values satisfy
%  \begin{equation}
%    \underline d_\theta\le\underline p_{\hat f}
%      \le\overline p_{\hat f}\le\overline d_\theta.
%    \label{eq-weak-duality}
%  \end{equation}
%\end{lemma}
%\begin{proof}
%  The input marginal equality is equivalent to
%  \[
%    \int_{\mathcal X\times\Xi}\eta(\bm x)\,\Pi(\mathrm d\bm x,\mathrm d\xi)
%      =\mathbb E_{\mathbb P}\left[\eta(\bm X)\right],
%    \qquad\forall\eta\in\mathcal F_{\mathrm b}(\mathcal X).
%  \]
%  Relaxing this equality and the rank inequalities gives the upper Lagrangian \citep[Section~5.1]{B.V04}
%  \[
%    \begin{aligned}
%      \mathcal L^+(\Pi,\eta,\bm\lambda,\bm\gamma)
%      ={}&\int_{\mathcal X\times\Xi}\mathbf1_{\{\xi\le-\hat g(\bm x)\}}\,\Pi(\mathrm d\bm x,\mathrm d\xi)
%        +\mathbb E_{\mathbb P}\left[\eta(\bm X)\right]
%        -\int_{\mathcal X\times\Xi}\eta(\bm x)\,\Pi(\mathrm d\bm x,\mathrm d\xi)\\
%      &+\sum_{k=1}^{K}\lambda_k\left(u_k-\int_{\mathcal X\times\Xi}
%          \mathbf1_{\{\xi<\xi_{(k)}^{\mathrm{te}}\}}\,\Pi(\mathrm d\bm x,\mathrm d\xi)\right)
%        +\sum_{k=1}^{K}\gamma_k\left(\int_{\mathcal X\times\Xi}
%          \mathbf1_{\{\xi\le\xi_{(k)}^{\mathrm{te}}\}}\,\Pi(\mathrm d\bm x,\mathrm d\xi)-l_k\right)\\
%      ={}&\mathbb E_{\mathbb P}\left[\eta(\bm X)\right]
%        +\bm\lambda^\top\bm u-\bm\gamma^\top\bm l
%        +\int_{\mathcal X\times\Xi}\left[
%          \mathbf1_{\{\xi\le-\hat g(\bm x)\}}-\eta(\bm x)-\zeta_{\bm\lambda,\bm\gamma}(\xi)
%        \right]\,\Pi(\mathrm d\bm x,\mathrm d\xi).
%    \end{aligned}
%  \]
%  Its supremum over finite nonnegative measures is finite exactly when the integrand is nonpositive at every $(\bm x,\xi)$. A positive value at any point permits an arbitrarily large point mass there; otherwise the supremum integral is zero, attained by the zero measure. The functional upper dual therefore minimizes $\mathbb E_{\mathbb P}[\eta(\bm X)]+\bm\lambda^\top\bm u-\bm\gamma^\top\bm l$ subject to $\eta(\bm x)\ge\mathbf1_{\{\xi\le-\hat g(\bm x)\}}-\zeta_{\bm\lambda,\bm\gamma}(\xi)$ at every pair. For the lower problem, reversing the rank penalties gives
%  \[
%    \mathcal L^-
%      =\mathbb E_{\mathbb P}\left[\eta(\bm X)\right]
%        -\bm\lambda^\top\bm u+\bm\gamma^\top\bm l
%        +\int_{\mathcal X\times\Xi}\left[
%          \mathbf1_{\{\xi\le-\hat g(\bm x)\}}-\eta(\bm x)+\zeta_{\bm\lambda,\bm\gamma}(\xi)
%        \right]\,\Pi(\mathrm d\bm x,\mathrm d\xi).
%  \]
%  Its infimum over finite nonnegative measures is finite exactly when the integrand is nonnegative everywhere, in which case the infimum integral is zero. The lower functional dual maximizes $\mathbb E_{\mathbb P}[\eta(\bm X)]-\bm\lambda^\top\bm u+\bm\gamma^\top\bm l$ subject to $\eta(\bm x)\le\mathbf1_{\{\xi\le-\hat g(\bm x)\}}+\zeta_{\bm\lambda,\bm\gamma}(\xi)$ at every pair.
%
%  For fixed rank multipliers, the best marginal functions are the pointwise envelopes. For $\bm x\in\mathcal X_i$, they satisfy the identities
%  \begin{equation}
%    \begin{aligned}
%      \eta^+_{\bm\lambda,\bm\gamma}(\bm x)
%        &:=\sup_{\xi\in\Xi}\left\{\mathbf1_{\{\xi\le-\hat g(\bm x)\}}
%          -\zeta_{\bm\lambda,\bm\gamma}(\xi)\right\}
%          =\max_{1\le j\le J}\left\{\overline h_{ij}-\zeta_j\right\},\\
%      \eta^-_{\bm\lambda,\bm\gamma}(\bm x)
%        &:=\inf_{\xi\in\Xi}\left\{\mathbf1_{\{\xi\le-\hat g(\bm x)\}}
%          +\zeta_{\bm\lambda,\bm\gamma}(\xi)\right\}
%          =\min_{1\le j\le J}\left\{\underline h_{ij}+\zeta_j\right\}.
%    \end{aligned}
%    \label{eq-dual-envelope-identities}
%  \end{equation}
%  To verify them, note that outside $\Xi_{2i+1}$ the two cell coefficients equal the actual failure indicator. In the overlapping interval, the lower coefficient zero is attained by $\xi=\left(\xi_{(i+1)}^{\mathrm{te}}-\hat g(\bm x)\right)/2$ if $i<K$, and by $\xi=-\hat g(\bm x)+1$ if $i=K$. The upper coefficient one is attained by $\xi=-\hat g(\bm x)$ whenever $-\hat g(\bm x)>\xi_{(i)}^{\mathrm{te}}$. The only remaining case is $-\hat g(\bm x)=\xi_{(i)}^{\mathrm{te}}$ with $i\ge1$. Then the overlapping interval is safe, but
%  \[
%    \zeta_{2i}=\zeta_{2i+1}-\gamma_i\le\zeta_{2i+1},
%    \qquad
%    1-\zeta_{2i}\ge1-\zeta_{2i+1}.
%  \]
%  Choosing the left endpoint therefore attains an upper-envelope value at least as large as the candidate from the overlapping interval. Together with \eqref{eq-cell-costs}, this proves both identities, including threshold equalities.
%
%  The envelopes are bounded and constant on the Borel cells $\mathcal X_i$, so they belong to $\mathcal F_{\mathrm b}(\mathcal X)$. They respectively minimize and maximize the expectation over feasible marginal functions. Substituting them into the functional dual objectives gives \eqref{eq-upper-measure-dual} and \eqref{eq-lower-measure-dual}. Introducing epigraph variables $z_i$ for the upper maxima and hypograph variables $z_i$ for the lower minima gives the equivalent LPs
%  \begin{equation}
%    \begin{aligned}
%      \overline d_\theta
%      =\inf_{\bm\lambda,\bm\gamma\ge\bm0,\ \bm z\in\mathbb R^{K+1}}\quad&
%        \bm\lambda^\top\bm u-\bm\gamma^\top\bm l+\sum_{i=0}^{K}\mathbb P(\mathcal X_i) z_i\\
%      \text{s.t.}\quad&z_i+\zeta_j\ge\overline h_{ij},
%        \qquad i=0,\ldots,K,\quad j=1,\ldots,J,
%    \end{aligned}
%    \label{eq-finite-upper-dual}
%  \end{equation}
%  and
%  \begin{equation}
%    \begin{aligned}
%      \underline d_\theta
%      =\sup_{\bm\lambda,\bm\gamma\ge\bm0,\ \bm z\in\mathbb R^{K+1}}\quad&
%        -\bm\lambda^\top\bm u+\bm\gamma^\top\bm l+\sum_{i=0}^{K}\mathbb P(\mathcal X_i) z_i\\
%      \text{s.t.}\quad&z_i-\zeta_j\le\underline h_{ij},
%        \qquad i=0,\ldots,K,\quad j=1,\ldots,J.
%    \end{aligned}
%    \label{eq-finite-lower-dual}
%  \end{equation}
%  Each $z_i$ can be chosen as its corresponding finite envelope, including on input cells with $\mathbb P(\mathcal X_i)=0$. Thus introducing these variables preserves both optimal values.
%
%  For a feasible joint law, the signs of the rank penalties give $\mathcal L^-\le p_{\hat f}(\Pi)\le\mathcal L^+$. Feasible lower dual objectives therefore bound the primal objectives from below, and feasible upper dual objectives bound them from above. Taking the respective optima over the common nonempty primal feasible set proves \eqref{eq-weak-duality}.
%\end{proof}
%
%The reduced duals use $2K$ nonnegative multipliers. Their LP forms \eqref{eq-finite-upper-dual} and \eqref{eq-finite-lower-dual} introduce $K+1$ free variables $z_i$ and $(K+1)(2K+1)$ pairwise inequalities.
%
%\Needspace{5\baselineskip}
%It remains to realize the cellwise objectives by feasible joint laws.
%
%\begin{theorem}[Exact finite reformulations and strong duality]
%  \label{thm-finite-dual-reformulations}
%  For $\mathcal Q=\mathcal Q^{\mathrm{DW}}(\theta)$ and every $\theta\ge0$, the failure probability bounds admit the exact representations
%  \begin{equation}
%    \overline p_{\hat f}
%      =\max_{\bm\pi\in\mathcal C_\theta}\sum_{i=0}^{K}\sum_{j=1}^{J}\overline h_{ij}\pi_{ij},
%    \qquad
%    \underline p_{\hat f}
%      =\min_{\bm\pi\in\mathcal C_\theta}\sum_{i=0}^{K}\sum_{j=1}^{J}\underline h_{ij}\pi_{ij}.
%    \label{eq-finite-primal-values}
%  \end{equation}
%  Their Lagrangian dual values in Lemma~\ref{lem-weak-duality} satisfy
%  \begin{equation}
%    \overline p_{\hat f}=\overline d_\theta,
%    \qquad
%    \underline p_{\hat f}=\underline d_\theta.
%    \label{eq-strong-duality}
%  \end{equation}
%  The finite primal and dual optima, the measure primal optima, and the functional dual optima are attained.
%\end{theorem}
%\begin{proof}
%  Any feasible joint law induces a matrix in $\mathcal C_\theta$. Integrating \eqref{eq-cell-costs} therefore bounds the upper measure problem by the finite maximization from above and the lower measure problem by the finite minimization from below.
%
%  For the reverse inequalities, fix $\bm\pi\in\mathcal C_\theta$. Draw $\bm x$ from $\mathbb P$. On $\mathcal X_i$ with $\mathbb P(\mathcal X_i)>0$, select a provisional residual cell $j$ with probability $\pi_{ij}/\mathbb P(\mathcal X_i)$; the row equality makes these probabilities sum to one. For $j\ne2i+1$, choose a fixed representative of $\Xi_j$, where the actual failure indicator equals both coefficients. In the overlapping cell, use the following assignments to define joint laws $\Pi^+$ and $\Pi^-$.
%
%  For the upper problem, assign $\xi=-\hat g(\bm x)$. This causes failure and remains in $\Xi_{2i+1}$ unless $-\hat g(\bm x)=\xi_{(i)}^{\mathrm{te}}$. At equality, which requires $i\ge1$, the assignment transfers the provisional mass to the left endpoint. For any $\widetilde\xi\in\left(\xi_{(i)}^{\mathrm{te}},\xi_{(i+1)}^{\mathrm{te}}\right)$ and every observed threshold,
%  \[
%    \mathbf1_{\{\xi_{(i)}^{\mathrm{te}}<\xi_{(k)}^{\mathrm{te}}\}}
%      =\mathbf1_{\{\widetilde\xi<\xi_{(k)}^{\mathrm{te}}\}},
%    \qquad
%    \mathbf1_{\{\xi_{(i)}^{\mathrm{te}}\le\xi_{(k)}^{\mathrm{te}}\}}
%      \ge\mathbf1_{\{\widetilde\xi\le\xi_{(k)}^{\mathrm{te}}\}}.
%  \]
%  Thus every open upper probability is unchanged and every closed lower probability can only increase. The upper assignment preserves all rank inequalities and attains the coefficient $\overline h_{i,2i+1}=1$.
%
%  For the lower problem, assign $\xi=\left(\xi_{(i+1)}^{\mathrm{te}}-\hat g(\bm x)\right)/2$ if $i<K$, and $\xi=-\hat g(\bm x)+1$ if $i=K$. These values remain in $\Xi_{2i+1}$ and strictly exceed the failure threshold, including at threshold equality. Hence the lower assignment retains the cell matrix and attains $\underline h_{i,2i+1}=0$.
%
%  If $\mathbb P(\mathcal X_i)=0$, the corresponding row is zero; set $\xi=\xi_{(1)}^{\mathrm{te}}$ on that input cell. All assignments are Borel measurable, and their finite mixtures define probability kernels. Combining them with $\mathbb P$ gives joint laws with the full prescribed input marginal, including randomization at input atoms. Both laws are feasible and satisfy
%  \[
%    p_{\hat f}(\Pi^+)=\sum_{i=0}^{K}\sum_{j=1}^{J}\overline h_{ij}\pi_{ij},
%    \qquad
%    p_{\hat f}(\Pi^-)=\sum_{i=0}^{K}\sum_{j=1}^{J}\underline h_{ij}\pi_{ij}.
%  \]
%  The lower law retains every provisional cell mass. The upper law may change the cell matrix at threshold equalities, but it preserves feasibility and the assigned upper objective. These identities establish the reverse inequalities required for \eqref{eq-finite-primal-values}.
%
%  For every $\theta\ge0$, the empirical product law induces the feasible matrix
%  \[
%    \pi_{i,2k}^{\mathrm{emp}}:=\mathbb P(\mathcal X_i)\frac{c_k-c_{k-1}}{N^{\mathrm{te}}},
%    \qquad i=0,\ldots,K,\quad k=1,\ldots,K,
%  \]
%  with all odd columns zero. Its open and closed rank probabilities are $c_{k-1}/N^{\mathrm{te}}$ and $c_k/N^{\mathrm{te}}$, respectively, so feasibility includes $\theta=0$. The nonnegative entries of a matrix in $\mathcal C_\theta$ sum to one, and the constraints are closed. Thus $\mathcal C_\theta$ is nonempty and compact. The finite objectives lie in $[0,1]$ and attain their optima, which the joint-law constructions realize in the measure problems.
%
%  The finite primal LPs in \eqref{eq-finite-primal-values} have row multipliers $z_i$ and rank multipliers $\lambda_k,\gamma_k$. Their dual constraints are precisely $z_i+\zeta_j\ge\overline h_{ij}$ and $z_i-\zeta_j\le\underline h_{ij}$. Finite LP strong duality therefore gives \eqref{eq-strong-duality} and dual attainment \citep[Section~5.2.3]{B.V04}. By \eqref{eq-dual-envelope-identities}, optimal rank multipliers also yield optimal bounded marginal functions for the original functional duals. Feasibility and finite optimal values suffice; strict rank inequalities are unnecessary, including at $\theta=0$.
%\end{proof}
%
%A finite optimizer specifies the cell masses; the proof supplies their within-cell realization.
%
%\Needspace{5\baselineskip}
%The ordered cells allow the extrema to be evaluated explicitly.
%
%\begin{corollary}[Closed-form bounds and optimal certificates]
%  \label{thm-closed-form-bounds}
%  Under the conditions of Theorem~\ref{thm-finite-dual-reformulations},
%  \begin{align}
%    \overline p_{\hat f}
%      &=\min\left\{1,\min_{k=1,\ldots,K}\left[u_k+\sum_{i=k}^{K}\mathbb P(\mathcal X_i)\right]\right\},
%      \label{eq-dw-closed-upper}\\
%    \underline p_{\hat f}
%      &=\max\left\{0,\max_{k=1,\ldots,K}\left[l_k+\sum_{i=k}^{K}\mathbb P(\mathcal X_i)-1\right]\right\}.
%      \label{eq-dw-closed-lower}
%  \end{align}
%  Each dual optimum admits a certificate with at most one nonzero rank multiplier, of value one.
%\end{corollary}
%\begin{proof}
%  For a candidate residual law $\mathbb Q\in\mathcal M_1(\Xi)$, define the finite transportation set
%  \[
%    \mathcal T(\mathbb Q):=\left\{\bm\pi\in\mathbb R_+^{(K+1)\times J}:
%    \begin{array}{ll}
%      \displaystyle\sum_{j=1}^{J}\pi_{ij}=\mathbb P(\mathcal X_i),&i=0,\ldots,K,\\[1mm]
%      \displaystyle\sum_{i=0}^{K}\pi_{ij}=\mathbb Q(\Xi_j),&j=1,\ldots,J
%    \end{array}\right\}.
%  \]
%  It contains the product matrix $\pi_{ij}=\mathbb P(\mathcal X_i)\mathbb Q(\Xi_j)$. Every residual cell is nonempty, so any nonnegative column allocation with total mass one is realized by a residual law. Proposition~\ref{prop-finite-rank-constraints} consequently gives
%  \[
%    \mathcal C_\theta
%      =\bigcup_{\mathbb Q\in\mathcal Q^{\mathrm{DW}}(\theta)}\mathcal T(\mathbb Q).
%  \]
%  The parameter $\mathbb Q$ fixes only residual partition masses in this finite matrix problem. The joint-law construction in Theorem~\ref{thm-finite-dual-reformulations} preserves feasibility and the assigned objective while allowing the residual marginal to change, including endpoint transfers in the upper problem.
%
%  For fixed cell masses $\mathbb Q(\Xi_j)$, the upper problem maximizes mass on pairs $j\le2i+1$; the lower problem maximizes safe mass on pairs $j>2i$ and subtracts it from one. Both are maximum-flow problems with source-to-row capacities $\mathbb P(\mathcal X_i)$, column-to-sink capacities $\mathbb Q(\Xi_j)$, and capacity two on each permitted row-to-column edge. A transportation matrix restricts to a feasible flow. Conversely, any flow can be completed to a transportation matrix because the remaining row and column masses are nonnegative with equal totals. Thus the maximum flow equals the maximum permitted mass.
%
%  A cut of capacity one is available, so a minimum cut crosses no edge of capacity two. Its source-side columns can be taken to be exactly the neighbors of its source-side rows. Since the neighborhoods are nested, adding rows whose neighbors are already included cannot increase the capacity. Apart from empty and full row sets, upper-problem cuts reduce to rows $i<k$ and columns $j<2k$, while safe-mass cuts reduce to rows $i\ge k$ and columns $j>2k$. For $k=1,\ldots,K$, their capacities are respectively
%  \[
%    \sum_{i=k}^{K}\mathbb P(\mathcal X_i)+\sum_{j<2k}\mathbb Q(\Xi_j),
%    \qquad
%    \sum_{i=0}^{k-1}\mathbb P(\mathcal X_i)+\sum_{j>2k}\mathbb Q(\Xi_j).
%  \]
%  Empty and full row sets give capacity one. The max-flow/min-cut theorem \citep{F.F56} therefore yields
%  \begin{align*}
%    \max_{\bm\pi\in\mathcal T(\mathbb Q)}\sum_{i,j}\overline h_{ij}\pi_{ij}
%      &=\min\left\{1,\min_k\left[\sum_{i=k}^{K}\mathbb P(\mathcal X_i)
%        +\sum_{j<2k}\mathbb Q(\Xi_j)\right]\right\},\\
%    \min_{\bm\pi\in\mathcal T(\mathbb Q)}\sum_{i,j}\underline h_{ij}\pi_{ij}
%      &=\max\left\{0,\max_k\left[\sum_{i=k}^{K}\mathbb P(\mathcal X_i)
%        +\sum_{j\le2k}\mathbb Q(\Xi_j)-1\right]\right\}.
%  \end{align*}
%
%  These values are nondecreasing in the respective open and closed cumulative residual probabilities. The rank constraints bound the former from above by $u_k$ and the latter from below by $l_k$. To attain all relevant limits simultaneously, choose laws $\mathbb Q^+$ and $\mathbb Q^-$ with cell probabilities
%  \begin{align*}
%    &\mathbb Q^+(\Xi_1):=u_1,\quad
%      \mathbb Q^+(\Xi_{2k}):=u_{k+1}-u_k,\quad
%      \mathbb Q^+(\Xi_{2k+1}):=0,\\
%    &\mathbb Q^-(\Xi_{2k-1}):=0,\quad
%      \mathbb Q^-(\Xi_{2k}):=l_k-l_{k-1},\quad
%      \mathbb Q^-(\Xi_J):=1-l_K,
%      \qquad k=1,\ldots,K.
%  \end{align*}
%  By Proposition~\ref{prop-finite-rank-constraints}, the prescribed masses are nonnegative and sum to one. Choosing a representative in each residual cell realizes both laws. Their cumulative probabilities satisfy
%  \begin{align*}
%    &\sum_{j<2k}\mathbb Q^+(\Xi_j)=u_k,
%      \qquad \sum_{j\le2k}\mathbb Q^+(\Xi_j)=u_{k+1}\ge l_k,\\
%    &\sum_{j\le2k}\mathbb Q^-(\Xi_j)=l_k,
%      \qquad \sum_{j<2k}\mathbb Q^-(\Xi_j)=l_{k-1}\le u_k.
%  \end{align*}
%  Both laws therefore belong to $\mathcal Q^{\mathrm{DW}}(\theta)$. The union representation of $\mathcal C_\theta$ and \eqref{eq-finite-primal-values} show that these allocations attain the upper and lower bounds in \eqref{eq-dw-closed-upper} and \eqref{eq-dw-closed-lower}. Zero masses and coincident cumulative limits are allowed, including at $\theta=0$.
%
%  Finally, let $\bm e_k$ be the $k$th unit vector. Certificates for the upper and lower LPs \eqref{eq-finite-upper-dual} and \eqref{eq-finite-lower-dual} are respectively
%  \begin{align*}
%    &\bm\lambda=\bm e_k,\qquad\bm\gamma=\bm0,\qquad
%      z_i=\mathbf1_{\{i\ge k\}},\\
%    &\bm\lambda=\bm0,\qquad\bm\gamma=\bm e_k,\qquad
%      z_i=\mathbf1_{\{i\ge k\}}-1.
%  \end{align*}
%  Their feasibility follows from
%  \[
%    \mathbf1_{\{i\ge k\}}+\mathbf1_{\{j<2k\}}\ge\overline h_{ij},
%    \qquad
%    \mathbf1_{\{i\ge k\}}-1+\mathbf1_{\{j\le2k\}}\le\underline h_{ij}.
%  \]
%  Their objectives are $u_k+\sum_{i=k}^{K}\mathbb P(\mathcal X_i)$ and $l_k+\sum_{i=k}^{K}\mathbb P(\mathcal X_i)-1$, the inner candidates in the closed-form expressions. Zero rank multipliers with $z_i=1$ or $z_i=0$ supply the outer candidates. Selecting the best candidate gives an optimal full dual solution, whose functional marginal multiplier follows from \eqref{eq-dual-envelope-identities}. The inclusion of $i=k$ in the tail sum retains inputs whose failure threshold equals the observed residual.
%\end{proof}
%
%The upper bound combines an allowed residual tail mass with the probability that the input threshold reaches the corresponding observed residual. The lower bound measures unavoidable overlap.
%
%For the response-threshold curves in Corollary~\ref{cor-simultaneous-threshold-bounds}, apply the same construction to $g-c$ and $\hat g-c$. Every residual, rank count, and endpoint $l_k,u_k$ is unchanged. The input partition now depends on $c$, and its cumulative probability at rank $k$ is $\mathbb P\left\{\hat g(\bm X)\le c-\xi_{(k)}^{\mathrm{te}}\right\}$. Thus Corollary~\ref{thm-closed-form-bounds} gives
%\begin{align}
%  \overline p_{\hat f}(c)
%    &=\min\left\{1,\min_{k=1,\ldots,K}\left[u_k
%      +\mathbb P\left\{\hat g(\bm X)\le c-\xi_{(k)}^{\mathrm{te}}\right\}\right]\right\},
%    \label{eq-threshold-closed-upper}\\
%  \underline p_{\hat f}(c)
%    &=\max\left\{0,\max_{k=1,\ldots,K}\left[l_k
%      +\mathbb P\left\{\hat g(\bm X)\le c-\xi_{(k)}^{\mathrm{te}}\right\}-1\right]\right\}.
%    \label{eq-threshold-closed-lower}
%\end{align}
%These identities hold for every $c\in\mathbb R$ and $\theta\ge0$, including $\theta=0$. When $\theta\ge\theta_{N^{\mathrm{te}}}(\delta)$, their coverage is simultaneous by \eqref{eq-threshold-curve-coverage}. An extremizing coupling may depend on $c$; the common feasible true joint law supplies the simultaneous guarantee.
%
%For $p_{\mathrm{tar}}\in(0,1)$, the upper curve satisfies $\overline p_{\hat f}(c)\le p_{\mathrm{tar}}$ exactly when some rank satisfies
%\[
%  u_k+\mathbb P\left\{\hat g(\bm X)\le c-\xi_{(k)}^{\mathrm{te}}\right\}\le p_{\mathrm{tar}}.
%\]
%Both terms must therefore be small at the same residual level.
%
%With a known residual law, the coupling problem is the classical fixed-marginal sum problem studied by \citet{M82} and \citet{F.N.S87}. Here the residual law also varies within a calibrated DW band, and the open and closed rank constraints retain atoms.
%
%Dependence uncertainty can remain even with a known residual marginal. For equally likely surrogate responses $0$ and $2$ and equally likely residuals $-1$ and $1$, failure occurs exactly when the response is $0$ and the residual is $-1$. Its probability ranges from $0$ to $1/2$ over couplings, whereas independence gives $1/4$. Further marginal calibration alone does not resolve this dependence uncertainty.
%
%\Needspace{7\baselineskip}
%Verified support gives information beyond the empirical residuals.
%
%\begin{remark}[Known interval support]
%  \label{rem-interval-support}
%  Suppose $\Xi=[a,b]$ with $-\infty<a\le b<\infty$ and the observed residuals lie in this interval. Then
%  \begin{align*}
%    \overline p_{\hat f}
%      &=\min\left\{\mathbb P\left\{\hat g(\bm X)\le-a\right\},
%        \min_{k=1,\ldots,K}\left[u_k+\mathbb P\left\{\hat g(\bm X)\le-\xi_{(k)}^{\mathrm{te}}\right\}\right]\right\},\\
%    \underline p_{\hat f}
%      &=\max\left\{\mathbb P\left\{\hat g(\bm X)\le-b\right\},
%        \max_{k=1,\ldots,K}\left[l_k+\mathbb P\left\{\hat g(\bm X)\le-\xi_{(k)}^{\mathrm{te}}\right\}-1\right]\right\}.
%  \end{align*}
%  Represent the support restriction on $\mathbb R$ by $\mathbb Q\{\xi<a\}\le0$ and $\mathbb Q\{\xi\le b\}\ge1$, and add $a,b$ to the ordered threshold list. At a newly added $a$, use upper and lower bounds zero; at a newly added $b$, use upper and lower bounds one. At coincident observed thresholds, tighten only the upper bound at $a$ and the lower bound at $b$. The enlarged sequences satisfy the compatibility conditions in Proposition~\ref{prop-finite-rank-constraints}, so the same finite-aggregation and ordered-cut arguments apply. The effective additional upper candidate is $\mathbb P\left\{\hat g(\bm X)\le-a\right\}$, and the effective additional lower candidate is $\mathbb P\left\{\hat g(\bm X)\le-b\right\}$; the other added candidates are redundant. If $a=b$, both bounds equal $\mathbb P\left\{\hat g(\bm X)\le-a\right\}$.
%
%  For the partition on $\Xi=[a,b]$, intersect the residual cells from the augmented threshold list with $\Xi$ and omit empty cells. Including the support endpoints retains their atoms and the corresponding changes in failure status. For a response threshold $c$, shift each surrogate-response threshold in the displayed event probabilities by $c$.
%\end{remark}
%
%Inputs with $\hat g(\bm X)\le-b$ must fail; those with $\hat g(\bm X)>-a$ are always safe. Thus $\mathbb P\left\{\hat g(\bm X)\le-a\right\}=0$ rules out failure.
%
%After zero failures in $N^{\mathrm{te}}$ independent tests, the one-sided exact binomial upper limit is $1-\delta^{1/N^{\mathrm{te}}}\approx\log(1/\delta)/N^{\mathrm{te}}$ \citep{C.P34}. At $c=0$, the simple DW case $K=1$ and $\mathbb P(\mathcal X_1)=0$ gives an upper bound $\beta_0(\theta)$ determined by the simultaneous-band calibration.
%
%Evaluation requires scalar DW inversion, sorting and counting the test residuals, and computing the $K$ input event probabilities $\mathbb P\left\{\hat g(\bm X)\le c-\xi_{(k)}^{\mathrm{te}}\right\}$ at each response threshold $c$. The endpoint table depends only on $N^{\mathrm{te}}$ and $\theta$ and can be precomputed. Given this table and the input event probabilities, the two scans cost $O(K)$ at each response threshold.
%
%For $\theta\ge\theta_{N^{\mathrm{te}}}(\delta)$, the exact expressions retain the coverage in Theorem~\ref{theorem-01} and Corollary~\ref{cor-simultaneous-threshold-bounds}. Independent auxiliary simulation of $\bm X\sim\mathbb P$ can supply simultaneous upper and lower confidence bounds on the input event probabilities. Substituting these bounds in the respective upper and lower expressions preserves conservative failure probability bounds. For one fixed $c$, simultaneous bounds at its $K$ evaluation thresholds suffice. To retain coverage of the entire response-threshold curve, the auxiliary bounds must hold for the surrogate-response CDF at every real threshold. They must hold conditionally on the test data, with their error probability included in the overall confidence budget.
%
%
% END PRESERVED BODY BEFORE NAMED RESULT TRANSITIONS

% BEGIN PRESERVED BODY BEFORE IMPLEMENTATION PROCEDURE (2026-11-26)
%\section{Problem Description}
%\label{sec-problem}
%
%Let $\bm X\sim\mathbb P$ be a random input taking values in $\mathcal X\subseteq\mathbb R^n$, and let $g:\mathcal X\to\mathbb R$ be a measurable limit-state function. The failure probability is defined as
%\[
%  p_f:=\Pr\{g(\bm X)\le0\}=\mathbb E_{\bm X\sim\mathbb P}\left[\mathbf 1_{\{g(\bm X)\le0\}}\right]=\int_{\mathcal X}\mathbf 1_{\{g(\bm x)\le0\}}\,\mathbb P(\mathrm d\bm x).
%\]
%
%Evaluating $g$ can require an expensive high-fidelity simulation. We use the training dataset
%\[
%  \mathcal D^{\mathrm{tr}}:=\left\{\left(\bm x_i^{\mathrm{tr}},g(\bm x_i^{\mathrm{tr}})\right)\right\}_{i=1}^{N^{\mathrm{tr}}}
%\]
%to fit a computationally efficient, measurable surrogate $\hat g:=\hat g_{\mathcal D^{\mathrm{tr}}}$. Model fitting and selection use only the training data and are completed before the test data are used. We condition on $\mathcal D^{\mathrm{tr}}$, treat $\hat g$ as fixed, and suppress this dependence in the notation.
%
%Define the signed residual function
%\[
%  \xi(\bm x):=g(\bm x)-\hat g(\bm x),\qquad \forall \bm x\in\mathcal X.
%\]
%For a new input $\bm X\sim\mathbb P$ independent of the training data, let $\mathbb Q^\star$ be the distribution of $\xi^\star:=\xi(\bm X)$. Let $\Xi\subseteq\mathbb R$ be a known Borel set with $\mathbb Q^\star(\Xi)=1$; the default choice is $\Xi=\mathbb R$. For any Borel set $\mathcal S$ in a Euclidean space, write $\mathcal B(\mathcal S)$ for its Borel sigma-algebra. Let $\mathcal M_+(\mathcal S)$ denote the finite nonnegative Borel measures concentrated on $\mathcal S$, and let $\mathcal M_1(\mathcal S)$ denote the probability measures.
%
%For a candidate residual law $\mathbb Q\in\mathcal M_1(\Xi)$, let $\Pi$ be a joint distribution on $\mathcal X\times\Xi$ with marginals $\Pi_{\bm X}=\mathbb P$ and $\Pi_\xi=\mathbb Q$. Its surrogate-based failure probability is
%\begin{equation}
%  p_{\hat f}(\Pi):=\Pr\{\hat g(\bm X)+\xi\le0\}=\int_{\mathcal X\times\Xi}\mathbf 1_{\{\hat g(\bm x)+\xi\le0\}}\,\Pi(\mathrm d\bm x,\mathrm d\xi).
%  \label{eq-surrogate-failure-probability}
%\end{equation}
%For fixed $\hat g$, this functional is linear in $\Pi$. In particular, the true joint distribution $\Pi^\star$ satisfies $p_{\hat f}(\Pi^\star)=p_f$.
%
%The residual marginal and its dependence on the input jointly determine the failure probability. We represent residual-marginal uncertainty by an ambiguity set $\mathcal Q$ calibrated to contain $\mathbb Q^\star$ with high confidence. For a given $\mathcal Q$, define the joint ambiguity set
%\[
%  \mathcal P(\mathcal Q) :=\left\{\Pi\in\mathcal M_1(\mathcal X\times\Xi):\Pi_{\bm X}=\mathbb P,\ \Pi_\xi\in\mathcal Q\right\}.
%\]
%The true residual is determined by the input, but its unknown functional relationship is not encoded in $\mathcal Q$. The set $\mathcal P(\mathcal Q)$ therefore allows all couplings with the specified marginals, including randomized conditional residual laws. The upper and lower failure probability bounds are
%\begin{equation}
%  \overline p_{\hat f}:=\sup_{\Pi\in\mathcal P(\mathcal Q)}p_{\hat f}(\Pi),
%  \label{eq-upper-bound}
%\end{equation}
%and
%\begin{equation}
%  \underline p_{\hat f}:=\inf_{\Pi\in\mathcal P(\mathcal Q)}p_{\hat f}(\Pi),
%  \label{eq-lower-bound}
%\end{equation}
%respectively. If $\mathbb Q^\star\in\mathcal Q$, then $\Pi^\star\in\mathcal P(\mathcal Q)$ and hence $\underline p_{\hat f}\le p_f\le\overline p_{\hat f}$.
%
%
%Section~\ref{sec-ambiguity-set} constructs the residual ambiguity set and establishes finite-sample coverage of the failure probability bounds, simultaneously across response thresholds. Section~\ref{sec-reformulation} then derives their Lagrangian duals, proves strong duality through a finite reformulation, and obtains closed-form expressions.
%
%
%\section{Ambiguity Set Design}
%\label{sec-ambiguity-set}
%
%Failure occurs when $\xi\le-\hat g(\bm X)$, so the residual cumulative distribution function (CDF) gives the marginal mass at or below each failure threshold. We use the D\"umbgen--Wellner (DW) construction \citep{L.J23} to control empirical CDF deviations simultaneously over all thresholds, with tolerances that vary across probability levels.
%
%Let $q$ and $p$ denote an empirical probability and a candidate probability. Following \citet{B.J79} and the CDF confidence bands of \citet{O95}, we measure their local discrepancy by the Bernoulli Kullback--Leibler divergence
%\begin{equation}
%  \phi(q,p):=q\log\frac{q}{p}+(1-q)\log\frac{1-q}{1-p},\qquad \forall q\in[0,1],\, p\in(0,1),
%  \label{eq-bernoulli-kl}
%\end{equation}
%where logarithms are natural and a term with zero numerator is interpreted as zero. For $q$ close to an interior $p$, the approximation $\phi(q,p)\approx(q-p)^2/[2p(1-p)]$ scales the deviation by the Bernoulli variance. Thus, the same absolute deviation receives a larger weight near the probability boundaries. The full divergence also covers zero and full empirical counts through $\phi(0,p)=-\log(1-p)$ and $\phi(1,p)=-\log p$. At the boundaries of the candidate probability, we set $\phi(0,0)=\phi(1,1)=0$, $\phi(q,0)=+\infty$ for $q>0$, and $\phi(q,1)=+\infty$ for $q<1$.
%
%To compare deviations simultaneously across probability levels, we use the additive correction of \citet{L.J23}. For $v\in(0,1)$, define
%\begin{equation*}
%  \rho(v):=\log\log\frac{e}{4v(1-v)},\qquad \omega(v):=\log\!\left(1+\rho(v)^2\right).
%\end{equation*}
%Before sampling, we fix the divergence parameter $s=1$, which gives \eqref{eq-bernoulli-kl}, and the correction parameter $\nu=1$. Set $\psi_0(v):=\rho(v)+\omega(v)$, with $\psi_0(0)=\psi_0(1)=+\infty$, and define the bivariate correction
%\begin{equation*}
%  \psi(q,p):=\min_{v\in[\min\{q,p\},\,\max\{q,p\}]}\psi_0(v).
%\end{equation*}
%The functions $\rho$ and $\omega$ vanish at $v=1/2$ and increase toward the probability boundaries. The minimum over the interval between $q$ and $p$ is therefore large only when both probabilities lie in the same tail. This correction balances fluctuations across probability scales in the simultaneous CDF constraint.
%
%We construct the empirical CDF from a test dataset
%\[
%  \mathcal D^{\mathrm{te}}:=\left\{\left(\bm x_j^{\mathrm{te}},g(\bm x_j^{\mathrm{te}})\right)\right\}_{j=1}^{N^{\mathrm{te}}}.
%\]
%The test inputs are independent draws from $\mathbb P$ and are independent of $\mathcal D^{\mathrm{tr}}$. We fix $N^{\mathrm{te}}\ge1$ before sampling and use these data to calibrate the residual ambiguity set.
%
%Define the signed test residuals
%\[
%  \xi_j^{\mathrm{te}}:=\xi(\bm x_j^{\mathrm{te}})=g(\bm x_j^{\mathrm{te}})-\hat g(\bm x_j^{\mathrm{te}}),\qquad j=1,\ldots,N^{\mathrm{te}}.
%\]
%Conditional on $\mathcal D^{\mathrm{tr}}$, these residuals are independent and identically distributed according to $\mathbb Q^\star$. Their empirical CDF is
%\begin{equation}
%  \widehat F^{\mathrm{te}}(t):=\frac{1}{N^{\mathrm{te}}}\sum_{j=1}^{N^{\mathrm{te}}}\mathbf 1_{\{\xi_j^{\mathrm{te}}\le t\}},\qquad \forall t\in\mathbb R.
%  \label{eq-empirical-cdf}
%\end{equation}
%For $\mathbb Q\in\mathcal M_1(\Xi)$, write $F_{\mathbb Q}(t):=\mathbb Q\left(({-\infty},t]\right)$. At a threshold $t$, $q=\widehat F^{\mathrm{te}}(t)$ is the mean of the Bernoulli indicators in \eqref{eq-empirical-cdf}, and $p=F_{\mathbb Q}(t)$ is the candidate probability. Thus $N^{\mathrm{te}}\phi(q,p)$ is the Bernoulli log-likelihood difference between the empirical and candidate probabilities. Combining it with the multiscale correction gives $N^{\mathrm{te}}\phi(q,p)-\psi(q,p)$.
%
%At the probability boundaries, the corrected discrepancy is $-\infty$ for $q=p\in\{0,1\}$ and $+\infty$ for $p=0<q$ or $p=1>q$. These conventions extend the constraint to every residual threshold.
%
%The following definition constrains the residual law through the corrected CDF discrepancy.
%
%\begin{definition}[DW-based CDF ambiguity set]
%  \label{definition-dw-ambiguity-set}
%  Given the empirical CDF $\widehat F^{\mathrm{te}}$ and a threshold $\theta\ge0$, define
%  \begin{equation} 
%    \label{eq-cdf-ambiguity-set}
%      \mathcal Q^{\mathrm{DW}}(\theta):=\left\{\mathbb Q\in\mathcal M_1(\Xi):\ 
%      \sup_{t\in\mathbb R}\left[ N^{\mathrm{te}}\phi\!\left(\widehat F^{\mathrm{te}}(t),F_{\mathbb Q}(t)\right) 
%       {}-\psi\!\left(\widehat F^{\mathrm{te}}(t),F_{\mathbb Q}(t)\right)\right]\le\theta\right\}. 
%  \end{equation} 
%\end{definition}
%
%We calibrate $\theta$ against a distribution-free reference law. Following \citet{L.J23}, let $V_1,\ldots,V_{N^{\mathrm{te}}}$ be independent standard uniform variables, independent of the training and test data. Their empirical CDF is
%\begin{equation*}
%  G_{N^{\mathrm{te}}}(u):=\frac{1}{N^{\mathrm{te}}}\sum_{j=1}^{N^{\mathrm{te}}}\mathbf 1_{\{V_j\le u\}},\qquad \forall u\in[0,1].
%\end{equation*}
%Define the uniform reference statistic
%\begin{equation}
%  T_{N^{\mathrm{te}}}^{\mathrm U}:=\sup_{u\in(0,1)}\left[N^{\mathrm{te}}\phi\!\left(G_{N^{\mathrm{te}}}(u),u\right)-\psi\!\left(G_{N^{\mathrm{te}}}(u),u\right)\right].
%  \label{eq-dw-reference-statistic}
%\end{equation}
%
%The following proposition calibrates $\theta$ to obtain finite-sample coverage of the residual law.
%
%\begin{proposition}[Finite-sample threshold calibration]
%  \label{proposition-01}
%  Under the sampling conditions stated above, let $\delta\in(0,1)$ and define
%  \begin{equation}
%    \theta_{N^{\mathrm{te}}}(\delta):=\inf\left\{\tau\ge0:\Pr\left\{T_{N^{\mathrm{te}}}^{\mathrm U}\le\tau\right\}\ge1-\delta\right\}.
%    \label{eq-dw-calibrated-threshold}
%  \end{equation}
%
%  Then, for every $\theta\ge\theta_{N^{\mathrm{te}}}(\delta)$, the DW-based CDF ambiguity set satisfies
%  \begin{equation}
%    \Pr\left\{\mathbb Q^\star\in\mathcal Q^{\mathrm{DW}}(\theta)\,\middle|\,\mathcal D^{\mathrm{tr}}\right\}\ge1-\delta.
%    \label{eq-cdf-coverage}
%  \end{equation}
%\end{proposition}
%\begin{proof}
%  The reference statistic is finite almost surely, and its distribution function is right-continuous. Hence \eqref{eq-dw-calibrated-threshold} gives
%  \[
%    \Pr\left\{T_{N^{\mathrm{te}}}^{\mathrm U}>\theta_{N^{\mathrm{te}}}(\delta)\right\}\le\delta.
%  \]
%  Conditioning on $\mathcal D^{\mathrm{tr}}$ fixes $\hat g$ and $\mathbb Q^\star$. Define the generalized quantile function $F_{\mathbb Q^\star}^{-1}(u):=\inf\{t\in\mathbb R:F_{\mathbb Q^\star}(t)\ge u\}$ for all $u\in(0,1)$. The sample $F_{\mathbb Q^\star}^{-1}(V_j)$ for $j=1,\ldots,N^{\mathrm{te}}$, has the same conditional joint law as the test residuals. Thus, given $\mathcal D^{\mathrm{tr}}$,
%  \[
%    \left(\widehat F^{\mathrm{te}}(t)\right)_{t\in\mathbb R}\overset{\mathrm d}{=}\left(G_{N^{\mathrm{te}}}\!\left(F_{\mathbb Q^\star}(t)\right)\right)_{t\in\mathbb R}.
%  \]
%  For $0<F_{\mathbb Q^\star}(t)<1$, the discrepancy at $\mathbb Q^\star$ is evaluated at a probability level included in \eqref{eq-dw-reference-statistic}. When $F_{\mathbb Q^\star}(t)$ is zero or one, the empirical CDF equals that boundary almost surely, giving discrepancy $-\infty$. The supremum in \eqref{eq-cdf-ambiguity-set}, evaluated at $\mathbb Q^\star$, is therefore bounded by $T_{N^{\mathrm{te}}}^{\mathrm U}$ under this coupling. Since $\mathbb Q^\star(\Xi)=1$, we obtain
%  \[
%    \Pr\left\{\mathbb Q^\star\notin\mathcal Q^{\mathrm{DW}}\!\left(\theta_{N^{\mathrm{te}}}(\delta)\right)\,\middle|\,\mathcal D^{\mathrm{tr}}\right\}
%    \le\Pr\left\{T_{N^{\mathrm{te}}}^{\mathrm U}>\theta_{N^{\mathrm{te}}}(\delta)\right\}\le\delta.
%  \]
%  For every test sample, $\theta\ge\theta_{N^{\mathrm{te}}}(\delta)$ implies $\mathcal Q^{\mathrm{DW}}\!\left(\theta_{N^{\mathrm{te}}}(\delta)\right)\subseteq\mathcal Q^{\mathrm{DW}}(\theta)$. This inclusion proves \eqref{eq-cdf-coverage} for every such $\theta$.
%\end{proof}
%
%Proposition~\ref{proposition-01} requires no continuity or moment assumptions on the residual law. For fixed $\phi$ and $\psi$, the reference law depends only on $N^{\mathrm{te}}$, so thresholds can be precomputed. \citet{L.J23} give finite-sample evaluation through uniform order-statistic boundary probabilities. A conservative quantile upper bound preserves \eqref{eq-cdf-coverage}.
%% If calibration is performed by simulation, quantile estimation error must be accounted for in the confidence guarantee. A threshold chosen from the test data also preserves \eqref{eq-cdf-coverage}, provided $\theta\ge\theta_{N^{\mathrm{te}}}(\delta)$ almost surely, by the same nesting argument.
%
%The following theorem transfers this coverage guarantee to the failure probability interval.
%
%\begin{theorem}[Finite-sample failure probability bounds]
%  \label{theorem-01}
%  Under the sampling conditions of Proposition~\ref{proposition-01}, let $\delta\in(0,1)$ and $\theta\ge\theta_{N^{\mathrm{te}}}(\delta)$, with the calibrated threshold defined in \eqref{eq-dw-calibrated-threshold}. Then the bounds defined in \eqref{eq-upper-bound} and \eqref{eq-lower-bound}, evaluated with $\mathcal Q=\mathcal Q^{\mathrm{DW}}(\theta)$, satisfy
%  \begin{equation}
%    \Pr\left\{\underline p_{\hat f}\le p_f\le\overline p_{\hat f}\,\middle|\,\mathcal D^{\mathrm{tr}}\right\}\ge1-\delta.
%    \label{eq-failure-probability-coverage}
%  \end{equation}
%\end{theorem}
%
%\begin{proof}
%  On the event $\mathbb Q^\star\in\mathcal Q^{\mathrm{DW}}(\theta)$, the true joint distribution $\Pi^\star$ belongs to $\mathcal P\left(\mathcal Q^{\mathrm{DW}}(\theta)\right)$. Hence
%  \[
%    \underline p_{\hat f}\le p_{\hat f}(\Pi^\star)=p_f\le\overline p_{\hat f}.
%  \]
%  By Proposition~\ref{proposition-01}, this event has conditional probability at least $1-\delta$ given $\mathcal D^{\mathrm{tr}}$, proving \eqref{eq-failure-probability-coverage}.
%\end{proof}
%
%Theorem~\ref{theorem-01} uses one residual-coverage event for both endpoints, so the confidence level applies to the entire interval. For fixed test data, $0\le\theta_1\le\theta_2$ implies $\mathcal Q^{\mathrm{DW}}(\theta_1)\subseteq\mathcal Q^{\mathrm{DW}}(\theta_2)$. The lower bound is nonincreasing and the upper bound nondecreasing in $\theta$, so higher confidence gives a wider or unchanged interval.
%
%The following corollary extends the guarantee to all response thresholds using the same residual ambiguity set.
%
%\begin{corollary}[Simultaneous bounds across response thresholds]
%  \label{cor-simultaneous-threshold-bounds}
%  Under the conditions of Theorem~\ref{theorem-01}, let $\mathcal Q=\mathcal Q^{\mathrm{DW}}(\theta)$. For $c\in\mathbb R$, write $p_f(c):=\Pr \left\{g(\bm X)\le c\right\}$ and define
%  \[
%    p_{\hat f}(\Pi;c):=\int_{\mathcal X\times\Xi}
%      \mathbf1_{\{\hat g(\bm x)+\xi\le c\}}\,\Pi(\mathrm d\bm x,\mathrm d\xi),
%  \]
%  with lower and upper bounds
%  \[
%    \underline p_{\hat f}(c):=\inf_{\Pi\in\mathcal P(\mathcal Q)}p_{\hat f}(\Pi;c),
%    \qquad
%    \overline p_{\hat f}(c):=\sup_{\Pi\in\mathcal P(\mathcal Q)}p_{\hat f}(\Pi;c).
%  \]
%  Then
%  \begin{equation}
%    \Pr\left\{\underline p_{\hat f}(c)\le p_f(c)\le\overline p_{\hat f}(c),
%      \ \forall c\in\mathbb R\,\middle|\,\mathcal D^{\mathrm{tr}}\right\}\ge1-\delta.
%    \label{eq-threshold-curve-coverage}
%  \end{equation}
%  At $c=0$, these definitions recover the original failure probability and its bounds.
%\end{corollary}
%\begin{proof}
%  On $\mathbb Q^\star\in\mathcal Q^{\mathrm{DW}}(\theta)$, the true joint law $\Pi^\star$ belongs to the same set $\mathcal P(\mathcal Q)$ for every $c$. Consequently,
%  \[ 
%    \underline p_{\hat f}(c) =\inf_{\Pi\in\mathcal P(\mathcal Q)} p_{\hat f}(\Pi;c) \le p_{\hat f}(\Pi^\star;c) =p_f(c) \le\sup_{\Pi\in\mathcal P(\mathcal Q)} p_{\hat f}(\Pi;c) =\overline p_{\hat f}(c), \qquad\forall c\in\mathbb R.
%  \]
%  Both bound curves are nondecreasing. Monotonicity and the CDF limits of $p_f$ imply that any violation is detected at a rational threshold or an atom of $g(\bm X)$, a countable set. The simultaneous event is therefore measurable. Proposition~\ref{proposition-01} gives the common residual-coverage event conditional probability at least $1-\delta$, proving \eqref{eq-threshold-curve-coverage}.
%\end{proof}
%
%Corollary~\ref{cor-simultaneous-threshold-bounds} permits response thresholds to be selected after observing the test data. Complementing the endpoints gives simultaneous bounds on $\Pr \left\{g(\bm X)>c\right\}$. The functions $g$, $\hat g$, and the law $\mathbb P$ remain fixed; fragility curves conditioned on hazard intensity require coverage for the corresponding conditional residual laws. The coverage argument also applies to any residual ambiguity set containing $\mathbb Q^\star$ with conditional probability at least $1-\delta$. It remains to evaluate the DW extrema in \eqref{eq-upper-bound} and \eqref{eq-lower-bound} and their response-threshold extensions.
%
%
%\section{Tractable Reformulation}
%\label{sec-reformulation}
%
%Throughout this section, the training data, test residuals, and $\theta\ge0$ are fixed. Input and residual cell probabilities give finite primal and dual problems and closed-form bounds. The main results use $\Xi=\mathbb R$; Remark~\ref{rem-interval-support} treats known interval support.
%
%For each empirical probability level $q_j:=j/N^{\mathrm{te}}$, define the scalar acceptance interval
%\begin{equation}
%  [\alpha_j(\theta),\beta_j(\theta)]
%  :=\left\{p\in[0,1]:
%  N^{\mathrm{te}}\phi(q_j,p)-\psi(q_j,p)\le\theta\right\},
%  \qquad j=0,\ldots,N^{\mathrm{te}}.
%  \label{eq-dw-rank-acceptance}
%\end{equation}
%Let $\xi_{(1)}^{\mathrm{te}}<\cdots<\xi_{(K)}^{\mathrm{te}}$ be the distinct ordered test residuals, and set
%\begin{equation}
%  c_0:=0,\quad
%  c_k:=\sum_{j=1}^{N^{\mathrm{te}}}\mathbf1_{\{\xi_j^{\mathrm{te}}\le\xi_{(k)}^{\mathrm{te}}\}},\quad
%  l_k:=\alpha_{c_k}(\theta),\quad
%  u_k:=\beta_{c_{k-1}}(\theta),\qquad k=1,\ldots,K.
%  \label{eq-dw-counts-and-levels}
%\end{equation}
%The counts retain all multiplicities, so $c_K=N^{\mathrm{te}}$ even when test residuals coincide. Set $\xi_{(0)}^{\mathrm{te}}:=-\infty$ and $\xi_{(K+1)}^{\mathrm{te}}:=+\infty$. Separating the observed values from the intervening intervals gives $J:=2K+1$ residual cells:
%\begin{equation}
%  \Xi_{2k}:=\left\{\xi_{(k)}^{\mathrm{te}}\right\}
%    \quad\left(k=1,\ldots,K\right),
%  \qquad
%  \Xi_{2k+1}:=\left(\xi_{(k)}^{\mathrm{te}},\xi_{(k+1)}^{\mathrm{te}}\right)
%    \quad\left(k=0,\ldots,K\right).
%  \label{eq-residual-cells}
%\end{equation}
%For any candidate residual law $\mathbb Q\in\mathcal M_1(\Xi)$, the open rank probability measures mass strictly below an observed residual, while the closed rank probability also includes mass at that value. For $k=1,\ldots,K$, the partition gives
%\[
%  \mathbb Q\left\{\xi<\xi_{(k)}^{\mathrm{te}}\right\}
%    =\sum_{j<2k}\mathbb Q(\Xi_j),
%  \qquad
%  \mathbb Q\left\{\xi\le\xi_{(k)}^{\mathrm{te}}\right\}
%    =\sum_{j\le2k}\mathbb Q(\Xi_j).
%\]
%The two probabilities differ by the atom mass $\mathbb Q(\Xi_{2k})$.
%
%For $i=0,\ldots,K$, the observed residual levels classify inputs by the location of the failure threshold $-\hat g(\bm x)$:
%\begin{equation}
%  \mathcal X_i:=\left\{\bm x\in\mathcal X:
%    \xi_{(i)}^{\mathrm{te}}\le-\hat g(\bm x)<\xi_{(i+1)}^{\mathrm{te}}\right\}.
%  \label{eq-input-cells}
%\end{equation}
%The input cell probabilities give the surrogate-response threshold probabilities through
%\[
%  \mathbb P\left\{\hat g(\bm X)\le-\xi_{(k)}^{\mathrm{te}}\right\}
%    =\sum_{i=k}^{K}\mathbb P(\mathcal X_i),
%  \qquad k=1,\ldots,K.
%\]
%The half-open convention assigns $-\hat g(\bm x)=\xi_{(k)}^{\mathrm{te}}$ to $\mathcal X_k$. The tail sum therefore includes threshold equalities even when the surrogate response has atoms. 
%The two partitions satisfy
%\[
%  \mathcal X=\bigsqcup_{i=0}^{K}\mathcal X_i,
%  \qquad
%  \Xi=\bigsqcup_{j=1}^{J}\Xi_j.
%\]
%Figure~\ref{fig-residual-input-cells} illustrates the two partitions for $K=3$.
%\begin{figure}[!htbp]
%  \centering
%  \includegraphics[width=0.5\columnwidth]{fig_residual_input_cells.png}
%  \caption{Illustration of the residual and input-threshold partitions for $K=3$. The residual partition consists of $K$ singleton cells and $K+1$ open intervals, whereas the input partition is induced by half-open intervals of the failure threshold $-\hat g(\bm x)$. Here, $w_i=\mathbb P(\mathcal X_i)$ denotes the probability mass of input cell $\mathcal X_i$.}
%  \label{fig-residual-input-cells}
%\end{figure}
%
%The following proposition expresses the DW ambiguity set through linear constraints on cumulative residual cell probabilities.
%
%\begin{proposition}[Finite representation of the DW ambiguity set]
%  \label{prop-finite-rank-constraints}
%  Each acceptance set in \eqref{eq-dw-rank-acceptance} is a nonempty closed interval containing $q_j$. The DW ambiguity set has the finite representation
%  \begin{equation}
%    \mathcal Q^{\mathrm{DW}}(\theta)
%    =\left\{\mathbb Q\in\mathcal M_1(\Xi):
%      \sum_{j<2k}\mathbb Q(\Xi_j)\le u_k,\quad
%      \sum_{j\le2k}\mathbb Q(\Xi_j)\ge l_k,\quad k=1,\ldots,K\right\}.
%    \label{eq-finite-rank-constraints}
%  \end{equation}
%  The sequences $l_k$ and $u_k$ are nondecreasing and satisfy
%  \[
%    0\le l_k,u_k\le1,\qquad l_{k-1}\le u_k,\qquad l_k\le u_{k+1},
%    \qquad k=1,\ldots,K,
%  \]
%  where $l_0:=0$ and $u_{K+1}:=1$.
%\end{proposition}
%\begin{proof}
%  For fixed $q$, the Bernoulli divergence increases as $p$ moves away from $q$, while $\psi(q,p)$ cannot increase as the interval between $q$ and $p$ expands. The corrected discrepancy therefore decreases toward $q$ from below and increases above $q$. It is continuous for $p\in(0,1)$ and equals $-\psi_0(q)\le0$ at $p=q$. The boundary conventions in Section~\ref{sec-ambiguity-set} give a nonempty closed acceptance interval also for $q\in\{0,1\}$. Scalar inversion of the DW constraint therefore gives the CDF band \citep[Supplement, Section~S.6]{L.J23}
%  \[
%    \alpha_{N^{\mathrm{te}}\widehat F^{\mathrm{te}}(t)}(\theta)
%      \le F_{\mathbb Q}(t)
%      \le\beta_{N^{\mathrm{te}}\widehat F^{\mathrm{te}}(t)}(\theta),
%    \qquad\forall t\in\mathbb R.
%  \]
%
%  For $k=1,\ldots,K-1$, the empirical CDF on $\left[\xi_{(k)}^{\mathrm{te}},\xi_{(k+1)}^{\mathrm{te}}\right)$ equals $c_k/N^{\mathrm{te}}$. Monotonicity makes the lower band constraint on this interval equivalent to $\mathbb Q\left\{\xi\le\xi_{(k)}^{\mathrm{te}}\right\}\ge\alpha_{c_k}(\theta)$. By continuity of probability from below, the upper band constraint is equivalent to $\mathbb Q\left\{\xi<\xi_{(k+1)}^{\mathrm{te}}\right\}\le\beta_{c_k}(\theta)$. The exterior intervals give $\mathbb Q\left\{\xi<\xi_{(1)}^{\mathrm{te}}\right\}\le\beta_0(\theta)$ and $\mathbb Q\left\{\xi\le\xi_{(K)}^{\mathrm{te}}\right\}\ge\alpha_{N^{\mathrm{te}}}(\theta)$; their other band constraints are automatic because $\alpha_0(\theta)=0$ and $\beta_{N^{\mathrm{te}}}(\theta)=1$. The partition identities above then yield \eqref{eq-finite-rank-constraints}.
%
%  With $p$ fixed, the corrected discrepancy decreases as $q$ approaches $p$ from either side. For empirical levels $q<q'$, an accepted value $p<q$ at level $q'$ remains accepted at level $q$; an accepted value $p>q'$ at level $q$ remains accepted at level $q'$. Since each acceptance interval contains its empirical level, both endpoints are nondecreasing. The inequalities $\alpha_{c_k}(\theta)\le c_k/N^{\mathrm{te}}\le\beta_{c_k}(\theta)$ give the stated compatibility conditions. Singleton intervals are allowed: if $q_j=1/2$ and $\theta=0$, the interval is $\{1/2\}$ because $\psi(1/2,p)=0$.
%\end{proof}
%
%Proposition~\ref{prop-finite-rank-constraints} restricts only the residual cell masses, leaving the distribution within each cell free. The scalar endpoints can be computed by one-dimensional bisection.
%
%For a feasible joint law, define its matrix of joint cell probabilities by $\pi_{ij}:=\Pi(\mathcal X_i\times\Xi_j)$. The fixed input marginal and the residual restrictions imply $\bm\pi=(\pi_{ij})\in\mathcal C_\theta$, where
%\begin{equation}
%  \mathcal C_\theta:=\left\{\bm\pi\in\mathbb R_+^{(K+1)\times J}:
%  \begin{array}{ll}
%    \displaystyle\sum_{j=1}^{J}\pi_{ij}=\mathbb P(\mathcal X_i),&i=0,\ldots,K,\\[1mm]
%    \displaystyle\sum_{i=0}^{K}\sum_{j<2k}\pi_{ij}\le u_k,&k=1,\ldots,K,\\[1mm]
%    \displaystyle\sum_{i=0}^{K}\sum_{j\le2k}\pi_{ij}\ge l_k,&k=1,\ldots,K
%  \end{array}\right\}.
%  \label{eq-cell-mass-polytope}
%\end{equation}
%To bound the failure objective, compare the indicator on each product cell:
%\begin{equation}
%  \underline h_{ij}:=\mathbf1_{\{j\le2i\}}
%  \le\mathbf1_{\{\xi\le-\hat g(\bm x)\}}
%  \le\overline h_{ij}:=\mathbf1_{\{j\le2i+1\}},
%  \qquad \forall (\bm x,\xi)\in\mathcal X_i\times\Xi_j.
%  \label{eq-cell-costs}
%\end{equation}
%For inputs in $\mathcal X_i$, the failure status can vary only within $\Xi_{2i+1}$. The other cells are uniformly failing or safe, as illustrated in Figure~\ref{fig-joint-partition} for $K=3$.
%\begin{figure}[!htbp]
%  \centering
%  \includegraphics[width=0.5\columnwidth]{fig_joint_partition.png}
%  \caption{Joint partition of the input and residual spaces for $K=3$. Each product cell has probability mass $\pi_{ij}=\Pi(\mathcal X_i\times\Xi_j)$. The blue, yellow, and red cells correspond to $(\underline h_{ij},\overline h_{ij})=(0,0)$, $(0,1)$, and $(1,1)$, respectively.}
%  \label{fig-joint-partition}
%\end{figure}
%
%The upper problem remains a linear program over joint measures, with the full input marginal fixed:
%\begin{equation}
%  \begin{aligned}
%    \overline p_{\hat f}
%    =\sup_{\Pi\in\mathcal M_+(\mathcal X\times\Xi)}\quad&
%      \int_{\mathcal X\times\Xi}\mathbf1_{\{\xi\le-\hat g(\bm x)\}}\,\Pi(\mathrm d\bm x,\mathrm d\xi)\\
%    \text{s.t.}\quad&
%      \int_{\mathcal X\times\Xi}\mathbf1_{\mathcal U}(\bm x)\,\Pi(\mathrm d\bm x,\mathrm d\xi)
%        =\mathbb P(\bm X\in\mathcal U),&&\forall\mathcal U\in\mathcal B(\mathcal X),\\
%    &\int_{\mathcal X\times\Xi}\mathbf1_{\{\xi<\xi_{(k)}^{\mathrm{te}}\}}\,\Pi(\mathrm d\bm x,\mathrm d\xi)\le u_k,
%        &&k=1,\ldots,K,\\
%    &\int_{\mathcal X\times\Xi}\mathbf1_{\{\xi\le\xi_{(k)}^{\mathrm{te}}\}}\,\Pi(\mathrm d\bm x,\mathrm d\xi)\ge l_k,
%        &&k=1,\ldots,K.
%  \end{aligned}
%  \label{eq-measure-upper}
%\end{equation}
%The lower problem minimizes the same integral under the same constraints. The marginal equalities fix both the input law and total mass one. The product of $\mathbb P$ and the empirical residual law is feasible for every $\theta\ge0$.
%
%The following lemma gives finite-dimensional Lagrangian duals and establishes weak duality.
%
%\begin{lemma}[Finite-dimensional Lagrangian duals and weak duality]
%  \label{lem-weak-duality}
%  Write $\bm u=(u_1,\ldots,u_K)^\top$ and $\bm l=(l_1,\ldots,l_K)^\top$. The nonnegative multipliers $\bm\lambda,\bm\gamma\in\mathbb R_+^K$ correspond to the upper and lower rank constraints. The input marginal equality has an unrestricted multiplier $\eta\in\mathcal F_{\mathrm b}(\mathcal X)$, the space of bounded Borel measurable real-valued functions. Define the residual penalty
%  \begin{equation}
%    \zeta_{\bm\lambda,\bm\gamma}(\xi)
%    :=\sum_{k=1}^{K}\lambda_k\mathbf1_{\{\xi<\xi_{(k)}^{\mathrm{te}}\}}
%      -\sum_{k=1}^{K}\gamma_k\mathbf1_{\{\xi\le\xi_{(k)}^{\mathrm{te}}\}},
%    \label{eq-dual-step-function}
%  \end{equation}
%  whose value on $\Xi_j$ is
%  \begin{equation}
%    \zeta_j:=\sum_{k=1}^{K}\lambda_k\mathbf1_{\{j<2k\}}
%      -\sum_{k=1}^{K}\gamma_k\mathbf1_{\{j\le2k\}},
%    \qquad j=1,\ldots,J.
%    \label{eq-dual-cell-values}
%  \end{equation}
%  Eliminating the input marginal multiplier gives the upper and lower Lagrangian dual problems
%  \begin{align}
%    \overline d_\theta
%      &:=\inf_{\bm\lambda,\bm\gamma\ge\bm0}\left\{
%        \bm\lambda^\top\bm u-\bm\gamma^\top\bm l
%        +\sum_{i=0}^{K}\mathbb P(\mathcal X_i)\max_{1\le j\le J}
%          \left\{\overline h_{ij}-\zeta_j\right\}\right\},
%      \label{eq-upper-measure-dual}\\
%    \underline d_\theta
%      &:=\sup_{\bm\lambda,\bm\gamma\ge\bm0}\left\{
%        -\bm\lambda^\top\bm u+\bm\gamma^\top\bm l
%        +\sum_{i=0}^{K}\mathbb P(\mathcal X_i)\min_{1\le j\le J}
%          \left\{\underline h_{ij}+\zeta_j\right\}\right\}.
%      \label{eq-lower-measure-dual}
%  \end{align}
%  Their optimal values satisfy
%  \begin{equation}
%    \underline d_\theta\le\underline p_{\hat f}
%      \le\overline p_{\hat f}\le\overline d_\theta.
%    \label{eq-weak-duality}
%  \end{equation}
%\end{lemma}
%\begin{proof}
%  The input marginal equality is equivalent to
%  \[
%    \int_{\mathcal X\times\Xi}\eta(\bm x)\,\Pi(\mathrm d\bm x,\mathrm d\xi)
%      =\mathbb E_{\mathbb P}\left[\eta(\bm X)\right],
%    \qquad\forall\eta\in\mathcal F_{\mathrm b}(\mathcal X).
%  \]
%  Relaxing this equality and the rank inequalities gives the upper Lagrangian \citep[Section~5.1]{B.V04}
%  \[
%    \begin{aligned}
%      \mathcal L^+(\Pi,\eta,\bm\lambda,\bm\gamma)
%      ={}&\int_{\mathcal X\times\Xi}\mathbf1_{\{\xi\le-\hat g(\bm x)\}}\,\Pi(\mathrm d\bm x,\mathrm d\xi)
%        +\mathbb E_{\mathbb P}\left[\eta(\bm X)\right]
%        -\int_{\mathcal X\times\Xi}\eta(\bm x)\,\Pi(\mathrm d\bm x,\mathrm d\xi)\\
%      &+\sum_{k=1}^{K}\lambda_k\left(u_k-\int_{\mathcal X\times\Xi}
%          \mathbf1_{\{\xi<\xi_{(k)}^{\mathrm{te}}\}}\,\Pi(\mathrm d\bm x,\mathrm d\xi)\right)
%        +\sum_{k=1}^{K}\gamma_k\left(\int_{\mathcal X\times\Xi}
%          \mathbf1_{\{\xi\le\xi_{(k)}^{\mathrm{te}}\}}\,\Pi(\mathrm d\bm x,\mathrm d\xi)-l_k\right)\\
%      ={}&\mathbb E_{\mathbb P}\left[\eta(\bm X)\right]
%        +\bm\lambda^\top\bm u-\bm\gamma^\top\bm l
%        +\int_{\mathcal X\times\Xi}\left[
%          \mathbf1_{\{\xi\le-\hat g(\bm x)\}}-\eta(\bm x)-\zeta_{\bm\lambda,\bm\gamma}(\xi)
%        \right]\,\Pi(\mathrm d\bm x,\mathrm d\xi).
%    \end{aligned}
%  \]
%  Its supremum over finite nonnegative measures is finite exactly when the integrand is nonpositive at every $(\bm x,\xi)$. A positive value at any point permits an arbitrarily large point mass there; otherwise the supremum integral is zero, attained by the zero measure. The functional upper dual therefore minimizes $\mathbb E_{\mathbb P}[\eta(\bm X)]+\bm\lambda^\top\bm u-\bm\gamma^\top\bm l$ subject to $\eta(\bm x)\ge\mathbf1_{\{\xi\le-\hat g(\bm x)\}}-\zeta_{\bm\lambda,\bm\gamma}(\xi)$ at every pair. For the lower problem, reversing the rank penalties gives
%  \[
%    \mathcal L^-
%      =\mathbb E_{\mathbb P}\left[\eta(\bm X)\right]
%        -\bm\lambda^\top\bm u+\bm\gamma^\top\bm l
%        +\int_{\mathcal X\times\Xi}\left[
%          \mathbf1_{\{\xi\le-\hat g(\bm x)\}}-\eta(\bm x)+\zeta_{\bm\lambda,\bm\gamma}(\xi)
%        \right]\,\Pi(\mathrm d\bm x,\mathrm d\xi).
%  \]
%  Its infimum over finite nonnegative measures is finite exactly when the integrand is nonnegative everywhere, in which case the infimum integral is zero. The lower functional dual maximizes $\mathbb E_{\mathbb P}[\eta(\bm X)]-\bm\lambda^\top\bm u+\bm\gamma^\top\bm l$ subject to $\eta(\bm x)\le\mathbf1_{\{\xi\le-\hat g(\bm x)\}}+\zeta_{\bm\lambda,\bm\gamma}(\xi)$ at every pair.
%
%  For fixed rank multipliers, the best marginal functions are the pointwise envelopes. For $\bm x\in\mathcal X_i$, they satisfy the identities
%  \begin{equation}
%    \begin{aligned}
%      \eta^+_{\bm\lambda,\bm\gamma}(\bm x)
%        &:=\sup_{\xi\in\Xi}\left\{\mathbf1_{\{\xi\le-\hat g(\bm x)\}}
%          -\zeta_{\bm\lambda,\bm\gamma}(\xi)\right\}
%          =\max_{1\le j\le J}\left\{\overline h_{ij}-\zeta_j\right\},\\
%      \eta^-_{\bm\lambda,\bm\gamma}(\bm x)
%        &:=\inf_{\xi\in\Xi}\left\{\mathbf1_{\{\xi\le-\hat g(\bm x)\}}
%          +\zeta_{\bm\lambda,\bm\gamma}(\xi)\right\}
%          =\min_{1\le j\le J}\left\{\underline h_{ij}+\zeta_j\right\}.
%    \end{aligned}
%    \label{eq-dual-envelope-identities}
%  \end{equation}
%  To verify them, note that outside $\Xi_{2i+1}$ the two cell coefficients equal the actual failure indicator. In the overlapping interval, the lower coefficient zero is attained by $\xi=\left(\xi_{(i+1)}^{\mathrm{te}}-\hat g(\bm x)\right)/2$ if $i<K$, and by $\xi=-\hat g(\bm x)+1$ if $i=K$. The upper coefficient one is attained by $\xi=-\hat g(\bm x)$ whenever $-\hat g(\bm x)>\xi_{(i)}^{\mathrm{te}}$. The only remaining case is $-\hat g(\bm x)=\xi_{(i)}^{\mathrm{te}}$ with $i\ge1$. Then the overlapping interval is safe, but
%  \[
%    \zeta_{2i}=\zeta_{2i+1}-\gamma_i\le\zeta_{2i+1},
%    \qquad
%    1-\zeta_{2i}\ge1-\zeta_{2i+1}.
%  \]
%  Choosing the left endpoint therefore attains an upper-envelope value at least as large as the candidate from the overlapping interval. Together with \eqref{eq-cell-costs}, this proves both identities, including threshold equalities.
%
%  The envelopes are bounded and constant on the Borel cells $\mathcal X_i$, so they belong to $\mathcal F_{\mathrm b}(\mathcal X)$. They respectively minimize and maximize the expectation over feasible marginal functions. Substituting them into the functional dual objectives gives \eqref{eq-upper-measure-dual} and \eqref{eq-lower-measure-dual}. Introducing epigraph variables $z_i$ for the upper maxima and hypograph variables $z_i$ for the lower minima gives the equivalent LPs
%  \begin{equation}
%    \begin{aligned}
%      \overline d_\theta
%      =\inf_{\bm\lambda,\bm\gamma\ge\bm0,\ \bm z\in\mathbb R^{K+1}}\quad&
%        \bm\lambda^\top\bm u-\bm\gamma^\top\bm l+\sum_{i=0}^{K}\mathbb P(\mathcal X_i) z_i\\
%      \text{s.t.}\quad&z_i+\zeta_j\ge\overline h_{ij},
%        \qquad i=0,\ldots,K,\quad j=1,\ldots,J,
%    \end{aligned}
%    \label{eq-finite-upper-dual}
%  \end{equation}
%  and
%  \begin{equation}
%    \begin{aligned}
%      \underline d_\theta
%      =\sup_{\bm\lambda,\bm\gamma\ge\bm0,\ \bm z\in\mathbb R^{K+1}}\quad&
%        -\bm\lambda^\top\bm u+\bm\gamma^\top\bm l+\sum_{i=0}^{K}\mathbb P(\mathcal X_i) z_i\\
%      \text{s.t.}\quad&z_i-\zeta_j\le\underline h_{ij},
%        \qquad i=0,\ldots,K,\quad j=1,\ldots,J.
%    \end{aligned}
%    \label{eq-finite-lower-dual}
%  \end{equation}
%  Each $z_i$ can be chosen as its corresponding finite envelope, including on input cells with $\mathbb P(\mathcal X_i)=0$. Thus introducing these variables preserves both optimal values.
%
%  For a feasible joint law, the signs of the rank penalties give $\mathcal L^-\le p_{\hat f}(\Pi)\le\mathcal L^+$. Feasible lower dual objectives therefore bound the primal objectives from below, and feasible upper dual objectives bound them from above. Taking the respective optima over the common nonempty primal feasible set proves \eqref{eq-weak-duality}.
%\end{proof}
%
%Lemma~\ref{lem-weak-duality} eliminates the function-valued marginal multiplier, leaving $2K$ nonnegative multipliers. The equivalent LPs \eqref{eq-finite-upper-dual} and \eqref{eq-finite-lower-dual} introduce $K+1$ free variables $z_i$ and $(K+1)(2K+1)$ pairwise inequalities.
%
%The following theorem proves that the finite cell allocations and their duals recover the original failure probability bounds.
%
%\begin{theorem}[Exact finite reformulations and strong duality]
%  \label{thm-finite-dual-reformulations}
%  For $\mathcal Q=\mathcal Q^{\mathrm{DW}}(\theta)$ and every $\theta\ge0$, the failure probability bounds admit the exact representations
%  \begin{equation}
%    \overline p_{\hat f}
%      =\max_{\bm\pi\in\mathcal C_\theta}\sum_{i=0}^{K}\sum_{j=1}^{J}\overline h_{ij}\pi_{ij},
%    \qquad
%    \underline p_{\hat f}
%      =\min_{\bm\pi\in\mathcal C_\theta}\sum_{i=0}^{K}\sum_{j=1}^{J}\underline h_{ij}\pi_{ij}.
%    \label{eq-finite-primal-values}
%  \end{equation}
%  Their Lagrangian dual values in Lemma~\ref{lem-weak-duality} satisfy
%  \begin{equation}
%    \overline p_{\hat f}=\overline d_\theta,
%    \qquad
%    \underline p_{\hat f}=\underline d_\theta.
%    \label{eq-strong-duality}
%  \end{equation}
%  The finite primal and dual optima, the measure primal optima, and the functional dual optima are attained.
%\end{theorem}
%\begin{proof}
%  Any feasible joint law induces a matrix in $\mathcal C_\theta$. Integrating \eqref{eq-cell-costs} therefore bounds the upper measure problem by the finite maximization from above and the lower measure problem by the finite minimization from below.
%
%  For the reverse inequalities, fix $\bm\pi\in\mathcal C_\theta$. Draw $\bm x$ from $\mathbb P$. On $\mathcal X_i$ with $\mathbb P(\mathcal X_i)>0$, select a provisional residual cell $j$ with probability $\pi_{ij}/\mathbb P(\mathcal X_i)$; the row equality makes these probabilities sum to one. For $j\ne2i+1$, choose a fixed representative of $\Xi_j$, where the actual failure indicator equals both coefficients. In the overlapping cell, use the following assignments to define joint laws $\Pi^+$ and $\Pi^-$.
%
%  For the upper problem, assign $\xi=-\hat g(\bm x)$. This causes failure and remains in $\Xi_{2i+1}$ unless $-\hat g(\bm x)=\xi_{(i)}^{\mathrm{te}}$. At equality, which requires $i\ge1$, the assignment transfers the provisional mass to the left endpoint. For any $\widetilde\xi\in\left(\xi_{(i)}^{\mathrm{te}},\xi_{(i+1)}^{\mathrm{te}}\right)$ and every observed threshold,
%  \[
%    \mathbf1_{\{\xi_{(i)}^{\mathrm{te}}<\xi_{(k)}^{\mathrm{te}}\}}
%      =\mathbf1_{\{\widetilde\xi<\xi_{(k)}^{\mathrm{te}}\}},
%    \qquad
%    \mathbf1_{\{\xi_{(i)}^{\mathrm{te}}\le\xi_{(k)}^{\mathrm{te}}\}}
%      \ge\mathbf1_{\{\widetilde\xi\le\xi_{(k)}^{\mathrm{te}}\}}.
%  \]
%  Thus every open upper probability is unchanged and every closed lower probability can only increase. The upper assignment preserves all rank inequalities and attains the coefficient $\overline h_{i,2i+1}=1$.
%
%  For the lower problem, assign $\xi=\left(\xi_{(i+1)}^{\mathrm{te}}-\hat g(\bm x)\right)/2$ if $i<K$, and $\xi=-\hat g(\bm x)+1$ if $i=K$. These values remain in $\Xi_{2i+1}$ and strictly exceed the failure threshold, including at threshold equality. Hence the lower assignment retains the cell matrix and attains $\underline h_{i,2i+1}=0$.
%
%  If $\mathbb P(\mathcal X_i)=0$, the corresponding row is zero; set $\xi=\xi_{(1)}^{\mathrm{te}}$ on that input cell. All assignments are Borel measurable, and their finite mixtures define probability kernels. Combining them with $\mathbb P$ gives joint laws with the full prescribed input marginal, including randomization at input atoms. Both laws are feasible and satisfy
%  \[
%    p_{\hat f}(\Pi^+)=\sum_{i=0}^{K}\sum_{j=1}^{J}\overline h_{ij}\pi_{ij},
%    \qquad
%    p_{\hat f}(\Pi^-)=\sum_{i=0}^{K}\sum_{j=1}^{J}\underline h_{ij}\pi_{ij}.
%  \]
%  The lower law retains every provisional cell mass. The upper law may change the cell matrix at threshold equalities, but it preserves feasibility and the assigned upper objective. These identities establish the reverse inequalities required for \eqref{eq-finite-primal-values}.
%
%  For every $\theta\ge0$, the empirical product law induces the feasible matrix
%  \[
%    \pi_{i,2k}^{\mathrm{emp}}:=\mathbb P(\mathcal X_i)\frac{c_k-c_{k-1}}{N^{\mathrm{te}}},
%    \qquad i=0,\ldots,K,\quad k=1,\ldots,K,
%  \]
%  with all odd columns zero. Its open and closed rank probabilities are $c_{k-1}/N^{\mathrm{te}}$ and $c_k/N^{\mathrm{te}}$, respectively, so feasibility includes $\theta=0$. The nonnegative entries of a matrix in $\mathcal C_\theta$ sum to one, and the constraints are closed. Thus $\mathcal C_\theta$ is nonempty and compact. The finite objectives lie in $[0,1]$ and attain their optima, which the joint-law constructions realize in the measure problems.
%
%  The finite primal LPs in \eqref{eq-finite-primal-values} have row multipliers $z_i$ and rank multipliers $\lambda_k,\gamma_k$. Their dual constraints are precisely $z_i+\zeta_j\ge\overline h_{ij}$ and $z_i-\zeta_j\le\underline h_{ij}$. Finite LP strong duality therefore gives \eqref{eq-strong-duality} and dual attainment \citep[Section~5.2.3]{B.V04}. By \eqref{eq-dual-envelope-identities}, optimal rank multipliers also yield optimal bounded marginal functions for the original functional duals. Feasibility and finite optimal values suffice; strict rank inequalities are unnecessary, including at $\theta=0$.
%\end{proof}
%
%Theorem~\ref{thm-finite-dual-reformulations} reduces the optimization to joint cell masses while retaining the full input law through the within-cell construction.
%
%The following corollary evaluates both bounds in closed form and identifies optimal dual certificates.
%
%\begin{corollary}[Closed-form bounds and optimal certificates]
%  \label{thm-closed-form-bounds}
%  Under the conditions of Theorem~\ref{thm-finite-dual-reformulations},
%  \begin{align}
%    \overline p_{\hat f}
%      &=\min\left\{1,\min_{k=1,\ldots,K}\left[u_k+\sum_{i=k}^{K}\mathbb P(\mathcal X_i)\right]\right\},
%      \label{eq-dw-closed-upper}\\
%    \underline p_{\hat f}
%      &=\max\left\{0,\max_{k=1,\ldots,K}\left[l_k+\sum_{i=k}^{K}\mathbb P(\mathcal X_i)-1\right]\right\}.
%      \label{eq-dw-closed-lower}
%  \end{align}
%  Each dual optimum admits a certificate with at most one nonzero rank multiplier, of value one.
%\end{corollary}
%\begin{proof}
%  For a candidate residual law $\mathbb Q\in\mathcal M_1(\Xi)$, define the finite transportation set
%  \[
%    \mathcal T(\mathbb Q):=\left\{\bm\pi\in\mathbb R_+^{(K+1)\times J}:
%    \begin{array}{ll}
%      \displaystyle\sum_{j=1}^{J}\pi_{ij}=\mathbb P(\mathcal X_i),&i=0,\ldots,K,\\[1mm]
%      \displaystyle\sum_{i=0}^{K}\pi_{ij}=\mathbb Q(\Xi_j),&j=1,\ldots,J
%    \end{array}\right\}.
%  \]
%  It contains the product matrix $\pi_{ij}=\mathbb P(\mathcal X_i)\mathbb Q(\Xi_j)$. Every residual cell is nonempty, so any nonnegative column allocation with total mass one is realized by a residual law. Proposition~\ref{prop-finite-rank-constraints} consequently gives
%  \[
%    \mathcal C_\theta
%      =\bigcup_{\mathbb Q\in\mathcal Q^{\mathrm{DW}}(\theta)}\mathcal T(\mathbb Q).
%  \]
%  The parameter $\mathbb Q$ fixes only residual partition masses in this finite matrix problem. The joint-law construction in Theorem~\ref{thm-finite-dual-reformulations} preserves feasibility and the assigned objective while allowing the residual marginal to change, including endpoint transfers in the upper problem.
%
%  For fixed cell masses $\mathbb Q(\Xi_j)$, the upper problem maximizes mass on pairs $j\le2i+1$; the lower problem maximizes safe mass on pairs $j>2i$ and subtracts it from one. Both are maximum-flow problems with source-to-row capacities $\mathbb P(\mathcal X_i)$, column-to-sink capacities $\mathbb Q(\Xi_j)$, and capacity two on each permitted row-to-column edge. A transportation matrix restricts to a feasible flow. Conversely, any flow can be completed to a transportation matrix because the remaining row and column masses are nonnegative with equal totals. Thus the maximum flow equals the maximum permitted mass.
%
%  A cut of capacity one is available, so a minimum cut crosses no edge of capacity two. Its source-side columns can be taken to be exactly the neighbors of its source-side rows. Since the neighborhoods are nested, adding rows whose neighbors are already included cannot increase the capacity. Apart from empty and full row sets, upper-problem cuts reduce to rows $i<k$ and columns $j<2k$, while safe-mass cuts reduce to rows $i\ge k$ and columns $j>2k$. For $k=1,\ldots,K$, their capacities are respectively
%  \[
%    \sum_{i=k}^{K}\mathbb P(\mathcal X_i)+\sum_{j<2k}\mathbb Q(\Xi_j),
%    \qquad
%    \sum_{i=0}^{k-1}\mathbb P(\mathcal X_i)+\sum_{j>2k}\mathbb Q(\Xi_j).
%  \]
%  Empty and full row sets give capacity one. The max-flow/min-cut theorem \citep{F.F56} therefore yields
%  \begin{align*}
%    \max_{\bm\pi\in\mathcal T(\mathbb Q)}\sum_{i,j}\overline h_{ij}\pi_{ij}
%      &=\min\left\{1,\min_k\left[\sum_{i=k}^{K}\mathbb P(\mathcal X_i)
%        +\sum_{j<2k}\mathbb Q(\Xi_j)\right]\right\},\\
%    \min_{\bm\pi\in\mathcal T(\mathbb Q)}\sum_{i,j}\underline h_{ij}\pi_{ij}
%      &=\max\left\{0,\max_k\left[\sum_{i=k}^{K}\mathbb P(\mathcal X_i)
%        +\sum_{j\le2k}\mathbb Q(\Xi_j)-1\right]\right\}.
%  \end{align*}
%
%  These values are nondecreasing in the respective open and closed cumulative residual probabilities. The rank constraints bound the former from above by $u_k$ and the latter from below by $l_k$. To attain all relevant limits simultaneously, choose laws $\mathbb Q^+$ and $\mathbb Q^-$ with cell probabilities
%  \begin{align*}
%    &\mathbb Q^+(\Xi_1):=u_1,\quad
%      \mathbb Q^+(\Xi_{2k}):=u_{k+1}-u_k,\quad
%      \mathbb Q^+(\Xi_{2k+1}):=0,\\
%    &\mathbb Q^-(\Xi_{2k-1}):=0,\quad
%      \mathbb Q^-(\Xi_{2k}):=l_k-l_{k-1},\quad
%      \mathbb Q^-(\Xi_J):=1-l_K,
%      \qquad k=1,\ldots,K.
%  \end{align*}
%  By Proposition~\ref{prop-finite-rank-constraints}, the prescribed masses are nonnegative and sum to one. Choosing a representative in each residual cell realizes both laws. Their cumulative probabilities satisfy
%  \begin{align*}
%    &\sum_{j<2k}\mathbb Q^+(\Xi_j)=u_k,
%      \qquad \sum_{j\le2k}\mathbb Q^+(\Xi_j)=u_{k+1}\ge l_k,\\
%    &\sum_{j\le2k}\mathbb Q^-(\Xi_j)=l_k,
%      \qquad \sum_{j<2k}\mathbb Q^-(\Xi_j)=l_{k-1}\le u_k.
%  \end{align*}
%  Both laws therefore belong to $\mathcal Q^{\mathrm{DW}}(\theta)$. The union representation of $\mathcal C_\theta$ and \eqref{eq-finite-primal-values} show that these allocations attain the upper and lower bounds in \eqref{eq-dw-closed-upper} and \eqref{eq-dw-closed-lower}. Zero masses and coincident cumulative limits are allowed, including at $\theta=0$.
%
%  Finally, let $\bm e_k$ be the $k$th unit vector. Certificates for the upper and lower LPs \eqref{eq-finite-upper-dual} and \eqref{eq-finite-lower-dual} are respectively
%  \begin{align*}
%    &\bm\lambda=\bm e_k,\qquad\bm\gamma=\bm0,\qquad
%      z_i=\mathbf1_{\{i\ge k\}},\\
%    &\bm\lambda=\bm0,\qquad\bm\gamma=\bm e_k,\qquad
%      z_i=\mathbf1_{\{i\ge k\}}-1.
%  \end{align*}
%  Their feasibility follows from
%  \[
%    \mathbf1_{\{i\ge k\}}+\mathbf1_{\{j<2k\}}\ge\overline h_{ij},
%    \qquad
%    \mathbf1_{\{i\ge k\}}-1+\mathbf1_{\{j\le2k\}}\le\underline h_{ij}.
%  \]
%  Their objectives are $u_k+\sum_{i=k}^{K}\mathbb P(\mathcal X_i)$ and $l_k+\sum_{i=k}^{K}\mathbb P(\mathcal X_i)-1$, the inner candidates in the closed-form expressions. Zero rank multipliers with $z_i=1$ or $z_i=0$ supply the outer candidates. Selecting the best candidate gives an optimal full dual solution, whose functional marginal multiplier follows from \eqref{eq-dual-envelope-identities}. The inclusion of $i=k$ in the tail sum retains inputs whose failure threshold equals the observed residual.
%\end{proof}
%
%Corollary~\ref{thm-closed-form-bounds} extends the classical fixed-marginal sum problem studied by \citet{M82} and \citet{F.N.S87} to the setting where the residual law varies within a calibrated DW band, with open and closed rank constraints accounting for possible atoms. Even when both marginals are known, dependence uncertainty may lead to substantially different failure probabilities. For equally likely surrogate responses $0$ and $2$ and equally likely residuals $-1$ and $1$, the failure probability ranges from $0$ to $1/2$ over all couplings, whereas independence gives $1/4$. Thus, further calibration of the residual marginal alone cannot eliminate dependence uncertainty.
%
%For the response-threshold curves in Corollary~\ref{cor-simultaneous-threshold-bounds}, applying the same construction to $g-c$ and $\hat g-c$ leaves the residuals, rank counts, and endpoints $l_k,u_k$ unchanged. Only the input partition depends on $c$. Hence, Corollary~\ref{thm-closed-form-bounds} yields
%\begin{align}
%  \overline p_{\hat f}(c)
%    &=\min\left\{1,\min_{k=1,\ldots,K}\left[u_k
%      +\mathbb P\left\{\hat g(\bm X)\le
%      c-\xi_{(k)}^{\mathrm{te}}\right\}\right]\right\},
%    \label{eq-threshold-closed-upper}\\
%  \underline p_{\hat f}(c)
%    &=\max\left\{0,\max_{k=1,\ldots,K}\left[l_k
%      +\mathbb P\left\{\hat g(\bm X)\le
%      c-\xi_{(k)}^{\mathrm{te}}\right\}-1\right]\right\}.
%    \label{eq-threshold-closed-lower}
%\end{align}
%These identities hold for every $c\in\mathbb R$ and $\theta\ge0$. When $\theta\ge\theta_{N^{\mathrm{te}}}(\delta)$, the coverage is simultaneous by \eqref{eq-threshold-curve-coverage}, even though the extremizing couplings may vary with $c$. In particular, for a target probability $p_{\mathrm{tar}}\in(0,1)$, the condition $\overline p_{\hat f}(c)\le p_{\mathrm{tar}}$ holds if and only if
%\[
%  u_k+\mathbb P\left\{\hat g(\bm X)\le
%  c-\xi_{(k)}^{\mathrm{te}}\right\}
%  \le p_{\mathrm{tar}}
%\]
%for at least one $k\in\{1,\ldots,K\}$. 
%
%The following remark incorporates known residual support into the bounds.
%
%\begin{remark}[Known interval support]
%  \label{rem-interval-support}
%  Suppose $\Xi=[a,b]$ with $-\infty<a\le b<\infty$ and the observed residuals lie in this interval. Then
%  \begin{align*}
%    \overline p_{\hat f}
%      &=\min\left\{\mathbb P\left\{\hat g(\bm X)\le-a\right\},
%        \min_{k=1,\ldots,K}\left[u_k+\mathbb P\left\{\hat g(\bm X)\le-\xi_{(k)}^{\mathrm{te}}\right\}\right]\right\},\\
%    \underline p_{\hat f}
%      &=\max\left\{\mathbb P\left\{\hat g(\bm X)\le-b\right\},
%        \max_{k=1,\ldots,K}\left[l_k+\mathbb P\left\{\hat g(\bm X)\le-\xi_{(k)}^{\mathrm{te}}\right\}-1\right]\right\}.
%  \end{align*}
%  Represent the support restriction on $\mathbb R$ by $\mathbb Q\{\xi<a\}\le0$ and $\mathbb Q\{\xi\le b\}\ge1$, and add $a,b$ to the ordered threshold list. At a newly added $a$, use upper and lower bounds zero; at a newly added $b$, use upper and lower bounds one. At coincident observed thresholds, tighten only the upper bound at $a$ and the lower bound at $b$. The enlarged sequences satisfy the compatibility conditions in Proposition~\ref{prop-finite-rank-constraints}, so the same finite-aggregation and ordered-cut arguments apply. The effective additional upper candidate is $\mathbb P\left\{\hat g(\bm X)\le-a\right\}$, and the effective additional lower candidate is $\mathbb P\left\{\hat g(\bm X)\le-b\right\}$; the other added candidates are redundant. If $a=b$, both bounds equal $\mathbb P\left\{\hat g(\bm X)\le-a\right\}$. For the partition on $\Xi=[a,b]$, intersect the residual cells from the augmented threshold list with $\Xi$ and omit empty cells. Including the support endpoints retains their atoms and the corresponding changes in failure status. For a response threshold $c$, shift each surrogate-response threshold in the displayed event probabilities by $c$.
%\end{remark}
%
%Remark~\ref{rem-interval-support} identifies inputs that must fail, $\hat g(\bm X)\le-b$, and those that are always safe, $\hat g(\bm X)>-a$. Thus $\mathbb P\left\{\hat g(\bm X)\le-a\right\}=0$ rules out failure.
%
%% After zero failures in $N^{\mathrm{te}}$ independent tests, the one-sided exact binomial upper limit is $1-\delta^{1/N^{\mathrm{te}}}\approx\log(1/\delta)/N^{\mathrm{te}}$ \citep{C.P34}. At $c=0$, the simple DW case $K=1$ and $\mathbb P(\mathcal X_1)=0$ gives an upper bound $\beta_0(\theta)$ determined by the simultaneous-band calibration.
%
%% Evaluation requires scalar DW inversion, sorting and counting the test residuals, and computing the $K$ input event probabilities $\mathbb P\left\{\hat g(\bm X)\le c-\xi_{(k)}^{\mathrm{te}}\right\}$ at each response threshold $c$. The endpoint table depends only on $N^{\mathrm{te}}$ and $\theta$ and can be precomputed. Given this table and the input event probabilities, the two scans cost $O(K)$ at each response threshold.
%
%% For $\theta\ge\theta_{N^{\mathrm{te}}}(\delta)$, the exact expressions retain the coverage in Theorem~\ref{theorem-01} and Corollary~\ref{cor-simultaneous-threshold-bounds}. Independent auxiliary simulation of $\bm X\sim\mathbb P$ can supply simultaneous upper and lower confidence bounds on the input event probabilities. Substituting these bounds in the respective upper and lower expressions preserves conservative failure probability bounds. For one fixed $c$, simultaneous bounds at its $K$ evaluation thresholds suffice. To retain coverage of the entire response-threshold curve, the auxiliary bounds must hold for the surrogate-response CDF at every real threshold. They must hold conditionally on the test data, with their error probability included in the overall confidence budget.
%
%\section{Implementation Procedure}
%
%
% END PRESERVED BODY BEFORE IMPLEMENTATION PROCEDURE
